\documentclass[11pt,a4paper]{article}
\makeatletter\def\input@path{{sections/paper/}{sections/paper/}}\makeatother
\input{sections/common.tex}
\fancyhead[R]{\small v43}
\hypersetup{pdftitle={一般型射影曲面のLagrangian torus fibrationと負Kähler--Einstein極限},pdfauthor={}}
\title{一般型射影曲面の\\Lagrangian torus fibrationと\\負Kähler--Einstein極限\\[0.5em]\large 相対SNC構成、分岐底空間、作用適合matroidal atlas}
\author{}
\date{2026年11月22日\quad 第43稿・現行改訂43}
\begin{document}
\maketitle
\begin{abstract}
任意の固定次数 $d\ge5$ と本文の一般性条件を満たす全平面退化に対し，
actual 負 Kähler--Einstein 形式 $\operatorname{Ric}(\omega_t)=-\omega_t$ に関する
大域 Lagrangian torus fibration
$\pi_t:X_t\to B_d(T)$ を構成する．
底空間はコンパクトな三価分岐二次元複体で，正則 fiber は $T^2$，
特異 fibers は nodal，branch，component core の模型を持つ．
標識付き収縮 $p_T:B_d(T)\to\mathcal C_d$ と合成した同じ射影 $f_t=p_T\pi_t$ は，
$R_t=\log(-\log|t|^2)$ に対して
\[
 \sup_{x,y\in X_t}
 \left|R_t^{-1}d_{g_t}(x,y)-d_\square(f_t(x),f_t(y))\right|\longrightarrow0
\]
を満たす．ここで $\mathcal C_d=\{w\in[0,1]^d:\max w_i=1,\,
|\operatorname{supp}w|\le3\}$ には unit-square 長さ距離を入れる．
従って正規化 GH 極限を actual fibration の射影によって実現し，
六次以上で $p_T$ が潰す各内部木上の全 fiber 合併も一様に収縮する．
六次以上の極限直径は $2\sqrt5$ である．
同じ射影による確率 KE 体積の極限は，$d$ 個の成分頂点に重み $1/d$ を持つ原子測度である．
全 cubical 計量空間を保持した測度収束と，その有限な測度の台を区別する．

構成では，共通 normal moment と reduced 葉を保つ相対平均化により
source を含む neck と closed triple corners を接着し，
positive-residue primitive の compact core filling で全曲面を覆う．
Source の中央全点径は actual 尺度で $O(1)$，全共通 moment での
残余境界への退出費用も $O(1)$ である．
後者は source 軸を含む oriented arm の二側 tensor 比較と直接経路による．
距離下界は，source を除かない正則 normal frame 上の scalar Schwarz 障壁と
全曲面上の clipped-depth map から導く．上界は，共通 moment の接続と
高々 $60$ 本の square gallery から得る．
Component の cut 深度を $N_j$ とすると，cusp の動径経路による接続費用は
$C(1+\log(N_j+1))$ であり，残余基底を先に選ぶことで全 fiber 合併にも同じ予算を使える．
関連 KE 原稿の未検証の大域比較仮説を用いず，明記した局所解析入力から
同じ $f_t$ の全域評価を直接証明する．

一般の symplectic 構成としては，正の weighted Kähler 表示を持ち，各成分に二つ以上の交点がある
任意 genus の SNC divisor に対する有限相対 compactification を与える．
$k\ge3$ の次数分割 $d=\sum e_i$ から，
$k$ 個の cores と $d\sum_{i<j}e_i e_j$ 個の nodes を持つ fibration を得て，
$[\omega]=aH$，$a>0$ の指定 Kähler 形式への適用，基底の整数位相，
順序付き小解消の vertical Kähler slice を記述する．
全平面の五次・六次では，focus singularities を許す大域作用適合 matroidal atlas を構成する．
この選択で五次の各 triangle に五 nodes，六次の各 triangle に三 nodes，各 bigon に二 nodes が入り，
全 branch の正規化長さは $a/(\pi(d-1))$ である．
混合五次分割 $(1,1,1,2)$ は局所作用適合 cores を持つが，
四 puncture neck の分母剛性と法円作用の分岐係数により大域 atlas を持たない．
比例曲率 SNC cores と branch graph を保ち追加の branch singularities を入れない範囲では，
作用適合 matroidal atlas が可能な次数分割は全平面の五次・六次に限られる．
Actual KE 構成の五次では，全 K2 end collars を保ち，source 近傍を anchor から
離して選ぶと，同じ Liouville core に局所作用適合 $K_4$ chart を構成できる．
両端の actual periods の整合は，residue weights の有限線形等式と source 配分の整数式に帰着する．
等重みでは対称な六つの四角形系と配分 $(3,1,1)$ がこれらを満たす．
五次では，初期 source 葉の全 tail を小面積に選び，source wall を含む全 moment の
面積指定と両端の作用接着を行うことで，この atlas と KE 計量実現を
同じ全域射影により同時に達成する．六次の同時実現は未証明である．
七次以上では，上記の actual KE 極限を実現する fibration の基底の各 core 頂点近傍も位相的に非 matroidal である．
\end{abstract}
\clearpage
\tableofcontents
\clearpage
\section{Introduction}

滑らかな射影曲面を Lagrangian torus の族として記述する問題では，
正則部分の action coordinates と，退化に現れる divisor の組合せ論を
大域的に結び付ける必要がある．
本稿の中心結果は，一般位置の全平面退化において，actual 負 Kähler--Einstein 形式に関する
大域 Lagrangian fibration と，その射影に標識付き基底収縮を合成した
同じ写像によって，径の尺度での計量極限を実現することである．
この構成を支える相対幾何として，正の重み付き SNC divisor を持つ射影曲面の有限相対 compactification を構成し，
それを一般型超曲面の退化へ適用する．
曲線成分には任意の genus を許し，divisor 上の marked Morse 葉層，
core 近傍の primitive，交点の product germs を同時に制御する．
この相対構成から，次数 $d\geq5$ の滑らかな超曲面
$X_d\subset\PP^3$ に対し，hyperplane class の正の倍数 $aH$ に属する
任意の Kähler form の下で，全正則 fibers が Lagrangian $T^2$ となる fibration を得る．
さらに，$d=e_1+\cdots+e_k$ という次数分割を選ぶことで，core の個数と型，
node 数，および底空間の homotopy type を変更できる．
五次・六次の全平面分割では，focus affine singularities を残した大域作用適合 matroidal atlas まで構成する．
また混合五次分割では，個別の core charts の存在と，任意の branch denominators を許した大域接続の不可能性を示す．
全平面五次で任意の奇数分母を許す局所 chart の整数構造と，
共有 branch に沿う大域的な算術制約は，companion 論文
\cite{QuinticArithmetic}で扱う．

固定した一般位置の全平面退化に対し，actual 負 Kähler--Einstein 形式で
同じ相対構成を計量的に制御する．$R_t=\log(-\log|t|^2)$ とすると，
標識付き基底収縮と fibration の合成 $f_t=p_T\pi_t$ が
\[
 \sup_{x,y\in X_t}
 \left|R_t^{-1}d_{g_t}(x,y)-d_\square(f_t(x),f_t(y))\right|\to0
\]
を満たすことを，任意の固定 $d\ge5$ で証明する
（定理\ref{ke:global-realization}，定理\ref{ke:global-metric-realization}）．
これは全曲面の ambient 距離についての主張であり，nodal，branch，core fibers と，
六次以上で基底収縮が潰す内部木上の全 fiber 合併を含む．
一般の symplectic 存在構成と，この特定退化族に対する KE 計量実現は，
仮定と証明を分けて記述する．関連する局所 KE 崩壊および
Lagrangian 構成との比較は第\ref{intro:ke-collapse}節で行う．

\paragraph{三つの存在主張の量化．}\label{intro:existence-scopes}
本稿で用いる選択の自由度は，次の三つで異なる．
\begin{enumerate}
\item \emph{指定 Kähler 形式での一般存在．}
滑らかな $X_d$ と $[\omega]=aH$ の Kähler 形式を先に固定し，
実際の punctured curves 上の所定の型の marked Morse 葉層の相対 isotopy class
（全平面分割では marked pants trees）を先に指定して
fibration を構成する（定理\ref{main:global}，定理\ref{part:global}）．
SNC 相対定理では，正の weighted divisor 表示と各成分の二つ以上の交点を仮定し，
同じ Kähler class の代表 $\omega'$ と normal holonomy を構成中に選ぶ
（定理\ref{snc:main}）．
\item \emph{Actual KE 計量による実現．}
次数と一般性条件を満たす退化族を先に固定し，一つの木系 $T$ と
全ての十分小さい複素 $t\ne0$ に対する射影 $\pi_t$ を選ぶ．
各 $\pi_t$ は連続であり，$f_t=p_T\pi_t$ の全域 distortion が零へ向かう．
この選択は $t$ 方向に連続又は滑らかな fibration 族を与えるとは限らず，
任意の先指定木系・critical values への適合も主張しない．
\item \emph{作用適合 matroidal atlas．}
五次・六次では，各指定 Kähler 形式に対して atlas を持つ fibration を選べる
（定理\ref{main:action-atlas}）．しかし，これと KE 計量実現を
\emph{同じ} $\pi_t$ で同時に達成することは未証明である
（注意\ref{main:atlas-ke-scope}）．七次以上では，実際に GH 極限を実現する
$\pi_t$ の基底の各 core 近傍が位相的に非 matroidal となる
（系\ref{ke:global-core-nonmatroidal}）．
\end{enumerate}

基底は，自由な辺を持たないコンパクトな分岐二次元複体である．
「境界がない」とはこの自由辺の不存在を意味し，分岐辺の内点には三枚の sheets が接する．
その fiber は $S^1\times(S^1\vee S^1)$ であり，
core 頂点の fiber は affine complement の logarithmic Weinstein skeleton である．
全 fiber の飽和と overlap 上の partition 一致から閉同値関係を作り，
大域商の Hausdorff 性・全射性・連結 fibers を得る
（補題\ref{glue:closed-whole-fiber-quotient}）．
Cubical 比較写像 $f_t=p_T\pi_t$ の全 fiber も連結だが，
六次以上では内部木上の torus の合併に三次元部分が現れ得る
（系\ref{topo:cubical-composite-connected}）．
従って $f_t$ 自体を Lagrangian torus fibration とは呼ばない．
全平面構成の fiber Euler 関数は node と core values だけに支持され，
各 node の寄与は $1$，各 core の寄与は $\binom{d-3}{2}$ である．
有限層化に対する $\chi_c$ の加法性から，
$\chi(X_d)=d\binom d2+d\binom{d-3}{2}$ と，
$p_T$ で同一視された全 fiber 合併の Euler 数を同時に得る
（命題\ref{inv:euler-pushforward}，系\ref{inv:cubical-fiber-euler}）．
この確認は数値から全空間を同定するものではなく，
実際に構成した全域写像の各 fiber を用いる．
正則な底空間には action--angle theorem による integral affine structure がある．
core と分岐辺を含めた matroidal charts の存在，その action lattice との整合，
および既知の algebraic filling との同定は，それぞれ追加の問題として扱う．

この設定は almost toric geometry の射程とも整合する．
Leung--Symington の閉 almost toric 四次元多様体の分類
\cite[Theorem 2.1, Table 1]{LeungSymington}は，通常の曲面を基底とし，
elliptic および focus--focus singularities を持つ fibrations を対象とする．
基底の境界がない場合には K3，Enriques 曲面，torus bundles が現れる．
adjunction により $K_{X_d}=(d-4)H$ を持つ一般型超曲面については，
本稿で許す branch と core の singularities が構成の本質を担う．
action--angle coordinates，nodal trade，almost toric geometry の基本事項は
Evans の講義録\cite{Evans}に従う．

より以前にSmith\cite{IvanSmithTorus}は，
\emph{homologically essential Lagrangian torus fibers}と
\emph{tame singularities}を仮定した閉symplectic四次元多様体について，
K3または一定のtorus bundlesへの位相的制限を証明している．
この定理も今回の一般型曲面に標準的なalmost toric
singularitiesだけを使うことの困難さを示すが，
\emph{本稿の分岐fiberとcore skeletonがSmithのtame条件を満たすか，
またregular torusのhomology classがessentialかを
追加検証せずに彼の定理を適用しない}．
我々のsingular fibersは，分岐辺に沿う
$S^1\times(S^1\vee S^1)$と，
nodal points上のpinched torusに加えて，
複数のstrataを持つcore Weinstein skeletonを含む．
従って通常のalmost toric fibrationと同じ局所型を
仮定したものではない．

\subsection{Hypersurface の pants 分解と一般型曲面}

代数的 hypersurface の大域的な torus fibration を考える上で，
Mikhalkin の pair-of-pants 分解\cite{Mikhalkin}は基本的な出発点である．
同論文の Theorems 1 および $1'$ は，tropical polyhedral complex への
stratified fibration と，高次元の pair of pants による分解を与える．
基底の homotopy type は幾何種数に等しい個数の球面の wedge である．
そこでの symplectic な結論は，各最大次元 cell に少なくとも一つ
Lagrangian fiber が存在するというものである．
全ての regular fibers の Lagrangian 性を得るためには，
接着に用いる form と local fibrations の整合をさらに示す必要がある．

一般型 hypersurface の Calabi--Yau 的な部分領域は，
Leung--Wan\cite{LeungWan}によって研究されている．
同論文は holomorphic volume forms を局在させる分解を構成し，
Mikhalkin の構成で既に Lagrangian である fibers について
asymptotically special となる性質を示す．
一般型における SYZ 的構造を調べる動機として，
この解析的な結果と，全 regular fibers を Lagrangian にする存在問題の両方が重要である．

Zhang\cite[Theorem 1.1]{Zhang}は，次数 $d$ の超曲面と微分同相な四次元多様体を，
$d(d-4)^2$ 個の四次元 pants と K3 に由来する領域などに分解している．
その pants 領域の境界には，torus を一般 fiber とする graph への fibration が構成される．
参照した preprint の Figure~1.1 と Theorem~1.1(ii) における特異 fiber は
$\Theta\times S^1$ であり，$\Theta$ は二つの三価頂点を持つ theta graph である．
これは本稿の一つの四価頂点を持つ $(S^1\vee S^1)\times S^1$ と同相ではない．
同論文の Introduction には，
五次曲面の五平面退化から五個の pants を得る記述もある．
本稿の五次の場合は，この既存の幾何学的構図を踏まえている．
一方，torus に対する積分 $\int_T\omega=0$ と，
各 tangent space 上での条件 $\omega|_T=0$ は異なるため，
本稿では後者を各 piece とその overlaps で直接確認する．

Golla--Martelli\cite{GollaMartelli}は，四次元多様体の pants 分解を
微分位相的な構成法として展開し，corner を含む貼合せと mapping-class data を記述している．
本稿でも，分岐底空間の cells と boundary faces の構造を理解する際にこの記述を用いる．
symplectic な構成では，これに fixed collars と実際の action periods の整合が加わる．
境界を固定した core の置換を検討する際には，corner に追加の torus 領域を置く操作が
基本群と boundary topology を変え得ることも考慮する必要がある．

\subsection{SNC divisor の有限相対 compactification}

Evans--Mauri\cite{EM}は，contact link 上の Lagrangian fibration を
Liouville flow によって completion 上へ延長する方法を与えた．
一般位置の四直線を $\PP^2$ から除いた四次元 pants に対しては，
$K_4$ の cone を基底とし，torus，$S^1\times(S^1\vee S^1)$，
core skeleton を fibers とする模型を構成している
\cite[Lemmas 2.16, 2.22; \S3.1]{EM}．
同論文は一般型曲面への応用も明示的に動機として挙げている．
また Remark 2.23 は，曲線上の Lagrangian leaves が一次元であるため，
circle-bundle lift に必要な curvature の vanishing condition が自動的であることを指摘する．
この点に genus 零の制約はなく，本稿の positive-genus 構成も同じ lifting principle を用いる．

本稿の相対問題では，completion の内部への cone extension に加え，
所定の divisor における有限時刻の circle collapse を制御する．
Lemma 2.16 の completion は contact form の正の関数倍を許すため，
固定した finite collar，normal moment，divisor 上の foliation，
core 近傍の primitive を同時に指定するには追加の構成が必要である．
同じ contact 構造を保つ正の rescaling でも既存の葉の Lagrangian 性は一般に変わる
（注意\ref{ke:contact-rescaling-collar-scope}）．
一方，\cite[\S\S2.3--2.4]{EM}の Hermitian metric の変更は選んだ Kähler class を，
exact symplectic deformation は de Rham class を保つので，
固定 cohomology class 自体を本稿固有の特徴とはしない．
ここでの相対条件は，Kähler representative と指定 product germs，
marked full leaves，選択する holonomy と periods，有限 endpoint の同時制御である．
SNC symplectic neighborhoods とその relative regularizations には
Tehrani--McLean--Zinger\cite[Definitions 2.11--2.12, Theorem 2.13]{TMZ}の枠組みを用いる．
同定理は一般には symplectic form の変形を許すため，本稿の固定 $\omega'$ に対する
charts は，product crossing germs を準備した後の相対 tubular Moser で構成する．
正の weighted divisor に適合する primitive の存在は
Borman--Sheridan--Varolgunes\cite[Propositions 2.7, 2.10]{BSV}に関係する．
同結果の重みは正の実数でよく，$H^1(M)=0$ を要求しないが，先指定した normal connection，
core primitive germ，full Morse leaves の同時保存までは結論していない．
同論文の period-one 規約とは $r_a=2\pi\mu_a$，$\kappa_a=2\pi w_a$ で対応する．
normal directions の対称化は Pascaleff\cite[\S4.2]{Pascaleff}の構成と関係する．
Kähler class の実数係数の表示から global logarithmic potential を作る段階では，
compact Kähler manifold の $\dd\dd^c$-lemma\cite[\S5]{DGMS}を用いる．

具体的には，$M$ を射影曲面，$D=\bigcup_aD_a$ を三重交点のない SNC divisor とし，
各 $D_a$ が少なくとも二つの intersection points を持つとする．
$D_a$ の genus，self-intersection，同じ二成分の複数交点には制限を加えず，
\[
 [\omega]=2\pi\sum_a w_a[D_a],\qquad w_a>0
\]
を仮定する．
定理\ref{snc:main}は，同じ class の Kähler representative $\omega'$ に対し，
$D$ を outer graph の逆像とする closed-cone fibration を与える．
各 punctured component 上には，あらかじめ指定した無 extrema の
marked Morse foliation を用いる．
full saddle leaves と finite endpoint を保ち，
交点の product germs および有限個の divisor-compatible flat patches も指定できる．
個々の divisor components の ampleness や $H^1(M;\RR)=0$ は仮定しない．

positive genus における period の制御には，connection holonomy の自由度を用いる．
固定 residues を持つ二つの collar primitives の差は，
plumbing の homology exact sequence により $H^1(D;\RR)$ に由来する．
これは dual graph の periods のほかに，各 $D_a$ の $2g_a$ 個の periods を含む．
crossing disks で零の closed one-form $\eta_a$ により normal connection を変え，
対応する coupling forms を Moser map $\Phi_s$ で固定すると，
\[
 [\Phi_s^*\lambda_{a,s}-\lambda_{a,0}]
       =-s w_a[\eta_a]
\]
となる（補題\ref{snc:holonomy}）．
これで component periods を消し，残る dual-graph periods を basic forms で補正する．
この Moser 操作は auxiliary normal-bundle charts を変えるものであり，
actual collar 上の Kähler form $\omega'$ は固定される．
最後に bounded primitive の cutoff を用いて，
core 近傍の primitive germ を保ったまま end の model と接続する．
さらに connection と coupling charts を固定し，closed basic correction と exact cutoff
に限る場合の必要十分条件を命題\ref{snc:fixed-holonomy}で与える．
Reeb 商 $q:D\to\Gamma$ に対し，障害は
\[
 H^1(D;\RR)/q^*H^1(\Gamma;\RR)
 \cong\bigoplus_a\operatorname{Hom}(L_a,\RR),\qquad
 \operatorname{rank}L_a=g_a
\]
であり，各成分から $g_a$ 本ずつ選んだ正則葉の period の消滅で判定できる．
曲面一枚での leaf-period と basic representative の関係は
Izosimov--Khesin--Mousavi\cite[Lemma 4.8, \S5]{IKM}の既知結果である．
本稿では整数格子の証明と SNC 全体の exact sequence を記述し，有限時刻の全葉 collapse と
core に相対的な接着へ適用する．任意の障害の消去には $\sum_a g_a$ 個の独立な holonomy
調整方向が必要十分である（系\ref{snc:minimal-holonomy-controls}）．
非零障害を消す connection 変更は一般に全 full saddle collars を支持外にできない
（注意\ref{snc:saddle-support-obstruction}）．これらは指定した補正方法の条件であり，
任意の Lagrangian fibration の存在を分類するものではない．

固定した元の form が交点で nonorthogonal な場合には，
commuting normal-circle actions の product model をそのまま要求できない．
これは\cite[Proposition 2.5]{BSV}の orthogonality 条件であり，
本稿では Kähler class 内の明示的な potential deformation によって product germs を作る．
さらに global residues と core の primitive germ の関係を記述し，
collar の area balance だけから global exactness を推論しない．
$\PP^1\times\PP^1$ の toric boundary，
二点で交わる二本の $(1,1)$ curves，
九点で交わる二本の smooth plane cubics は，
nonample components，multiple intersections，positive genus をそれぞれ検査する例になる．

\subsection{実際の退化族上での接着}

Fernández de Bobadilla--Pełka\cite[Proposition 5.6, Theorem 6.14]{FBP}は，
Kähler SNC 退化の近傍で ivy-like な Lagrangian torus fibrations を構成し，
一般型を含む広い範囲の退化を扱う．
曲線 stratum 上の Morse saddle から branch が生じる点は本稿にも共通する．
同論文が大域的に構成するのは，maximal および submaximal strata に対応する，
nearby fiber の境界・角を持つ余次元零の部分領域である．
component の内部に対応する残余領域を埋めることは追加の問題であり，
同論文の Example 5.8 もこの差を説明している．
本稿では，この残余領域に core を含む finite compactification を入れ，
double-curve necks および triple-point models と有限の overlaps 上で接着する．
同論文の Theorem~6.14(b) の fibration の定義域はこの部分領域だが，
同定理(c) の GH 収束は構成した Kähler 計量を入れた $W_z$ 全体を扱い，
退化写像が proper なら $W=X$ と取れる．従って同論文の GH 結論まで
部分領域だけのものとは解釈しない．その計量と本稿の actual 負 KE 計量の区別は
第\ref{intro:ke-collapse}節で述べる．

central fiber から symplectic parallel transport で fibration を移す考え方には，
Ruan の Calabi--Yau hypersurface の構成\cite{Ruan}がある．
Ruan の五次は $\PP^4$ 内の Calabi--Yau 三次元多様体であり，
本稿の $\PP^3$ 内の一般型五次曲面とは異なる次元の問題を扱う．
また Mak--Matessi--Ruddat--Zharkov の共同研究について，
Ruddat の二〇二六年六月の講演概要\cite{RuddatAnnouncement2026}では，
reflexive polytope に対応する toric Fano manifold 内の compact Calabi--Yau hypersurface 上で，
余次元二の discriminant を持つ Lagrangian torus fibration の構成が発表されている．
概要は Evans--Mauri の方法に従う Liouville flow と帰納的な mixed local models を挙げる．
これは関連する大域的構成の動向として引用し，本稿での比較は公表概要の射程にとどめる．
本稿の一般型曲面の SNC cores と整数分母の問題には，指定した局所模型の接着条件が加わる．
さらに Zharkov の二〇二六年四月の講演概要\cite{ZharkovSoftCells2026}は，
phase-tropical および ober-tropical な超曲面の比較から，
permutahedral skeleton を soft cells に置き換える局所模型を述べている．
元の超曲面との symplectic な比較と Lagrangian skeleton fibers が発表の対象であり，
pants の局所構成を比較する上で関連する．
これも公開された講演概要として引用し，指定した SNC data に相対的な接着定理を
その概要だけから導いたり，完成論文の証明を確認したものとして扱ったりしない．

四次曲面における transport と moment 保存の記述は
Evans\cite[Chapter 6]{EvansNotes}にも詳しい．
Chakravarthy--Groman\cite[Theorem 1.6]{CGK3}は，
symplectic Kulikov model の component 上の almost toric fibrations を
nearby K3 へ接着する大域定理を証明している．
その component hypotheses には anticanonical cycles が含まれるため，
多数の punctures と saddles を持つ本稿の曲線 strata には別の collar 構成が必要になる．
同論文の改訂版 Lemma 2.13 が扱うのは，必要な cylinders を潰した
central fiber の商への transport の連続性である．
本稿では smooth strata 内の compact saturated subsets を運び，
double curves と triple points の中心部分を正確な $xy=t$，$uvw=t$ の models で覆う．

位相型と linear monodromy だけでは symplectic gluing は決まらない．
曲面上の Morse systems には weighted または measured Reeb graph の情報があり
\cite{DMT,IKM}，compact focus--focus fiber の近傍には
semi-global invariant がある\cite[Theorem 2.1]{VuNgoc}．
一重 focus の十分小さい飽和近傍では $H_1\cong\ZZ$ であり，
正則 torus の normal vanishing cycle は消え，残余 cycle が generator となる．
従って既に存在する全近傍の symplectic map が指定 primitive に関して exact であることは，
その一つの残余 period の一致と同値である（補題\ref{ke:focus-neighborhood-period}）．
Node 上の primitive germ だけを固定してもこの period は変えられるが，
一つの外側正則 fibre 上の primitive まで固定すれば障害は零となる
（例\ref{ke:focus-local-primitive-shift}）．これは constant-zero Taylor 不変量と，
\cite[\S7]{VuNgoc}の exact 版に現れる作用の定数項との区別を具体化する．
そこで，本稿は actual family 上の同じ Hamiltonian circle action と
normal moment $\mu$ を用い，reduced surface 上の functions $h_\mu$ を
marked foliation に相対的に補間する．
$\{\mu,h_\mu\}=0$ によって Lagrangian 性を保ち，
既に構成した nodal fiber はその compact fiber 全体と近傍を transport する．
その際に葉の分割と Eliasson 座標・cycles を同時に運べば Taylor 不変量が保存され，
exact 移送では作用の定数項も保存される（命題\ref{ke:focus-whole-neighborhood-transport}）．
この方法は測度や period の不変量を無視するものではなく，
同じ実際の form 上で leaves と action differentials が一致するように作るものである．
正則 torus band の既定 base 同定については，normal moment の一致に加え，
第二作用が整数 shear と定数 translation を除いて一致することが必要十分である
（命題\ref{neck:regular-period-matching}）．固定 primitive に関する exact 同定では
その translation も零を要求する．これは action--angle theorem の局所的な帰結であり，
独立な第二作用の自動一致や，任意の先指定 collar map の大域延長を主張しない．
circle-invariant nodal fibrations の reduction による記述には
Groman--Varolgunes\cite[\S7.3]{GV}を用いる．

component 上の局所点 blow-up には，
Chakravarthy--Groman\cite[Theorem 1.1]{CGblowup}の almost toric comparison と，
Fleming\cite[Proposition 4.6, Construction 4.7]{Fleming}の Kähler blow-up を用いる．
これらは regular elliptic-boundary patches 内で適用し，
core と crossing germs を固定する．
容量は例外曲線の実面積 $e$ で測り，実際の複素 blow-up 上の radial K\"ahler form と
半径 $\sqrt{e/\pi}$ の symplectic blow-up の相対同定を直接構成する
（補題\ref{glue:area-normalized-blowup}）．先に飽和した affine-convex patch を取り，
その内部へ球と容量 triangle が入るよう $e$ を小さくする．比較写像の恒等域はその
飽和 patch の外であり，球の外全体ではない．
その後の divisor-relative Moser と normal-circle conjugacy は，
曲線の genus に依存しない．
positive genus の場合も，同じ実際の marked curve 上の isotopies を用い，
normal-bundle の degree の一致に加えて line-bundle の smoothing relation を保持する．

\subsection{次数分割，偏極，底空間の整数位相}

まず平面退化の場合，定理\ref{main:global}は $d$ 個の cores，
$d\binom d2$ 個の nodal fibers，
$\binom d2(d-4)$ 本の branch edges を持つ fibration を構成する．
$d=5$ の底空間は四単体の二次骨格である．
一般の $d$ では，同じ実際の double curves 上の marked pants trees から
有限複体 $B_d(T)$ を作り，
\[
 B_d(T)\simeq\bigvee^{\binom{d-1}{3}}S^2
\]
への整数鎖ホモトピーと $H_2$ の具体的な基底を与える．
球面の wedge という型は Mikhalkin の結果と整合するが，
ここでは指定された branch edges，bigons，triangles に対する鎖データを記述する．

任意の順序付き分割 $d=e_1+\cdots+e_k$，$k\geq3$ に対しても，
一般の smooth hypersurfaces $S_i=\{F_i=0\}$ を用い，
$\prod_iF_i+tG_d=0$ を順序付き小解消することができる．
定理\ref{part:global}は，この退化から $k$ 個の cores を持つ fibration を与える．
第 $i$ core は
$S_i\setminus\bigcup_{j\ne i}(S_i\cap S_j)$ の logarithmic skeleton に由来し，
node 総数は $d\sum_{i<j}e_i e_j$，
branch edge 数は $(d-4)\sum_{i<j}e_i e_j$ である．
定理\ref{part:homotopy}はその底空間について
\[
 B_{\mathbf e}(T)\simeq
 \bigvee^{\,p_g(X_d)-\sum_i p_g(S_i)}S^2
\]
を示し，Reeb graphs の cycles，同じ triple-point type の複数交点，
および tetrahedron boundaries から整数基底を構成する．
例えば次数六の分割 $(2,2,2)$ と $(1,1,4)$ は，
同じ symplectic 曲面上に $b_2(B)=10$ と $9$ の底空間を与える．
従って，core の選択は底空間の homotopy type にも反映される．

偏極の範囲はこの自由度とは別に正確に記述できる．
小解消の vertical class を
$\xi=aH+\sum_i\beta_i[D_i]$ とし，
$\delta_\ell=\beta_{\ell+1}-\beta_\ell$，
$E_\ell=e_1+\cdots+e_\ell$ と置く．
定理\ref{part:cone}は，その Kähler slice が
\[
 \delta_\ell>0\quad(1\leq\ell<k),\qquad
 \sum_{\ell=1}^{k-1}E_\ell\delta_\ell<a
\]
という sharp な開単体であることを示す．
平面分割では $E_\ell=\ell$ となり，定理\ref{pol:cone}を回収する．
この slice 全域の記述には Bézout と Nakai--Moishezon を用いる．
大域 fibration の局所 blow-up construction に使うのは，
一様な patches を確保した後の十分小さい ray である．

\subsection{core の topology と matroidal charts}

core の homotopy type と embedded Liouville skeleton の同定は別の問題である．
Hattori\cite{Hattori}は一般位置の hyperplane complement の homotopy type を決定し，
関連する arrangement の記述は Papadima--Suciu\cite[\S1.1]{PapadimaSuciu}にもある．
Ruddat--Sibilla--Treumann--Zaslow\cite{RSTZ}の combinatorial skeleton は，
Zhou\cite{Zhou}によって，適切な exact symplectic structure の
Liouville skeleton として実現されている．
本稿では，指定された compactification の collar を保つことが必要なため，
実際の logarithmic potential から得られる skeleton を用いる．
特定の RSTZ model との相対同定は追加条件として残る．
Matessi\cite{MatessiPants,Matessi}の Lagrangian pants は半次元の部分多様体を作る構成であり，
symplectic pants 全体の torus fibration とは区別される．
Abouzaid--Auroux--Katzarkov\cite{AAK}の対象も，
hypersurface に付随する toric blow-up や conic bundle であり，
compact hypersurface 自身への本稿の gluing theorem は別途必要である．

一般型曲面内の部分的な almost toric geometry にも先行研究がある．
Evans--Urz\'ua\cite{EvansUrzua}は，symplectic form を変形した一般型曲面で，
rational homology balls の非有界な列の symplectic embeddings を構成した．
Brendel--Hauber--Schmitz\cite{BHS}は，$B_{dpq}$ の skeleton 近傍からの
partial embeddings によって移される小さい Lagrangian tori を
ambient symplectomorphism に関して区別し，一般型の五次曲面や
simply connected Godeaux surface にも，適切な symplectic form の下で応用している．
ここでは displacement energy germ が重要な不変量である．
本稿が指定する K\"ahler class，compact な底の全体，core と neck の同時接着は
別途確認する条件であり，局所 torus の非同値性から fibration 全体の非同値性を推論しない．
この区別は，構成後の torus fibers や atlas の同値性を研究する際にも維持する．

torus fibers の Floer 不変量を比較する方向では，Venugopalan--Woodward
\cite[Definitions~1.1--1.2, Theorem~1.3]{VenugopalanWoodward2026}が，
compact almost toric 四次元多様体の torus fiber に対する Maslov index 二の disk counts を
tropical graphs で表す公式を与えている．同定理は elementary polyhedral decomposition と，
potential に寄与する全 disks の tropical graphs の collisions が generic interior points にある
という仮定を伴う．本稿の分岐底複体と core fibers はこの almost toric 設定に含まれず，
公式を本稿の fibers に用いるには別の比較定理が必要である．
また非単調な場合の potential は Maurer--Cartan space 上の関数を同型まで考えるものであり，
選択に依存しない一つの Laurent polynomial とみなして atlas の非同値性を判定しない．

matroidal charts との比較には，Bergman complex の位相を使う．
Ardila--Klivans\cite{ArdilaKlivans}は lattice of flats による subdivision と
phylogenetic-tree spaces との関係を示し，
Feichtner--Sturmfels\cite[Theorem 4.1]{FeichtnerSturmfels}は
nested-set complexes による unimodular triangulation を与えた．
Petersen graph と $K_{3,3}$ は後者の Examples 1.1--1.2 に既に現れる．
Shaw\cite[\S2.4]{Shaw}の matroidal fans は，
本稿で比較する smooth tropical local models の範疇を与える．
これらの既知の模型を含め，定理\ref{trop:cubic-classification}では，
二価点を除いたリンクが三価である rank-three matroidal fans を分類する．
連結 matroid から現れるのは
$K_4,K_{3,3},\mathrm{Petersen},\mathrm{Heawood}$ であり，
非連結 matroids と Euclidean factors を含めてもリンクの第一 Betti 数は八以下である．
体上での realizability はこの上界の仮定ではない．

曲線配置の tropicalization を matroid の外側で扱う比較対象として，
Aksnes\cite[Theorem 4.15]{Aksnes}は，平面内の transverse very affine arrangement の
補集合を arroid fan によって記述する．
本稿の Reeb 展開を施した contact link と，divisorial tropicalization の support が
同相であるとは限らない．また，この結果から actual action periods の一致は従わない．
従って，matroidal charts の障害の後に比較する，より広い fan の範疇として区別する．

一方，平面退化の各 core link は
$b_1=\binom{d-2}{2}$ を持つ．
従って $d\geq7$ では，その core の局所位相だけで matroidal chart を排除できる
（定理\ref{trop:degree-classification}）．
次数分割の場合も
\[
 b_1(\Gamma_i)=1+\frac{e_i(d-e_i)(d-4)}2
\]
であり，同じ障害が全ての partitions に及ぶ
（命題\ref{trop:partition-obstruction}）．
六次では，全ての core が matroidal な局所位相を持つ候補は全平面退化に限られる．
その全 core links が Petersen となる labelled split systems は正確に六個で，
$K_6$ の one-factorizations と対応する（定理\ref{topo:petersen}）．

五次では Betti 数だけの判定をさらに精密化する。
divisor の各 punctured component の内部 Reeb graph は、core link の誘導部分グラフになる。
この block と個々の triple-point labels を保持すると、$(1,1,3)$ と $(1,2,2)$ の
全 core の同時 matroidal 化は不可能である一方、$(1,1,1,2)$ では四つの標識付き
abstract graph systems が存在する（定理\ref{q5:classification}）。
その条件は三つの quartet splits の parity で記述できる。
Reeb graph の実現には good orientation が必要であり、
Michalak\cite[Definition 2.2, Theorem 5.6]{MichalakReeb}の実現問題との関係を明示した上で、
本稿で使う境界付き曲面には saddle pants と cylinders による直接の構成を与える。
従って有限グラフの個数を、embedded marked isotopy classes の個数とは同一視しない。


\subsection{局所作用座標から大域 atlas へ}

Petersen link が得られても，標準 fan の格子と実際の作用格子を合わせる必要がある．
critical slopes の rationality，annular area の Farey 条件，branch lattice の saturation は，
そのための局所的な条件である（命題~\ref{trop:farey-condition}，\ref{trop:branch-denominator}）．
固定 annular area は $\operatorname{area}/(2\pi w)=1/m$，$m$ は正の奇整数，という
離散条件を持つ（命題~\ref{per:odd-reciprocal}）．
一方，marked isotopy class 内で foliation の代表を選べるなら，五直線補集合で各 critical
basin の面積を $2\pi w_b$ にし，内部 annulus の面積を $2\pi w_a$ にできる．
これにより十五面全ての作用格子が標準 Petersen fan の格子と一致する
（定理~\ref{per:local-realization}）．
これは Evans--Mauri の Liouville cone construction\cite{EM}と，標準 tropical lattice
\cite{GKM}を，実際の作用積分によって接続する構成である．
四直線の $K_4$ core については，有向四周期に沿う basin areas を用い，六面の作用座標と
$H_1(W_4;\ZZ)\cong N_4^*$ の整数同型を与える（命題~\ref{per:k4-realization}）．
Petersen の全 meridians に一つの ambient-linear extension を要求すると半整数の障害があるが，
この追加条件は一つの core の作用適合 chart の存在とは異なる．

大域 atlas では，さらに chart transitions を検査する必要がある．
Shaw\cite[Definition~3.1]{ShawSurfaces}，Jell--Rau--Shaw
\cite[Definitions~2.1, 3.1, 5.1]{JellRauShaw}，L\'opez de Medrano--Rinc\'on--Shaw
\cite[Definitions~2.2--2.4]{LRSChern}の枠組みでは，transitions は ambient integer affine maps
の制限である．face ごとの unimodularity と branch saturation だけでは足りない．
本稿は三価 star の唯一の線形関係を用いて延長条件を記述し，両 chart で同じ Hamiltonian
circle の作用が ambient-affine ならその条件が成立することを示す
（補題~\ref{atl:branch-extension}，命題~\ref{atl:common-hamiltonian}）．
focus--focus points では作用格子の monodromy が非自明である\cite[Remark~3.1]{VuNgoc}ため，
atlas はその有限集合を除いた底空間に与え，除いた点には実際の focus affine singularities を残す．

Evans--Mauri の掲載版 Propositions~6.11--6.12 と Theorem~6.14 は，
Assumption~1 の下で toric neighborhoods と双対 monodromy を比較する\cite{EM}．
この仮定は各一次元 stratum が $\PP^1$ で零次元 stratum と正確に二点で交わることを要求し，
各直線に $m-1\ge3$ 個の交点を持つ本稿の一般位置 $m$ 直線配置には直接適用できない．
本稿の三穴・四穴 sphere とその saddle branches には，critical actions を直線に保つ別の接着が必要である．
実際，作用適合 matroidal chart が存在すれば，各 incident sheet の critical action は
normal moment の有理勾配の affine function でなければならない
（命題\ref{atl:critical-action-obstruction}）．
固定した exact symplectic neck 上で，normal action，両端の open germs，toric collars，
saddle 点の局所 germ，marked tree を全て保ちながら，critical leaf の regular arc だけを動かして
この作用を非直線にできる（例\ref{atl:fixed-form-bent-branch}）．
その点には任意の分母を許しても作用適合 matroidal chart はなく，
Lagrangian 性と atlas の存在が別の条件であることを示す．
このため，actual neck の全空間上で一つの primitive を取り，三重点の toric reference actions
と両側の作用の加法定数を同時に固定する．
固定した正負の normal moments における symplectic transport はこの primitive を exact 差まで保つ．
境界作用に対する Stokes の公式から reduced total area の affine 性を得て，
critical basins と内部 annuli の面積を指定する parameter に相対的な Moser construction を行う．
これにより有限 $t\ne0$ でも critical action functions を直線に保てる
（系~\ref{ng8:central-criterion}）．
measured Reeb graphs\cite{DMT,IKM}との関係では，既に固定した二つの systems を
無条件に同一視するのではなく，同じ actual reduced forms 上で必要な critical-area data を実現している．

この接着を用いると，五次と六次の全平面退化で大域 atlas を得る
（定理~\ref{ga8:global}）．$b_i=\beta_i/(2\pi)$ とし，
\[
 c=\frac{a}{2\pi(d-1)},\qquad w_{i,j}=c+b_j-b_i
\]
という weights を選ぶ．五次では $f_i(j)=(i+j)/2\bmod5$ の対称な四周期系を用い，
各共有曲線の五つの blow-up points を三つの basins へ $(3,1,1)$ と配分する．
六次では各曲線の四 basins へ一つずつ，内部 annulus へ二つを配分する．
面積変化と normal framing の変化が同じ整数条件を与え，両端の affine actions が一致する．
五次の各 triangle には五つの nodes，六次の各 triangle と bigon にはそれぞれ三つと二つの nodes が入る．
全 branch lengths は $a/(\pi(d-1))$ となる．
五次の対称な標識付き recipes と六次の Petersen systems はそれぞれ六個あるが，
この数を非同型な大域 atlases の個数とは同一視しない．

\subsection{混合分割での局所存在と大域障害}

分割 $(1,1,1,2)$ では，quadric の三本の $(1,1)$ curves の補集合に
$K_{3,3}$ chart を構成する（定理~\ref{quad:realization}）．
標準 $K_{3,3}$ fan が二つの tropical lines の積であることは，parallel connection に対する
Shaw--Werner\cite[Theorem~7.2]{ShawWerner}の記述に従う．
ここで必要な追加の計算は，自己交点二の actual normal framing と非零 critical periods である．
各 component の一つの basin をその weight 分だけ広げると，内部 annulus の面積比が一となり，
九面の作用行列は全て unimodular となる．
二直線と conic を境界に持つ三つの plane cores にも，同様に非零 periods を用いた
$K_4$ charts がある（命題~\ref{mixper:realization}）．
従って，四つの標識付き quartet systems の全てで，各 core を個別に実現できる．

分岐における作用の不整合は，tropical divisor の言葉で表せる．
単位 facet weights を持つ三価 star の primitive transverse lifts を
$u_1,u_2,u_3$，分岐方向を $v$ とし，
$u_1+u_2+u_3=\kappa v$ と書くと，sheet ごとに整数アフィンな作用関数 $h$ の係数は
\[
 \mathfrak d(h)=\sum_{a=1}^3\dd h_a(u_a)-\kappa\dd h(v)
\]
となる．これは Allermann--Rau の tropical principal divisor の係数そのものである
\cite[Definition~3.4, Remark~3.5]{AllermannRau}．
同文献の abstract cycles に対する構成
\cite[Definition~6.1, Construction~6.4]{AllermannRau}と整合するが，
本稿では各 divisor component に固有の normal action を用いるため，
全 core の係数を初めから一つの大域 principal divisor と解釈しない．
共有 neck の normal circles が互いに逆向きであることから，
対応する二端の係数が相殺するという必要条件を直接用いる．

分母の自由度を除くのは，局所 fan の分類だけではなく，
共有する四 puncture 曲線の annulus と nodal framing の整合性である．
一つの focus node が与える integral shear は almost toric geometry の標準的な事実である
\cite[\S4.3, equation~(4.10)]{SymingtonFour}．
本稿では，両側の annulus の unimodularity と四 punctures の整数 cycle shears を同時に用い，
$q_Aq_B\mid2$ を導く（命題~\ref{q9:rigidity}）．SNC の node budget を合わせると，
二端の分母 $q_A,q_B$ に対して内部 annulus の node 数が
\[
 n_{\rm int}=\frac{2}{q_Aq_B}
\]
に等しいことを示す．三価 branch の saturation により分母は奇数であるから，
node 数の整数性は $q_A=q_B=1$ および $n_{\rm int}=2$ を強制する．
この量子化条件を混合五次の三本の conic necks に適用すると，
高分母を許した場合にも core coefficients の総和は $-6$ となる．
従って現在の SNC cores と branch graph を保つ範囲では，
混合五次の大域作用適合 matroidal atlas は分母の選択によらず存在しない．
これは regular curves と annuli の周囲の議論であり，
focus values 自体に matroidal chart を延長する主張ではない．

三本の plane--plane branches の分母を任意のまま残しても，line components の寄与は相殺される．
各 mixed plane core の係数総和は $-1$，quadric core は $-3$ なので，全体は $-6$ となり，
対向端の相殺条件に反する（定理~\ref{hd9:mixed-no-go}）．
これと位相障害を合わせると，現在の比例曲率 SNC 構成で cores と branch graph を保つ
大域作用適合 matroidal atlas が可能な次数分割は，分母を制限せずに，全平面の五次と六次に限られる
（系~\ref{hd9:classification}）．

同じ四 puncture の剛性は，六次の全三十 branches の分母が一であり，
十五 bigons の各々に二 nodes が入ることも強制する．これは一つの weight recipe に依存しない．
一方，四直線の局所 $K_4$ core では，正の basin areas と六つの unimodular action matrices を
同時に指定することで，任意の奇数 $q\ge3$ に対して分母 $(q,1,1,1)$ を持つ chart を得る
（\cite{QuinticArithmetic}）．この例は，局所の奇数分母条件と，実際の共有 neck に沿う
大域的な剛性との違いを具体的に示す．

\subsection{負 Kähler--Einstein 退化との計量的な比較}\label{intro:ke-collapse}

定理\ref{ke:global-realization}は，Lagrangian fibration の射影に標識付き基底収縮を合成した写像によって，
負 KE 計量の径の尺度での Gromov--Hausdorff 極限を実現する．
Actual KE 形式に関する fibration $\pi_t:X_t\to B_d(T)$ と
標識付き収縮 $p_T$ の合成について，全曲面の ambient 距離に対する
一様な距離歪みが消えることを証明する．
正規化は $\operatorname{Ric}(\omega_t)=-\omega_t$，
$\bar g_t=R_t^{-2}g_t$，$R_t=\log(-\log|t|^2)$ である．
以下では既存結果の範囲，本稿が明記して用いる局所入力，
および全域化のために証明する評価を区別する．

本稿の fibration の存在と，負 Kähler--Einstein 計量の崩壊は異なる主張である．
全平面退化 $X_t=\{\prod_i L_i+tF_d=0\}$ に対し，
$\operatorname{Ric}(\omega_t)=-\omega_t$ と規格化すると
$[\omega_t]=2\pi(d-4)H$ なので，本稿の存在定理の偏極範囲に入る．
しかし，各 $t$ で $\omega_t$-Lagrangian な fibration が存在することは，
その fibers の KE 直径が小さいことや，退化 parameter に関する計量的一様性を含意しない．

Song\cite[Theorem~1.2]{SongNegativeKE}の stable degeneration の定理を
この族に適用すると，成分内部に基点を選んだ非崩壊極限は
$\PP^2$ から $d-1$ 本の直線を除いた $d$ 個の完備 KE 曲面となる．
従って元の尺度では，各成分の内部領域には四実次元幾何が残る．
固定成分の空間的 blow-down は companion 原稿の Part I
\cite[\texttt{cor:arrangement-cone}]{MultiscaleKECurrent}の対象である．
そこでの $C_{2,d-1}$ は，本稿の compact な極限 $\mathcal C_d$ の成分頂点における
metric tangent cone と labels を含めて一致する
（命題\ref{topo:component-cone-germ}）．さらに本稿の同じ $f_t$ の distortion を使うと，
元の尺度と径尺度の間に適切な $1\ll a_t\ll R_t$ を選び，
成分頂点上の基点に関する $(X_t,a_t^{-1}d_{g_t})$ の pointed GH 極限も
この錐となる（系\ref{ke:intermediate-component-cone}）．さらに同系では，
基点の target 値の成分頂点からのずれも尺度の選択に含めることで，
component core と変更 collar の全領域内の任意の基点列へ結論を拡張する．
さらに命題\ref{topo:cubical-local-models}と\ref{topo:cubical-moving-blowups}は，
cubical target の全局所模型と移動基点の拡大極限を分類する．
$\epsilon_t=\operatorname{dis}(f_t)$ とし，$1\ll a_t\ll R_t$，
$R_t\epsilon_t=o(a_t)$ を満たす尺度では，actual 曲面の全ての移動基点の
部分列極限も，有限基点を持つ $C_{2,d-1}$，$\mathbb R\times\mathcal Y_{d-2}$，
全角 $3\pi/2$ の錐，および $\mathbb R^2$ に尽くされ，各型を実現できる
（系\ref{ke:intermediate-moving-classification}）．
ここで $\mathcal Y_{d-2}$ は原点で貼った $d-2$ 本の半直線である．
これは全域歪みの閉球への移送と有限複体の幾何による帰結であり，
上の誤差条件を除いた任意の中間尺度や極限の交換，tensor 収束は主張しない．
固定補空間の Part I--III，退化族の局所一様解析と条件付き大域枠組の Part IV--V，
本稿の actual fibration による直接実現は，この対象と量化を区別して接続する．
Part II の標識付き残余複素曲線の period と，本稿の第二作用 period も別のデータである．

他方，Zhang\cite[Theorems~1.1, 2.4]{ZhangNegativeKE}は，
全空間が $\mathbb Q$-factorial な simple toroidal canonically polarized degenerationで，
深い stratumへ基点を動かした局所 torus collapse を証明した．
曲面の二重曲線では一つの円，三重点では二次元トーラスが潰れる．
これは compact な全基底への GH 収束とは異なる．
本稿の族には $xy=tz$ 型の非 $\mathbb Q$-factorial な全空間特異点があり，
小解消後も crepant exceptional curve 上で相対 canonical bundle の次数が零となるため，
同定理の相対 ample 仮定をそのまま使えない．
Mandel\cite{MandelNegativeKE}の負 KE neck の構成も，
我々の compact family との全域的な計量比較を自動的には与えない．

Fu--Hein--Jiang\cite[\S1.1, Main Theorem]{FuHeinJiangCusp}は，
孤立楕円特異点へ退化する六次曲面の actual 負 KE 計量を接着により解析し，
形成される complex hyperbolic cusp の先端に Tian--Yau gravitational instanton が
bubble として現れることを証明している．これは負 KE 退化の実際の計量を扱う
先行結果であり，全平面 SNC 退化の全曲面を同じ Lagrangian 射影で比較する
本稿の主張とは，退化模型と極限の尺度を区別して比較する．

Fernández de Bobadilla--Pełka\cite[Theorem 6.14(b),(c)]{FBP}では，
部分領域上の ivy-like fibration と，proper な場合の全 nearby fiber の GH 収束が得られる．
その計量は，選んだ Kähler 形式への exact な補正から作る $\omega_q^\varepsilon$ による．
Proposition~6.13 と Theorem~6.14(c) はこの計量を再尺度化して双対複体へ収束させる．
本稿との比較では，領域の全域性だけでなく，二形式・尺度・射影を同時に区別する．
Calabi--Yau の場合にも，これと Ricci-flat 形式は区別されている
\cite[Introduction, Remark 7.4]{FBP}．同じ cohomology class の形式を各固定 $t$ で
Moser 移送できても，その移送の計量的な一様性は別に必要である．
実際，$\phi_t^*\omega_t'=\omega_t$ から従うのは
$\phi_t^*g_t'=\omega_t(\,\cdot\,,\phi_t^*J\,\cdot\,)$ であり，
これを $g_t=\omega_t(\,\cdot\,,J\,\cdot\,)$ と比較する評価は自動ではない．
本稿では局所移送の変位，全域 scalar 下界と portal 上界を直接評価し，
指定された actual 負 KE 計量と同じ $f_t=p_T\pi_t$ を比較する．

局所 torus fibration により近い先行結果は，
Ruan\cite[Proposition~3.1, Lemma~3.3, Theorem~3.1]{RuanHMinimal}である．
同論文は homothetic toric potentials と一様 $C^\infty$ quasi-isometry の下で，
potential の torus 平均，Hamiltonian-gradient flow，および陰関数定理を用い，
固定された縮小 polytope 上に H-minimal Lagrangian torus fibration を構成する．
Theorem~4.2 は simple toroidal canonical degeneration の最大退化点へ消える
minimal Lagrangian torus を与える．従って平均化と Moser 法の組合せ自体を
本稿の新規な方法とは扱わない．
本稿の定理\ref{ke:short-torus-averaging}は，$C^0$ の tensor 比較と短い作用軌道だけから，
変位と ambient fiber diameter を直接評価する弱い入力・弱い結論の定量化である．
H-minimal 性，special Lagrangian 性，変形後の intrinsic fiber metric は結論しない．

別研究\cite{KECubicalCurrent}の diameter-scale 候補は
\[
 R_t=\log(-\log|t|^2),\qquad
 (X_t,R_t^{-2}g_t)\longrightarrow(\mathcal C_d,d_\square).
\]
この収束は同稿の粗い比較仮定 $\mathrm{(CP)}_d$ と，
増大する積領域の一様比較・整合性を課す $\mathrm{(PBE)}_d$ の下でのみ主張されている．
さらに我々の fibration による収縮と解釈するには，適切に選びパラメータ化した
$\pi_t:X_t\to B_d(T)$ と系\ref{topo:KE-cubical-comparison}の $p_T$ について，
KE 側の近似写像 $\Phi_t$ が
\[
 \sup_{x\in X_t}d_\square\bigl(\Phi_t(x),p_T\pi_t(x)\bigr)\longrightarrow0
\]
となることを追加で示すのが一つの十分条件である．
命題\ref{topo:ke-map-compatibility}は，この条件から正則・特異 fibers の
\emph{ambient} 直径収縮と距離歪みの評価を導く初等的な比較原理である．
本稿はこの条件付き定理を引用して収束を結論する経路を採らず，
同じ $f_t=p_T\pi_t$ の全域上下界を直接証明する．
従って $\mathrm{(CP)}_d,\mathrm{(PBE)}_d$ 自体の証明を主張しない．
$d=5$ では $p_T$ は同相であり，$d\ge6$ で潰される内部木上の全 fiber 合併は
命題\ref{ke:collapsed-tree-full-fiber-collapse}で制御する．

三重点の積座標では，三枚の cubical squares を一つの toric hexagon によって
全閉領域上で同時に実現できる（命題\ref{ke:triple-hexagon}）．
定理\ref{ke:triple-moving-fibration}は，
$\ell_a=-\log|\zeta_a|^2>m_t/2$ の全 phase を含む領域で
actual KE metric と cusp Hermitian model の一様 tensor 比較を仮定し，
$\min\ell_a>m_t$ における Lagrangian fibration の
ambient fiber diameter $O(m_t^{-1})$ と，
raw logarithmic-depth map との誤差 $O((1+\log m_t)/R_t)$ を与える．
例えば $m_t=R_t$ と選べる．この tensor 比較は，体積密度の上下界だけからは従わない．
関連 KE 原稿\cite{KELocalCurrent}の
\texttt{coarse:prop:density-counterexample} は，滑らかな負 KE 計量に限っても
その一般的な推論が成立しないことを示す．
固定された affine simplex の compact subset は径の尺度では一つの triple vertex に縮むので，
既知の固定 compacta 上の局所収束だけで上の動く領域の仮定を埋めることもできない
（命題\ref{ke:fixed-simplex-scale}）．
この解析的比較は complete simplex model と境界体積比較によって検証できる．
系\ref{ke:actual-closed-corner}は任意の固定 $d\ge5$ で，三枚の閉 square 全体に一つの
actual Lagrangian torus bundle を構成し，$m_t=R_t$ のとき fiber 自身の intrinsic diameter も
$O(R_t^{-1})$ と評価する．
さらに定理\ref{ke:closed-corner-gh}は，同じ Moser flow と同じ fibers を用いて，
三枚の閉 square 自身の内在距離への projection distortion を
$O((\log R_t)/R_t)$ と評価する．
系\ref{ke:closed-corner-protected-gh}により，$\min\ell_i>R_t/2$ の保護領域内の
全経路を許してもこの評価は保たれる．
模型側の入力は，関連 KE 主稿の係数 $1$ の rank-one Schwarz 評価，
全正深度の二重 cusp 漸近，および clipped depth map による全経路の長さ下界である．
これに有限 square gallery の持ち上げを合わせ，model metric の Lie derivative 評価で
actual Lagrangian bundle へ移す．
標的の距離は閉 corner の内在距離として定義されるが，命題\ref{topo:coordinate-retraction}の非伸長 retraction により，全複体 $\mathcal C_d$ の距離の制限と等しい．
この局所距離評価だけから全曲面の shortcut の制御は得られない．

その相対接着に向け，命題\ref{ke:relative-fiber-preserving}は，
元の phase tori が既に actual Lagrangian である飽和 collar 上で
averaging の chain primitive による Moser flow が fiber 分割を保つことを示す．
Metric の torus 不変性は不要である．
Actual neck との orbit partition と quantitative buffer は後述の別の構成で検証し，
この命題だけから既存の大域 fibration への接着を結論しない．

同じ normal circle を持つ neck と triple collar を扱うため，
定理\ref{ke:staged-neck-averaging}では normal average $\eta=A_1\omega$ を先に作り，
その後に残る円を平均する．第一の map を $\phi$，第二を $\psi$ とすると，
$\psi$ は normal circle と可換であるから，$\Theta=\phi\psi$ の各 torus は
neck と共通の actual action $\phi\rho_1\phi^{-1}$ を含む．
単一の primitive $\omega=d\lambda$ を固定すれば
$\Theta^*\lambda=A_2A_1\lambda+dq$ であり，全整数 cycle の作用周期が一致する．
命題\ref{ke:staged-hamiltonian-comparison}は，平均順序や整数基底の違いを，
共通の buffer 上の Hamiltonian isotopy によって比較する．
Exact symplectomorphism であることだけを Hamiltonian 性の根拠にはせず，
primitive を exact 差まで保つ path homotopy を明示する．
系\ref{ke:staged-closed-corner-gh}は，この構成にも
$O((\log R_t)/R_t)$ の局所 GH 誤差を与える．
従来の一段階 bundle と同一の fiber 分割を主張するものではない．

Source-free な固定 compact base 上では，関連 KE 主稿\cite{KELocalCurrent}の
\texttt{pair:prop:bundle-conic} が射影的 normal-bundle coordinates から一つの正則円作用を与え，
\texttt{pair:cor:global-primitive} と \texttt{pair:cor:actual-common-moment} が
global primitive と円平均後の共通 moment interval を与える．
後者は固定 end での KE 収束と cusp 漸近を使う．
本稿の定理\ref{ke:fixed-base-neck-moser}は，capped density の大域下界と
Chern--Lu 評価を組み合わせ，固定残余基底上の moving normal strip で
actual 二側比較・短円・共通 Moser source と像を与える．
命題\ref{ke:fixed-base-moment-connectors}では同一 normal moment 上の残余点を
有界長で接続し，その level 全体の正規化 ambient 径 $O(R_t^{-1})$ を得る．
固定残余基底では source を含む一つの作用で actual neck と閉じた pinched torus も
構成する（定理\ref{ke:projected-source-neck}）．
補題\ref{ke:source-reduction-pinched-fiber}では，固定点での smooth exact transport と
canonical moment の零への規格化を確認し，位相的な零 reduction と ambient の
滑らかな Lagrangian strata を区別する．Normal orbit は primitive な消滅cycleであり，
閉じた source fibre は focus--focus 点を一つ持つ pinched torus となる．
Moment に依存する残余関数 $A(\mu,p)$ を用いた写像の非退化条件は
$d_pA(0,q)\ne0$ である（系\ref{ke:source-focus-transversality}）．
この写像の条件と葉の分割の型を分け，後続の相対補間では source の submersion germ を保つ．
Focus の局所型と monodromy 自体は既知の結果\cite[\S\S2--4]{ZungFocus}であり，
ここでは actual KE neck の共通作用・buffer・全閉 fibre へ適用する条件を実証する．

三重点への接続には，節\ref{ke:neck-corner-transition-section}で
固定比に分離した三つの arm を用いる．射影 Morse 座標と一つの共通 triple 座標の
対数比は，全 phase・全固定階の cusp 微分に関して指数的に小さい．
積を保つ smooth 補間は非正則になり得るが，平均形式と凸 path の tame 性を保つ．
Canonical primitive の base-only closed 差と exact 差を分離することで，
normal period の厳密な一致も得る．第二 period は別の条件である．
さらに節\ref{ke:relative-corner-section}で比較 Hamiltonian 自体を評価し，
三つの端に指定した normal-first 平均を一つの actual corner bundle へ相対統合する．
定理\ref{ke:slow-cut-closed-corner-gh}は
$m_t\to\infty$，$\log m_t=o(R_t)$ の遅い cut に対しても，
同じ corner 射影の局所 GH 近似を保証する．
節\ref{ke:neck-corner-exhaustion-section}では，固定残余 compact，normal cut，
共通 moment 幅，退化 parameter の閾値の順に選ぶ．Centered normal flow の自然性と
moment-dependent reduced foliation により，moving end の全軌道を共通領域に収める．
系\ref{ke:neck-corner-saturated-union}は，closed corner と切詰めた neck の
全ての実際の交わりで全 fiber が一致することを示し，それらの和集合上の
actual KE--Lagrangian 部分 fibration を与える．
補題\ref{ke:fixed-marked-tree-exhaustion}は，source の小円板での一価な isotopy と，
固定 compact の外に支持を持つ end 補正を用い，全 stage と全ての小さい複素 $t$ で
同じ marked pants 木系を選ぶ．これは全 torus cycle の parameter monodromy の
自明化や，異なる $t$ の nodal neighborhoods の semi-global symplectic 同値を意味しない．
後述の定理\ref{ke:actual-global-fibration}は，これを component cores へ相対接着して
全曲面を覆う．その同じ fibration の全曲面距離比較は，
後述の global clipped-depth map と共通 moment の接続で証明する．

Source の計量的な問題については，節\ref{ke:source-connectors-section}で
固定中央 rescaling の volume density を $L_t^{-2}$ と両側から比較する．
Canonical potential は一様に有界で，中央 moment は $O(L_t^{-1})$ にある．
Ricci 下界と Laplacian の障壁により actual inradius $O(L_t^{-1/2})$ を得て，
中央点を oriented arm の断面へ有界長で接続する．
系\ref{ke:fixed-base-central-band}は source-free な中央 moment 帯の全点を，
命題\ref{ke:source-central-band-reduction}は source 近傍の中央帯を固定 core へ，
それぞれ有界長で結ぶ．さらに固定 source tube 上の正の canonical potential の
局所対数勾配から，中央点の actual 球が normal 深度 $\kappa L_t$ 以上に留まることを示す．
距離函数の complex Sobolev norm に Dinh--Marinescu--Vu の
Moser--Trudinger 不等式\cite[Theorem~1.2]{DMVMT}を適用し，
logarithmic cover の周期数 $\asymp L_t$ を用いると，この広い中央帯で
actual 体積の Gaussian tail を得る．Bishop--Gromov 比較と合わせて，
定理\ref{ke:source-central-diameter}は固定点を含む中央 source の全点径を
$O(1)$ に抑える．固定残余基底では中央 moment 帯と閉じた pinched torus
全体にも適用できる（系\ref{ke:fixed-base-source-central-band}）．
用いた Moser--Trudinger 不等式自体は既存結果であり，ここで検証するのは
退化する実計量の距離函数への一様な適用，周期因子と actual 球の領域条件である．
積分測度には固定した補助 Euclidean potential を用い，actual potential の
一様 Hölder 連続性は仮定しない．
Potential の一様連続度による局所径評価\cite[Proposition~1.4]{GGZDiameter}は
同じ $O(1)$ の径を与える別の十分条件だが，その連続度はまだ得ていない．
本稿の全点径評価はその modulus を使わない．

Source の oriented arm では，normal covector の Schwarz 不等式に
固定 normal 境界・source-free な残余境界・moving normal 境界の上界を与える．
最後の境界では正の potential により orientation を保つ固定半径の actual 球を取り，
局所 Schwarz 評価を用いる．最大値原理で normal covector を全 arm へ制御し，
射影下界と体積上界を合わせて source 軸を含む二側 tensor 比較を得る
（命題\ref{ke:source-arm-uniform-tensor}）．
固定 normal 座標での残余水平経路と，最後の normal 深度補正により，
全共通 moment で同じ moment の残余境界へ $O(1)$ で退出できる
（定理\ref{ke:source-all-moment-exit}）．
中央点には短い gate 退出を先に使うので，この直接証明は上記 Gaussian 評価にも
依存しない（補題\ref{ke:source-direct-central-exit}）．
Gaussian 体積評価は独立な補助結果として保持する．
固定残余基底における同一 moment の全 fibers の合併も，
系\ref{ke:source-inclusive-same-moment-collapse}により収縮する．

成分 core については節\ref{ke:component-core-transfer-section}で，
fixed-gauge smooth KE 収束から canonical primitive の差を明示し，
固定 compact から actual KE 形式への小さい symplectic 移送を構成する．
境界の点を固定する相対移送には $H^1$ の制限写像で記述される flux 条件があり，
normal fiber periods だけではこれを検出できない．
一般位置の $m=d-1$ 直線補空間の compact core $K$ と $S=\partial K$ については，
Cohen--Suciu\cite[Proposition~2.5 and \S6.9]{CohenSuciuBoundary}の既知の境界 cohomology を用い，
\[
 H^1(S;\RR)/\operatorname{im}H^1(K;\RR)
 \cong H^2(K,S;\RR),\qquad
 \dim H^2(K,S;\RR)=\binom{d-2}{2}
\]
と計算できる．命題\ref{ke:component-flux-space}はこれを有限個の boundary periods による
相対 Moser の判定へ適用する．系\ref{ke:nonextendible-small-core-germ}では，
同じ actual KE 二形式を保つ小さい collar germs でも各非零障害方向を実現する．
平均形式も parameter 方向に Moser 移送することで exact collar germ も得る．
一方，定理\ref{ke:actual-global-fibration}の構成では，actual link の葉を先に固定し，
同じ actual 二形式の primitive を相対不変化して内部へ延ばすため，
点固定の相対 germ を追加入力とする必要はない．


開いた二次元セルの固定 compact 部分では，さらに明示的な product model と局所 GH 極限を得る．
関連 KE 原稿\cite{KELocalCurrent} の \texttt{simplex:cor:actual-dominant-product} は，
三重点の独立深度 $a,b\in[\tfrac12L_t^\eta,2L_t^{1-\eta}]$，$0<\eta<1/2$ において，
actual metric と product-cusp metric の全固定階の差を $O(L_t^{-\eta})$ と評価する．
証明では complete simplex と product domains の包含によって model volume ratio を $1$ に近づけ，
全深部三重点の actual semi-flat 比較と合わせ，
Savin\cite[Theorem~1.3]{SavinSmall} の小摂動正則性から tensor 比較を得る．
本稿の定理\ref{ke:open-cell-collapse}はこれを用い，
$[\eta,1-\eta]^2$ 上の actual KE--Lagrangian torus fibers の intrinsic diameter を
$O(L_t^{-\eta})$ とし，Moser 像領域の内在距離について平坦長方形への GH 収束を証明する．
Projection distortion も $O(L_t^{-\eta})$ である．
全曲面内の ambient distance は領域外の近道を許すため，この局所定理だけから
大域 GH 収束や既存 fibration の全 fibers に関する収縮を結論しない．

\paragraph{固定成分 core と target の距離．}
節\ref{ke:component-core-geometry-section}では complete KE 形式に対し，
任意の固定 compact 上で形式と primitive を保つ smooth 正則化と，
正 residue の有限型 Weinstein core を構成する．既存の相対 SNC 定理を用い，
その小さい飽和近傍を Lagrangian に filling する．

全域への相対構成は節\ref{ke:component-link-cover-section}--
\ref{ke:actual-core-gluing-section}で行う．Period $1$ の正の meridian に対し，
$M_a+2\pi w_a$ が既存の向き付き canonical normal moment と厳密に一致することを，
volume-frame ratio の整数 winding と exact averaging transport から示す．
$N_j\iota_j\gg1$ の深さで normal moment faces を取り，corner では二つの
moment cuts を smooth monotone minimum で丸める．
固定段階でも，さらに深い tube 上で同じ centered flow の逆像 buffer を使い，
fixed-shell moment の上界を加えると，既存の部分 fibration が中央の全特異 locus の
固定近傍を覆うことが分かる．その補集合は各 component の固定 compact に入る．

Rochon--Zhang の全 cusp 微分評価\cite[\S5, (5.3)--(5.6)]{RochonZhang}と
三重点の全固定階比較から，actual generators の $C^1$ smallness を導く．
Actual symplectic circle に対する centered scalar average は primitive を exact に補正し，
common normal action に不変な cutoffs で貼り合わせても
$Z^{\rm ad}\Psi>0$ を保つ．これにより既存 core を変えず，closed link 全体に外向きな
Liouville field を得る．Saddle を動かす相対葉族補間は，
各 critical pants と全 regular annuli に一つの vector field を作ることで証明する．
この補間は非零 moment band に置くので source の全 pinched tori は保たれる．

境界から内部への cone extension 自体は Evans--Mauri の既存結果
\cite[Lemma~2.16]{EM}に基づく．本稿では contact form を替えずに compact domain 上で用い，
actual link の既存葉，normal/corner overlaps，source を避けた相対補間，
全 fiber 単位の切詰めを合わせて検証する．
従って全 $X_t$ を覆う actual KE--Lagrangian fibration が得られる．
内側の primitive は compactly supported exact correction により Morse--Smale に選び，
index 順の閉 filtration によって skeleton の弱い isotropic stratification を確保する
（補題\ref{ke:compact-relative-morse-smale}）．
境界の contact form と既存 normal/corner actions は変えない．
各 component の変更領域は深度 $CN_j$ 以下の固定 compact image にあり，
cusp の動径経路により接続費用を $C(1+\log(N_j+1))$ と評価できる
（命題\ref{ke:component-logarithmic-connectors}）．
従って十分遅い対角選択でその正規化 ambient 径を消せる．
この定理自体の結論は保持 corner の局所構成を伴う全域の存在であり，
距離歪みの消失には次の追加評価を用いる．
この相対構成全体の一般的新規性は未確認である．

Target の組合せ型は既知の cubical barycentric subdivision\cite{BabsonBilleraChan}であり，
その構成自体の新規性は主張しない．命題\ref{topo:coordinate-retraction}は，unit-square 長さ距離に対する座標部分複体の等長性を与え，
閉 corner の target 距離を全体距離へ同定する．
定理\ref{topo:diameter-six-labels}は有限幾何のみで，六次以上の target の
厳密直径 $2\sqrt5$ と全最大対を決定する．
これらは actual 曲面の全域 shortcut を制御する結果とは区別する．

\paragraph{全曲面での距離下界と同じ基底値．}
Source を含む固定 residual domain 上に正則 normal frame $x$ を取り，
$s=-\log|x|^2$，$U=|d\log s|_{g_t}^2$ とする．
Bochner 計算から $\Delta\log U\ge U-1$ が得られる．
これは Schwarz 型評価\cite{YauSchwarz}の rank-one 版であり，
本稿では実計量と複素 Laplacian の規格化を含めて係数を直接計算する．
固定外側 shell，残余基底の境界，moving central shell の三つの境界を検証し，
\[
 B=1+\varepsilon_j+A_j\mathfrak n_j/s+C_js/(\kappa L_t)
\]
を障壁として用いる．二つの cut を十分離すことで，source を除くことなく
active な $|d\log\ell_i|_{g_t}$ を $1+o(1)$ に抑える
（定理\ref{ke:sharp-global-one-active}）．
体積密度の二側比較だけから sharp tensor bound を推論してはいない．

Global clipped-depth map $\Theta_t$ は，高々三つの非零 coordinates と
少なくとも一つの plateau を持つ．二つの active coordinates では全深部 triple comparison，
一つの場合は上の scalar estimate を用いる．
得られる長さ下界は全曲面の任意の経路に直接成立し，領域の出入りの回数に依存しない．
その後で，保持 corner では旧 softmax と短い Moser 移送を比較する．
Neck の任意の点には，同じ canonical moment を持つ旧 corner 境界上の参照 torus がある．
Source と branch を含む同一 moment 集合の接続評価と，既に証明した
$\Theta_t$ の Lipschitz 評価により，その参照 torus での値比較を全 neck へ移す
（補題\ref{ke:same-moment-corner-reference}）．
Core の像の小ささと合わせて
$\sup_{X_t}d_\square(\Theta_t,p_T\pi_t)=o(1)$ を得る．
ここで用いる接続評価は局所解析から独立に証明してあり，全域の距離上界を先取りしない．
Clipped depth が個々の Lagrangian fiber 上で一定であるとも仮定しない．

\paragraph{基底の接着と全曲面での距離上界．}
旧 corner の境界 torus と新しい cellular branch locus は同じ strata ではない．
命題\ref{ke:global-base-parameterization}は保持 corner を少し内側へ写し，
その外周に neck strips と component sectors を配置する．
旧 softmax 値との差を $o(1)$ に保ち，内部木の全枝に一つの canonical moment を使う．
この同じ moment 上の source を含む接続定数 $C_j^{\rm neck}$ は，
component cut $N_j$ より先に固定できる．Neck の component 側の尾部も深度
$CN_j$ 以下に入ることを別に示す．従って
$C_j^{\rm neck}\le1+\log(N_j+1)$ となるよう cut を選べば，
全接続に一つの予算 $H_j=C(1+\log(N_j+1))$ を使える．
Target の最短路を高々 $60$ 本の square 線分に分け，共通 moment の portals で繋ぐと
同じ $f_t$ の全曲面距離上界を得る．その誤差は
$C_d\{b_j^{-1}+(1+\log R_t+\log(N_j+1))/R_t\}$ 以下である．
同一 moment の参照点を使った下界でも，残余接続費用をこの同じ cut に吸収できる．
式\eqref{ke:eq:global-realization-explicit-cut-budget}は，上下両側を合わせた distortion を
固定族の一つの係数と，$b_j,\varepsilon_j,\mathfrak m_j,K_j,N_j$ の明示式で抑える．
これは許容 cut の下での評価であり，残余定数の増大率や最適収束速度を主張しない．
定理\ref{ke:global-metric-realization}では有限個の stage 定数を先に選び，
共通の遅い対角化で上下両側の誤差を零にする．
これらの source 評価，scalar 障壁，基底の相対 parameter 化と全域距離比較が，
本稿で既存の局所理論に加えて証明する部分である．
全体の先行研究に対する新規性の確定と包括的な独立検証は，公表前の課題である．

\subsection{他の tropical 構造との比較と残る問題}

Cartwright の tropical complex は，退化の dual complex に intersection numbers を加えて，
線形関数の層を指定する構造である\cite[Definitions~2.5, 2.8; Construction~3.2]{CartwrightComplexes}．
作用周期から作る本稿の atlas と同一視することはできない．
実際，ordered small resolution の最初の component における local intersection matrix は，
全平面の場合 $J_{d-1}-dI$ で負定値となり，tropical complex の local positivity condition を満たさない
（命題~\ref{atl:cartwright-distinction}）．
同論文 Remark~3.4 も，Gross--Siebert 型の affine structure との特異集合の違いを区別している．
本稿でも，focus points を残した atlas に compact smooth tropical manifolds の大域定理を
追加の検証なしに適用しない．

また，標準 fan の chart と特定の algebraic filling との同定は別の条件である．
Kapranov\cite[Theorem~4.3.3]{KapranovChow}の blow-up 記述によれば，
$\overline M_{0,5}$ の boundary は十本の $(-1)$-curves の Petersen 配置である．
五直線 complement の模型は common-circle determinants と boundary homology によって，
この標準 $M_{0,5}$ filling と区別される
（命題~\ref{trop:divisor-recovery}，系~\ref{trop:not-standard-M05}）．
二つの boundary の $b_1$ はそれぞれ十と十一である．
従って作用適合 Petersen chart の存在から，その filling の置換を結論してはいけない．

全平面五次の高分母 chart の同時接着とその支持 graph の問題は，
companion 論文\cite{QuinticArithmetic}で，整数データの必要条件と
幾何学的な実現条件を区別して検討する．そこでは非空の高分母支持が葉を持たず，
最小次数三以上・高辺八本以上の三つの五頂点 graph 型，二十六標識に限られることを示す．
非空支持は三点連結であり，単位分母の辺は matching をなす．
特に一つの core に接する二本の branches の分母が一なら，全分母が一となる．一方，
残る整数接着の存在・一般の分母剛性は未解決である．
同 companion 論文では，全ての pair products が正である四ray整数データの正値性を，
完全 matching との同値性で判定する．
単位枝の node 配分は $\{5,0,0\}$ または $\{4,1,0\}$ に限られ，
二単位枝を持つ $G_{5,14}$ では後者だけが残る．
これらは局所の正面積と整数接着に関する必要条件であり，残る三型の大域実現を主張しない．
本稿の構成に続く幾何的な問題として，atlases の同値関係と変形空間，
および core を替える相対 surgery が残る．
四 puncture の内部 cones の格子 index を $m^-,m^+$ に緩める場合も，
$q_Aq_B\mid(m^-+m^+)$ という必要条件が残る．
一般次数における非 matroidal links を，より広い tropical models で扱う可能性も残る．
hyperplane class 以外の Kähler classes，special Lagrangian 条件，本稿以外の退化族における canonical metric の collapse は，
本稿の存在定理から独立した課題として扱う．
Gross--Siebert\cite{GrossSiebert}の affine geometry からの退化再構成との比較にも，
本稿で得た実際の作用周期に加えて，必要な追加データを明示する必要がある．

\section{設定と主定理}\label{sec:main}

全ての多様体、形式、写像は、特に断らない限り滑らかである。
Hamiltonian vector fieldは$\iota_{X_H}\omega=-\mathrm dH$で定義し、normal circleの周期は$2\pi$とする。
従って標準複素直線上の正のmomentは$|z|^2/2$である。
$\dd^c$は$\dd^c\log|z|^2=\dd\theta$となる規約を用いる。
$H$はhyperplane classで、$\int_{\PP^1}H=1$と正規化する。

\begin{definition}\label{def:fibration}
本稿の\emph{分岐底空間上のLagrangian torus fibration}とは、コンパクトなsymplectic 4多様体$(X,\omega)$から有限2複体$B$への、連結なfibersを持つproper continuous surjection $\pi:X\to B$で、$B$の稠密な開2次元stratum上では滑らかなsubmersionとなり、そのfibersがLagrangian $T^2$となるものをいう。
特異fiberには、以下に指定する局所的な層構造を与え、その滑らかな各stratumはisotropicとする。
ここでstratificationは\cite[Definition 2.1]{EM}の弱い意味、すなわち閉部分集合のfiltrationの各差が滑らかな多様体となるという意味で用い、Whitney条件は含めない。
\emph{自由な辺がない}とは、2次元stratumの閉包に接する各開1次元stratumに、ちょうど一枚だけのsheetが接することがない、という意味である。
\end{definition}

この定義は、基底が通常の2次元多様体であるという条件を含まない。
本稿では分岐辺の横断的な断面は三叉グラフ$Y$である。
正則な底空間にはaction-angle theoremによるintegral affine structureがあるが、分岐を横断する単一のaffine chartは定義からは得られない。

各$i\ne j$に対し、$d-2$個の葉が$\{1,\ldots,d\}\setminus\{i,j\}$でラベル付けされた三価木$T_{ij}=T_{ji}$を選ぶ。
これは実際の多穴球面$C_{ij}^{\circ}$上のmarked pants分解で実現し、単なる抽象木の選択とは区別する。
$B_d(T)$は、各成分のcoreを一つの頂点に、各$T_{ij}$の内部頂点を$i,j$間の分岐辺に、内部辺を二角形に対応させ、各三重点を三角形に対応させた2複体である。
厳密なcellular定義と接着写像は底空間の節で与える。

\begin{theorem}[境界のない分岐底空間への構成]\label{main:global}
$X_d\subset\PP^3$を滑らかな次数$d\ge5$の曲面、$\omega$を$[\omega]=aH$, $a>0$を満たすKähler formとする。
任意に選んだ上記のmarked pants treesを用いて
\[
\pi_d:(X_d,\omega)\longrightarrow B_d(T)
\]
という定義\ref{def:fibration}のfibrationを構成できる。そのfiberは次の表の通りであり、他の型の特異fiberを生じない。
\begin{center}
\begin{tabularx}{\linewidth}{@{}lXr@{}}
\toprule
底空間のstratum & fiber & 個数または成分数\\
\midrule
正則な開2次元部分 & Lagrangian $T^2$ & ---\\
開分岐辺 & $S^1\times\figeight$ & $\binom d2(d-4)$\\
面内のnodal values & 一度pinchされた$T^2$ & $d\binom d2$\\
core頂点 & $\PP^2$の$d-1$直線補集合のlogarithmic Weinstein skeletonのsymplectic image & $d$\\
\bottomrule
\end{tabularx}
\end{center}
各分岐辺にはちょうど三枚のsheetが接する。
基底は$\binom d3$枚の三角形、$\binom d2(d-5)$枚の二角形を持ち、
\[
B_d(T)\simeq\Sk_2\Delta^{d-1}\simeq\bigvee^{\binom{d-1}{3}}S^2.
\]
$d=5$では$B_5\cong\Sk_2\Delta^4$であり、$d\ge6$では$B_d(T)$は$\Sk_2\Delta^{d-1}$と同相ではない。
\end{theorem}

coreの同定は、適切に選んだlogarithmic Weinstein構造のskeletonとの同定である。
特定のembedded skeletonを最初から指定して保存する、という主張ではない。
Morse型coreの等方的なstratificationには\cite[\S2.3]{EM}を用いる。
core fiberのEuler数は$\binom{d-3}{2}$であり、$d$とともに変化する。
従ってこの定理は、全ての次数に共通の一種類のpants fiberだけで構成するという結論も含まない。

\begin{theorem}[vertical Kähler slice]\label{main:cone}
後述の順序付き小解消$f:\widetilde{\cX}_d\to\Delta$の中心成分を$D_1,\ldots,D_d$とする。
$a>0$を固定し、$\sum_i[D_i]=0$による共通定数の不定性を除いて
\[
\xi=aH+\sum_{i=1}^d\beta_i[D_i]
\]
を考える。
底を十分縮小した全空間上のKähler classとして$\xi$が実現できるための必要十分条件は
\[
\beta_1<\beta_2<\cdots<\beta_d,
\qquad a+\sum_i\beta_i-d\beta_d>0
\]
である。$\delta_k=\beta_{k+1}-\beta_k$を用いると、この領域は
\[
\delta_k>0\quad(1\le k\le d-1),\qquad
\sum_{k=1}^{d-1}k\delta_k<a
\]
という開単体になる。
\end{theorem}

これは全相対Kähler coneの記述ではなく、指定したvertical affine sliceとの交わりの完全な記述である。
大域fibrationの証明には、この単体の中の十分小さいrayを用いる。
単体の全ての点で同じ局所surgeryが実行できると主張する必要はない。

\paragraph{証明の構造.}
まず明示的な平面退化と偏極を構成する。
次に、直線配置補集合の有限体積compactificationを、境界上のmarked Morse葉層に相対的に作る。
小さいdivisor-compatible blow-upによってfocus--focus fibersを挿入し、実際の中心成分へfibrationを移す。
その後、divisorから正の距離を置いた飽和領域を、同じ全空間のsymplectic connectionで平滑化へ運ぶ。
二重曲線ではmomentに依存するMorse葉層を補間し、三重点ではtoric moment cutsを使う。
この順序により、nodeのコンパクトなfiber全体を保ち、gluingの際にfiberを途中で切ることを避ける。

\paragraph{一般化と局所的な障害.}
上の平面模型に必要な相対 compactification は，正の重み付き Kähler divisor 表示を持つ
任意の genus の SNC divisor に対して構成する（定理\ref{snc:main}）．
これにより，$d=e_1+\cdots+e_k$，$k\geq3$ の次数分割に対し，
$k$ 個の core と $d\sum_{i<j}e_i e_j$ 個の nodes を持つ構成を得る
（定理\ref{part:global}）．その底空間は
\[
 B_{\boldsymbol e}\simeq
 \bigvee^{\,p_g(X_d)-\sum_i p_g(S_i)}S^2
\]
であり，整数鎖の強変形収縮を明示する（定理\ref{part:homotopy}）．
ここで $S_i\subset\PP^3$ は次数 $e_i$ の smooth surface である．

この構成の core を保つ場合，三価 Bergman link の分類から $d\geq7$ の全ての次数分割に
局所位相障害が生じる（定理~\ref{trop:cubic-classification}，命題~\ref{trop:partition-obstruction}）．
六次では全平面分割だけが残り，六つの同時 Petersen systems がある．
五次では全平面と $(1,1,1,2)$ だけが全 core の matroidal local topology を持ち得る．
後者の標識付き抽象 graph systems は四個である（定理~\ref{q5:classification}）．

\paragraph{作用適合 atlas の存在と分類.}
作用適合 matroidal atlas の transitions は，ambient integer affine maps の制限とする．
正則 faces の余接線格子は actual fibration の作用格子に等しい．
focus points を通常の matroidal charts から除き，その punctured neighborhoods に
実際の focus affine singularities を残す（定義~\ref{atl:definition}）．
branch denominator $q$ は，primitive branch tangent に対する正 normal action の
微分の絶対値であり，以下の構成は全 branch で $q=1$ を満たす．

\begin{theorem}[五次・六次の大域作用適合 atlas]\label{main:action-atlas}
$d=5,6$，$X_d\subset\PP^3$ smooth とし，$\omega$ を
$[\omega]=aH$，$a>0$ を満たす Kähler form とする．
定理~\ref{main:global} の全平面 fibration を，focus affine singularities を許す
大域作用適合 matroidal atlas を持つように選べる．
core charts は五次では標準 $K_4$，六次では標準 Petersen fans である．
五次の十 triangles には各五つ，六次の二十 triangles と十五 bigons にはそれぞれ三つと二つの
nodal values が入る．全 branch の正規化された affine length は $a/(\pi(d-1))$ である．
\end{theorem}

\begin{remark}[作用 atlas と KE 計量実現の選択]\label{main:atlas-ke-scope}
定理\ref{main:action-atlas}は，指定された Kähler form ごとに fibration を選ぶ存在定理である．
各小さい $t\ne0$ の actual 負 KE form もこの対象となる．
五次ではさらに，系\ref{ke:quintic-simultaneous-main}と
定理\ref{ke:quintic-simultaneous-realization}により，作用 atlas を持つ同じ射影を
全曲面 ambient 距離の GH 実現にも使える．等重みの対称四角形系を選べば，
各 triangle に五つの nodal values が入る．
その証明では actual primitive の両端切片と整数 framing を有限式で合わせ，
初期 source 葉の全 tail を小面積に選び，source wall を含む全 moment の面積輸送を行う．
同じ canonical moment，保持 corner の全葉，中央 nodal neighborhoods と
全 overlap の whole-fiber partition を保つため，
命題\ref{ke:moment-relative-replacement}から同じ射影の距離比較を得る．
任意に先指定した source fibers・大きい保持 annuli・critical values への適合は要求しない．
六次でも，定理\ref{ke:sextic-simultaneous-realization}が標準等重み Petersen 系について
同じ射影による atlas・GH 実現・確率体積極限を与える．
その追加内容は，内部二 source block の累積位置を含む二端 cut，
各側三 sources の separator による全 moment 実現，中央 source 全近傍を保つ
parameter の固定，及び平行移動まで一致する三面の整数アフィン接続である．
\end{remark}

定理\ref{main:action-atlas}の証明は，局所作用座標，整数 framing，critical-area を指定した actual neck 接着を
まとめる定理~\ref{ga8:global} で与える．
個別の core charts を並べるだけでは，作用の加法定数と atlas transitions は決まらない．
本稿は固定 nonzero moments 上の exact symplectic transport と相対 Moser により
この接着条件を満たす．

混合分割 $(1,1,1,2)$ では，三つの $K_4$ と一つの $K_{3,3}$ の各 core に個別の
作用適合 chart がある（系~\ref{mixper:individual}）．
しかし，四 puncture neck の両端の annulus を比較すると，正の分母 $q_A,q_B$ は
$q_Aq_B\mid2$ を満たし，三価 branch の奇数性により両方とも一となる
（命題~\ref{q9:rigidity}）．この剛性と法円作用の分岐係数の総和 $-6$ により，
混合五次の大域接続は任意の分母を許しても不可能である
（定理~\ref{hd9:mixed-no-go}）．

\begin{theorem}[現行 SNC 構成の matroidal atlas の分類]\label{main:atlas-classification}
$d\ge5$，$d=\sum_{i=1}^ke_i$，$k\ge3$ とする．本稿の比例曲率 SNC cores，
既存の branch graph，標準 coarse matroidal cones を保ち，追加の branch singularities を
入れない範囲で，大域作用適合 matroidal atlas を持つ fibration を構成できる次数分割は，
全平面五次と全平面六次に限る．branch denominators に一という仮定を課さない．
\end{theorem}

証明は系~\ref{hd9:classification}で与える．全平面六次の atlas では，
三十本全ての branch denominators が自動的に一となり，各 bigon は二 nodes を持つ
（系~\ref{q9:sextic-rigidity}，注意~\ref{q9:node-budget}）．
一方，四直線の一つの core では任意の奇数 $q\ge3$ に対し分母 $(q,1,1,1)$ の chart が
存在する（\cite{QuinticArithmetic}）．従って分母剛性は大域接着を使う条件である．
分類の範囲は指定した SNC 構成であり，core surgery，fan の変更，追加の branch singularities
を許す全ての fibrations を分類する主張ではない．

\paragraph{五次の整数アフィン接着との関係．}
全平面五次で任意の奇数分母を許す chart の同時接着については，局所の整数行列，
共有 puncture の整数 clutching，node 配分，および対向する分岐係数の相殺を同時に調べる必要がある．
これらの算術条件，高分母支持 graph に対する障害，および局所整数族の分類は，
companion 論文\cite{QuinticArithmetic}にまとめる．
定理~\ref{main:action-atlas}で構成する分母一の大域 atlas の存在証明は，本稿内で完結する．

% !TeX root = ../../surface-lagrangian-torus-fibration.tex
% 2026-11-22 / v2.4: 精密主定理，解析入力，証明経路とcutの量化順序．
\section{負 KE 極限を実現する大域 fibration}
\label{ke:global-realization-section}

本節では，負 Kähler--Einstein 計量の Gromov--Hausdorff 極限を，
actual Lagrangian fibration の射影によって実現する主定理を述べる．
後続の節で全域 fibration を構成し，その同じ射影に対する全曲面の距離の上下界を証明する．
最後の統合と対角選択は定理\ref{ke:global-metric-realization}に置く．

\subsection{固定した退化族と計量の規格化}

整数 $d\ge5$ を固定し，まず\eqref{deg:family}の全平面退化を考える．
具体的には，一般位置の平面 $\{L_i=0\}$ と次数 $d$ の式 $F_d$ による
\[
 X_t=\left\{\prod_{i=1}^dL_i+tF_d=0\right\}\subset\PP^3
\]
を扱う．ここで，$F_d$ が三重点で非零，二重線上での零点が
相異なる単純零点であり，十分小さい $t\ne0$ に対して $X_t$ が滑らかであることを仮定する．
族とその幾何学的補助データは固定し，以下で定数の次数一様性や族一様性は要求しない．

$X_t$ 上の負 KE form とその実計量を $\omega_t,g_t$ と書き，
\[
 \operatorname{Ric}(\omega_t)=-\omega_t,\qquad
 [\omega_t]=2\pi(d-4)H,\qquad
 L_t=-\log|t|^2,\qquad R_t=\log L_t
\]
と規格化する．全曲面上の Riemannian distance を $d_{g_t}$ とし，
\[
 \bar g_t=R_t^{-2}g_t,\qquad \bar d_t=R_t^{-1}d_{g_t}
\]
を用いる．$\bar d_t$ は全 $X_t$ 内の経路を許した距離であり，
三重点の内在距離や保護領域内の距離ではない．
極限複体は系\ref{topo:KE-cubical-comparison}の
\[
 \mathcal C_d=
 \{w\in[0,1]^d:\max_i w_i=1,\ |\operatorname{supp}(w)|\le3\}
\]
であり，unit Euclidean squares の貼合せによる全体の長さ距離 $d_\square$ を入れる．
この規格化は，関連 KE 原稿\cite{KELocalCurrent,KECubicalCurrent}および
定理\ref{ke:closed-corner-gh}の規格化と同じである．

\subsection{大域的な幾何学的実現の主定理}

\begin{theorem}[負 KE 極限の Lagrangian fibration による実現]
\label{ke:global-realization}
上の各固定次数と退化族について，実際の marked pants 分解で実現できる
一つの木系 $T$ と，すべての十分小さい $t\ne0$ に対する
\[
 \pi_t:(X_t,\omega_t)\longrightarrow B_d(T)
\]
であって，次の二条件を満たすものが存在する．
\begin{enumerate}
\item $\pi_t$ は定義\ref{def:fibration}の大域 fibration である．
すなわち全 $X_t$ を覆う連結 fiber の連続全射で，正則 fiber は
\emph{actual form $\omega_t$ に関する} Lagrangian $T^2$ である．
特異 fiber には定理\ref{ke:actual-global-fibration}の nodal，branch，core 模型と等方的 strata を用い，
接着の際に whole fibers を失わない．
\item 系\ref{topo:KE-cubical-comparison}の labelled map
$p_T:B_d(T)\to\mathcal C_d$ に対し，$f_t=p_T\circ\pi_t$ と置くと，
\begin{equation}\label{ke:eq:global-realization-distortion}
 \varepsilon_t:=
 \sup_{x,y\in X_t}
 \left|R_t^{-1}d_{g_t}(x,y)
              -d_\square(f_t(x),f_t(y))\right|
 \longrightarrow0
 \qquad(t\longrightarrow0)
\end{equation}
が成り立つ．
\end{enumerate}
木系の組合せ型と component labels は固定するが，nodal values，
collar の位置，作用座標の表示，および基底の parameter 化には $t$ 依存を許す．
任意に既定の $T$ や任意に固定した critical values に対する実現は，
この定理の結論に含めない．
連続性は各固定 $t$ の写像 $\pi_t:X_t\to B_d(T)$ についての条件であり，
$t$ 方向に連続又は滑らかな fibration 族を与えるという主張ではない．
\end{theorem}

第一条件は定理\ref{ke:actual-global-fibration}，第二条件は
定理\ref{ke:global-ambient-lower}と定理\ref{ke:global-ambient-upper}で証明する．
命題\ref{ke:global-base-parameterization}が両評価を同じ $f_t$ に揃え，
定理\ref{ke:global-metric-realization}が共通の遅い対角選択を与える．
次数範囲は各固定 $d\ge5$ であり，正則部分の action--angle structure を超える
全基底の matroidal atlas や special Lagrangian 性は結論に含めない．

\begin{remark}[Atlas との同時実現に十分な値比較]\label{ke:atlas-ke-value-criterion}
$d=5,6$ とし，本定理の $f_t=p_T\pi_t$ に加えて，一つの固定木系 $\widehat T$ に対し，
atlas を持つ別の actual Lagrangian fibration
$\widehat\pi_t:X_t\to B_d(\widehat T)$ を選べたとする．
対応する labelled 収縮を用いて $\widehat f_t=p_{\widehat T}\widehat\pi_t$ と置く．
もし
\[
 \delta_t:=\sup_{x\in X_t}d_\square(f_t(x),\widehat f_t(x))\longrightarrow0
\]
を証明できれば，命題\ref{topo:ke-map-compatibility}の直接の適用により
\[
 \operatorname{dis}(\widehat f_t)
 \le\operatorname{dis}(f_t)+2\delta_t\longrightarrow0
\]
となる．これは新たな比較原理ではなく，既存の一様基底値比較をこの二つの射影へ適用したものである．
命題\ref{ke:moment-relative-replacement}は，この値比較の具体的な十分条件を与える．
同じ保持 corner の全葉，各 component の core 領域，同じ canonical normal
moment と anchor cuts を保ち，marked product/cone 商と全 overlap の
whole-fiber 接着を持つ相対変更なら，
$\delta_t\le C(1+\log(N_{j(t)}+1))/R_t$ となる．
Reduced 方向の変更を微分まで小さくする必要はない．
定理\ref{ng8:componentwise-critical-area}と系\ref{ng8:p0-small-area-flexibility}により，
相対面積実現の有限判定と source を保持する微小変更は得られた．
Actual core chart には命題\ref{ke:actual-anchor-fan-criterion}の
anchor period 条件を使える．五次の固定 anchor では，
系\ref{ke:k4-protected-anchor-selection}が保持 tail の period による有限判定を与える．
未解決なのは，actual データにおけるその条件の検証，所定の atlas targets に対する
面積配分の可解性と，反対側 neck の切片・整数 framing を含む全 open overlaps の一致である．
定理\ref{main:action-atlas}の各固定 $t$ における存在だけではこの条件は得られず，
一般の symplectic conjugacy だけで critical-action の障害が消えるとも主張しない．
\end{remark}

$f_t$ は全射なので，\eqref{ke:eq:global-realization-distortion}により，
その graph correspondence，または命題\ref{topo:ke-map-compatibility}により
\begin{align*}
 d_{\mathrm{GH}}\bigl((X_t,\bar d_t),(\mathcal C_d,d_\square)\bigr)
 &\le\tfrac12\varepsilon_t,\\
 \sup_{c\in\mathcal C_d}
 \operatorname{diam}_{\bar d_t}f_t^{-1}(c)
 &\le\varepsilon_t,\\
 \sup_{b\in B_d(T)}
 \operatorname{diam}_{\bar d_t}\pi_t^{-1}(b)
 &\le\varepsilon_t
\end{align*}
を同時に得る．ここで直径は全空間の距離を制限した ambient diameter である．
特異 fiber もこの上限に含み，元の尺度での収縮や fiber 内在直径の収縮までは要求しない．
$d=5$ では $p_T$ は同相である．$d\ge6$ では，
一つの $p_T$-fiber 上にある $\pi_t$-fibers の集合全体が縮むことが必要であり，
各 torus の点ごとの収縮だけでは十分でない．

\begin{remark}[径を一にする規格化との関係]
\label{ke:global-realization-scale}
$D_d=\operatorname{diam}(\mathcal C_d,d_\square)>0$ とする．
\eqref{ke:eq:global-realization-distortion}と全射性から
\[
 \left|R_t^{-1}\operatorname{diam}(X_t,g_t)-D_d\right|
 \le\varepsilon_t
\]
である．従って径を一に規格化した計量の極限は
$(\mathcal C_d,D_d^{-1}d_\square)$ となる．
定理\ref{topo:diameter-six-labels}により $d\ge6$ では $D_d=2\sqrt5$ であり，
五次では既存の有限計量計算\cite[\texttt{thm:intro-diameter-C5}]{KECubicalCurrent}から
$D_5=\sqrt{13}$ である．これらの有限計量結果と主定理から，
$R_t^{-1}\operatorname{diam}(X_t,g_t)$ の極限も定まる
（系\ref{ke:global-diameter-limit}）．
元の負 KE 尺度の極限，$R_t$ による規格化，および径一の規格化は区別する．
\end{remark}

\subsection{解析入力の出典と適用範囲}
\label{ke:global-realization-transfer}
\label{ke:proof-inputs}

主定理の証明では，退化族の局所解析，固定成分の cusp 幾何，および本稿の
actual fibration の構成を次のように接続する．退化族に関する局所入力の正本は
\cite[Part~IV]{KELocalCurrent}であり，以下に英字で記す定理ラベルは同稿を指す．
各固定 $d$ と族に対する定数を使い，固定 compact 上の収束から
任意の moving 領域における一様性を推論することはない．

\begin{enumerate}
\item \emph{Global volume 下界と oriented 領域の volume 上界．}
\texttt{coarse:cor:degree-d-volume-lower} は，全曲面上の滑らかな capped density に
対する actual KE 体積の下界を与える．
\texttt{coarse:lem:pair-arm-density} と
\texttt{coarse:lem:complete-volume-upper} は，source を含む oriented arm で
比較する model density と上界を与える．後者の完全計量比較で必要な
log volume ratio の上方有界性は，各固定 $t$ の閉包と model の人工境界での発散から
検証する条件であり，未構成の actual metric buffer を仮定しない．
これらを用い，本稿で別に証明する
補題\ref{ke:capped-projective-lower}の固定残余領域上の射影計量下界，
補題\ref{ke:source-tube-potential}の正の potential と
補題\ref{ke:positive-poisson-log-gradient}の局所対数勾配，
命題\ref{ke:source-arm-uniform-tensor}の三境界 Schwarz 評価を組み合わせる．
体積密度の二側有界性だけを tensor 二側比較に置き換えてはいない．

\item \emph{全深部 triple の tensor 比較と complete simplex の幾何．}
\texttt{coarse:thm:full-triple-tensor} は，三重点の全 phase と全深度比を含む
$\min_a\ell_a\ge m$ 上で，$L_t\ge L_*$，$m_*\le m<L_t/3$ を満たす全 cut に一様な
$O(m^{-1})$ の tensor 比較を全固定微分階で与える．
その局所的な核は \texttt{coarse:thm:boundary-to-triple} である．
同定理は固定 face の volume 下界と，各固定 $t$ の radial 閉包近傍での
smooth 負 KE 計量という入力を分離し，全平面退化への適用では
上の capped 下界でこれを満たす．
Volume ratio が $1$ に近づく評価と
Savin\cite[Theorem~1.3]{SavinSmall}の小摂動正則性を組み合わせる段階が，
有界な volume ratio のみからの推論との相違である．
\texttt{simplex:lem:rank-one-schwarz}，
\texttt{simplex:prop:dominant-product}，
\texttt{simplex:cor:active-cogram}，
\texttt{simplex:thm:closed-corner-model} は，同じ complete simplex 上の
長さと active covectors の比較を与える．
本稿の定理\ref{ke:slow-cut-closed-corner-gh}はこれらを
actual Lagrangian な closed corner と有限 cut の距離評価へ移す．
全曲面の source 又は二重曲線の全領域が triple 比較の適用域に含まれるとはしない．

\item \emph{固定 residual compact と固定 end の normal moment．}
\texttt{pair:prop:bundle-conic}，
\texttt{pair:cor:global-primitive}，
\texttt{pair:thm:fixed-end-moment}，
\texttt{pair:cor:actual-common-moment} は，source-free な固定残余基底の
正則 bundle 座標，canonical primitive，固定 end における moment の比較を与える．
固定 compact と end を先に取り，その後 $t\to0$ とする順序を用いる．
Moving な全軌道を含む actual flow buffer，source の全 moment 接続，
複数 end の同時接着は，本稿の
節\ref{ke:staged-neck-section}，
\ref{ke:neck-corner-exhaustion-section}および
定理\ref{ke:source-all-moment-exit}で追加して証明する．
\texttt{pair:lem:diagonal-fixed-end} の遅い対角選択は，
これらの評価を証明した後に使う原理である．

\item \emph{固定成分上の smooth convergence と SNC cusp 評価．}
Song\cite[Theorem~1.2]{SongNegativeKE}による成分極限の同定と
fixed-gauge convergence は，上記
\texttt{pair:thm:fixed-end-moment} の証明で対象族へ適用される．
成分は $\PP^2$ から $d-1$ 本の直線を除いたもので，
log canonical bundle は $\mathcal O_{\PP^2}(d-4)$ である．
Rochon--Zhang\cite[\S5, (5.3)--(5.6), Theorem~5.1]{RochonZhang}の
bounded cusp derivatives を持つ potential を用い，
本稿の節\ref{ke:component-core-geometry-section}と
\ref{ke:component-core-transfer-section}で正の residue，
固定 compact の正則化と小さい actual 移送を得る．
Divisor が smooth な場合だけの polyhomogeneous expansion は使わず，
$t$ とともに動く任意の compact 上での一様収束も仮定しない．
\end{enumerate}

一般の Moser 同定が Lagrangian 性を運ぶことと，その移送が KE 距離を保つことは別である．
本稿では下界と上界を actual $g_t$ に対して直接証明する．
\cite{KECubicalCurrent}の
$\mathrm{(G)}_d+\mathrm{(CP)}_d+\mathrm{(PBE)}_d$ による条件付き GH 定理，
又は $\mathrm{(TC)}_5$ を主定理の入力にはしない．
既存の近似写像 $\Phi_t$ との一様な値比較から
$f_t=p_T\pi_t$ へ距離歪みを移す命題\ref{topo:ke-map-compatibility}は，
別の有効な十分条件として保持する．
本稿の直接証明を，これらの外部比較仮説自体の検証へ読み替えることはできない．

\subsection{同じ射影を保つ六段階の証明}
\label{ke:proof-route}

構成を確認した後に，同じ射影の距離評価を追うための経路を以下にまとめる．
局所 GH 評価は第3段階，全曲面の ambient 距離比較は第5段階である．

\begin{enumerate}
\item \emph{Neck と closed corner の全葉を一致させる．}
系\ref{ke:neck-corner-saturated-union}は，
共通の向き付き normal action と canonical moment を保ち，
全ての実際の overlap で reduced leaves と whole fibers を一致させる．
補題\ref{ke:fixed-marked-tree-exhaustion}により，
実際の marked pants 分解から得る木系を全段階と全ての小さい複素 $t$ で固定する．

\item \emph{Actual core を埋め，全曲面の fibration を作る．}
定理\ref{ke:actual-global-fibration}は，closed moment links，
primitive の相対不変化と nontrapping，saddle を動かす葉族補間，
補間後の全 fiber 単位の trim を用いる．
有限な closed whole-fiber quotient の Hausdorff 性を先に確保し，
全 $X_t$ の被覆，連結 fibers，actual Lagrangian 性と特異模型を得る．
Source の pinched tori と保持 corner の葉はこの相対構成で保たれる．

\item \emph{後続の接続に使う局所費用を評価する．}
定理\ref{ke:source-all-moment-exit}は，source の全共通 moment で，
同じ canonical moment の残余境界まで actual 長さ $O(1)$ の退出を与える．
Source 軸を含む arm tensor 比較と中央点からの短い直接退出による証明なので，
Gaussian 可積分性を必要としない．
定理\ref{ke:slow-cut-closed-corner-gh}で closed corner の距離を評価し，
命題\ref{ke:component-logarithmic-connectors}で各 component の
core・変更 collar・neck 尾部の費用を $C(1+\log(N_j+1))$ に抑える．

\item \emph{一つの基底 parameter 化を選ぶ．}
命題\ref{ke:global-base-parameterization}により，第2段階の同じ fiber partition を
$B_d(T)$ と同定し，$f_{j,t}=p_T\pi_{j,t}$ とする．
保持 corner の像，neck strips，component sectors をこの写像で比較し，
内部木の全枝にも一つの canonical moment を用いる．

\item \emph{この同じ写像の全域下界と上界を証明する．}
下界は定理\ref{ke:sharp-global-one-active}の source を除かない scalar 障壁から，
系\ref{ke:global-depth-lipschitz}の global clipped-depth map，
命題\ref{ke:global-depth-base-comparison}の一様な基底値比較を経て，
定理\ref{ke:global-ambient-lower}へ至る．
長さ評価は全曲面の任意の経路に成立し，領域の出入りの回数に依存しない．
上界は補題\ref{ke:coherent-corner-portals}の共通 moment の接続を
高々 $3\binom{\min\{d,6\}}3\le60$ 本の square gallery に沿って結ぶ
定理\ref{ke:global-ambient-upper}で得る．
下界の値比較で用いる接続は第3段階で証明してあり，
全域距離上界を先取りしない．

\item \emph{有限段階の誤差を一つの対角選択にまとめる．}
定理\ref{ke:global-metric-realization}は，
式\eqref{ke:eq:global-realization-explicit-cut-budget}の両側評価に
共通の遅い選択 $j=j(t)$ を適用し，
\eqref{ke:eq:global-realization-distortion}を結論する．
各 torus だけでなく，$p_T$ が潰す木上の全 fiber 合併を同じ評価に含める．
\end{enumerate}

\subsection{切断尺度と退化 parameter の量化順序}
\label{ke:cut-selection-order}

固定族の下で，各有限段階 $j$ のデータを次の順に選ぶ．
これは任意に先指定した moving 領域への一様評価ではなく，
既に証明した局所評価を同じ全域構成で用いるための選択である．
\begin{enumerate}
\item $b_j\to\infty$ と $\varepsilon_j\downarrow0$ を取り，
旧 corner cut $m_j=6b_j$，固定 residual compact，既存の normal cut $n_j$，
共通 moment 幅 $\iota_j$ を順に選ぶ．
固定 residual domain の接続定数 $C_j^{\rm neck}$ は，次の $N_j$ より先に固定する．
\item Component cut $N_j$ を，既存の下限，
$N_j\iota_j\gg1$，および
\[
 C_j^{\rm neck}\le1+\log(N_j+1)
\]
を満たすように選ぶ．これにより core と変更 collar を作り，
neck と component の全接続に一つの対数的予算を使える．
\item Scalar 障壁専用の補助 cut $\mathfrak n_j$ と境界定数を先に決め，
その後に lower clipping depth $\mathfrak m_j$ と upper factor $K_j$ を選ぶ．
$\mathfrak m_j$ は固定段階の core の深度上界も支配させる．
$\mathfrak n_j$ と既存の $n_j$，$\mathfrak m_j$ と旧 corner cut $m_j$ は別の量である．
正確な下限は定理\ref{ke:sharp-global-one-active}と
命題\ref{ke:global-depth-base-comparison}に置く．
\item 最後に，全ての構成条件と誤差評価に加え，
\[
 R_t\ge j\{1+\log\mathfrak m_j+\log K_j+\log(N_j+1)\}
\]
を保証する閾値 $r_j>0$ を選ぶ．
$r_{j+1}<r_j/2$ として，$j(t)=\max\{j:|t|<r_j\}$ と置く．
\end{enumerate}
残余定数の増大率や固定 compact の収束閾値は有効に評価していないので，
この量化順序は普遍的な収束速度を与えない．また，各 $\pi_t$ の連続性と
$t$ 方向の連続又は滑らかな選択は区別する．

Source の marking には命題\ref{ke:source-marking}の条件が必要である．
指定 nodal fiber が actual source core と交わり，候補近似 $\Phi_t$ が raw map と
一様に $o(1)$ で近いとする．このとき
$\sup_{X_t}d_\square(\Phi_t,p_T\pi_t)\to0$ を得るには，
その nodal value $b_{ij,t}$ が $p_T(b_{ij,t})\to e_i+e_j$ を満たす必要がある．
系\ref{ke:neck-base-value-control}の構成はこの位置条件を満たす．
任意の先指定 critical values 又は任意の第二 periods の実現を，
以上の六段階に追加して主張するものではない．

\input{sections/paper/degeneration_polarization.tex}
\section[SNC 因子と connection holonomy による有限相対 compactification]{SNC 因子と connection holonomy による\\有限相対 compactification}
\label{snc:section}
% 2026-11-21 / v1.2: 固定 connection の全 leaf-period 条件と最小 holonomy 自由度．

直線配置に対する構成の本質は，個々の成分が ample であることではない．
必要なのは，Kähler class が divisor 全体の正の重み付き和で表されることである．
有理曲線の場合の graph periods に加え，正の genus を持つ成分上の periods は
normal connection の holonomy を選ぶことで解消できる．
本節では，任意の genus，負の自己交点数，同じ二成分の複数交点，
および二価の成分を含む形でこの構成を証明する．

\subsection{仮定，marking，結論}

$M$ を連結な滑らかな射影複素曲面，$D=\bigcup_{a=1}^{r}D_a$ を
reduced simple normal crossings divisor とする．
各 $D_a$ は genus $g_a$ の滑らかな曲線とし，三重交点はない．
異なる二成分が複数の点で交わることを許す．
各交点を個別にラベル付けし，その集合を $E$ とする．
$e=|E|$，$n_a=|D_a\cap\bigcup_{b\ne a}D_b|$ と置く．
$n_a$ は隣接する成分の数ではなく，\emph{交点の数}である．
以下では
\begin{equation}\label{snc:valence}
 n_a\geq2\qquad(1\leq a\leq r)
\end{equation}
を仮定する．$G$ は頂点が $D_a$，辺が個々の交点である dual multigraph とする．

Kähler form $\omega$ と正の実数 $w_a$ を，
\begin{equation}\label{snc:positive-class}
 [\omega]=2\pi\sum_{a=1}^{r}w_a[D_a],\qquad w_a>0
\end{equation}
となるように固定する．重みは存在だけでなくデータの一部として固定する．
normal-circle の周期は $2\pi$，Hamiltonian convention は
$\iota_V\omega=-\dd\mu$，normal moment は $\mu\geq0$ とする．
angular connection の curvature convention は
\begin{equation}\label{snc:curvature-convention}
 \int_{D_a}\dd\alpha_a=-2\pi D_a^2
\end{equation}
である．

各 $D_a$ から交点の小円板を除いた genus $g_a$，$n_a$ 境界成分の曲面を $P_a$ とする．
$P_a$ 上に，境界 collar で radial circles を葉とし，
内部には $2g_a+n_a-2$ 個の simple saddles だけを持つ実数値 Morse function の
connected-level-set 葉層 $\mathcal F_a$ を指定する．
各 connected critical leaf は一つの saddle を含む．
$n_a\geq2$ なら，このような Morse function は存在する．
例えば genus 零の cobordism に split--merge saddle pair を $g_a$ 組挿入すればよい．
$g_a=0,n_a=2$ の場合は annulus 上の正則な円周葉層である．
marking は，交点のラベルと，実際の $P_a$ 上の葉層の相対 isotopy class を含む．
その Reeb graph を $T_a$ と書く．$b_1(T_a)=g_a$ であり，有理成分の場合に限り tree である．
各交点 $p\in D_a\cap D_b$ に対応する $T_a,T_b$ の葉を同一視し，
そこでできる二価頂点を必要に応じて除いた graph を $\Gamma$ とする．
全成分が有理で二価なら，$\Gamma$ は円周である．

\begin{theorem}[SNC 因子の有限相対構成]
\label{snc:main}
仮定 \eqref{snc:valence}--\eqref{snc:positive-class} の下で，
$\omega$ と同じ cohomology class の Kähler form $\omega'$ と，
連結 fiber を持つ proper continuous surjection
\[
 \pi:(M,\omega')\longrightarrow\operatorname{Cone}_{\rm cl}(\Gamma)
\]
を次の性質を持つように構成できる．
\begin{enumerate}
\item cone point の fiber は $W=M\setminus D$ の，適切に選んだ
  logarithmic Weinstein skeleton $\Sigma$ である．
\item $W\setminus\Sigma$ の正則 fiber は Lagrangian $\TT^2$ であり，
  $\Gamma$ の三価頂点から延びる開放射辺上の fiber は
  $S^1\times\figeight$ である．他の内部 singular-fiber locus はない．
\item outer graph $\Gamma$ の逆像は正確に $D$ である．
  $D_a$ 上の connected fibers は指定した $\mathcal F_a$ の葉と一致し，
  crossing collars では radial circles，交点では一点となる．
\item 各交点の十分小さい holomorphic charts で，両 divisor を座標軸とする
  正の diagonal product Kähler germs をあらかじめ指定できる．
  collar の normal-circle action germs はそれらの product charts と一致する．
  各成分の normal connection の holonomy は構成中に選択する．
  collar の葉層は，その選択後の normal charts を用いて，固定した $D_a$ 上の
  $\mathcal F_a$ を lift したものである．
\item 各 $D_a$ の有限個の指定点が $\mathcal F_a$ の互いに異なる正則葉上にあり，
  crossing collars の外にあるなら，それらの周囲に divisor-compatible な
  holomorphic flat Kähler balls を同時に用意できる．
  これらの点は regular elliptic-boundary patches に含まれる．
\item $\Sigma$ を選んだ後の end の変更は，その近傍の Liouville structure を変えず，
  新しい trapping を作らない．全 singular fibers は，Evans--Mauri の意味で
  isotropic strata を持つ．特定の skeleton や Whitney 正則性は指定しない．
\end{enumerate}
個々の $D_a$ の ampleness，$D_a^2$ の正値性，
$H^1(M;\RR)=0$，または $Q=(D_a\cdot D_b)$ の非退化性を追加仮定する必要はない．
\end{theorem}

\begin{corollary}[多直線配置の marked な有限 compactification]
\label{comp:finite}
$M=\PP^2$，$D=L_1\cup\cdots\cup L_m$ を一般位置の $m\geq4$ 本の直線の和とし，
$c>0$ と各 $L_a$ 上の marked Morse 葉層を固定する．
ある Kähler form $\omega'\in cH$ に対して，定理\ref{snc:main}の構成が得られる．
各成分の Reeb graph は $m-1$ 個の labelled leaves と $m-3$ 個の三価頂点を持つ
tree であり，$\Gamma$ は対応する葉を交点ごとに貼り合わせた graph である．
crossing germs，full saddle collars，および有限個の指定された flat patches も
同時に保持できる．
\end{corollary}
\begin{proof}
$2\pi\sum_a w_a=c$ となる正の weights を選べば
$[\omega]=2\pi\sum_a w_a[L_a]=cH$ である．
全 $g_a=0$，$n_a=m-1\geq3$ なので，定理を適用できる．
\end{proof}

形式を $\omega'$ に替える点には実質的な意味がある．
固定した $\omega$ と固定した $D$ に対して，交点が非直交であるまま
commuting normal-circle actions を要求することはできない．
この限界と，primitive の重みを変更する場合の限界を後で明示する．

\begin{lemma}[連結性と affine 性]
\label{snc:affine}
\eqref{snc:positive-class} から $D$ は連結であり，$W$ は affine である．
\end{lemma}
\begin{proof}
$D$ が二つの非空な disjoint unions に分かれるなら，
対応する正の実 divisor を $A,B$ として $A\cdot B=0$ となる．
$[\omega]=2\pi(A+B)$ と Kähler positivity より
$A^2=[\omega]\cdot A/(2\pi)>0$，$B^2>0$ である．
これは Hodge index theorem に反する．従って $D$ は連結である．

Kähler cone は開いているので，$w$ を十分近い正の有理ベクトル $v$ に替えても
$2\pi\sum v_a[D_a]$ は Kähler class である．
分母を払うと，同じ support $D$ を持つ effective integral divisor
$A_D$ で，$\mathcal O_M(A_D)$ が正の curvature metric を持つものを得る．
従って $A_D$ は ample である．十分大きい倍を very ample とし，
対応する section の零集合を hyperplane section として実現すると，
その補集合 $W$ は affine となる．
\end{proof}

\subsection{重みを固定した global logarithmic primitive}

各 $\mathcal O_M(D_a)$ に任意の smooth Hermitian metric を取り，
canonical section を $s_a$ とする．
$-\dd\dd^c\log\|s_a\|^2$ は $W$ 上の式から $M$ 上の smooth closed
$(1,1)$-form $\Theta_a$ に延び，$[\Theta_a]=2\pi[D_a]$ である．
個々の $\Theta_a$ が正である必要はない．
compact Kähler manifold の $\dd\dd^c$-lemma より，smooth real function $u$ を用いて
\[
 \omega=\sum_a w_a\Theta_a+\dd\dd^c u
\]
と書ける\cite[Lemma (5.11), pp.~266--267; Remark (5.14)]{DGMS}．従って
\begin{equation}\label{snc:log-potential}
 \psi_0=-\sum_a w_a\log\|s_a\|^2+u,\qquad
 \lambda_0=\dd^c\psi_0,\qquad
 \dd\lambda_0=\omega|_W
\end{equation}
である．重みが全て正なので $\psi_0$ は exhaustion であり，
$\dd\dd^c\psi_0=\omega>0$ だから strictly plurisubharmonic である．
$D_a$ の正の meridian $\mathfrak m_a$ に対して
\begin{equation}\label{snc:residues}
 \lim_{\rho\to0}\int_{\mathfrak m_a(\rho)}\lambda_0=-2\pi w_a.
\end{equation}

SNC chart では $\psi_0$ は $-\sum w_a\log|z_a|^2$ と smooth function の和である．
そのため critical locus は compact であり，compact subset of $W$ 内の
小さい potential perturbation で Morse にできる．
この変更は $M$ 上の Kähler class を変えず，divisor の近傍にも影響しない．
全 critical points を含む固定 regular sublevel を compact Weinstein domain として取り，
補題\ref{ke:compact-relative-morse-smale}を適用する．これは
定理\ref{snc:main}に依存しない compact Weinstein domain の補題であり，
この時点の二形式と境界 collar の primitive を変えずに，内部に支持された
exact correction $dQ$ によって Morse--Smale にできる．この補正を外側で零に延ばし，
以後 $\lambda_0$ にはこの $dQ$ を含める．End の logarithmic primitive と全 periods
は変わらず，同補題により skeleton は次元順の閉 isotropic filtration を持つ．
その logarithmic Weinstein skeleton を $\Sigma$ とする．
$\Sigma$ は $W$ の homotopy type を持つ connected compact set であり，
この閉 filtration は Evans--Mauri\cite[Definitions~2.1 and 2.5]{EM}の意味での
isotropic strata を与える．Whitney 条件は用いない．

次に，全 critical values より高い end の中で，交点の Kähler form を product にする．
同時に，指定された divisor 上の有限個の正則点を中心とする互いに素な charts を取り，
divisor が一つの座標軸となる holomorphic flat balls を作る．
holomorphic chart の potential を $\phi=Q+O(r^3)$ と書き，
所望の正の diagonal quadratic form $Q'$ に対して
\[
 \phi'=(1-\chi(\log r))\phi+\chi(\log r)Q'
\]
を用いる．transition を長い logarithmic annulus に取ると
$h(r)=\chi(\log r)$ に対する $r|h'(r)|+r^2|h''(r)|$ を小さくでき，
$Q'-\phi=O(r^2)$ なので cutoff error は正の二形式の凸結合に対して任意に小さくなる．
従って Kähler positivity を保ち，cohomology class も変えない．
補正の微分は $O(r)$ であり，logarithmic normal differential は end で発散するので，
chart を縮めれば新しい critical point は生じない．
変更領域の potential value を元の全 critical values より高く保つため，
$\Sigma$ も変わらない．
以下ではこれまでの変更後の form と primitive を再び $\omega',\lambda_0$ と書く．
この変更は指定した $D_a$ 上の foliation を動かさない．
正則葉上の点では，normal-circle action と葉の局所 transverse coordinate が
regular elliptic-boundary model を与える．
ambient form が holomorphic flat であることと，
fibration を symplectic coordinates でその標準形にすることは区別して用いる．

crossing product charts を固定する relative symplectic tubular charts を選ぶ．
normal first derivative は complex-linear に取る．
normal coordinate の変換は $z\mapsto a z+O(|z|^2)$，$a\neq0$ であるから，
angular pole の差を引いた remainder は有界である．
交点では二つの charts が初めから product として一致する．
この regularization と以下の connection の変更は，
標準的な SNC symplectic neighborhoods の枠組みに属する
\cite[\S2.4, Appendix A]{EM}\cite{TMZ}．
ただし\cite[Theorem 2.13]{TMZ}は一般には symplectic form の変形を許す．
ここで固定した $\omega'$ に対する charts は，既に product とした crossing germs に
相対的な tubular Moser により得る．実際，normal derivative を complex-linear かつ
symplectic に選べば，tubular pullback と coupling form は zero section 上で二形式として一致する．
その差の radial homotopy primitive は二次以上で消え，固定 product charts では零に取れる．
十分薄い近傍で両形式の線分は非退化なので，Moser field は zero section と normal first jet，
および crossing charts を固定する．従ってこの段階で actual $\omega'$ をさらに変形しない．

\subsection{任意の自己交点数に対する curvature balance}

$Q_{aa}=D_a^2$，$Q_{ab}=|D_a\cap D_b|$ $(a\neq b)$ とし，
$A_a=\int_{D_a}\omega'$ とする．仮定から
\begin{equation}\label{snc:area-balance}
 A_a=2\pi\sum_b Q_{ab}w_b
 =2\pi\left(w_aD_a^2+\sum_{p\in D_a\cap D_b}w_b\right).
\end{equation}
最後の和は全ての交点 $p$ について取り，同じ $b$ が複数回現れ得る．

各 crossing cap $p\in D_a\cap D_b$ で
$\omega'|_{D_a}=\dd I\wedge\dd\theta_b$，$p=\{I=0\}$ と書く．
$P_a$ 全体には
\[
 \kappa_a=-\omega'|_{D_a}/w_a
\]
を指定する．cap では normal angular frame $\theta_a$ を用い，
\begin{equation}\label{snc:connection-cap}
 \alpha_a=\dd\theta_a+A_{a,p}(I)\dd\theta_b,\qquad
 A_{a,p}=0\text{ near }I=0,\qquad
 A_{a,p}=(w_b-I)/w_a\text{ near }P_a
\end{equation}
と指定する．この curvature の全積分は
\begin{equation}\label{snc:curvature-integral}
 \int_{D_a}\kappa_a
 =\frac{2\pi\sum_{p\in D_a\cap D_b}w_b-A_a}{w_a}
 =-2\pi D_a^2.
\end{equation}
従って，これは $N_{D_a/M}$ の connection として実現できる．
$D_a^2\leq0$ でも，cap 内の curvature が必要な分を補償するため問題はない．

$P_a$ 上で $\beta_a=0$，cap では
\[
 \beta_a=(I-w_b+w_aA_{a,p}(I))\dd\theta_b
\]
とする．$P_a$ との overlap では $\beta_a=0$ であり，
$\dd\beta_a=\omega'|_{D_a}+w_a\kappa_a$ である．
normal bundle の coupling model は
\begin{equation}\label{snc:coupling}
 \Omega_a=\omega'|_{D_a}+\dd(\mu_a\alpha_a),\qquad
 \lambda_a=(\mu_a-w_a)\alpha_a+\beta_a.
\end{equation}
$\dd\lambda_a=\Omega_a$ であり，十分薄い collar では $\Omega_a$ は symplectic である．
交点では両側の式が
\[
 (\mu_a-w_a)\dd\theta_a+(\mu_b-w_b)\dd\theta_b
\]
に一致するので，これらの primitive は一つの
$\lambda_{\rm col}$ に貼り合わさる．

ここで connection を選ぶ自由は，実際の $\omega'$ の collar を失わせない．
二つの connections の差を $\pi^*\eta_a$ と書き，
$\alpha_{a,s}=\alpha_{a,0}+s\pi^*\eta_a$ とする．
対応する coupling forms の相対 Moser field を
\[
 \iota_{Y_s}\Omega_{a,s}=-\mu_a\pi^*\eta_a
\]
で定める．crossing product charts では $\eta_a=0$ と取れる．
$Y_s$ は zero section とそこでの first jet を固定し，
normal-circle に不変で，$Y_s\mu_a=0$ である．
よって縮めた collar 上で，connection の変更を equivariant，moment-preserving な
symplectomorphism に吸収できる．以下の式はこの同定後の実際の collar 上の式とみなす．

\begin{lemma}[full Morse 葉層を保つ Liouville field]
\label{snc:collar-dynamics}
$\iota_Z\omega'=\lambda_{\rm col}$ とする．$P_a$ 上では
\[
 Z\mu_a=\mu_a-w_a,\qquad Z^{\rm horizontal}=0.
\]
cap では $A=A_{a,p}$ と略記すると
\begin{equation}\label{snc:cap-dynamics}
 Z\mu=\mu-w_a,\qquad
 ZI=\frac{I-w_b+w_aA(I)}{1+\mu A'(I)}.
\end{equation}
従って normal circles と $\mathcal F_a$ の lift からなる collar 葉層を，
saddle fiber 全体も含めて $Z$ の流れが保つ．
\end{lemma}
\begin{proof}
$P_a$ 上では
$\Omega_a=(1-\mu_a/w_a)\omega'|_{D_a}+\dd\mu_a\wedge\alpha_a$
である．よって $Z=(\mu_a-w_a)\partial_{\mu_a}$ である．
cap では
\[
 \lambda_{\rm col}=(\mu-w_a)\dd\theta_a+
     (I-w_b+\mu A(I))\dd\theta_b
\]
を \eqref{snc:coupling} に代入して \eqref{snc:cap-dynamics} を得る．
normal radius を小さくすれば $1+\mu A'(I)>0$ である．
cap の radial circle leaves と $P_a$ 上の全 connected critical leaves が保存される．
$g_a=0,n_a=2$ では $P_a$ は annulus であり，同じ計算がその正則円周葉層を保つ．
\end{proof}

\subsection{plumbing，component holonomy，graph periods}

$N$ を $D$ の小さい plumbing neighborhood，$S=\partial N$ とする．
$N$ は $D$ に retract する．曲面を交点で貼り合わせる Mayer--Vietoris 計算により
\[
 H_2(N;\ZZ)\simeq\ZZ^r,\qquad H_3(N;\ZZ)=0,
\]
かつ
\begin{equation}\label{snc:divisor-homology}
 0\longrightarrow\bigoplus_aH_1(D_a;\ZZ)
 \longrightarrow H_1(D;\ZZ)
 \longrightarrow H_1(G;\ZZ)\longrightarrow0
\end{equation}
を得る．右端は free なのでこの列は非標準的に分裂し，$H_1(N;\ZZ)$ も free である．
複数交点は $G$ の別々の辺として数える．

\begin{proposition}[meridian と divisor の完全列]
\label{snc:integer-sequence}
整数係数で
\begin{equation}\label{snc:integer-exact}
 0\longrightarrow\operatorname{coker}(Q:\ZZ^r\to\ZZ^r)
 \longrightarrow H_1(S;\ZZ)
 \longrightarrow H_1(D;\ZZ)\longrightarrow0
\end{equation}
が成り立つ．左側の subgroup は normal meridians で生成され，
その関係式は $Q$ の行で与えられる．
特に $Q$ が特異でもよく，cokernel の torsion も保持される．
有理成分だけの場合には $H_1(D;\ZZ)=H_1(G;\ZZ)$ となる．
\end{proposition}
\begin{proof}
Poincaré--Lefschetz duality と universal coefficient theorem より
$H_2(N,S;\ZZ)\simeq H^2(N;\ZZ)\simeq\ZZ^r$，
$H_1(N,S;\ZZ)\simeq H^3(N;\ZZ)=0$ である．
$H_2(N)\to H_2(N,S)$ は intersection pairing，すなわち行列 $Q$ である．
対の完全列の該当部分の cokernel を取れば \eqref{snc:integer-exact} を得る．
normal disk の境界が meridian なので，生成元と関係式の説明も従う．
\end{proof}

各 $T_a$ は $n_a$ 個の一価頂点と $2g_a+n_a-2$ 個の三価頂点を持つ．
従って $b_1(T_a)=g_a$ であり，これらを交点ごとに貼り合わせると
\begin{equation}\label{snc:graph-counts}
 b_1(\Gamma)=\sum_a g_a+e-r+1,\qquad
 |V_3(\Gamma)|=\sum_a(2g_a+n_a-2)=2\sum_a g_a+2e-2r.
\end{equation}
$D\to\Gamma$ は $H_1$ の全射を誘導するが，正の genus の成分があると
一般には同型でない．component 上の $2g_a$ 個の periods 全てを消すには，
graph から引き戻した basic forms だけでは不足し得る．
この不足は次の holonomy の自由度で解消する．

\begin{lemma}[Reeb 商で消える正則葉の cycles]
\label{snc:reeb-leaf-kernel}
Cap へ radial に延長した Reeb graph を，辺を延ばす前の graph と同じ位相型で同定し，
商写像を $q_a:D_a\to T_a$，$q:D\to\Gamma$ とする．
整数格子
\[
 L_a=\ker\bigl((q_a)_*:H_1(D_a;\ZZ)\longrightarrow H_1(T_a;\ZZ)\bigr)
\]
は rank $g_a$ の primitive Lagrangian sublattice である．
$T_a$ の maximal subtree の外の $g_a$ 本の辺から正則葉を一つずつ選ぶと，
その向き付き円周 $\ell_{a,1},\ldots,\ell_{a,g_a}$ が $L_a$ の基底になる．
また，次は完全列である：
\begin{equation}\label{snc:reeb-homology-sequence}
 0\longrightarrow\bigoplus_aL_a\longrightarrow H_1(D;\ZZ)
 \xrightarrow{q_*}H_1(\Gamma;\ZZ)\longrightarrow0.
\end{equation}
従って実係数では
\begin{equation}\label{snc:leaf-obstruction-space}
 \mathcal O_{\mathcal F}(D):=
 H^1(D;\RR)/q^*H^1(\Gamma;\RR)
 \cong\bigoplus_a\operatorname{Hom}(L_a,\RR),
 \qquad \dim\mathcal O_{\mathcal F}(D)=\sum_a g_a.
\end{equation}
$q^*H^1(\Gamma;\RR)$ の各 class は，regular annuli 内に支持され，
全 critical leaves と crossing points の近傍で零の smooth closed basic form で表せる．
\end{lemma}

\begin{proof}
閉曲面の simple Morse function に対する，level periods の消滅と smooth basic
representative の対応は既知である\cite[Lemma~4.8, p.~2421; \S5, pp.~2424--2425]{IKM}．
ここでは本節の marked pants decomposition に対し，整数格子と SNC 接続を直接確認する．
選んだ $g_a$ 本の正則円周で $D_a$ を切ると，pants と cap disks を貼る graph は tree
になるため，切断曲面は連結で genus 零である．切断後の各円周の両岸を結ぶ arcs を
戻せば，これらの円周に intersection pairing で双対な cycles が得られる．
従って円周の classes は独立で primitive であり，互いに素なので交点数は全て零である．
Rank は $g_a$ なので Lagrangian sublattice を張る．
一方，graph の各 loop は pants 内の arcs で閉曲面上へ lift できるから $(q_a)_*$ は全射である．
$b_1(T_a)=g_a$ と rank を比較すれば，選んだ格子がその kernel 全体である．

$H_1(D)$ の列\eqref{snc:divisor-homology}と，$T_a$ を交点ごとに貼る graph の列
\[
 0\longrightarrow\bigoplus_aH_1(T_a;\ZZ)
 \longrightarrow H_1(\Gamma;\ZZ)\longrightarrow H_1(G;\ZZ)\longrightarrow0
\]
を比較する．右端の写像は恒等であり，左端は全射なので，kernel の列から
\eqref{snc:reeb-homology-sequence}を得る．実係数で双対化すれば
\eqref{snc:leaf-obstruction-space}が従う．

最後に，$\Gamma$ 上の cellular $1$-cochain で class を表す．各辺の内部に短い interval
を取り，その cochain の値を積分に持つ bump form $g(u)\dd u$ を置く．
全 vertices と指定された有限個の regular values を避けて支持を選べる．
葉の商へ引き戻し，支持の外で零に延ばせば，crossings と全 critical leaves の近傍でも
滑らかな closed basic form が得られる．同じ式を normal direction に一定に延ばせば，
plumbing collar 上でも全 fibers の tangent spaces を消す closed form となる．
\end{proof}

\begin{lemma}[connection holonomy による periods の制御]
\label{snc:holonomy}
$\eta_a$ を $D_a$ 上の closed real one-form で，全 crossing disks の近傍で $0$ とする．
connection を $\alpha_{a,s}=\alpha_a+s\pi^*\eta_a$ に替え，
$\beta_a$ は変えずに $\Omega_{a,s},\lambda_{a,s}$ を \eqref{snc:coupling} で定める．
十分薄い collar 上に，crossing germs と zero section を固定する Moser isotopy
$\Phi_s$ が存在し，$\Phi_s^*\Omega_{a,s}=\Omega_{a,0}$，$\mu_a\circ\Phi_s=\mu_a$ である．
その primitive の class の変化は
\begin{equation}\label{snc:holonomy-variation}
 [\Phi_s^*\lambda_{a,s}-\lambda_{a,0}]
     =-s w_a[\pi^*\eta_a].
\end{equation}
この操作は $D_a$ 上の marked foliation を固定し，新しい normal charts 上に
その同じ foliation を lift した collar fibration を保つ．
\end{lemma}
\begin{proof}
$\eta_a$ が closed なので curvature と $\beta_a$ は変わらない．
$\iota_{Y_s}\Omega_{a,s}=-\mu_a\pi^*\eta_a$ と置くと，
$Y_s\mu_a=0$，$\alpha_{a,s}(Y_s)=0$ であり，horizontal part は $O(\mu_a)$ である．
従って十分薄い collar では時間 $0\leq s\leq1$ の flow が存在し，
zero section，そこでの normal first derivative，および crossing germs を固定する．
Cartan の公式から
\begin{align*}
 \frac{\dd}{\dd s}\Phi_s^*\lambda_{a,s}
 &=\Phi_s^*\left((\mu_a-w_a)\pi^*\eta_a
                +\iota_{Y_s}\Omega_{a,s}
                +\dd(\lambda_{a,s}(Y_s))\right)\\
 &=-w_a\Phi_s^*(\pi^*\eta_a)
      +\dd\bigl(\Phi_s^*(\lambda_{a,s}(Y_s))\bigr).
\end{align*}
$\Phi_s$ は identity と isotopic なので，cohomology を取って積分すると
\eqref{snc:holonomy-variation} を得る．
全 crossing disks で identity なので，他成分の model とそのまま貼り合わさる．
zero section 上も identity であり，$P_a$ 上では引き続き
$\kappa_a=-\omega'|_{D_a}/w_a$，$\beta_a=0$ である．
従って新しい connection の coupling charts では
$Z=(\mu_a-w_a)\partial_{\mu_a}$ であり，fixed marked foliation の lift を保つ．
この chart を $\Phi_1$ で actual collar に引き戻せばよい．
$\Omega_{a,s}$ は補助的な normal-bundle model であり，この引き戻し後の形式は
最初の $\Omega_{a,0}$ である．従って actual manifold 上の Kähler form $\omega'$ 自体は
この holonomy 補正では変更しない．
\end{proof}

$H^1(D_a;\RR)$ の任意の class は crossing disks と，
指定した flat points を中心とする有限個の小円板の近傍で $0$ の closed form で表せる．
実際，disks 上での局所 primitive を cutoff して引けばよい．
これを他成分側で $0$ として plumbing に延長すると，
\eqref{snc:divisor-homology} の component part を任意に制御できる．
大きい holonomy の補正が必要な場合でも，その補正を選んだ後に collar を縮めれば，
Moser isotopy の存在に支障はない．

$\lambda_{\rm col}$ と $\lambda_0$ は同じ $\omega'$ の primitive であり，
同じ重み $w_a$ を用いたので meridian residues が一致している．
その差の meridian integral は radius に依存せず，radius $0$ の極限では $0$ になる．
従って全 meridian periods は正確に $0$ である．
\eqref{snc:integer-exact} を実数係数で双対化すると，差の class は
$H^1(D;\RR)$ に由来する．torsion は real closed one-form の period では検出されない．
まず補題\ref{snc:holonomy}で各 $D_a$ 上の全 periods を消す．
その後に残る class は \eqref{snc:divisor-homology} により $H^1(G;\RR)$ から来る．
以下では，この holonomy 補正後の実際の collar primitive を再び
$\lambda_{\rm col}$ と書く．

$G$ の spanning tree の外の各辺に一つの regular cap arm を対応させ，
そこで支持された basic form $g(I)\dd I$ の積分を選ぶ．
これら $e-r+1$ 個の forms は残る graph periods を任意に補正できる．
その和を $\eta$ とすると
\begin{equation}\label{snc:exact-correction}
 \lambda_{\rm col}-\eta-\lambda_0=\dd f
\end{equation}
となる．$\eta$ は全 collar fibers の tangent spaces を消す．
その補正は Liouville field の angular components だけを変え，
\eqref{snc:cap-dynamics} は変えない．
各 cap arm の端では $\eta=0$ とし，normal direction に同じ式を延長するため，
$\eta$ は end 全体の smooth closed form である．

normal first derivatives を complex-linear に選んだことから，
両 primitive の angular poles を差し引いた remainder は有界である．
holonomy の Moser isotopies も normal first derivatives を固定するので，
この有界性を保つ．従って \eqref{snc:exact-correction} の $\dd f$ は有界であり，
$f$ の定数を一つ固定すれば $f$ も有界である．
後者は，finite plumbing cover 内で base point から各点まで
一様に有限な長さの経路を選び，既に全 period が消えていることから従う．
この議論で $H^1(M;\RR)$ の消滅は使わない．

\subsection{core を保つ cutoff}

十分高い level で $0$ から $1$ へ移る cutoff $\chi(\psi_0)$ を取り，
\[
 \lambda=\lambda_0+\dd\bigl(\chi(\psi_0)f\bigr)
\]
とする．$\dd\lambda=\omega'$ であり，core の近傍では $\lambda_0$，
十分 end 側では $\lambda_{\rm col}-\eta$ と一致する．
$\iota_{X_H}\omega'=-\dd H$ とすると
\begin{equation}\label{snc:cutoff-estimate}
 Z\psi_0=Z_0\psi_0-\chi(\psi_0)X_f\psi_0.
\end{equation}
$\chi'$ を含む項は $X_{\psi_0}\psi_0=0$ により消える．
normal distance を $\rho_a$ と書けば，SNC charts と $\omega'$ の正値性から
\[
 Z_0\psi_0\geq C\sum_a\rho_a^{-2},\qquad
 |X_f\psi_0|\leq C'\left(1+\sum_a\rho_a^{-1}\right)
\]
となる．よって cutoff を十分 end 側に置くと，変更領域の全てで $Z\psi_0>0$ となる．
新しい zero，trapping，または低い critical region への正時間の戻りは生じない．
したがって nonescaping set は正確に元の $\Sigma$ のままである．

adapted Liouville primitive の一般的存在は
\cite[Proposition 2.10]{BSV}にも与えられている．
ここで加えたのは，指定した full Morse 葉層を保つ curvature と holonomy の選択，
component および graph の全 periods の補正，core の germ を固定した cutoff，
および後述の finite endpoint の制御である．

\subsection{actual contact link と finite endpoint}

$P_a$ と transition caps では $\mu_a=\eps_a<w_a/2$ とし，
crossing product charts の内部だけでこれらの hypersurfaces の corners を丸める．
全ての cuts は collar fibers で飽和させる．
$Z\mu_a=\mu_a-w_a<0$ なので，得られる link $S_\infty$ は $Z$ に横断的であり，
$\alpha=\lambda|_{S_\infty}$ は contact form である．
normal-circle lift と toric orbits はそこで
\[
 S_\infty\longrightarrow\Gamma_S\simeq\Gamma
\]
という $\dd\alpha$-Lagrangian fibration を与える．
full saddle fiber は $S^1\times\figeight$ である．

$\Sigma$ 外の内側の全ての点は正時間の flow により $S_\infty$ に達する．
低い sublevel では logarithmic Weinstein structure，
高い end では \eqref{snc:cutoff-estimate} の nontrapping を使う．
従って内側の $\Sigma$ の補集合は
$((-\infty,0]\times S_\infty,\dd(e^t\alpha))$ と同定され，
$(t,x)\mapsto(e^t,f(x))$ が cone fibration を与える．
$\Sigma$ 全体を cone point に送る延長は，Evans--Mauri の cone-extension argument
\cite[Lemma 2.16]{EM}によって連続である．
これは completion の議論であり，有限体積側への延長は以下で別途確認する．

normal moment は
\begin{equation}\label{snc:time}
 \mu_a(t)=w_a+e^t(\mu_a(0)-w_a),\qquad
 T_a=\log\frac{w_a}{w_a-\mu_a(0)}
\end{equation}
を満たす．$P_a$ 上では曲線方向の点が固定されるので，
$T_a$ に normal circle を潰せば指定された葉に到達する．

transition cap $I\in[a,b]$，$a>0$ では，少し大きい interval に buffer を取り，
$0\leq\mu\leq\eps_a<w_a/2$ において $|ZI|\leq C$ と抑える．
\[
 T_a\leq\frac{2\eps_a}{w_a},\qquad
 |I(t)-I(0)|\leq\frac{2C\eps_a}{w_a}\quad(0\leq t\leq T_a).
\]
$\eps_a<a w_a/(4C)$ とし，上側の buffer からも退出しないように縮めれば，
$\mu=0$ に達するまで $I\geq a/2>0$ である．$C=0$ ではこの制約は不要である．
従って別の divisor に先に到達することはない．
全ての caps が有限個であるため，一つの十分薄い collar で同時に処理できる．

さらに $J=I-w_b+\mu A(I)$ と置くと $ZJ=J$ であり，
\begin{equation}\label{snc:endpoint}
 I_{\rm end}
 =\frac{w_aI_0-\mu_0w_b+w_a\mu_0A(I_0)}{w_a-\mu_0},\qquad
 \frac{\partial I_{\rm end}}{\partial I_0}
 =\frac{w_a(1+\mu_0A'(I_0))}{w_a-\mu_0}>0.
\end{equation}
第二式では transition cap の link 上で $\mu_0=\eps_a$ が一定であることを使う．
$P_a$ 側では $I_{\rm end}=I_0$ である．

product crossing では二つの actual normal moments が同時に存在する．
ここでのみ最初の collapse time を $\min(T_a,T_b)$ と書く．
$R=(w_a-\mu_a)/(w_b-\mu_b)$ は flow invariant であり，endpoint は
\[
 \begin{cases}
 (0,w_b-w_a/R),&R\geq w_a/w_b,\\
 (w_a-w_bR,0),&R\leq w_a/w_b
 \end{cases}
\]
となる．丸めた link arc 上で $R$ は厳密に単調なので，
その quotient は二本の divisor arms の quotient と同相である．
$R=w_a/w_b$ では両 circle が同時に潰れ，交点に到達する．

従って endpoint map $\Gamma_S\to\Gamma_D$ は，$P_a$ 上で恒等，
transition arms 上で厳密単調，crossing で同相である．
同じ flow による写像なので overlaps で一致し，graph 全体の同相を与える．
最初の到達時刻 $T(q)$ は正の連続関数である．
saddle の近傍では一定で，product crossing では二つの連続関数の最小値である．
base は $-\infty<t\leq T(q)$ に cone point を加えたものとなり，
\[
 (t,q)\longmapsto(e^{t-T(q)},q)
\]
で $\operatorname{Cone}_{\rm cl}(\Gamma)$ と同相になる．
この同相による閉端での連続性も actual collar 座標で確認できる．
$r=e^{t-T(q)}$ と置くと，component face と transition cap では
$r=1-\mu_a/w_a$，product crossing では
\[
 r=\max\{1-\mu_a/w_a,\,1-\mu_b/w_b\}
\]
である．従って $D$ への任意の接近で $r\to1$ となる．
Graph coordinate は component 上では固定葉の quotient，transition cap では
\eqref{snc:endpoint}の厳密単調な endpoint coordinate により連続に延びる．
Product crossing では不変量 $R$ を用い，交点への接近で $R\to w_a/w_b$ となるため，
同じ crossing value へ収束する．Normal charts の同定は zero section で恒等なので，
この延長は実際の $D$ 上の指定された全葉に一致する．
collapse は全て fiber ごとに行われるので，新しい fiber の切断や連結成分は生じない．
outer graph の逆像は正確に $D$ となり，boundary fibers の記述も従う．
以上で定理\ref{snc:main}を証明した．

\subsection{固定形式と固定 residue に関する限界}

\begin{proposition}[fixed form の orthogonality 障害]
\label{snc:orthogonality}
固定した $(M,\omega)$ と固定した SNC divisor $D$ に対して，
交点近傍で commuting normal-circle actions を持つ product toric collar が存在するなら，
各交点で $D_a$ と $D_b$ は $\omega$-orthogonal である．
従って定理\ref{snc:main}で $\omega'$ を任意の元の $\omega$ に置き換えることはできない．
\end{proposition}
\begin{proof}
交点の tangent space において，第一の circle の固定部分空間は $TD_a$，
非自明な weight space はその normal plane である．
$\omega$ の invariance により，この二つの weight spaces の mixed pairing は $0$ である．
product action では第二の divisor の tangent plane が第一の normal plane なので，
$TD_a$ と $TD_b$ は symplectically orthogonal である．
これは \cite[Proposition 2.5]{BSV} の局所線形代数でもある．

具体例として，$\PP^2$ の Fubini--Study form と
$D=\{x_1=0\}\cup\{x_2-x_1=0\}\cup\{x_0=0\}$ を取る．
各成分の valence は二で，正の weights による \eqref{snc:positive-class} を満たせる．
しかし $[1:0:0]$ の affine coordinates では最初の二直線の tangent spaces は
$\CC(0,1)$ と $\CC(1,1)$ であり，標準 Hermitian form に関して直交しない．
\end{proof}

元の $\omega$ 上に fibration を戻したければ，fixed-class Kähler path に対する
Moser isotopy を用いればよい．このとき compactifying divisor は $D$ と isotopic な
symplectic divisor へ移る．固定した holomorphic $D$ を同時に保存するとは限らない．
元の form が必要な holomorphic product germs を既に持つ場合には，
potential の Morse perturbation を $D$ から離して行えるため，
Moser を divisor の近傍で恒等に選ぶことができる．

\begin{proposition}[global residues と core germ]
\label{snc:residue-change}
$W=M\setminus D$ とし，closed one-form $\eta$ の residue を
$\operatorname{Res}_a(\eta)=-(2\pi)^{-1}\int_{\mathfrak m_a}\eta$ と定義する．
次は完全列である：
\begin{equation}\label{snc:global-residue-exact}
 H^1(M;\RR)\longrightarrow H^1(W;\RR)
 \xrightarrow{\operatorname{Res}}\RR^r
 \xrightarrow{(c_a)\mapsto\sum c_a[D_a]} H^2(M;\RR).
\end{equation}
従って，同じ $\omega|_W$ の primitive の重みを $w$ から $w'$ に変更できるための
必要十分条件は $\sum_a(w_a'-w_a)[D_a]=0$ である．
ただし $\Sigma\hookrightarrow W$ が homotopy equivalence である core の近傍で
primitive の germ を固定したまま変更できる重みは，$w'=w$ に限られる．
\end{proposition}
\begin{proof}
相対 cohomology の低次数において
$H^2(M,W;\RR)\simeq\RR^r$ であり，その生成元は各成分の relative Thom class である．
成分同士の交点は実 codimension 四なので，この次数に新たな生成元や関係は生じない．
対 $(M,W)$ の完全列を，Thom class と meridian integral で表示すると
\eqref{snc:global-residue-exact} を得る．
同じ二形式の primitive の差は closed one-form であるから，最初の主張が従う．
relative de Rham class と prescribed wrapping numbers の関係は
\cite[\S1.1, Proposition 2.10]{BSV}でも明示されている．

差の closed form が $\Sigma$ の近傍で $0$ なら，
$\Sigma\hookrightarrow W$ が $H_1$ の同型を誘導するため，その全 periods は $0$ である．
従って差は exact であり，全 residues も $0$ となる．
これは primitive の germ の固定に関する障害であって，
異なる weights による skeleton の集合としての同一性を一般に否定するものではない．
\end{proof}

特に，面積の線形条件 \eqref{snc:area-balance} だけで
$[\omega]=2\pi\sum w_a[D_a]$ という global 条件を代用してはいけない．
前者は collar の curvature balance であり，後者は complement 上の global exactness を保証する．
また，graph-period correction は weights を固定した後の操作である．
residues が異なる差を basic graph forms と bounded exact cutoff だけで消すことはできない．

\subsection{valence 条件と具体例}

$P_a$ の Euler 数は $2-2g_a-n_a$ である．
境界で radial な実数値 Morse function に extrema を許さず saddles だけを置くなら，
その数は $2g_a+n_a-2$ でなければならない．
$n_a\geq2$ なら上述の構成により実現でき，$n_a=2$ を除外する必要はない．
$n_a=0$ なら compactness により最大と最小が必要である．
$n_a=1$ なら唯一の境界成分上の値が一定なので，非定数の関数は少なくとも一つの
内部 extremum を持つ．従って $n_a\geq2$ は，ここで指定した
実数値 Morse function による無 extrema の collar 模型の条件として sharp である．
circle-valued foliations や elliptic strata を許す別の模型は，この主張の対象ではない．

\begin{corollary}[graph と core の数値]
\label{snc:counts}
定理\ref{snc:main}では
\[
 b_1(\Gamma)=\sum_a g_a+e-r+1,\qquad
 \#\{\text{開放射分岐辺}\}=2\sum_a g_a+2e-2r,
\]
\[
 \chi(\Sigma)=\chi(M)-2r+2\sum_a g_a+e.
\]
\end{corollary}
\begin{proof}
最初の二式は \eqref{snc:graph-counts} である．
三重交点がないので $\chi(D)=\sum_a(2-2g_a)-e$ である．
$W$ は有限 CW 型を持つ実四次元の orientable manifold なので
$\chi(W)=\chi_c(W)=\chi(M)-\chi(D)$ である．
$\Sigma$ が $W$ の homotopy type を持つことから最後の式を得る．
\end{proof}

\paragraph{$\PP^1\times\PP^1$ の toric boundary.}
二本ずつの horizontal fibers と vertical fibers の和を $D$ とする．
$r=e=4$，各成分の valence は二で，$G$ と $\Gamma$ は円周である．
各成分の自己交点は $0$ であり，個々の成分は ample ではない．
しかし両 ruling の coefficient が正なら，その class は Kähler であり，
四成分全てに正の weights を配ることができる．
$W=(\CC^*)^2$ で，core は $T^2$ の homotopy type を持つ．
この例は nonample components と valence 二を同時に含む．

\paragraph{二本の $(1,1)$ 曲線.}
$M=\PP^1\times\PP^1$ 上で二本の smooth $(1,1)$ curves を一般に取り，
二点で横断的に交わらせる．$r=e=2$，各 valence は二で，
$G$ は二頂点を二辺で結ぶ multigraph，$\Gamma$ は円周である．
任意の正の $w_1,w_2$ に対して weighted class は Kähler であり，
\[
 Q=\begin{pmatrix}2&2\\2&2\end{pmatrix},\qquad
 H_1(S;\ZZ)\simeq\ZZ^2\oplus\ZZ/2.
\]
実際，$Q$ の Smith normal form は $\operatorname{diag}(2,0)$ であり，
\eqref{snc:integer-exact} の右端 $H_1(G;\ZZ)=\ZZ$ は free なので列は非標準的に分裂する．
$\chi(W)=2$，$b_1(W)=1$ である．後者は \eqref{snc:global-residue-exact} と，
二成分の class がともに $(1,1)$ であることから従う．
$W$ は complex dimension 二の Stein manifold なので二次元 CW 型を持ち，
$\chi(W)=1-b_1(W)+b_2(W)$ より $b_2(W)=2$ となる．
この core を simple focus--focus fiber と同定することはできない．
少なくともその Euler 数は，simple nodal fiber の $1$ と異なる．

\paragraph{三本の $(1,1)$ 曲線.}
同じ曲面上に三本の smooth $(1,1)$ curves を一般に取り，三重交点を避ける．
各 pair は二点で横断的に交わるため $r=3$，$e=6$，全 valence は四である．
従って $b_1(\Gamma)=4$，trivalent vertices は六個，core の Euler 数は四である．
これは実際に branching を持ち，直線配置の単純グラフには含まれない例である．
正の weights の任意の選択に対して weighted class は Kähler である．

\paragraph{負の自己交点を含む例.}
Hirzebruch surface $\mathbb F_n$ の toric boundary を
$S_0+F_0+S_\infty+F_\infty$ とする．
$S_0^2=-n$，$S_\infty=S_0+nF$，$F_i^2=0$ であり，全 valence は二である．
weights を $u_0,u_\infty,v_0,v_\infty>0$ とすると，weighted divisor class は
\[
 (u_0+u_\infty)S_0+(nu_\infty+v_0+v_\infty)F.
\]
$v_0+v_\infty>nu_0$ ならこの class は ample であり，定理が適用される．
cap の curvature balance は負の $S_0^2$ を含めて
\eqref{snc:curvature-integral} により同時に解かれる．


\paragraph{二本の smooth plane cubics.}
$M=\PP^2$ で，smooth cubics $D_1,D_2$ を九点で横断的に交わるように取る．
$g_1=g_2=1$，$r=2$，$e=9$，$n_1=n_2=9$ であり，
任意の正の weights が Kähler class を与える．
各 $P_a$ の Morse foliation は九個の saddles を持ち，
\[
 b_1(\Gamma)=10,\qquad |V_3(\Gamma)|=18,\qquad \chi(\Sigma)=12.
\]
meridians の整数表示は
\[
 H_1(W;\ZZ)=\langle \mathfrak m_1,\mathfrak m_2\mid
          3\mathfrak m_1+3\mathfrak m_2=0\rangle
     \simeq\ZZ\oplus\ZZ/3.
\]
これは pair の homology exact sequence と，line が各 cubic と三点で交わることから従う．
この例では $H^1(D;\RR)$ の dimension は十二であるのに対して
$H^1(\Gamma;\RR)$ の dimension は十であり，
foliation を固定した basic forms だけでは任意の period 差を補正できないことが分かる．
補題\ref{snc:holonomy}は，二つの elliptic components の holonomy を選んでこれを解消する．

\subsection{固定 connection と全 leaf-period 条件}

以下で固定するのは unitary holonomy の値だけでなく，normal connections と
actual collar への coupling charts，およびその full fiber partition である．
許す補正は closed foliation-basic form と，core を保つ exact cutoff である．
この範囲で，先の period 補正の必要十分条件を記述できる．

\begin{proposition}[固定 coupling charts に対する相対 primitive 接着]
\label{snc:fixed-holonomy}
本節の logarithmic models の下で，$\lambda_0$ と $\lambda_{\rm col}$ は
同じ二形式 $\omega'$ の primitive であり，同じ residues $-2\pi w_a$ を持つとする．
同じ angular poles を引いた remainders が有界であることも，この model の条件に含める．
Plumbing retraction の境界への制限を $\rho:S\to D$ と書くと，
命題\ref{snc:integer-sequence}により一意な $c\in H^1(D;\RR)$ があり，
$[\lambda_{\rm col}-\lambda_0]=\rho^*c$ となる．次は同値である．
\begin{enumerate}
\item Collar の closed basic form $\eta$ と有界関数 $f$ により
  $\lambda_{\rm col}-\eta-\lambda_0=\dd f$ とできる．
  $\eta$ は全 critical leaves と crossing germs の近傍で零に選べる．
\item $c\in q^*H^1(\Gamma;\RR)$，すなわち $[c]=0$ in $\mathcal O_{\mathcal F}(D)$．
\item 補題\ref{snc:reeb-leaf-kernel}の各 $\ell_{a,k}$ に対して
  \begin{equation}\label{snc:fixed-leaf-period-test}
   \int_{\widetilde\ell_{a,k}}(\lambda_{\rm col}-\lambda_0)=0
   \qquad(1\le k\le g_a)
  \end{equation}
  が成り立つ．ここで $\widetilde\ell_{a,k}$ は collar 上の任意の circle-bundle lift である．
\end{enumerate}
必要なら固定 charts を十分薄い共通 subcollar に制限する．これらの条件の下で，
core の primitive germ，その固定 collar の全葉，normal moments，
full saddle collars と crossing germs を保つ本節の exact-cutoff 構成を実行できる．
Finite endpoint で得られる $D$ 上の全葉も元の指定と一致する．
\end{proposition}

\begin{proof}
Meridian periods が零なので，\eqref{snc:fixed-leaf-period-test}は lift の選び方に依存せず，
$c([\ell_{a,k}])$ に等しい．補題\ref{snc:reeb-leaf-kernel}より (2) と (3) は同値である．
各 $\widetilde\ell_{a,k}$ は一つの full torus fiber 内に取れるので，
basic $\eta$ と $\dd f$ の積分は零である．従って (1) なら (3) である．
逆に (2) の下では同補題の edge-bump representatives を collar に延ばして $\eta$ を取る．
差は exact となり，同じ angular poles を引いた remainder の有界性と
有限 plumbing cover の経路による議論から，$\dd f$ と $f$ を有界に取れる．
従って (1) が成立する．

Core へ接ぐ cutoff の nontrapping 評価は\eqref{snc:cutoff-estimate}と同じである．
ここでは一般の basic $\eta$ を $P_a$ の regular annuli にも置くので，
curve 上の点を固定することと，全葉を保つことを区別する．
$\lambda_{\rm col}-\eta$ の Liouville field を $Z$ とすると
\[
 Z\mu_a=\mu_a-w_a,\qquad \alpha_a(Z)=0,\qquad
 (1-\mu_a/w_a)\iota_{Z^{\rm horizontal}}\omega'|_{D_a}=-\eta.
\]
$\eta$ は葉の tangent を消すから，horizontal part は葉の tangent である．
さらに $\mu_a\le\eps_a<w_a/2$ では滑らかに有界である．
従って有限時刻\eqref{snc:time}までの base flow は同じ全正則葉上の diffeomorphism を与え，
normal circle を潰す endpoint も同じ全葉である．
Critical collars では $\eta=0$ なので full saddle leaves をそのまま保つ．
Cap の basic form $g(I)\dd I$ は angular components だけを変えるため，
\eqref{snc:cap-dynamics}と\eqref{snc:endpoint}も変わらない．
Crossing germs も支持の外であり，既存の閉端への連続性が使える．
\end{proof}

\begin{corollary}[必要な holonomy 調整の独立方向]
\label{snc:minimal-holonomy-controls}
固定された葉層に対し，$\sum_a g_a$ 本の closed real one-forms を用いる connection
変更と basic correction で，任意の $c\in H^1(D;\RR)$ の障害を消せる．
Holonomy 変更から $\mathcal O_{\mathcal F}(D)$ への線形写像の rank は正確に $\sum_a g_a$ である．
従って全ての可能な障害を制御する線形 family には，これだけの独立方向が必要十分である．
ここで holonomy 変更後の lift は新しい coupling charts を用いる．
\end{corollary}

\begin{proof}
各 $a$ で，$\int_{\ell_{a,j}}\eta_{a,k}=\delta_{jk}$ となる closed forms を選ぶ．
補題\ref{snc:reeb-leaf-kernel}の primitive lattice とその双対 cycles により選択でき，
局所 primitive の cutoff により全 crossing disks と指定 flat patches の近傍で零にできる．
$p_{a,k}=c([\ell_{a,k}])$ とし，connection を
\[
 \alpha_a\longmapsto\alpha_a+
       \sum_{k=1}^{g_a}(p_{a,k}/w_a)\pi^*\eta_{a,k}
\]
と変える．補題\ref{snc:holonomy}により障害の各座標は
$p_{a,k}-w_a(p_{a,k}/w_a)=0$ となる．残りは $q^*H^1(\Gamma)$ の class なので，
命題\ref{snc:fixed-holonomy}の basic correction で処理できる．
選んだ $\sum_a g_a$ 方向と障害座標に関する行列は，重み $-w_a$ の非零対角行列であり，
rank の主張も従う．
\end{proof}

\begin{remark}[Connection 補正の支持の制約]\label{snc:saddle-support-obstruction}
Closed form $\eta_a$ が全ての full critical-leaf collars と crossing collars の近傍で零なら，
$[\eta_a]|_{L_a}=0$ である．実際，regular annulus 上の葉の period は Stokes の定理により一定で，
端の critical collar または cap で零になる．従って，このような支持の connection 変更は
$\mathcal O_{\mathcal F}(D)$ を変えない．非零障害を消す際には，一般に全ての full saddle collars を
同時に固定することはできない．Saddle 点だけの小さい disks を固定する条件はこれより弱い．
Holonomy 補正後も $D$ 上の指定 full leaves は保たれるが，actual collar での lift は新しい
coupling charts による．障害消去後の basic graph correction は，これとは異なり
全 critical-leaf collars を支持外にできる．
\end{remark}

二本の smooth plane cubics の例では，$\dim H^1(D)=12$，$b_1(\Gamma)=10$ であり，
固定 coupling charts の障害は二つの elliptic components 上の正則葉の period 一つずつである．
各成分の二つの periods 全てを holonomy で消すことは十分だが，必要な独立調整は計二つで足りる．
全成分が有理なら $\mathcal O_{\mathcal F}(D)=0$ であり，この段階には holonomy 変更を要しない．
一方，命題\ref{ke:component-flux-space}の点固定 collar germ の障害空間とは異なる quotient
であり，その消滅を意味しない．また，任意の fibration の存在に対する必要条件や，
全ての foliation-preserving Liouville fields の分類を主張しているわけではない．

Unitary holonomy の位相値だけを固定する場合にも注意が必要である．
$\eta_a$ の全 periods が $2\pi\ZZ$ に入る変更は holonomy phases を変えないが，
\eqref{snc:holonomy-variation}の real primitive class は $-w_a[\eta_a]$ だけ動く．
この変更は normal gauge を変えるので，上の固定 coupling charts の条件には含めない．
位相値だけを固定する問題では，この整数 period の自由を別途扱う必要がある．

\subsection{実際の成分へ移す際に保持されるデータ}
\label{snc:output-data}

定理\ref{snc:main}は固定した Kähler class 内の一つの model として，
次の相対データを与える．
\begin{enumerate}
\item crossing で固定した holomorphic product germs と，
それに一致する commuting normal-circle action germs．
\item 各 labelled $D_a$ 上の実際の marked Morse 葉層と，
その connected critical leaf 全体を含む full saddle collar．
正の genus の場合の holonomy はこの構成の一部として選ぶ．
\item core から離れた有限個の divisor-compatible holomorphic flat balls と，
regular elliptic-boundary patches．局所 nodal surgery を用いる場合には，
その surgery に必要な symplectic chart との同定を別途行う．
\item core と局所 surgery の compact support を避けてさらに縮められる divisor collars．
normal moments は各 zero section 上で $0$ と正規化される．
\item core 近傍の primitive germ を固定した component および graph の period 補正．
この補正だけから，別個に選ばれた symplectic pieces の action data の一致を
結論するものではない．
\end{enumerate}
actual component への移送では，固定 crossing germs を持つ Kähler path，
divisor-relative Moser，normal-circle action germs の同定，
および compact saturated truncation 全体の transport を用いる．
これらは本節の有限 compactification の後で確認すべき別の相対条件である．
正則 band で先指定した第二作用まで合わせる条件は命題\ref{neck:regular-period-matching}，
同じ actual form 上で全葉を選んで接続する構成は第\ref{sec:gluing}節に記す．
作用適合 atlas を得るには，さらに critical action の有理直線性
（命題\ref{atl:critical-action-obstruction}）を満たすように葉族を選ぶ必要がある．
例\ref{atl:fixed-form-bent-branch}は，normal action と両端 germs の固定だけから
その直線性が従わないことを示す．

\input{sections/paper/neck_models.tex}
\input{sections/paper/global_gluing.tex}
\section{任意の pants 木から得られる分岐基底}\label{topo:section}

本節の結果は有限複体についての定理であり，シンプレクティックな接着の成立を仮定しない．
したがって，前節の大域構成とは独立に用いることができる．
以下，$d\geq5$，$[d]=\{1,\ldots,d\}$ とする．

\begin{definition}\label{topo:definition}
各非順序対 $ij=\{i,j\}\subset[d]$ に対し，葉が $[d]\setminus\{i,j\}$ でラベル付けされた
有限三価木 $T_{ij}=T_{ji}$ を選ぶ．内部頂点の集合を $P_{ij}$，内部辺の集合を $A_{ij}$ とし，
葉 $k$ に隣接する内部頂点を $p_{ij}(k)$ と書く．
分岐複体 $B(T)$ を次のセルで定める．
\begin{enumerate}
\item 頂点は $v_1,\ldots,v_d$ とする．
\item $p\in P_{ij}$ ごとに $v_i$ と $v_j$ を結ぶ辺 $e_{ij,p}$ を置く．
\item 内部辺 $pq\in A_{ij}$ ごとに，$e_{ij,p}$ と $e_{ij,q}$ に沿う二角形を一枚付ける．
\item $i<j<k$ ごとに，$e_{ij,p_{ij}(k)}$，$e_{jk,p_{jk}(i)}$，
$e_{ik,p_{ik}(j)}$ を境界とする三角形を一枚付ける．
\end{enumerate}
二角形は二本の異なる辺に沿う閉円板である．この記述は有限正則 CW 複体を定め，
細分すれば有限単体的多面体になる．
\end{definition}

三価木の辺数の公式から
\[
 |P_{ij}|=d-4,\qquad |A_{ij}|=d-5
\]
である．葉と葉に接する開辺を除く部分 $T_{ij}^{\rm int}$ は，内部頂点を全て含む木であり，
$d=5$ では一点である．

\begin{proposition}\label{topo:counts}
$B(T)$ の各辺の内部にはちょうど三枚の面が接する．セル数は
\[
 V=d,\qquad E=\binom d2(d-4),\qquad
 F_\triangle=\binom d3,\qquad F_2=\binom d2(d-5).
\]
頂点 $v_i$ のリンク $G_i$ は，完全グラフ $K_{d-1}$ の各頂点 $j$ を
$T_{ij}^{\rm int}$ で展開した連結三価グラフであり，
\[
 |V(G_i)|=(d-1)(d-4),\quad
 |E(G_i)|=\frac32(d-1)(d-4),\quad
 b_1(G_i)=\binom{d-2}{2}.
\]
\end{proposition}

\begin{proof}
$e_{ij,p}$ に接する面は，$T_{ij}$ における $p$ の三つの flags に対応する．
内部辺の flag は二角形，葉 $k$ の flag は三角形 $ijk$ を与えるため，面の数は三である．
リンクの頂点は $(j,p)$，$j\ne i$，$p\in P_{ij}$ である．
その辺は，各 $T_{ij}^{\rm int}$ の内部辺と，$j,k\ne i$ に対応する
$(j,p_{ij}(k))$ と $(k,p_{ik}(j))$ を結ぶ辺である．
内部木を全て縮約すると $K_{d-1}$ になり，連結性と第一 Betti 数が従う．
三価性から $2|E(G_i)|=3|V(G_i)|$ も従う．
\end{proof}

\subsection{縮約写像と整数係数の鎖収縮}

$S_d=\operatorname{Sk}_2\Delta^{d-1}$ と置く．同じ $ij$ の平行辺を $S_d$ の辺 $ij$ に送り，
二角形をその辺に潰し，三角形を対応する三角形に写す写像
\[
 q:B(T)\longrightarrow S_d
\]
がある．より具体的には，$ij$ の全ての辺と二角形の和集合を
\[
 \bigl(T_{ij}^{\rm int}\times[0,1]\bigr)/
 \bigl(T_{ij}^{\rm int}\times\{0\}\sim v_i,\ 
 T_{ij}^{\rm int}\times\{1\}\sim v_j\bigr)
\]
と表示し，第二因子への射影を使う．三角形上の写像は，この境界のパラメータ表示と整合するように選ぶ．
各辺の内点の逆像は $T_{ij}^{\rm int}$，それ以外の点の逆像は一点である．
したがって $q$ は cell-like な写像，すなわち全ての点逆像が可縮な compact 多面体である写像である．

\begin{theorem}\label{topo:sdr}
任意の木の組 $T=(T_{ij})$ に対し，$q$ はホモトピー同値であり，
\[
 B(T)\simeq\bigvee^{\binom{d-1}{3}}S^2.
\]
さらに，セル鎖写像 $q_*$ は整数係数で明示的な鎖切断と鎖ホモトピーを持つ．
\end{theorem}

\begin{proof}
$i<j$ に対して $e_{ij,p}$ を $v_i$ から $v_j$ に向ける．各内部木に根
$r_{ij}\in P_{ij}$ を選び，内部辺を根から遠ざかる向きに向ける．
その向きが $p\to u$ である二角形を $b_{ij,pu}$ と書き，
\[
 \partial b_{ij,pu}=e_{ij,u}-e_{ij,p}
\]
となる向きを与える．根から $p$ への経路に沿う二角形鎖を $D_{ij}(p)$ とすると
\[
 D_{ij}(r_{ij})=0,\qquad
 \partial D_{ij}(p)=e_{ij,p}-e_{ij,r_{ij}}.
\]
三角形 $ijk$ の向きを $i<j<k$ の順序で定め，その鎖を $\tau_{ijk}$ とする．
補正三角形
\begin{align*}
 \widehat\tau_{ijk}={}&\tau_{ijk}
   -D_{ij}(p_{ij}(k))-D_{jk}(p_{jk}(i))
   +D_{ik}(p_{ik}(j))
\end{align*}
は
\[
 \partial\widehat\tau_{ijk}
 =e_{ij,r_{ij}}+e_{jk,r_{jk}}-e_{ik,r_{ik}}
\]
を満たす．したがって，$S_d$ の頂点を同じ頂点に，辺 $ij$ を
$e_{ij,r_{ij}}$ に，三角形 $ijk$ を $\widehat\tau_{ijk}$ に送る写像
$s:C_*(S_d;\ZZ)\to C_*(B(T);\ZZ)$ は鎖写像であり，$q_*s=1$ である．

次数一の生成元に対し $H(e_{ij,p})=D_{ij}(p)$ と定め，それ以外の次数で
$H=0$ と置く．木の経路の一意性から，内部辺 $p\to u$ に対して
$D_{ij}(u)-D_{ij}(p)=b_{ij,pu}$ である．これと補正三角形の式を使うと，
全てのセルについて
\[
 1-sq_* =\partial H+H\partial
\]
が成り立つ．また $Hs=0$，$q_*H=0$，$H^2=0$ である．
これは整数係数の鎖複体の明示的な strong deformation retraction である．

基本群について，二角形の関係は同じ $ij$ の全ての辺を相対端点ホモトピーで同一視する．
残る三角形の関係は完全グラフの全てのループを消すので，$B(T)$ は単連結である．
$S_d$ も単連結であるため，いま証明した整数ホモロジー同値と
単連結 CW 複体に対するホモロジー版 Whitehead 定理から，$q$ はホモトピー同値である
\cite[Corollary 4.33]{Hatcher}．
最後に，$S_d$ は単連結な二次元複体であり，$H_2(S_d;\ZZ)$ は自由群である．
その基底を Hurewicz の定理で球面からの写像として実現し，再び Whitehead の定理を使えば，
主張した球面の wedge を得る．球面数は
$\chi(S_d)-1=d-\binom d2+\binom d3-1=\binom{d-1}{3}$ である．
\end{proof}

\begin{corollary}[整数二サイクルの基底]\label{topo:cycles}
$2\leq i<j<k\leq d$ に対し
\[
 Z_{ijk}=\widehat\tau_{ijk}-\widehat\tau_{1jk}
              +\widehat\tau_{1ik}-\widehat\tau_{1ij}
\]
と置くと，これらは $H_2(B(T);\ZZ)$ の基底を与える．
\end{corollary}

\begin{proof}
標準複体では，$[ijk]-[1jk]+[1ik]-[1ij]$ は四面体 $[1ijk]$ の境界である．
任意の二サイクルから，頂点 $1$ を含まない各三角形の係数だけこれらを差し引くと，
残りは $1$ を含む三角形だけに支持される．そのような二サイクルは零である．
実際，辺 $ij$，$i,j>1$ はその支持の中では三角形 $1ij$ にだけ現れる．
同じ理由で基底の一次独立性も分かる．鎖切断 $s$ で移せばよい．
\end{proof}

この証明は有限個の次数の行列の階数計算を一般次数の証明に代用しない．
各木について根と経路を列挙するだけで，整数係数の切断・ホモトピー・二サイクルを出力できる．

\subsection{通常の双対複体との差}

\begin{proposition}\label{topo:nonhomeomorphic}
$d=5$ では $q:B(T)\to S_5$ は同相写像である．$d\geq6$ では
$B(T)$ と $S_d$ は同相でない．
\end{proposition}

\begin{proof}
$d=5$ では各内部木は一点であり，二角形が存在しない．
$d\geq6$ では点 $x$ の局所ホモロジー
$H_2(Y,Y\setminus\{x\};\ZZ)$ の階数で区別できる．
三枚の面が接する $B(T)$ の辺内ではこの階数は二である．一方 $S_d$ では，
面内・辺内・頂点における階数はそれぞれ
\[
 1,\qquad d-3,\qquad \binom{d-2}{2}
\]
である．これは点のリンクの第一 Betti 数から従う．$d\geq6$ ではどれも二でない．
\end{proof}

したがって，木による三価化は通常の細分ではない．
自由な辺がないこと，双対複体とホモトピー同値であること，双対複体と同相であることは
区別しなければならない．また，この議論自体は integral-affine atlas を構成しない．

\begin{corollary}[KE 退化の cubical candidate への位相的射]
\label{topo:KE-cubical-comparison}
$d\ge5$ に対し，別の研究で用いる有限立方体複体
\[
 \mathcal C_d=
 \left\{w\in[0,1]^d:\max_i w_i=1,\ 
             |\operatorname{supp}(w)|\le3\right\}
\]
を標準 square の貼合せで定める．正規化
\[
 \rho_d:\mathcal C_d\longrightarrow S_d,\qquad
 w\longmapsto\frac{w}{\sum_iw_i}
\]
は同相であり，
\[
 p_T=\rho_d^{-1}\circ q:B(T)\longrightarrow\mathcal C_d
\]
は全ての点逆像が可縮な連続全射で，整数ホモロジーの同型を誘導する
ホモトピー同値である．$d=5$ では $p_T$ は同相，
$d\ge6$ では同相でない．
\end{corollary}
\begin{proof}
$\sum_iw_i\ge1$ であり，$S_d$ の各点 $x$ は
$x/\max_i x_i\in\mathcal C_d$ を一意に定めるので $\rho_d$ は同相である．
したがって $p_T$ の点逆像は $q$ の点逆像と一致し，
定理\ref{topo:sdr} の鎖ホモトピーとホモトピー同値がそのまま適用できる．
同相性の区別は命題\ref{topo:nonhomeomorphic}から従う．
\end{proof}

\begin{lemma}[連結 fiber を持つ閉写像の合成]
\label{topo:connected-preimages}
$X$ を compact Hausdorff 空間，$Y$ を Hausdorff 空間とし，
$p:X\to Y$ を連結 fiber を持つ連続全射とする．
任意の連結部分集合 $A\subset Y$ に対し，$p^{-1}(A)$ は連結である．
従って compact Hausdorff 空間の間の連結 fiber の連続全射を合成しても，
連結 fiber の連続全射となる．
\end{lemma}
\begin{proof}
$p$ は閉写像である．$p^{-1}(A)$ 内の相対閉集合は
$C\cap p^{-1}(A)$，$C\subset X$ 閉，と書け，その像は $p(C)\cap A$ である．
従って制限 $p_A:p^{-1}(A)\to A$ も閉写像である．
$p^{-1}(A)=U\sqcup V$ が非自明な分離なら，連結な各 fiber は
$U,V$ の一方に全体として入る．よって $p(U),p(V)$ は
$A$ を覆う非空な互いに素な相対閉集合となり，連結性に反する．
合成の点逆像は，後の写像の連結な点逆像を前の写像で引き戻したものなので，
最初の主張を適用すればよい．
\end{proof}

\begin{corollary}[Cubical 射の全 fiber の連結性]
\label{topo:cubical-composite-connected}
本稿で構成する $\pi:X\to B_d(T)$ に対し，$f=p_T\pi:X\to\mathcal C_d$ は
proper な連結 fiber の連続全射である．特に
\[
 f^{-1}(w)=\pi^{-1}\bigl(p_T^{-1}(w)\bigr)
\]
は，$p_T^{-1}(w)$ が内部木である場合も連結である．
\end{corollary}
\begin{proof}
$\pi$ の主張は補題\ref{glue:closed-whole-fiber-quotient}の適用，
$p_T$ の主張は系\ref{topo:KE-cubical-comparison}による．
前補題を用い，properness は compactness から従う．
\end{proof}

$p_T$ の点逆像が可縮であることから $f$ の点逆像の可縮性は結論しない．
$d\ge6$ で内部木が辺を持つ場合，その開辺上の正則 torus の合併は
局所的に interval と $T^2$ の積であり，$f$ の fiber は三次元部分を含み得る．
Lagrangian torus fibration と呼ぶのは $\pi$ であり，
$f$ はその全 fiber 合併を扱う計量比較写像である．

\begin{remark}[比較の射程]
系\ref{topo:KE-cubical-comparison}は候補極限複体と Lagrangian 分岐基底の
\emph{labelled topology} の比較である．
$q$ は一般に距離を保存せず，$p_T$ から
KE 計量の GH 収束，action-compatible integral-affine atlas，
あるいは special Lagrangian 構造は導けない．
$d\ge6$ の非同相性は，両基底をそのまま同じ位相空間として
同一視することに対する具体的な障害である．
\end{remark}


\paragraph{既知の cubical subdivision との関係．}
\label{topo:cubical-barycentric-comparison}
$\mathcal C_d$ の各閉 cell は，空でない集合の組
$A\subseteq B\subseteq\{1,\ldots,d\}$，$|B|\le3$ に対して
\[
 Q_{A,B}=\{w_i=1\ (i\in A),\quad 0\le w_i\le1\ (i\in B\setminus A),
 \quad w_i=0\ (i\notin B)\}
\]
と書ける．その face poset は，単体の二骨格の非空 face poset の
閉区間 $[A,B]$ を包含で順序付けたものに等しい．
従って組合せ型は既知の cubical barycentric subdivision，すなわち
Babson--Billera--Chan\cite[1996年preprint, \S2.3, Proposition~2.2, Corollary~2.4]{BabsonBilleraChan}
の barycentric cover の特殊例である．
この組合せ的構成自体の新規性を主張しない．以下では各二次元 cell を
\emph{unit Euclidean square} とする長さ距離を固定し，
その非伸長 retraction と直径を別に証明する．距離結果の先行研究との
網羅的な照合は未了である．

\subsection{座標部分複体の非伸長 retraction}

以後，各 unit square の Euclidean 長さを貼り合わせた距離を $d_\square$ とし，
$Q_{p;qr}=e_p+[0,1]e_q+[0,1]e_r$ と書く．
非空な $I\subset\{1,\ldots,d\}$ に対して
\[
 \mathcal C_I=\{w\in\mathcal C_d:\operatorname{supp}(w)\subset I\}
\]
とし，その cubical cells の Euclidean 長さから定まる内在距離を $d_I$ と書く．
$|I|=1$ の場合は一点，$|I|=2$ の場合は二本の unit intervals の和である．
$|I|\ge3$ の場合は各 unit square を貼り合わせた有限長さ空間となる．
この定義だけでは $d_I$ は $d_\square$ の制限を意味しないが，
以下の直接構成によって両者は等しい．

\begin{proposition}[座標部分複体の非伸長 retraction]
\label{topo:coordinate-retraction}
$a\in I$ を一つ選び，$r_{I,a}:\mathcal C_d\to\mathbb R^d$ を
\[
 \bigl(r_{I,a}(w)\bigr)_b=
 \begin{cases}
  \min\!\left\{1,
   \left(w_a^2+\displaystyle\sum_{\ell\notin I}w_\ell^2\right)^{1/2}
       \right\},&b=a,\\
  w_b,&b\in I\setminus\{a\},\\
  0,&b\notin I
 \end{cases}
\]
で定める．これは $(\mathcal C_d,d_\square)$ から
$(\mathcal C_I,d_I)$ への $1$-Lipschitz retraction である．特に
\begin{equation}\label{topo:eq:coordinate-distance}
 d_I(x,y)=d_\square(x,y)\qquad(x,y\in\mathcal C_I).
\end{equation}
さらに，各 unit square $Q\subset\mathcal C_d$ 自身も
非伸長 retraction の像であり，$Q$ 上で $d_\square$ は
その Euclidean 距離に等しい．
\end{proposition}
\begin{proof}
外側の全座標と $a$ 座標を一つの非負座標にまとめるので，
非零座標の個数は増えず，像の support は $I$ に含まれる．
入力の値 $1$ を取る座標が $I\setminus\{a\}$ にあればそのまま残り，
それ以外なら像の $a$ 座標が $1$ となる．
したがって像は $\mathcal C_I$ に入る．
$w\in\mathcal C_I$ では $w_a\le1$ より $r_{I,a}(w)=w$ である．

写像 $u\mapsto\min\{1,\|u\|_2\}$ は Euclidean 距離に関して
$1$-Lipschitz である．これは逆三角不等式と，実数の clipping が
$1$-Lipschitz であることから従い，$u=0$ と $\|u\|_2=1$ でも成立する．
その微分が存在する点では勾配の長さは $1$ 以下であり，
clipping の外側では零である．変更しない他の座標はこの $u$ と直交するから
\[
 \|r_{I,a}(w)-r_{I,a}(w')\|_2^2
 \le\sum_{b\in I\setminus\{a\}}|w_b-w'_b|^2
       +|w_a-w'_a|^2+\sum_{\ell\notin I}|w_\ell-w'_\ell|^2
 =\|w-w'\|_2^2.
\]

ここでは target の内在距離も確認する必要がある．
一つの source square の固定座標を $p$ とする．
$p\in I\setminus\{a\}$ なら像でも $p$ 座標が $1$ であり，
$p=a$ または $p\notin I$ なら像では $a$ 座標が $1$ である．
残る可変座標は高々二つで，その label は元の square と $I,a$ から固定される．
従って各 source square は一つの target cell（square，辺，または点）に写る．
上の Euclidean 評価から，その square 内の任意の曲線の長さは増えない．
有限個の cell 内線分をつないだ経路に適用して長さを足し，
piecewise Euclidean 距離の定義に従って下限を取れば
\[
 d_I(r_{I,a}(x),r_{I,a}(y))\le d_\square(x,y)
\]
を得る．従って任意の rectifiable curve についても，距離空間の
$1$-Lipschitz 写像の一般的性質として長さは増加しない．
包含写像による逆向きの不等式と retraction 性から
\eqref{topo:eq:coordinate-distance} が従う．

最後に，$Q=\{w_p=1,\ 0\le w_q,w_r\le1,\ w_b=0
\text{ for }b\notin\{p,q,r\}\}$ に対して
\[
 P_Q(w)=e_p+w_qe_q+w_re_r
\]
を取る．これは $Q$ 上恒等な Euclidean 非伸長写像で，
target の $Q$ 自身は凸な Euclidean square なので
$P_Q:(\mathcal C_d,d_\square)\to Q$ も非伸長である．
従って $Q$ 上の全体距離と Euclidean 距離は等しい．
\end{proof}

\begin{corollary}[六つ以下の label による最短路]
\label{topo:six-label-gallery}
$d\ge3$ とする．任意の二点 $x,y\in\mathcal C_d$ は，
高々六つの label を持つ座標部分複体内の最短路で結べる．
その最短路は
\[
 N_d^*=3\binom{\min\{d,6\}}{3}
\]
本以下の unit-square 内線分からなる gallery として取れる．
特に $d=5$ では $N_d^*=30$，$d\ge6$ では $N_d^*=60$ である．
また $D_d=\operatorname{diam}(\mathcal C_d,d_\square)$ とすれば
$D_d=D_6$ for $d\ge6$．
\end{corollary}
\begin{proof}
$\operatorname{supp}(x)\cup\operatorname{supp}(y)$ は高々六元である．
これを含み，$3\le|I|\le\min\{d,6\}$ となる $I$ を選ぶ．
有限連結 cubical length space $\mathcal C_I$ は compact であり，最短路を持つ．
命題\ref{topo:coordinate-retraction}によって，その最短路は全体でも最短である．

全体の最短路が同じ unit square の二点を通れば，その間の長さは
両端間の Euclidean 距離に等しい．Euclidean 空間の三角不等式の等号条件より，
その区間は square 内の線分である．従って各 square の逆像は閉区間であり，
同じ square を出て再訪することはない．有限個のこれらの区間を使って
parameter interval を順に覆えば，各 square を高々一度用いる gallery を得る．
$\mathcal C_I$ の square の個数は $|I|\binom{|I|-1}{2}=3\binom{|I|}{3}$
なので所望の上界を得る．

$d\ge6$ のとき，任意の二点間距離は label の置換により
$\mathcal C_6$ 内の二点間距離に等しいから $D_d\le D_6$．
逆向きは $\mathcal C_6\subset\mathcal C_d$ の等長性から従う．
\end{proof}


\subsection{成分頂点の計量錐と固定補空間の無限遠}

\begin{proposition}[成分頂点近傍の厳密な錐模型]
\label{topo:component-cone-germ}
$i\in[d]$ を固定し，
\[
 C_{2,d-1}=\{u\in\mathbb R_{\ge0}^{[d]\setminus\{i\}}:
                   |\operatorname{supp}(u)|\le2\}
\]
に各座標 quadrant の Euclidean 長さを貼った距離 $d_C$ を入れる．
任意の $0<r<1/2$ に対し，
\[
 \iota_i(u)=e_i+\sum_{k\ne i}u_ke_k
\]
は閉球 $\overline B_{d_C}(0,r)$ から
$\overline B_{d_\square}(e_i,r)$ への等長全単射である．
両閉球の距離は，それぞれ錐全体と $\mathcal C_d$ 全体の距離の制限とする．
従って $e_i$ における一意な metric tangent cone は $C_{2,d-1}$ であり，
その link は各辺の長さが $\pi/2$ の完全グラフ $K_{d-1}$ である．
\end{proposition}
\begin{proof}
$\mathcal C_d$ の Euclidean 座標表示は曲線の長さを増加させない．
$w_i<1$ なら別の座標 $w_k=1$ が存在するため，
$d_\square(e_i,w)\ge\|e_i-w\|_2\ge1$ である．
一方，閉 star $S_i=\{w\in\mathcal C_d:w_i=1\}$ 内では，
$u=w-e_i$ に沿う動径線分と Euclidean 下界により
$d_\square(e_i,w)=\|u\|_2$ となる．錐側も $d_C(0,u)=\|u\|_2$ なので，
両閉球は集合として上の写像で一致する．

錐の半径 $r$ への動径射影は，各 quadrant では Euclidean 閉球への射影であり，
共通の face 上で一致する．従って piecewise Euclidean 長さに関して非伸長である．
二端点の norm が $r$ 以下なら，錐内の距離の下限は同じ閉球内の経路で取れ，
その値は原点経由の $2r$ 以下である．これらの経路を $S_i$ へ写せるので，
$d_\square(\iota_i(u),\iota_i(v))\le d_C(u,v)$ である．

逆に $S_i$ 内の経路は同じ長さの錐内経路になる．$S_i$ の外へ出る経路は，
退出時と再入時に，$w_i=w_k=1$ を満たすいずれかの gate を通る．
二端点の各可変座標は $r$ 以下だから，gate までの片側の長さは少なくとも $1-r$ である．
従ってこの経路の全長は $2(1-r)>2r\ge d_C(u,v)$ であり，短縮を与えない．
全経路で下限を取り逆向きの不等式を得る．
閉球の等長は任意に小さい $r$ で成り立つため tangent cone の主張も従う．
座標軸が link の頂点，二軸の quadrant が長さ $\pi/2$ の辺を与える．
\end{proof}

全平面退化の成分を $H_i=\{L_i=0\}\subset\PP^3$ と書くと，
\[
 U_i=H_i\setminus\bigcup_{k\ne i}(H_i\cap H_k),\qquad
 K_{H_i}+\sum_{k\ne i}(H_i\cap H_k)=(d-4)H|_{H_i}.
\]
すなわち $U_i$ は $\PP^2$ から一般位置の $d-1$ 本の直線を除いた固定補空間である．
固定空間の漸近錐を扱う companion 原稿
\cite[\texttt{cor:arrangement-cone}]{MultiscaleKECurrent}の $(n,m)=(2,d-1)$ は，
同じ $\operatorname{Ric}=-1$ の規格化で上の $C_{2,d-1}$ を与える．
軸の label は二重線 $H_i\cap H_k$，quadrant の label は三重点 $H_i\cap H_k\cap H_l$ に対応する．
従って固定成分の無限遠模型と compact な cubical 極限の成分頂点の接錐は，
計量と labels を保って一致する．この有限複体の命題だけから，
actual 退化族の任意の moving 領域の tensor 比較や中間尺度の収束は従わない．
後者について，本稿の大域 distortion を別に使う帰結を
系\ref{ke:intermediate-component-cone}で述べる．

\subsection{全ての局所錐と移動基点の拡大極限}

原点を共有する $m$ 本の半直線を，その path metric で貼った空間を
$\mathcal Y_m$ と書く．また $\mathcal K_{3\pi/2}$ は，三つの Euclidean quadrants
を循環的に貼った全角 $3\pi/2$ の Euclidean cone とする．
以下の積には Euclidean product distance を入れる．

\begin{proposition}[Cubical target の四種類の計量的局所模型]
\label{topo:cubical-local-models}
任意の $q\in\mathcal C_d$ について，十分小さい半径の閉球は，次の空間の
指定された中心の閉球と厳密に等長である．距離はいずれも全空間の ambient 距離の制限である．
\begin{enumerate}
\item $q=e_i$ のときは $(C_{2,d-1},0)$．
\item $|\operatorname{supp}q|=2$ のときは
      $(\mathbb R\times\mathcal Y_{d-2},(0,0))$．
      これには $q=e_i+e_j$ も含む．
\item $q=e_i+e_j+e_k$ のときは $(\mathcal K_{3\pi/2},0)$．
\item 残る点，すなわち support が三元で，その三座標が全て $1$ ではない点では
      $(\mathbb R^2,0)$．
\end{enumerate}
従ってこれらが各点の一意な metric tangent cone である．
\end{proposition}
\begin{proof}
まず局所的な貼合せが ambient 距離も決めることを確認する．有限個の closed squares
のうち $q$ を含まないもの，および incident square の $q$ を含まない faces は，
$q$ から正の Euclidean 距離だけ離れている．その全てを避ける $a>0$ を取ると，
$q$ の Euclidean $a$ 近傍は incident tangent sectors の $a$ 近傍と，
各 sector の距離と貼合せを保って同一視できる．Sector 上の半径 $r$ への動径射影は
Euclidean 閉球への射影であり，共通 faces 上で一致するため長さを増加させない．
よって $r<a/2$ の二端点は，局所錐の半径 $r$ の球内で長さ $2r$ 以下の経路で結べる．
一方，$a$ 近傍の外を通る経路の長さは $2(a-r)>2r$ 以上である．
従って小球上の全点対距離には近傍外からの短縮がなく，局所錐との厳密な等長を得る．

成分頂点では命題\ref{topo:component-cone-germ}を用いる．Support が $\{i,j\}$ の点は
長さ $2$ の折れ辺 $E_{ij}(u)$，$0<u<2$ 上にある．ここで
\[
 E_{ij}(u)=\begin{cases}e_i+u e_j,&0\le u\le1,\\
                      (2-u)e_i+e_j,&1\le u\le2.\end{cases}
\]
各 $k\notin\{i,j\}$ について $E_{ij}(u)+v e_k$，$v\ge0$ が一枚の局所半平面を与える．
$u=1$ では二つの quadrants が直線境界で貼られ，一つの半平面になる．
全ての $d-2$ 枚が $v=0$ の同じ直線だけを共有するので，第二の積模型を得る．

三座標が全て $1$ の点では，その一つを $1$ に保ち残り二つを減らす三つの
quadrants が現れる．三本の負向き座標軸を境界として循環的に貼られるので，
link は長さ $3\pi/2$ の circle である．Support が三元の残りの点では，
$1$ の座標が一つなら一枚の square の内部であり，二つなら二枚の half-planes の
共通境界にある．いずれも局所錐は Euclidean plane となる．
\end{proof}

\begin{proposition}[移動基点の全ての拡大極限]
\label{topo:cubical-moving-blowups}
$\lambda_\nu\to\infty$ と任意の $q_\nu\in\mathcal C_d$ を取る．
Pointed metric spaces $(\mathcal C_d,\lambda_\nu d_\square,q_\nu)$ は
部分列を取れば pointed GH 収束し，その全ての可能な極限は，等長の違いを除いて
\[
 (C_{2,d-1},z),\qquad
 (\mathbb R\times\mathcal Y_{d-2},(0,z)),\qquad
 (\mathcal K_{3\pi/2},z),\qquad
 (\mathbb R^2,0)
\]
である．各 $z$ は対応する錐又は $\mathcal Y_{d-2}$ の任意の有限の点を取る．
任意に指定した $\lambda_\nu\to\infty$ に対し，これらの各基点付き空間を
実現する $q_\nu$ を選べる．
\end{proposition}
\begin{proof}
Compactness により $q_\nu\to q$ としてよい．前命題の $q$ における固定した
局所錐 $Z$ で $q_\nu$ を表し，その座標を $\lambda_\nu$ 倍した点を $z_\nu$ とする．
前命題の full-ambient 等長を半径 $\rho>0$ の閉球上で固定する．任意の固定 $A$ に対し，
$d_Z(0,z_\nu)/\lambda_\nu\to0$ なので，十分大きい $\nu$ では
$d_Z(0,z_\nu)+A<\lambda_\nu\rho$ である．従って二つの中心を動かした半径 $A$ の球は，
この等長を rescale した写像の定義域と像に含まれる．全点対距離も既に保存されているから，
その閉球は $(Z,z_\nu)$ の閉球と厳密に等長である．
以下ではこの四つの $Z$ の移動基点だけを調べる．

$Z=\mathbb R^2$ では平行移動を使う．
$Z=\mathbb R\times\mathcal Y_{d-2}$ では第一因子の平行移動を除き，
第二因子の半直線を部分列で固定する．その座標が有限の値へ収束すれば同じ積の
その点を基点とする極限を得る．無限大へ向かえば，各固定球は半直線の頂点を避け，
二次元 Euclidean 球となる．

$Z=\mathcal K_{3\pi/2}$ では，$z_\nu$ が有界なら properness により収束部分列を取れる．
原点からの距離が無限大へ向かえば，各固定球は Euclidean 球になる．実際，
単位半径の circle は compact で，全ての点が平坦な近傍を持つため，共通の
正の flat-ball radius がある．錐の dilation によってこの半径は基点の半径に比例する．

$Z=C_{2,d-1}$ でも有界な $z_\nu$ は有限の基点を与える．それ以外では support を
部分列で固定し，さらに各非零座標が有限値へ収束するか無限大へ向かうよう選ぶ．
正座標は高々二つである．二座標が無限大へ向かう場合は，固定球の局所模型は
$\mathbb R^2$ である．一座標だけが無限大へ向かう場合は，その座標を平行移動し，
残りの座標は高々一つだけ正であることから
$\mathbb R\times\mathcal Y_{d-2}$ を得る．
ここで単に closed star の内在距離を使っているのではない．発散する各座標の
最小値を $m_\nu$ とすると，半径 $A$ の二端点を結び，その正座標の star を
出る経路は長さ $2(m_\nu-A)>2A$ を要する．Star 内の積座標による閉球の経路は
長さ $2A$ 以下であり，外部の経路は短縮を与えない．これは固定配置錐の
positive-star 評価\cite[\texttt{lem:star-ball-isometry}]{MultiscaleKECurrent}と同じ機構である．

最後に各型を実現する．対応する前命題の点 $q$ を一つ固定し，その局所模型内で
指定した有限基点を $1/\lambda_\nu$ 倍した点を $q_\nu$ とする．
平坦な型は任意の一枚の square の内点を固定すればよい．固定球の等長により
所要の pointed 収束を得る．
\end{proof}

\subsection{六次以上の cubical target の厳密直径}

\begin{lemma}[三 label corner の singleton からの距離]
\label{topo:corner-singleton-radius}
$J=\{i,j,k\}$ とし，$x\in\mathcal C_J$ とする．すると
\[
 d_J(e_i,x)\le\sqrt5.
\]
等号は $x=e_j+e_k$ の場合に限る．特に $x_i>0$ なら不等号は真である．
\end{lemma}
\begin{proof}
$x_i=1$ なら，同じ square 内の線分を使って距離は $\sqrt2$ 以下である．
それ以外では $j,k$ を交換して $x_j=1$ とできる．
$u=x_i$，$v=x_k$ と置く．$w_i=1$ の square と $w_j=1$ の square を
共有辺で開いて $2\times1$ rectangle にすると，両点の差は $(2-u,v)$ となる．
この rectangle 内の線分を戻して
\[
 d_J(e_i,x)^2\le(2-u)^2+v^2\le5
\]
を得る．最後の等号には $u=0$，$v=1$ が必要である．

$x=e_j+e_k$ の下界も直接確認できる．各 square を singleton--triple
対角線で二分し，頂点値
\[
\begin{array}{c|ccccc}
z&e_i&e_j,e_k&e_i+e_j,e_i+e_k&e_j+e_k&e_i+e_j+e_k\\\hline
\sqrt5\,h(z)&0&4&2&5&3
\end{array}
\]
を affine に延長する．各三角形の二つの直交辺方向の増分は，
符号を除いて $1/\sqrt5$ と $2/\sqrt5$ である．
従って $h$ は $1$-Lipschitz であり，
$h(e_j+e_k)-h(e_i)=5/\sqrt5=\sqrt5$ が所望の下界を与える．
\end{proof}

\begin{theorem}[六 label からの直径の安定値]
\label{topo:diameter-six-labels}
任意の $d\ge6$ に対し
\[
 \operatorname{diam}(\mathcal C_d,d_\square)=2\sqrt5=\sqrt{20}.
\]
距離が $2\sqrt5$ に等しい点対は，互いに素な三元集合 $A,B$ による
\[
 v_A=\sum_{i\in A}e_i,\qquad v_B=\sum_{i\in B}e_i
\]
の対に限る．従ってこの直径と等号点対の記述は，KE の解析的な仮定と独立である．
\end{theorem}
\begin{proof}
まず $S=\operatorname{supp}(x)$，$T=\operatorname{supp}(y)$ とし，
$|S\cup T|\le5$ の場合を考える．
$S\cap T\ne\varnothing$ なら $i\in S\cap T$ を選び，
$S,T$ をそれぞれ含む三元集合 $J,K$ を取る．
$S\cap T=\varnothing$ なら少なくとも一方の大きさは二以下なので，
必要なら $x,y$ を交換して $|S|\le2$ とする．
$i\in T$ を選び，$J\supset S\cup\{i\}$，$K\supset T$ を三元集合に補う．
どちらの場合も $i\in J\cap K$ であり，少なくとも一方の実際の support に $i$ が属する．
補題\ref{topo:corner-singleton-radius}と包含を用いると
\[
 d_\square(x,y)
 \le d_J(x,e_i)+d_K(e_i,y)<2\sqrt5.
\]
片方の評価が真の不等式になることが，ここで必要である．
従って残るのは，両点の support がそれぞれ三元で互いに素な場合である．
座標を置換し，値 $1$ の座標をそれぞれ $1,4$ として
\[
 x=e_1+a e_2+b e_3\in Q_{1;23},\qquad
 y=e_4+c e_5+e e_6\in Q_{4;56},\qquad
 0\le a,b,c,e\le1
\]
と書く．

\emph{四通りの展開による上界．}
始点 square を平面の $[0,1]^2$ とし，終点 square を
$[-3,-2]\times[-1,0]$ に置く．三回の left step と一回の down step の
四つの順序は，次の edge-adjacent galleries として全て存在する：
\begin{align*}
 LLLD:&\quad Q_{1;23},\ Q_{1;34},\ Q_{4;13},\ Q_{4;36},\ Q_{4;56},\\
 LLDL:&\quad Q_{1;23},\ Q_{1;34},\ Q_{4;13},\ Q_{4;15},\ Q_{4;56},\\
 LDLL:&\quad Q_{1;23},\ Q_{1;34},\ Q_{1;45},\ Q_{4;15},\ Q_{4;56},\\
 DLLL:&\quad Q_{1;23},\ Q_{1;25},\ Q_{1;45},\ Q_{4;15},\ Q_{4;56}.
\end{align*}
各共有辺に沿う反転で展開すると，始点と終点の座標はそれぞれ
$(a,b)$ と $(-2-e,-c)$ となる．
この二点を結ぶ Euclidean segment が直線 $u=0,-1,-2$ と $v=0$ を横切る順序を取れば，
上の四つの gallery の少なくとも一つに segment 全体が含まれる．
格子頂点を通る場合は，二つの横切りのどちらを先に取ってもよい．
各 square 上の展開は等長なので
\begin{equation}\label{topo:eq:disjoint-square-upper}
 d_\square(x,y)^2
 \le (a+e+2)^2+(b+c)^2
 \le 4^2+2^2=20.
\end{equation}
最後の等号には $a=b=c=e=1$ が必要である．
したがって全体の直径は $2\sqrt5$ 以下で，等号点対は定理の候補に限られる．

\emph{八種類の三角形による下界．}
$d=6$，$A=\{1,2,3\}$，$B=\{4,5,6\}$ とする．
空でない $S\subset A\cup B$，$|S|\le3$ に対応する頂点に対し，
$a(S)=|S\cap A|$，$b(S)=|S\cap B|$ と置く．
各 square を singleton と triple の頂点を結ぶ対角線で二分し，
次の頂点値を各三角形上 affine に延長する：
\[
\begin{array}{c|rrrrrrrrr}
 (a,b)&(3,0)&(2,0)&(1,0)&(2,1)&(1,1)&(1,2)&(0,1)&(0,2)&(0,3)\\\hline
 \sqrt5\,h&-5&-3&-2&-1&0&1&2&3&5.
\end{array}
\]
三角形は singleton から pair，triple への subset chain であり，
加える label の $A/B$ 型は八通りである．
対応する $\sqrt5\,h$ の二つの直交辺方向の増分は
\[
\begin{array}{c|rrrrrrrr}
 \text{型}&AAA&AAB&ABA&ABB&BAA&BAB&BBA&BBB\\\hline
 \text{増分}&(-1,-2)&(-1,2)&(2,-1)&(2,1)&(-2,-1)&(-2,1)&(1,-2)&(1,2).
\end{array}
\]
従って全三角形上で $|\nabla h|^2=(1^2+2^2)/5=1$ であり，
$h$ は貼合せの長さ距離に関して $1$-Lipschitz である．
しかし
\[
 h(v_B)-h(v_A)=\frac{10}{\sqrt5}=2\sqrt5
\]
なので $d_\square(v_A,v_B)\ge2\sqrt5$ となる．
命題\ref{topo:coordinate-retraction}の等長性により，この下界は
$d>6$ の中の任意の六 label にも保たれる．
上界と座標置換対称性を合わせると，直径と全等号点対の記述を得る．
\end{proof}

\begin{corollary}[六次以上の直径対の線形安定性]
\label{topo:diameter-six-stability}
$d\ge6$，$0\le\delta<\sqrt{20}-\sqrt{13}$ とし，
$d_\square(x,y)\ge\sqrt{20}-\delta$ とする．
すると $A=\operatorname{supp}(x)$ と $B=\operatorname{supp}(y)$ は
互いに素な三元集合であり，
\[
 d_\square(x,v_A)+d_\square(y,v_B)
 \le\sqrt{20}\,\delta-\frac{\delta^2}{2}
 \le\sqrt{20}\,\delta.
\]
\end{corollary}
\begin{proof}
命題\ref{topo:coordinate-retraction}と，既存の五次の直径定理
\cite[\texttt{thm:intro-diameter-C5}, \texttt{sec:finite-diameter}]{KECubicalCurrent}
により，support の和が五元以下なら距離は $\sqrt{13}$ 以下となるので，
両 support は互いに素な三元集合である．上の座標を用い，
$u=2-a-e$，$v=2-b-c$ とすると $0\le u,v\le2$ であり，
\[
 20-\bigl((a+e+2)^2+(b+c)^2\bigr)
 =(8u-u^2)+(4v-v^2)\ge2(u+v).
\]
\eqref{topo:eq:disjoint-square-upper}から
$u+v\le(20-d_\square(x,y)^2)/2$ となる．
各点から対応する頂点までの square 内の線分の長さは，
それぞれ $2-a-b$，$2-c-e$ 以下なので，その和は $u+v$ 以下である．
距離の仮定を代入すれば結論を得る．
\end{proof}

\begin{remark}[五次の直径の使用範囲]
定理\ref{topo:diameter-six-labels}の直径と全等号点対の証明は，
上の展開・calibration と座標 retraction だけで自己完結し，五次の厳密直径を使わない．
既存の五次の $\sqrt{13}$ を使用するのは，
系\ref{topo:diameter-six-stability}の明示的な gap
$\sqrt{20}-\sqrt{13}$ の箇所に限る．
Actual KE 曲面の直径漸近を得るには，全曲面の distortion が消えることを
別に証明しなければならない．
\end{remark}

\subsection{分岐基底の立方体距離への定量的収縮}

上の $p_T$ は同相でなくても，基底上に\emph{人工的な}長さ距離を
入れると KE 側の候補立方体複体へ収束させられる．
これは収縮する方向を正確に特定する位相幾何学的な模型であり，
実際の Lagrangian fibration が誘導する action metric の収束ではない．

\begin{theorem}[有限分岐基底の立方体計量への収縮]
\label{topo:metric-collapse}
$d\ge5$ と三価木の系 $T=(T_{ij})$ を固定する．
$\mathcal C_d$ に各 unit square の Euclidean length metric
$d_\square$ を入れ，$p_T:B(T)\to\mathcal C_d$ を
系\ref{topo:KE-cubical-comparison} の全射とする．
$B(T)$ の，セルと $p_T$ に適合する有限細分上に
任意の補助 piecewise Euclidean metric $h$ を固定する．
$\varepsilon\in(0,1]$ に対して，各細分セルの接ベクトル $v$ に
\[
 |v|_\varepsilon^2
 =|d p_T(v)|_\square^2+\varepsilon^2|v|_h^2
\]
を与え，これを貼った内在的長さ距離を $d_\varepsilon$ とする．
このとき $C=C(d,T,h)<\infty$ が存在し，任意の
$x,y\in B(T)$ について
\[
 0\le d_\varepsilon(x,y)
       -d_\square(p_Tx,p_Ty)\le C\varepsilon .
\]
従って
\[
 d_{\mathrm{GH}}\bigl((B(T),d_\varepsilon),
                  (\mathcal C_d,d_\square)\bigr)
 \le \frac{C\varepsilon}{2}\longrightarrow0 .
\]
$d\ge6$ では各 $\varepsilon>0$ の分岐基底は
$\mathcal C_d$ と同相ではないにもかかわらず，この収束は成立する．
\end{theorem}

\begin{proof}
$\rho_d\circ p_T=q$ は各三角形セル $\tau_{ijk}\subset B(T)$
から $S_d$ の三角形 $[ijk]$ への homeomorphism である．
$\rho_d^{-1}$ は $x\mapsto x/\max_i x_i$ なので，
$[ijk]$ を三つの最大座標領域に分ければ smooth かつ Lipschitz である．
したがって $p_T|_{\tau_{ijk}}$ の逆写像は
$\mathcal C_d$ の対応する三つの unit squares 上で
区分的 Lipschitz となり，補助 metric $h$ に関する
一様定数 $L=L(d,T,h)$ を持つ．
$p_T^{-1}(z)$ は，$z$ が support 二の辺の内部にあるとき，
対応する $T_{ij}^{\rm int}$ の有限木（longitudinal 座標を固定した
各二角形の横断部分）であり，それ以外は一点である．
従って全ての fiber の内在的な $h$-直径は一様に
$D=D(d,T,h)$ 以下である．Fiber 内の任意の接続曲線の
$d_\varepsilon$-長は高々 $\varepsilon D$ にできる．

各 unit square は ambient $\mathbb R^d$ 内の凸集合である．
$\mathcal C_d$ の長さ最短路が同じ square を離れて再び戻るなら，
その間を square 内の線分で置き換えると短くなる
（元の任意の曲線の長さは ambient Euclidean chord 以上）．
従って最短路は，少なくとも一つの選択では，
有限個の square の gallery として記述でき，
系\ref{topo:six-label-gallery}により，その訪問数は
$N_d=3\binom{\min\{d,6\}}{3}$ 以下に取れる．
Edge に沿う部分もそれを含む任意の square に割り当てればよい．

$p_Tx$ と $p_Ty$ を結ぶこの gallery の各 square 内の線分を，
対応する $\tau_{ijk}$ の $p_T$ による逆像へ持ち上げる．
その $h$-長は元の線分長の $L$ 倍以下なので，
\[
 L_\varepsilon(\widetilde\gamma_{\rm square})
 \le L_\square(\gamma_{\rm square})
      +\varepsilon L\,L_\square(\gamma_{\rm square}).
\]
異なる square に移る点では，持ち上げた端点同士は
同じ $p_T$-fiber にある．この二点をその木の中で
$\varepsilon D$ 以下の費用で接続する．
最初と最後にも同様の fiber connector を使う．
全誤差は
\[
 \varepsilon\bigl(L\operatorname{diam}(\mathcal C_d)
                          +(N_d+1)D\bigr)
\]
以下である．逆向きの不等式は
$|dp_T(v)|_\square\le |v|_\varepsilon$ から従う．
$p_T$ は全射なので，その graph correspondence の distortionは
$C\varepsilon$ 以下となり，GH評価を得る．
\end{proof}

\begin{proposition}[収縮速度の次数による鋭い差]\label{topo:metric-collapse-sharp}
定理\ref{topo:metric-collapse}と同じ$p_T,h,d_\varepsilon$を固定し，
$p_T$に付随するgraph correspondenceのdistortionを
\[
 \mathcal D_\varepsilon=
 \sup_{x,y\in B(T)}
 \left|d_\varepsilon(x,y)-d_\square(p_Tx,p_Ty)\right|
\]
とする．$\varepsilon\downarrow0$に対し，
\[
 \mathcal D_\varepsilon=
 \begin{cases}
  \Theta_{T,h}(\varepsilon^2),&d=5,\\
  \Theta_{d,T,h}(\varepsilon),&d\ge6.
 \end{cases}
\]
従って五次では
$d_{\mathrm{GH}}((B(T),d_\varepsilon),(\mathcal C_5,d_\square))
=O_h(\varepsilon^2)$ に改善できる．
$d\ge6$ での下界は\emph{このgraph correspondenceのdistortion}に
対するものであり，別のcorrespondenceによって定義される
$d_{\mathrm{GH}}$の同じ次数の下界を主張しない．
\end{proposition}
\begin{proof}
任意の$h$-rectifiable曲線$\gamma$について積分の三角不等式を用いると，
\[
 L_\varepsilon(\gamma)
 =\int\sqrt{|(p_T\gamma)'|_\square^2
                  +\varepsilon^2|\gamma'|_h^2}
 \ge\sqrt{L_\square(p_T\gamma)^2+
                          \varepsilon^2L_h(\gamma)^2}.
\]
従って全$x,y$に対し
\begin{equation}\label{topo:metric-lower-mixed}
 d_\varepsilon(x,y)
 \ge\sqrt{d_\square(p_Tx,p_Ty)^2+\varepsilon^2d_h(x,y)^2}.
\end{equation}
$d\ge6$なら$d-4\ge2$であり，少なくとも一つの内部木
$T_{ij}^{\rm int}$ は正の長さを持つ．同じ$p_T$-fiberの相異なる
$x,y$を固定すると，\eqref{topo:metric-lower-mixed}により
$\mathcal D_\varepsilon\ge\varepsilon d_h(x,y)>0$である．
上界は定理\ref{topo:metric-collapse}による．

$d=5$では全ての内部木は一点であり，$p_T$は同相である．
有限細分上で$p_T$の逆はpiecewise Lipschitzなので，ある$L<\infty$に対し
$|v|_h\le L|dp_T(v)|_\square$が全細分セルで成り立つ．
$\mathcal C_5$内の測地線をこの同相の逆で持ち上げれば
\[
 d_\varepsilon(x,y)\le
   \sqrt{1+\varepsilon^2L^2}\,d_\square(p_Tx,p_Ty)
\]
となる．この上界と$d_\square$の直径の有限性から
$\mathcal D_\varepsilon=O(\varepsilon^2)$を得る．
また任意の相異なる$x,y$を固定すると，
\eqref{topo:metric-lower-mixed}と
$d_\square(p_Tx,p_Ty)>0,\ d_h(x,y)>0$により
$\mathcal D_\varepsilon\ge c_{x,y,h}\varepsilon^2$
が十分小さい$\varepsilon$に対して成立する．
\end{proof}

\begin{proposition}[有限squareに対するthree-input比較原理]
\label{topo:finite-portal-criterion}
$T$を固定し，$p_T:B(T)\to\mathcal C_d$を上の写像とする．
$B(T)$上の任意の長さ距離$d_t$に対し，非負の定数
$\alpha_t,\sigma_t,\lambda_t,\kappa_t$が次を満たすと仮定する．
\begin{enumerate}
\item[\textup{(N)}]
任意の$x,y\in B(T)$に対し
$d_t(x,y)\ge d_\square(p_Tx,p_Ty)-\alpha_t$
（大域no-shortcut）．
\item[\textup{(S)}]
任意のunit square $Q\subset\mathcal C_d$と任意の線分
$\gamma\subset Q$に対し，$p_T$の下で$\gamma$に一致する
rectifiable lift $\widetilde\gamma$が存在し，
\[
 L_{d_t}(\widetilde\gamma)\le
 (1+\sigma_t)L_\square(\gamma)+\lambda_t
\]
（各squareの水平持ち上げ）．
\item[\textup{(F)}]
全ての$z\in\mathcal C_d$で
$\operatorname{diam}_{d_t}p_T^{-1}(z)\le\kappa_t$
（分岐木fiber・portalの接続費用）．
\end{enumerate}
$N_d=3\binom{\min\{d,6\}}{3}$, $D_d=\operatorname{diam}(\mathcal C_d,d_\square)$
とすると
\[
 \sup_{x,y\in B(T)}
 |d_t(x,y)-d_\square(p_Tx,p_Ty)|
 \le\max\{\alpha_t,\,
       \sigma_tD_d+N_d\lambda_t+(N_d+1)\kappa_t\}.
\]
右辺が零へ行けば
$(B(T),d_t)\to(\mathcal C_d,d_\square)$ が
Gromov--Hausdorff収束する．
\end{proposition}
\begin{proof}
下界は\textup{(N)}．上界には$p_Tx$と$p_Ty$を結ぶ
系\ref{topo:six-label-gallery}の最短路を取り，$N_d$個以下の
square内線分からなるgalleryを選ぶ．
各線分を\textup{(S)}で持ち上げ，隣り合う線分の持ち上げ端点が
属する同一fiberの二点を，\textup{(F)}で評価される ambient 距離を
任意の正誤差で近似する経路で接続し，最初と最後も$x,y$へ接続する．
これらの経路が fiber 内に留まる必要はない．近似の誤差を零へ送れば，$d_t(x,y)$ は
\[
 (1+\sigma_t)d_\square(p_Tx,p_Ty)
   +N_d\lambda_t+(N_d+1)\kappa_t
\]
以下である．$p_T$は全射なので，graph correspondenceの
distortion boundからGH収束が従う．
\end{proof}

\begin{remark}[端点だけを用いる有限 gallery の上界]
\label{topo:finite-endpoint-comparison}
上界には target 線分へ厳密に写る lift を必要としない．一般の距離空間
$(X,d_X)$，写像 $F:X\to\mathcal C_d$ と非負定数 $a,\lambda,\kappa$ に対し，
各 $x$ の anchor $q_x$ を
$d_\square(F(x),q_x)\le a$ となるように選ぶ．各 unit square $Q$ と $z\in Q$ に
portal $p_Q(z)\in X$ を選び，
\[
 \begin{aligned}
 d_X(p_Q(z),p_Q(z'))&\le |z-z'|+\lambda &&(z,z'\in Q),\\
 d_X(p_Q(z),p_{Q'}(z))&\le\kappa &&(z\in Q\cap Q'),\\
 d_X(x,p_Q(q_x))&\le\kappa &&(q_x\in Q)
 \end{aligned}
\]
が成立するとする．点の選択に連続性は不要であり，一つの比較に必要な有限個の
端点に対して整合的に選べれば足りる．このとき，前命題と同じ $N_d$ に対し
\[
 d_X(x,y)\le d_\square(F(x),F(y))
       +N_d\lambda+(N_d+1)\kappa+2a .
\]
実際，$q_x,q_y$ を結ぶ $M\le N_d$ 本の最短 gallery に第一の不等式を使い，
$M-1$ 個の junction と初め・終わりに残る二つの不等式を使う．
最後に anchor の誤差 $2a$ を加えればよい．$q_x=q_y$ の場合は一つの square の
定値線分を用いる．全て ambient 距離に対する三角不等式なので，
portal 間の経路を fiber 又は局所領域の中へ閉じ込める仮定も不要である．
\end{remark}

\begin{example}[各squareの局所モデルだけでは不十分]
\label{topo:portal-counterexample}
$Q=[0,1]^2$ の二つのコピー$Q_1,Q_2$を，原点だけで貼り合わせた
長さ空間$Y=Q_1\vee_{\{0\}}Q_2$を考える．
両コピー上でidentityとなる折り畳み写像$p:Y\to Q$を取る．
この写像は非伸長なので，命題\ref{topo:finite-portal-criterion}の
\textup{(N)}は$\alpha=0$で成立する．
各squareの直線はどちらのコピーにも等長に持ち上がるので，
\textup{(S)}も$\sigma=\lambda=0$で成立する．
しかし二つの右上頂点は同じ$p$-fiberにあり，
$Y$での距離は$2\sqrt2$である．
従って \textup{(F)} は破れ，
\[
 \operatorname{diam}Y=2\sqrt2,\quad
 \operatorname{diam}Q=\sqrt2,\quad
 d_{\mathrm{GH}}(Y,Q)\ge\frac{\sqrt2}{2}.
\]
特に有限個の局所的なsquare isometriesと大域no-shortcutだけから
global GH収束は導けない．
\end{example}

\begin{example}[飽和した局所領域の外側を通る薄い shortcut]
\label{topo:protected-thin-bypass}
局所の protected distance，square の持上げ上界，小さい全 fibers だけでは，
全域 no-shortcut は従わない．この点を一般の compact length space で確認する．
Target を二枚の unit square からなる Euclidean rectangle
\[
 Y=[0,1]_s\times[-1,1]_u
\]
とし，$0<\varepsilon<1/4$ に対して source は同じ集合に
\[
 g_\varepsilon=du^2+a_\varepsilon(u)^2ds^2
\]
を入れた $X_\varepsilon$ とする．$a_\varepsilon$ は
$\varepsilon\le a_\varepsilon\le1$ を満たす滑らかな関数で，
$u\le\varepsilon/3$ と $u\ge4\varepsilon/3$ では $1$，
$2\varepsilon/3\le u\le\varepsilon$ では $\varepsilon$ に等しいものを取る．
境界と corner を許した Riemannian length metric を用い，
$f_\varepsilon:X_\varepsilon\to Y$ は集合としての恒等写像とする．

全 fibers は一点なので，その ambient 径は零である．
$g_\varepsilon\le g_Y$ より，target の各 square 内線分は恒等写像で
長さを増やさず持ち上がる．従って\textup{(S)}の誤差と\textup{(F)}の誤差は共に零である．
また
\[
 W=[0,1]\times[-1,0],\qquad
 V_\varepsilon=[0,1]\times[-1,\varepsilon/3]
\]
は $f_\varepsilon$-飽和であり，その上では metric は Euclidean なので
\[
 d_W^{\rm int}(x,y)=d_{V_\varepsilon}^{\rm int}(x,y)
 =d_Y(f_\varepsilon(x),f_\varepsilon(y))
 \qquad(x,y\in W).
\]
しかし $x=(0,0)$，$y=(1,0)$ を，上端の $u=\varepsilon$ まで上がり，
そこで水平移動して下りる経路で結ぶと
\[
 d_{X_\varepsilon}(x,y)\le2\varepsilon+\varepsilon=3\varepsilon,
 \qquad d_Y(f_\varepsilon(x),f_\varepsilon(y))=1.
\]
従って全域 no-shortcut の誤差は零へ行かない．

これは actual KE 退化の反例ではなく，局所的な protected 比較と
全域の\textup{(S),(F)}から\textup{(N)}を形式的に推論できないことを示す例である．
局所領域 $W$ 自体の飽和性も満たすが，上半分には同じ protected 比較を与えていない．
従って，飽和した\emph{全域被覆}上の局所距離比較と十分な overlap 制御を
さらに仮定した命題を否定する例とは扱わない．
本研究でも，全域被覆・全 overlap・actual 下界の確認を省略しない．
\end{example}

\begin{remark}[実際の幾何計量への接続に残るもの]
定理\ref{topo:metric-collapse}の人工計量では
\textup{(N)}が$\alpha=0$で成り立ち，
\textup{(S)}の誤差は$O(\varepsilon)$，
\textup{(F)}も$O(\varepsilon)$である．
しかしactual Lagrangian action metricやKE退化計量が
\textup{(N),(S),(F)}を満たすことは未証明である．
特に局所モデルの各squareに距離比較があるだけでは
分岐木fiber全体とportalを短く接続できるとは限らない．
逆に，$p_T$のgraph correspondenceが距離distortion
$\delta_t\to0$を持つなら，同一fiberの任意の二点$x,y$について
$d_t(x,y)\le\delta_t$なので，\textup{(F)}の$\kappa_t\to0$は
\emph{$p_T$を漸近写像に固定する限り必要}である．
例\ref{topo:portal-counterexample}は，
局所のisometric liftsと大域no-shortcutが揃っていても
この必要条件を得られないことを示す．
この一般原理の新規性は主張せず，対象研究の未解決な
幾何学的入力を明確化するために用いる．
\end{remark}

\begin{remark}[計量収縮の範囲と新規性]
これは $T$ を固定した有限 PL/CW 複体に対する直接構成で，
$\varepsilon$ は幾何から導いた parameter ではなく任意に入れた値である．
従って actual action-affine metric，KE計量，symplectic transport の
いずれについても新たな収束定理を主張しない．
この種の有限複体のcollapse自体の先行研究との新規性は未確認である．
\end{remark}

\subsection{KE 近似写像と Lagrangian fibration の適合性}

\begin{proposition}[近似写像の一様比較による fiber 収縮]
\label{topo:ke-map-compatibility}
$T$ を固定し，$p_T:B_d(T)\to\mathcal C_d$ を
系\ref{topo:KE-cubical-comparison}の全射とする．
コンパクト距離空間 $(X_t,\bar d_t)$，全射
$\pi_t:X_t\to B_d(T)$，写像 $\Phi_t:X_t\to\mathcal C_d$ が与えられ，
\[
 \sup_{x,y\in X_t}
 \bigl|\bar d_t(x,y)-d_\square(\Phi_t(x),\Phi_t(y))\bigr|\le a_t,
 \qquad
 \sup_{x\in X_t}d_\square(\Phi_t(x),p_T\pi_t(x))\le b_t
\]
を満たすとする．$f_t=p_T\circ\pi_t$ と置くと，
\[
 \operatorname{dis}(f_t)\le a_t+2b_t,\qquad
 d_{\mathrm{GH}}\bigl((X_t,\bar d_t),(\mathcal C_d,d_\square)\bigr)
 \le \frac{a_t+2b_t}{2},
\]
および
\[
 \sup_{c\in\mathcal C_d}\operatorname{diam}_{\bar d_t}f_t^{-1}(c)
 \le a_t+2b_t,\qquad
 \sup_{b\in B_d(T)}\operatorname{diam}_{\bar d_t}\pi_t^{-1}(b)
 \le a_t+2b_t
\]
が成り立つ．直径は全空間の距離を制限したものであり，
fiber 内の経路長距離ではない．特に $a_t,b_t\to0$ なら $f_t$ は GH 近似を与える．
\end{proposition}
\begin{proof}
任意の $x,y$ に対して三角不等式から
\[
 \bigl|d_\square(\Phi_t(x),\Phi_t(y))
       -d_\square(f_t(x),f_t(y))\bigr|\le2b_t
\]
なので，距離歪みは $a_t+2b_t$ 以下である．
$f_t$ は全射だから，その graph は両空間間の correspondence であり，
GH 距離の評価が従う．$f_t(x)=f_t(y)$ の場合に同じ評価を使うと
$\bar d_t(x,y)\le a_t+2b_t$ を得る．
各 $\pi_t$-fiber はある $f_t$-fiber に含まれるので最後の評価も従う．
\end{proof}

\begin{remark}[実計量への適合性と特異 fibers]
KE 退化では $\bar d_t=R_t^{-1}d_{g_t}$，
$R_t=\log(-\log|t|^2)$ を取る．
関連研究\cite{KECubicalCurrent}の $\mathrm{(CP)}_d+\mathrm{(PBE)}_d$ は
適切な $\Phi_t$ に対する $a_t\to0$ を条件付きで与えるが，
本稿の symplectic construction から $b_t\to0$ はまだ示されていない．
Moser transport の一様な距離制御も別途必要であり得る．
本命題で $a_t,b_t\to0$ が成り立てば，pinched torus，
$S^1\times\figeight$，core skeleton を含め，全 fiber の縮小を結論する．
一方，regular fibers の点ごとの縮小や，特異値集合の低次元性だけでは
この一様評価を得られない．また $d\ge6$ では，
$p_T$ の同じ木 fiber 上にある複数の $\pi_t$-fibers の集合全体も潰す必要がある．
この比較原理に一般的新規性は主張しない．
\end{remark}

\section{六次の場合の同時 Petersen 化}\label{topo:petersen-section}

$d=6$ のとき，各 $T_{ij}$ は四つの葉を二組に分けるだけで定まり，
各 $G_i$ は $K_5$ の各頂点を二個の三価頂点に展開したグラフである．
分割を独立に選ぶとリンクの同型型は一定ではないが，六つのリンクを同時に
Petersen グラフにできる．しかも，そのような分割系は完全に分類できる．

\begin{definition}
$K_6$ の完全マッチングは，六つのラベルを三組に分けたものである．
$K_6$ の one-factorization は，その十五本の辺を五個の完全マッチングに分解したものである．
\end{definition}

\begin{theorem}\label{topo:petersen}
ラベルを固定した $d=6$ の分割系で，全ての $G_i$ が Petersen グラフとなるものは
ちょうど六個である．それらは $K_6$ の六個の one-factorization と自然に対応する．
対応は，辺 $ij$ を含む完全マッチングの残りの二組を $T_{ij}$ の分割に取ることで与えられる．
ラベルの置換を許せば，この六個は一つの同型類をなす．
\end{theorem}

まず存在とその対称性を明示する．六つのラベルを
$\{\infty\}\cup\mathbb F_5$ とし，$c\in\mathbb F_5$ に対して
\[
 M_c=\bigl\{\{\infty,c\},\{c+1,c-1\},\{c+2,c-2\}\bigr\}
\]
と置く．五個の $M_c$ は one-factorization $\mathcal F$ をなす．
$i=\infty$ のリンクでは，$T_{\infty c}$ の二つの頂点を，葉 $c\pm1$ に接する
$A_c$ と，葉 $c\pm2$ に接する $B_c$ と書ける．その辺は
\[
 A_cB_c,\qquad A_cA_{c+1},\qquad B_cB_{c+2}\quad(c\in\mathbb F_5)
\]
であり，これは Petersen グラフの標準表示である．
平行移動 $c\mapsto c+a$ は $\mathcal F$ を保ち，$\infty$ を固定する．
反転 $c\mapsto c^{-1}$（$0\leftrightarrow\infty$）も $\mathcal F$ を保つ．
実際，反転は $M_0$ を固定し，$M_1\leftrightarrow M_3$，
$M_2\leftrightarrow M_4$ と入れ替える．したがって，この分割系の対称群は
六つのラベル上で推移的に作用し，全てのリンクが Petersen グラフになる．

\begin{lemma}\label{topo:factorizations}
$K_6$ の完全マッチングは十五個あり，各完全マッチングはちょうど二個の
one-factorization に属する．この二個はその完全マッチング以外を共有せず，
その中の任意の一辺の両端を入れ替える置換で互いに移る．
したがって one-factorization は六個で，互いにラベルの置換で同型である．
\end{lemma}

\begin{proof}
完全マッチングの数は $5\cdot3\cdot1=15$ である．一個を固定し，その三組を
$\{A_0,A_1\}$，$\{B_0,B_1\}$，$\{C_0,C_1\}$ と書く．
これと辺を共有しない完全マッチングは
\[
 \{A_aB_b,\ B_{1-b}C_c,\ C_{1-c}A_{1-a}\},
 \qquad(a,b,c)\in\{0,1\}^3
\]
の八個である．二個が辺を共有するのは，添字の三成分のうち相異なるものが
ちょうど一個のときである．したがって，残りを四個の互いに素な完全マッチングに
分解する問題は，三次元 cube の四頂点の独立集合を選ぶ問題になる．
そのような独立集合は偶数パリティと奇数パリティの二個だけである．
例えば $000$ を含むものは，$111$ を含めると大きさ二以下になるので，
残る三個は $011,101,110$ に限られる．
一組の両端交換は一個の添字を反転し，この二個を入れ替える．
従って one-factorization の総数は $15\cdot2/5=6$ である．
固定マッチングのラベルの置換と両端交換を使えば，全てが同型であることも分かる．
\end{proof}

\begin{lemma}\label{topo:local-petersen}
一つのラベル $i$ を固定する．$G_i$ が Petersen グラフであるための必要十分条件は，
$j\ne i$ に対する五個の完全マッチング
\[
 N_{ij}=\{ij\}\cup\{T_{ij}\text{ の二組}\}
\]
が one-factorization をなすことである．
\end{lemma}

\begin{proof}
十分性は，補題\ref{topo:factorizations}の同型性と上の明示的構成から従う．
必要性のため，$G_i$ の五本の内部辺が与える完全マッチングを除く．
Petersen グラフの任意の完全マッチングの補集合は，二個の五角形である．
この事実は標準表示で直接確認できる．完全マッチングに属する spoke の数は奇数である．
三本の場合は，残る二つの添字が外側の五角形と内側の星の両方で隣接する必要があり，不可能である．
五本の場合は明らかで，一本の場合は残りの外側と内側の四頂点のマッチングが一意に決まり，
その補集合も二個の五角形になる．
Petersen グラフには三角形がないので，完全マッチングの辺はこの五角形の弦になれず，
全て二つの五角形の間を結ぶ．内部辺を縮約すると，これらの二個の五角形は
$K_5$ の辺を互いに素に分割する．
従って，一方の五角形に沿って残る五ラベルを巡回的に付け直せば，
分割系は上の $A_c,B_c$ の表示に一致する．この表示の $N_{ij}$ は
one-factorization である．
\end{proof}

\begin{proof}[定理\ref{topo:petersen}の証明]
各リンクが Petersen であると仮定し，補題\ref{topo:local-petersen}による
頂点 $i$ の one-factorization を $\mathcal F_i$ とする．
$\mathcal F_i$ と $\mathcal F_j$ は $N_{ij}$ を共有する．
ラベル $p$ を固定し，$\mathcal F=\mathcal F_p$ と置くと，
補題\ref{topo:factorizations}より，各 $j\ne p$ に対し
\[
 \mathcal F_j=\mathcal F\quad\text{又は}\quad
 \mathcal F_j=(pj)\mathcal F
\]
である．ここで $(pj)$ は両端交換の置換である．

ある $a$ について $\mathcal F_a=(pa)\mathcal F\ne\mathcal F$ だったとする．
$b\ne p,a$ を取る．もし $\mathcal F_b=\mathcal F$ なら，
$\mathcal F_a\cap\mathcal F_b$ の唯一の完全マッチングは $pa$ を含み，
$ab$ を含めないので，$N_{ab}$ の存在に矛盾する．
残る場合は $\mathcal F_b=(pb)\mathcal F$ である．
$N_{ab}$ を $(pa)$ で引き戻したものは，$\mathcal F$ に属し $pb$ を含むため，
ある残りの三ラベル $c,e,f$ を使って
\[
 N=\{pb,ac,ef\}
\]
と書ける．一方 $N_{ab}\in(pb)\mathcal F$ から
\[
 (pb)(pa)N=\{pa,bc,ef\}\in\mathcal F
\]
も従う．これは $N$ と異なるのに辺 $ef$ を共有するので矛盾である．
したがって全ての $\mathcal F_i$ は同一である．
逆向きは既に示した．六個という数とラベル置換の下での一意性は
補題\ref{topo:factorizations}から従う．
\end{proof}

\begin{corollary}[六次の同時 Petersen 系におけるラベル対称性の破れ]
\label{topo:petersen-symmetry}
$d=6$ とし，$\mathfrak T_6^{\rm Pet}$ を全ての頂点リンクが
Petersen となるラベル付き分割系の集合とする．
$|\mathfrak T_6^{\rm Pet}|=6$ であり，
$S_6$ はラベル置換によってこの集合上に推移的に作用する．
各 $T\in\mathfrak T_6^{\rm Pet}$ の安定化部分群の位数は $120$，
指数は $6$ である．
特に全 $S_6$ で不変な単一のラベル付き Petersen 分割系は存在しない．
\end{corollary}
\begin{proof}
定理\ref{topo:petersen}より全ての系は一つのラベル置換軌道に属し，
軌道の大きさは六である．従って軌道・安定化群公式から
$|\operatorname{Stab}_{S_6}(T)|=6!/6=120$ となる．
全 $S_6$ に固定された系は軌道が一点になるはずなので存在しない．
\end{proof}

\begin{remark}
KE候補$\mathcal C_6$の座標置換は全$S_6$で対称である一方，
各Petersen分岐基底の\emph{選択した木系}は
その対称性を保たない．従ってラベル付きPetersen構成の比較に
全$S_6$-equivarianceを課すなら，単一の$T$ではなく
六個の選択の置換を同時に管理する必要がある．
これは別の分岐基底やラベルを忘れた幾何構成の存在を排除しない．
\end{remark}

この同時 Petersen 化は，六次の場合の全ての core の\emph{位相的なリンク}を統一する．
それだけで格子付きファン，integral-affine atlas，あるいは対応する
シンプレクティックな core の同定を与えるわけではない．

\input{sections/paper/ke_fiber_compatibility.tex}
\input{sections/paper/ke_open_cell_collapse.tex}
\input{sections/paper/ke_closed_corner_gh.tex}
% !TeX root = ../../surface-lagrangian-torus-fibration.tex
% 2026-11-23 / v1.10: Source-inclusive 正規形の共通像・正逆評価・全 orbit の定義域．
\section{共通の normal circle に適合する二段階平均}
\label{ke:staged-neck-section}

三重点の torus 平均を，neck 全体で共通に用いる normal circle の平均と，
end collar 上の残る circle の平均に分ける．
この順序によって，完成後の end tori が一つの actual Hamiltonian circle を含むことを
厳密に保証できる．単一 primitive を用いると normal moment の加法定数と
全 marked cycle の作用周期も同時に固定される．
平均化と Moser 法そのものの先行研究は
\cite[Section~3]{RuanHMinimal}を参照されたい．
以下の二段階構成と定量的適用の包括的な新規性確認は未完了である．

本節の円作用の周期はすべて $1$ とし，moment の符号は
\[
 \iota_V\omega=-d\mu
\]
に統一する．従って moment は一周の primitive 積分と同じ規格化である．
周期 $2\pi$ の生成場による通常の action 座標を用いる場合は，
生成場・moment・周期積分の規格化を同時に変更する．
また，以下の flow は指定する内側領域から出発して共通の外側領域に留まると仮定する．
作用・形式が一致する開集合を出る trajectory の一致は，局所的な ODE の一意性からは従わない．

\subsection{一つの円の平均と primitive の正確な輸送}

$U$ を複素曲面，$\omega=d\lambda$ を K\"ahler form とし，
$\rho_s$ を $U$ 上の正則円作用とする（固定点を許す）．
生成場を $V$，平均と chain homotopy を
\[
 \mathcal A\xi=\int_0^1\rho_s^*\xi\,ds,\qquad
 K\xi=\int_0^1(1-s)\rho_s^*(\iota_V\xi)\,ds
\]
と書く．Centered kernel $1/2-s$ に置き換えた場合も，
以下の恒等式と結論は同じである．

\begin{lemma}[円平均の exact Moser transport]
\label{ke:circle-exact-transport}
$\eta=\mathcal A\omega$，
$\omega_s=(1-s)\eta+s\omega$，
$\lambda_s=(1-s)\mathcal A\lambda+s\lambda$ とする．
$\beta=K\omega$ および
\[
 \iota_{Y_s}\omega_s=\beta,\qquad \phi_0=\mathrm{id}
\]
による flow が全時刻で定義される $\rho$-飽和した内側領域を取る．この領域では
\begin{equation}\label{ke:eq:circle-exact-transport}
 \phi_s^*\omega_s=\eta,\qquad
 \phi_s^*\lambda_s=\mathcal A\lambda+dq_s,
 \quad
 q_s=\int_0^s\phi_\tau^*
       \bigl(\lambda_\tau(Y_\tau)-K\lambda\bigr)\,d\tau .
\end{equation}
$\mu=(\mathcal A\lambda)(V)$ は $\eta$ の Hamiltonian moment である．
$\phi=\phi_1$ によって移した円作用
$\widehat\rho_s=\phi\circ\rho_s\circ\phi^{-1}$ は actual $\omega$ に関して Hamiltonian で，
その moment は
\begin{equation}\label{ke:eq:circle-actual-moment}
 \widehat\mu=\mu\circ\phi^{-1}.
\end{equation}
この moment の値は，対応する actual 円軌道上の $\lambda$ の一周積分に等しい．
\end{lemma}
\begin{proof}
Cartan の公式と部分積分により
\[
 dK+Kd=\mathcal A-I.
\]
従って
\[
 d\beta=\eta-\omega,\qquad
 \beta=\mathcal A\lambda-\lambda-d(K\lambda).
\]
正則作用による平均とその凸結合は K\"ahler form なので非退化である．
Moser の微分式から最初の pullback identity を得る．
Primitive については
\[
 \partial_s\lambda_s+\mathcal L_{Y_s}\lambda_s
 =\lambda-\mathcal A\lambda+\beta
     +d\bigl(\lambda_s(Y_s)\bigr)
 =d\bigl(\lambda_s(Y_s)-K\lambda\bigr)
\]
であり，flow で引き戻して積分すれば
\eqref{ke:eq:circle-exact-transport}の第二式となる．

$\mathcal A\lambda$ は円作用に不変なので
\[
 \iota_V\eta=-d\bigl((\mathcal A\lambda)(V)\bigr)=-d\mu.
\]
これを $\phi$ で移すと\eqref{ke:eq:circle-actual-moment}を得る．
周期 $1$ の軌道 $s\mapsto\rho_s(x)$ 上で
$(\mathcal A\lambda)(V)=\mu(x)$ は一定である．
Exact differential の周期は零なので，
$\int_{\phi(\rho_\bullet(x))}\lambda=\mu(x)$ である．
\end{proof}

\begin{lemma}[Actual circle の短さによる双方向 Moser buffer]
\label{ke:short-circle-actual-buffer}
上記の正則円作用と K\"ahler form に対し，作用不変な開集合 $U$ 全体で
$|V|_{g_\omega}\le\varepsilon$ と仮定する．
この補題の評価には $\omega$ の exactness，作用不変な背景計量，
tensor の二側比較は不要である．
$w(a)=1-a$ または $w(a)=1/2-a$ とし，
\[
 \beta=\int_0^1 w(a)\rho_a^*(\iota_V\omega)\,da,
 \quad \eta=\mathcal A\omega,
 \quad \omega_s=(1-s)\eta+s\omega,
 \quad \iota_{Y_s}\omega_s=\beta
\]
と置く．$c_w=(\pi/2)\|w\|_{L^2(0,1)}$ とすれば，
\begin{equation}\label{ke:eq:circle-actual-speed}
 |Y_s|_{g_\omega}
 \le\frac{\varepsilon\|w\|_{L^2(0,1)}}{2\sqrt{s(1-s)}}
 \qquad (0<s<1).
\end{equation}
従って，任意の二つの時刻を結ぶ存在する trajectory の actual 長さは
\begin{equation}\label{ke:eq:circle-actual-length}
 \operatorname{Length}_{g_\omega}(\phi_{s,s_0}(x))
 \le c_w\varepsilon,
 \qquad
 c_w=\begin{cases}
 \pi/(2\sqrt3),&w(a)=1-a,\\
 \pi/(4\sqrt3),&w(a)=1/2-a
 \end{cases}
\end{equation}
である．ここで時間を逆向きに辿る場合も同じ評価を用いる．

$K\subset U$ を compact な $\rho$-飽和集合とする．
$K\subset D\Subset U$ なる開集合 $D$ が存在し，
\begin{equation}\label{ke:eq:circle-actual-buffer}
 \delta:=d_{g_\omega}^{U}(K,\partial D)>c_w\varepsilon
\end{equation}
を満たすとする．$d_{g_\omega}^{U}$ は $U$ 内の経路による actual 距離である．
このとき $K$ のある $\rho$-飽和開近傍からの forward flow は全時刻で存在する．
また，時刻 $1$ に $K$ の点を終点として指定する backward flow も時刻 $0$ まで存在する．
従って最大の time-one flow の定義域と像はともに $K$ を含む．
各円軌道 $\mathcal O\subset K$ に対して
\begin{equation}\label{ke:eq:circle-actual-diameter}
 \operatorname{diam}_{d_{g_\omega}^{U}}\phi_{1,0}(\mathcal O)
 \le (2c_w+1/2)\varepsilon.
\end{equation}
ここで直径は ambient 距離による．
像の円自身の intrinsic diameter，flow の $C^1$ 評価，
元の neck base projection の保存は結論に含めない．
\end{lemma}
\begin{proof}
$g_a=\rho_a^*g_\omega$ と置く．正則性により
$g_\eta=\int_0^1g_a\,da$ であり，$d\rho_a(V)=V$ と仮定から
$g_a(V,V)\le\varepsilon^2$ である．任意の接ベクトル $X$ に対し
\begin{align*}
 |\beta(X)|
 &\le\varepsilon\int_0^1 |w(a)|\,|X|_{g_a}\,da\\
 &\le\varepsilon\|w\|_{L^2(0,1)}
       \left(\int_0^1g_a(X,X)\,da\right)^{1/2}
 =b|X|_{g_\eta},\qquad b=\varepsilon\|w\|_{L^2(0,1)}.
\end{align*}
共通複素構造を $J$ とする．Moser 方程式から
$g_{\omega_s}(Y_s,Y_s)=\beta(JY_s)$ なので，
$r=|Y_s|_{g_\eta}$，$q=|Y_s|_{g_\omega}$ と置けば
\[
 (1-s)r^2+s q^2\le br.
\]
平方完成によって
\[
 s q^2\le br-(1-s)r^2\le\frac{b^2}{4(1-s)}
\]
となり，\eqref{ke:eq:circle-actual-speed}を得る．
$\int_0^1(s(1-s))^{-1/2}\,ds=\pi$ を用いて
\eqref{ke:eq:circle-actual-length}が従う．
端点で表示した上界は発散するが可積分である．
$Y_s$ 自身は K\"ahler 形式 $\omega_s$ の非退化性により
$[0,1]\times U$ 上で滑らかに定義されている．

任意の時刻に $K$ から出発する trajectory が初めて $D$ の境界へ達すれば，
そこまでの長さは $\delta$ 以上になる．
これは\eqref{ke:eq:circle-actual-length}に反する．
従って trajectory は $\overline D\Subset U$ に留まり，
compact 集合上の滑らかな時間依存 ODE の延長定理によって
順方向・逆方向とも $[0,1]$ 全体へ延びる．
Strict な余裕と $K$ の compact 性により，この結論は $K$ の開近傍にも成立する．
$K$ と作用群が compact なので，その近傍に含まれる
$\rho$-飽和開近傍を選べる．Backward flow の時刻 $0$ の点は forward flow の定義域に属するから，
最大 flow の像も $K$ を含む．

元の円軌道の任意の二点は高々半周で結べ，その長さは $\varepsilon/2$ 以下である．
両端にそれぞれ長さ $c_w\varepsilon$ 以下の Moser trajectory を付ければ
\eqref{ke:eq:circle-actual-diameter}となる．
\end{proof}

\begin{corollary}[同じ compact 領域を覆う actual 円作用]
\label{ke:short-circle-image-buffer}
補題\ref{ke:short-circle-actual-buffer}の仮定に加えて
\begin{equation}\label{ke:eq:circle-image-buffer}
 \delta>(2c_w+1/2)\varepsilon
\end{equation}
とする．このとき time-one Moser map $\phi$ の定義域には
$\rho$-飽和した開集合 $O\supset K$ を取ることができ，その像も $K$ を含む：
\[
 K\subset O,\qquad K\subset\phi(O).
\]
従って $\widehat\rho=\phi\rho\phi^{-1}$ は $K$ の開近傍上の円作用となる．
さらに $\omega=d\lambda$ の場合は補題\ref{ke:circle-exact-transport}の
actual Hamiltonian action と moment の結論がその近傍で成立する．
\end{corollary}
\begin{proof}
$F$ を，時刻 $0$ からの trajectory が全時刻で $D$ 内に留まる点の集合とする．
これは開集合であり，time-one map は $F$ 上の滑らかな埋込である．
その最大の $\rho$-飽和部分
\[
 O=\{y\in F:\rho_a(y)\in F\text{ for every }a\in\mathbb R/\mathbb Z\}
\]
も開である．実際，ある軌道が開集合 $F$ に含まれれば，軌道の compact 性と
作用の連続性によりその近くの全軌道も $F$ に含まれる．
補題の forward 結論と $K$ の飽和性から $K\subset O$ である．

$x\in K$ を取り，backward flow による時刻 $0$ の点を $y$ とする．
$d_{g_\omega}^{U}(x,y)\le c_w\varepsilon$ である．
任意の $z=\rho_a(y)$ は $y$ から高々半周で結べるので
\[
 d_{g_\omega}^{U}(x,z)\le(c_w+1/2)\varepsilon.
\]
$y$ から $z$ への短い軌道弧も $D$ を出られない．
もし出れば，$x$ から $\partial D$ までの経路の長さが
$(c_w+1/2)\varepsilon<\delta$ となるからである．
$z$ からの forward Moser trajectory が $D$ を出るなら，
先ほどの $x$ から $z$ までの経路と合わせて
$x$ から $\partial D$ まで長さ $(2c_w+1/2)\varepsilon$ 以下の経路ができる．
これは\eqref{ke:eq:circle-image-buffer}に反する．
従って全ての $z=\rho_a(y)$ は $F$ に属し，$y\in O$ となる．
$x=\phi(y)$ だから $K\subset\phi(O)$ が従う．
最後の主張は既存の exact transport identity を適用すればよい．
\end{proof}

\begin{remark}[Neck への適用範囲]
上記の二つの結果は，actual normal generator の一様な短さと
actual 距離の buffer から normal averaging を実行する条件である．
Rank-one 評価と volume 比較から短さを得る場合も，
source-free neck 全体でこれらの入力が一様であることは別の評価を要する．
固定残余基底上では定理\ref{ke:fixed-base-neck-moser}で検証し，
source を含む固定段階への拡張は補題\ref{ke:fixed-stage-source-neck}で扱う．
複数の triple end における同一作用，reduced leaves，markings の同時整合，
source・branch・core への接着や全曲面の GH 収束は，この補題だけの結論ではない．
それぞれ系\ref{ke:neck-corner-saturated-union}，
定理\ref{ke:actual-global-fibration}，および
定理\ref{ke:global-ambient-lower}，\ref{ke:global-ambient-upper}で証明する．
\end{remark}

\subsection{固定残余領域上の法線作用の平均エネルギー}

\begin{proposition}[法線方向の積分評価と短い円軌道]
\label{ke:neck-integrated-normal-energy}
Source と三重点から離れた固定座標で $uv=t$，$z\in D$ とし，
$|u|,|v|<r$ に対して $s=-\log(|u|^2/r^2)\in(0,M_t)$ と置く．
ここで $M_t=-\log(|t|^2/r^4)$ である．
固定した $Q\Subset D$ 上で，actual negative KE form $\omega_t$ の体積が
annulus--disc complete KE 積計量の体積と $a_0\le s\le M_t-a_0$ で
一様に二側比較可能とする（$a_0>0$ は固定）．
$V=2\pi\partial_{\arg u}$，$\rho$ をその周期 $1$ の作用とし，
\[
 E_t(s,z)=\int_0^1 |V|_{g_t}^2(\rho_a(u,z))\,da
\]
と置く．$2a_0\le m\le M_t/2$ ならば，全 $z\in Q$ について
\begin{equation}\label{ke:eq:normal-energy-mass}
 \int_m^{M_t-m} E_t(s,z)\,ds\le C/m .
\end{equation}
定数 $C$ は $t,m,z$ に依存しない．任意の $A>0$ について
$[m,M_t-m]$ 内で $E_t>A/m^2$ となる深度の Lebesgue 測度は
$Cm/A$ 以下である．それ以外の深度では actual 円軌道の長さは
$\sqrt A/m$ 以下である．
\end{proposition}

\begin{proof}
$\Omega=du/u\wedge dz$，$F_t=\log(\omega_t^2/|\Omega|^2)$ とし，
$\bar F_t$ を $u$ の角度について平均する．$|\Omega|^2$ は
$\Omega$ の定める正の体積形式で，その定数規格化は以下に影響しない．
$d^c=\tfrac{\sqrt{-1}}2(\bar\partial-\partial)$ の規格化では
$dd^cF_t=\omega_t$ である．$d(s)=\min(s,M_t-s)$ とすると，
積計量の体積密度は $z\in Q$ 上で
$(\pi/M_t)^2\csc^2(\pi s/M_t)\asymp d(s)^{-2}$ に比較可能だから
\[
 \bigl|\bar F_t(s,z)+2\log d(s)\bigr|\le C_0 .
\]
局所対数座標 $\xi=\log u$ で $s=-2\Re\xi+2\log r$ である．
$\mathcal A\omega_t=dd^c\bar F_t$ の $\xi\bar\xi$ 係数から
$E_t=8\pi^2\partial_s^2\bar F_t$ を得る．特に $\bar F_t$ は $s$ の凸関数である．

$s\in[m,M_t-m]$ で $s\pm m/2$ を取ると，三点の $d$ の値の比は
$1/2$ と $3/2$ の間にある（$d(s)$ を分母とする）．
従って前進・後退の差商は絶対値 $C_1/m$ 以下であり，
凸性から $|\partial_s\bar F_t|\le C_1/m$ を得る．
両端における導関数の差を取れば
\eqref{ke:eq:normal-energy-mass} が従う．
測度評価は Markov の不等式，長さ評価は
$\int_0^1 |V|_{g_t}(\rho_a x)\,da\le E_t(s,z)^{1/2}$ による．
\end{proof}

上の体積仮定は，一般位置全平面退化の固定 source-free pair 座標では，
関連 KE 主稿の \texttt{neck:prop:volume-envelope} から従う．
例えば $m=R_t$，$A=R_t$ とすれば，各固定残余点で中央 strip のうち
深度の測度 $O(1)$ の例外を除いて円軌道長は $O(R_t^{-1/2})$ である．
この例外集合は残余点と $t$ に依存する．

\begin{corollary}[鋭い係数を持つ同一残余点での法線接続上界]
\label{ke:neck-axial-upper}
命題\ref{ke:neck-integrated-normal-energy}の座標で $z\in Q$ を固定する．
$2a_0\le a\le b\le M_t-2a_0$ とする．ある一定の角度で深度 $a$ から $b$ へ
進む動径弧の actual 長さは
\[
 \int_a^b \frac{ds}{d(s)}+C
\]
以下である．従って両端円の任意の点を，同じ $z$ を保つ経路で
この上界に $\tfrac12(\sqrt{E_t(a,z)}+\sqrt{E_t(b,z)})$ を加えた
長さ以下で結べる．
\end{corollary}

\begin{proof}
まず $b\le M_t/2$ とする．$\bar F_t=-2\log s+O(1)$ で，
命題の差商評価に各点で $m=d(s)=s$ を用いると $|s\partial_s\bar F_t|\le C_1$
である．部分積分して
\[
 \int_a^b s\partial_s^2\bar F_t\,ds
 =[s\partial_s\bar F_t]_a^b-[\bar F_t]_a^b
 \le 2\log(b/a)+C_2 .
\]
$JV=4\pi\partial_s$ と $J$ 不変性から，
$g_t(\partial_s,\partial_s)$ の角度平均は
$E_t/(16\pi^2)=\tfrac12\partial_s^2\bar F_t$ である．
従って Cauchy--Schwarz により動径弧長の角度平均は
\[
 \int_a^b \sqrt{\tfrac12\partial_s^2\bar F_t}\,ds
 \le \sqrt{\log(b/a)\bigl(\log(b/a)+C_2/2\bigr)}
 \le \log(b/a)+C_2/4 .
\]
上側半区間では $s$ を $M_t-s$ に置き換える．
中点を横切る場合は角度平均を取った後で両半区間の上界を加えるので，
全区間に共通する一つの角度を選べる．
両端では円の短い方の弧を付加する．
\end{proof}

この接続上界は正規化前の距離に対し加法誤差 $O(1)$ を持つ．
両端の円エネルギーも有界な場合は，正規化した法線接続費用が
$R_t^{-1}\int_a^b ds/d(s)+O(R_t^{-1})$ 以下となる．
この評価だけでは異なる残余点間の portal 接続費用や，全曲面を通る経路の下界は得られない．
前者は命題\ref{ke:fixed-base-moment-connectors}と
系\ref{ke:source-inclusive-same-moment-collapse}，
後者は定理\ref{ke:global-ambient-lower}で別に評価する．

\begin{proposition}[固定残余基底上の moving deep moment strip]
\label{ke:neck-deep-moment-strip}
関連 KE 主稿の \texttt{pair:prop:bundle-conic} の固定 source-free compact
基底 $Q\Subset D^\circ$ 上で，双対 Hermitian norms を用いて
$uv=t$，$s=-\log(\|v\|^2/r^2)$，$M_t=-\log(|t|^2/r^4)$ とする．
$V$ は $v$ を正方向に回す周期 $1$ の生成場とし，
$\mu_t=(\mathcal A\lambda_t)(V)$ は同稿 \texttt{pair:cor:global-primitive}
の canonical moment とする．
定数 $m_0,c,C>0$ と $\vartheta\in(0,1/2)$ を固定できて，
$m_0\le m\le\vartheta M_t$ のすべての深度について，全 $Q$ 上で
\[
 \frac c m\le\mu_t\big|_{s=m}\le\frac C m,
 \qquad
 -\frac C m\le\mu_t\big|_{s=M_t-m}\le-\frac c m
\]
が成り立つ．従って任意の $m=m_t$ がこの範囲にあれば，
共通 interval $I_t=[-c/(2m_t),c/(2m_t)]$ の各 level は
$m_t<s<M_t-m_t$ 内で各 annulus に一つの円を持つ．
その商は $I_t\times Q$ であり，$Q$ が境界付き領域なら境界も保つ．
特に $\mu_t=0$ の level は両側深度が少なくとも $\vartheta M_t$
である領域に入る．
\end{proposition}

\begin{proof}
局所正則 frame で $\Omega=dv/v\wedge dz$ をとる．
$\bar F_t=\mathcal A\log(\omega_t^2/|\Omega|^2)$ とすると，
定数倍および base-only factor を除いた primitive の規格化から
$\mu_t=-2\pi\partial_s\bar F_t$ を得る．
Hermitian norm による $s$ と正則 frame の座標による深度の差は
固定 compact 上で有界な base-only 関数であり，$\partial_s$ は変わらない．
有限個の少し大きい積座標に命題\ref{ke:neck-integrated-normal-energy}で使った
体積包絡を適用して，$d(s)=\min(s,M_t-s)$ とすれば
ある固定 $a_0>0$ について
\[
 |\bar F_t(s,z)+2\log d(s)|\le C_0
 \qquad (a_0\le s\le M_t-a_0)
\]
を $Q$ 全体に一様に得る．Base coordinate の変更による加法項は
有限個の compact patches 上で有界であり，差商には現れない．

$2\log\Lambda>2C_0+1$ となる固定 $\Lambda>1$ を選び，
$\vartheta=1/(2\Lambda)$ とする．$2a_0\le m\le\vartheta M_t$ なら，
$m$ と $\Lambda m$ は左側半区間にある．凸性により
\[
 \partial_s\bar F_t(m,z)
 \le\frac{\bar F_t(\Lambda m,z)-\bar F_t(m,z)}{(\Lambda-1)m}
 \le -\frac{1}{(\Lambda-1)m}.
\]
逆向きの大きさの上界は命題\ref{ke:neck-integrated-normal-energy}の
差商評価 $|\partial_s\bar F_t|\le C_1/m$ による．
右側は $M_t-s$ を使って同様に求める．
$\partial^2_s\bar F_t>0$ なので $\mu_t$ は各 annulus の深度に対し
厳密に減少する．中間値定理と陰関数定理で，$I_t$ 全体にわたる
一意な深度断面が滑らかに貼り合わさる．半径を正に変えても
円束の標識は保たれるので商は $I_t\times Q$ になる．
最後は $m=\vartheta M_t$ を選ぶ．
\end{proof}

この interval は normal 深度を動かしつつ，すべての固定残余点で
同じ canonical moment を使うためのものである．
既存の固定 end の moment 漸近を moving cut に代入する必要はない．
レベル上で点ごとの円エネルギーを抑える評価と actual Moser buffer は，
定理\ref{ke:fixed-base-neck-moser}で別に与える．
Source を含む固定段階と三重点へ動く collars の同時選択には，
補題\ref{ke:fixed-stage-source-neck}と定理\ref{ke:neck-corner-diagonal-matching}を用いる．

\begin{corollary}[軌道平均エネルギーを用いる Moser 速度評価]
\label{ke:neck-averaged-speed}
補題\ref{ke:short-circle-actual-buffer} の Moser 方程式において，
点 $x$ で
$E(x)=\int_0^1 |V|_{g_\omega}^2(\rho_a x)\,da$ とすると，
\[
 |Y_\tau(x)|_{g_\omega}
 \le \frac{\|w\|_{L^\infty}\sqrt{E(x)}}{2\sqrt{\tau(1-\tau)}} ,
 \qquad 0<\tau<1 .
\]
作用不変な開集合上で $E\le e^2$ が一様に成り立てば，
補題\ref{ke:short-circle-actual-buffer} と
系\ref{ke:short-circle-image-buffer} の長さ・buffer の結論は
$\varepsilon=e$，$c_w=(\pi/2)\|w\|_{L^\infty}$ で成り立つ．
\end{corollary}

\begin{proof}
$g_a=\rho_a^*g_\omega$ とすると
\[
 |\beta(X)|\le
 \left(\int_0^1 |w(a)|^2 g_a(V,V)\,da\right)^{1/2}
 |X|_{g_\eta}\le \|w\|_{L^\infty}\sqrt{E(x)}\,|X|_{g_\eta}.
\]
この定数を補題\ref{ke:short-circle-actual-buffer} の平方完成に用いる．
また $E\le e^2$ 上で一周の長さが $e$ 以下なので，
二点間の短い方の軌道弧の長さは $e/2$ 以下である．
これを使って両向きの軌道追跡と飽和近傍の証明をそのまま適用できる．
\end{proof}

積分評価\eqref{ke:eq:normal-energy-mass}は所与の深度での点ごとの上界や，
内側集合から小エネルギー領域の境界までの actual 距離下界を与えない．
例えば残余方向を固定しても凸関数の二階微分は小さい区間に集中し得る．
被覆を得るには，所要の開 collar 上の一様な $E$ の上界と actual buffer，
または軌道ごとに同じ被覆を保証する変位評価が必要となる．
本稿では定理\ref{ke:fixed-base-neck-moser}で前者を固定残余基底上に構成し，
補題\ref{ke:fixed-stage-source-neck}で source を含む場合へ拡張する．

\begin{corollary}[縮小する actual buffer からの Moser 被覆]
\label{ke:neck-shrinking-buffer}
上記の annulus--disc 座標で，compact な飽和集合 $K_t$ を含む飽和開集合
$U_t\subset\{d(s)>m_t\}$ を取る．$g_t$ に関する全曲面の距離で
\[
 \delta_t=\min\bigl(1,\operatorname{dist}_{g_t}(K_t,X_t\setminus U_t)\bigr)
\]
と置く．$m_t\delta_t^2\to\infty$ ならば，十分小さい $t$ で
$K_t$ を source と像の双方に含む飽和 source を持つ normal Moser map が存在する．
その円軌道像の ambient 直径は $O((m_t\delta_t)^{-1})$ である．
\end{corollary}

\begin{proof}
$N_t=\{\operatorname{dist}_{g_t}(\cdot,K_t)<\delta_t/2\}$ と置く．
各 $q\in N_t$ の actual $\delta_t/4$-ball の閉包は $U_t$ に含まれる．
残余射影への localized Chern--Lu 評価を半径 $\delta_t/4$ で用いると，
中心で $p^*h_D\le C\delta_t^{-2}\omega_t$ を得る．
この半径依存性は，距離 cutoff の勾配二乗と Laplacian の下界を
$O(\delta_t^{-2})$ で制御する Ricci 下界による距離比較から従う．
関連 KE 主稿の \texttt{neck:lem:rank-one-circle} と体積上界により
$N_t$ 上で $|V|_{g_t}\le\varepsilon_t=C'/(m_t\delta_t)$ となる．

$D_t=\{\operatorname{dist}_{g_t}(\cdot,K_t)<\delta_t/4\}$ とする．
$\overline D_t$ から任意の位相まで周期の半分以下で円軌道を辿るとき，
$N_t$ の最初の退出までの長さは $\varepsilon_t/2$ 以下である．
これは $\delta_t/4$ より短いので，$\overline D_t$ を通る全軌道が $N_t$ に留まる．
従って $N_t$ に含まれる最大の飽和開集合に $\overline D_t$ が含まれ，
そこで同じ短円上界を使える．$K_t$ から $\partial D_t$ までの actual 距離は
$\delta_t/4$ 以上であり，$m_t\delta_t^2\to\infty$ により
$(2c_w+1/2)\varepsilon_t$ を上回る．系\ref{ke:short-circle-image-buffer}と
補題\ref{ke:short-circle-actual-buffer}の直径評価を適用する．
\end{proof}

\subsection{固定残余領域での actual tensor 比較と Moser 被覆}

\begin{lemma}[大域 capped density を用いた射影計量の下界]
\label{ke:capped-projective-lower}
固定次数 $d\ge5$ の一般位置全平面退化を考え，
$\Omega_{t,d}^{\rm low}$ と $h_{t,d}^{\rm low}$ を
関連 KE 主稿の \texttt{coarse:cor:degree-d-volume-lower} で定める．
$\varphi_t=\log(\omega_t^2/\Omega_{t,d}^{\rm low})$ とすると，
ある固定 $A,C>0$ に対し
\[
 \operatorname{tr}_{\omega_t}(\omega_{\mathrm{FS}}|_{X_t})
 \le C\exp(A\varphi_t)
\]
が全 $X_t$ で成り立つ．従って $\varphi_t$ が一様に上から有界な領域では
$g_t\ge c g_{\mathrm{FS}}|_{X_t}$ となる．
\end{lemma}

\begin{proof}
KE 側の capped density の証明で $A_0$ を十分大きく固定して，
$h_{t,d}^{\rm low}\ge c_0\omega_{\mathrm{FS}}|_{X_t}$ を得る．
実際，背景は $(d-4)\omega_{\mathrm{FS}}|_{X_t}$ で，
すべての cap の負部分と frame errors の和は
その $O_d(A_0^{-1})$ 倍に抑えられ，残りは半正である．
同じ大域体積下界から $\varphi_t\ge-C_0$ が全 $X_t$ で成り立つ．
$T_t=\operatorname{tr}_{\omega_t}(\omega_{\mathrm{FS}}|_{X_t})$ とすると，
正則包含 $X_t\to\PP^3$ の Chern--Lu 不等式は
$\Delta_{\omega_t}\log T_t\ge-C_1-C_2 T_t$ を与える．
ここでは射影空間の固定した curvature upper bound と
$\operatorname{Ric}(\omega_t)=-\omega_t$ を用いる．
$dd^c\varphi_t=\omega_t-h_{t,d}^{\rm low}$ なので
\[
 \Delta_{\omega_t}(\log T_t-A\varphi_t)
 \ge (Ac_0-C_2)T_t-(C_1+2A).
\]
$Ac_0>C_2$ と取り，compact $X_t$ 上の最大点で $T_t$ を上から抑える．
そこでの $\varphi_t\ge-C_0$ と合わせ，
$\log T_t-A\varphi_t\le C'$ を全域に得る．
\end{proof}

\begin{theorem}[固定 source-free base の moving strip での actual Moser 被覆]
\label{ke:fixed-base-neck-moser}
関連 KE 主稿の \texttt{pair:prop:bundle-conic} の正則束座標を用い，
source と三重点を避けた固定した $Q\Subset B$ 上で
$uv=t$，$s=-\log(\|v\|^2/r^2)$，
$M_t=-\log(|t|^2/r^4)$，$d(s)=\min(s,M_t-s)$ とする．
$Q$ は compact とする．固定近傍に依存する有限な $m_Q$ がある．
$|t|$ を十分小さくし，$m_Q\le m_t\le M_t/2$ に対して
$K_t=\{p\in Q,\ d(s)\ge m_t\}$ とする．
$K_t$ を source と像の双方に含む normal Moser map が存在し，
その軌道像の actual ambient 直径は $O(m_t^{-1})$ である．
さらに $K_t$ を含む深部開 strip 上で，局所積座標の complete
annulus--disc 計量 $g_{\rm prod}$ と $g_t$ は一様に二側比較可能である．
定数は $m_t,t$ に依存しない．従って固定した十分大きい cut にも使え，
$m_t\to\infty$ と選べば軌道の変位は零へ収束する．
\end{theorem}

\begin{proof}
まず $\overline Q$ の少し大きい固定残余近傍を有限個の局所積座標で覆う．
各 patch で outer annulus--disc を固定し，その内部に
三重に入れ子になった annulus--disc を，正の固定な normal radii と residual radii
で選ぶ．最も外側の入れ子領域（outer chart には compact に含まれる）では
outer chart の体積 envelope と capped density が一様に比較可能である．
これは \texttt{neck:prop:volume-envelope} の証明にある
ordinary-pair capped density の比較である．同命題の体積上界から
そこでは $\varphi_t\le C$ であり，補題\ref{ke:capped-projective-lower}
によって $g_t\ge c g_{\mathrm{FS}}$ を得る．

内側の入れ子領域から中間領域の外へ出る経路は，その最初の退出までに，
残余座標又は $u,v$ の一方を固定した正の幅だけ変化させる．
これらの座標は固定な ambient 座標上で有界な微分を持つ．
実際，正則束座標は $t$ に依存しない ambient 変換で作られ，
コンパクトな座標範囲では変換とその逆の微分が一様に有界である．
よって経路は比較が成立する部分だけで一様に正の actual 長さを持つ．
従って内側の各点の一様な半径の actual ball の閉包は中間領域に含まれる．

この actual ball 上で outer annulus--disc の complete product への
正則恒等写像に localized Chern--Lu 評価
（関連 KE 主稿 \texttt{coarse:lem:localized-chern-lu}）を適用する．
曲率 $-1$ の二曲面の積の holomorphic sectional curvature は
$-1/2$ 以下なので，中心で $g_t\ge c g_{\rm prod}$ を得る．
既存の体積上界と合わせれば複素次元 $2$ の行列の固有値比較により
$g_t\le C g_{\rm prod}$ も従う．有限被覆で定数を一様にできる．
固定 norm による $s$ は局所積座標の深度と有界な base-only 項だけ異なり，
その項の微分も一様に有界である．従って十分深い領域上で
\[
 |V|_{g_t}\le C/d(s),\qquad
 |dp|_{g_t}\le C,\qquad
 |d\log d(s)|_{g_t}\le C
\]
となる．最後の不等式は $s=M_t/2$ の継目では片側微分，
又は経路に沿う Lipschitz 評価として読む．$|dp|$ の標的には固定基底計量を用いる．

$Q\Subset B'\Subset B$ を固定して，Moser の外側領域を
$U_t=\{p\in B',\ d(s)>m_t/2\}$ とする．
$K_t$ から $U_t$ の境界へ向かう任意の経路は，最初の退出以前に
基底で固定した正の距離を移動するか，$\log d(s)$ を $\log 2$ 以上変える．
従って actual buffer は正の定数以上である．ここで局所比較のための有限被覆は
$\overline{B'}$ 上に選んでいる．同じ議論で $B'$ と $m_t/2$ を少し内側へ
縮めた比較的 compact な領域を $U_t$ 内に取れる．
$U_t$ 上の $|V|_{g_t}\le C_Q/m_t$ を用いる．上で得た正の buffer を
$\beta_Q$ とし，$m_Q$ を
$(2c_w+1/2)C_Q/m_Q<\beta_Q$ となるまで大きく取れば，
全ての $m_t\ge m_Q$ について系\ref{ke:short-circle-image-buffer}を適用できる．
この条件も有限被覆の tensor 比較も $m_t\to\infty$ 自体を必要としない．
従って固定 cut を含めて求める Moser map を得る．
直径は補題\ref{ke:short-circle-actual-buffer}の評価による．
\end{proof}

定理\ref{ke:fixed-base-neck-moser} は residual base が固定した場合を扱う．
Source や三重点へ近付く residual collars では，射影計量で測った入れ子領域間の幅や
capped density と complete envelopes の比較を別に用意する．
Source は定理\ref{ke:projected-source-neck}，三重点の全 end との整合は
定理\ref{ke:neck-corner-diagonal-matching}と系\ref{ke:neck-corner-saturated-union}で扱う．
固定 compact の結論と命題\ref{ke:neck-deep-moment-strip} を合わせると，
同じ固定基底上で moving canonical moment strip を actual form に移せる．
複数 end の接着では，これに加えて上記の正則束作用と triple 作用との
開 overlap 上での整合を用いる．

\begin{proposition}[同一 normal moment 上の残余方向を含む有界な接続]
\label{ke:fixed-base-moment-connectors}
命題\ref{ke:neck-deep-moment-strip}と定理\ref{ke:fixed-base-neck-moser}の
固定 source-free 基底で，$P$ を連結な compact 境界付き曲面とし，
その少し大きい固定近傍上で両結果を用いる．
固定近傍に依存する十分大きい有限な $m_P$ を取り，
$m_P\le m_t\le\vartheta M_t$ とし，
$I_t=[-c/(2m_t),c/(2m_t)]$ を共通 moment interval とする．
円平均後の canonical moment を $\mu_t$，同じ normal Moser map を
$\phi_t$ とし，$a\in I_t$ に対して
\[
 E_t(a)=\{p\in P,\ m_t\le s\le M_t-m_t,\ \mu_t=a\}
\]
と置く．$t,a,m_t$ に依存しない $C_P$ が存在して
\begin{equation}\label{ke:eq:fixed-base-moment-connectors}
 \operatorname{diam}_{d_{g_t}}E_t(a)\le C_P,
 \qquad
 \operatorname{diam}_{d_{g_t}}\phi_t(E_t(a))
 \le C_P+O(m_t^{-1})
\end{equation}
となる．従って actual normal moment の同じ level に属する
この neck 内の全 fibres の合併は，正規化距離 $R_t^{-1}d_{g_t}$ で
一様に $O(R_t^{-1})$ の直径を持つ．
接続経路は moment level や個々の fibre の中に留まる必要はなく，
結論は ambient diameter についてのものである．
\end{proposition}

\begin{proof}
$d(s)=\min(s,M_t-s)$ とする．
命題\ref{ke:neck-deep-moment-strip}の証明から，定数を共通に取り
\[
 |\mu_t(s,p)|\le C_0/d(s),\qquad
 c_0/s\le\mu_t(s,p)\le C_0/s
 \quad(m_0\le s\le\vartheta M_t)
\]
を得る．右端では $s$ を $M_t-s$ に置き換え，符号を逆にする．
各 $p$ で $s\mapsto\mu_t(s,p)$ は厳密に減少する．
$A=2C_0/\vartheta$ と取ると，$a>A/M_t$ の level はすべて左端側にあり，
\[
 c_0/a\le s(a,p)\le C_0/a.
\]
$a<-A/M_t$ なら $M_t-s(a,p)$ が同様に比較できる．
残る $|a|\le A/M_t$ では，左端側にある点に上の下界を適用すると
$s(a,p)\ge c_0M_t/A$ となる．右端側も同様なので，
ある固定 $\kappa>0$ に対して
\[
 \kappa M_t\le s(a,p)\le(1-\kappa)M_t
\]
を得る．以上の三場合のいずれでも，任意の二つの残余点について
\begin{equation}\label{ke:eq:moment-depth-connector}
 \left|\int_{s(a,p)}^{s(a,q)}\frac{ds}{d(s)}\right|\le C_1
\end{equation}
が一様に成り立つ．この議論は $s(a,p)$ の残余微分を評価しない．
用いた $c_0,C_0,\vartheta$ と深度の下限 $m_0$ は固定近傍だけに依存するので，
$m_P$ を $m_0$ と前定理の Moser 下限以上に取れば，固定した有限 cut にも
同じ証明が適用される．特に固定 stage を先に選び，その後 $t\to0$ とする用途に使える．

固定した円束の smooth connection と $P$ 上の固定補助計量を取る．
連結 compact 曲面 $P$ の任意の二点は，この補助計量で長さが一様に有界な
区分的 smooth な経路で結べる．定理\ref{ke:fixed-base-neck-moser}の
二側比較を有限個の局所積座標に適用すると，深度 $s$ を固定して
この経路を水平に持ち上げた actual 長さも一様に有界である．
実際，局所正則 frame と固定 Hermitian norm の深度の差，および
connection form は固定基底上で一様有界であり，normal 方向の寄与は
$d(s)^{-1}$ 倍される．同じ二側比較から，固定残余点・固定角度の
動径弧の長さは $C\int ds/d(s)$ 以下である．

そこで $E_t(a)$ の二点を，第一点の深度を保つ水平経路，
第二点の残余点での動径経路，最後の短い円弧によって結ぶ．
動径費用は\eqref{ke:eq:moment-depth-connector}，
円弧の費用は $|V|_{g_t}=O(m_t^{-1})$ から有界である．
全経路は深部 neck 内に取り，最初の不等式を得る．
最後に両端へ $\phi_t$ の trajectory を付ける．
各 trajectory の actual 長さは $O(m_t^{-1})$ なので第二式が従う．
これは同一 $a$ 上の任意の二点に一様であり，reduced foliation の
葉の選択や葉内での長さに依存しない．
\end{proof}

\begin{corollary}[異なる moment を含む固定基底の中央帯]
\label{ke:fixed-base-central-band}
命題\ref{ke:fixed-base-moment-connectors}と同じ固定連結 source-free 基底 $P$，
canonical moment $\mu_t$，normal Moser map $\phi_t$ を用いる．
$A>0$ を固定し，前命題の定数を用いて
\[
 m_P\le m_t\le\min\{\vartheta,c/(2A)\}M_t
\]
とする．これは固定した十分大きい cut を含み，$m_t=o(M_t)$ の選択でも満たされる．
十分小さい $t$ について
\[
 E_{t,A}=\{p\in P,\ m_t\le s\le M_t-m_t,\ |\mu_t|\le A/M_t\}
\]
と置くと，$t,m_t$ に依存しない $C_{P,A}$ が存在して
\[
 \operatorname{diam}_{d_{g_t}}E_{t,A}\le C_{P,A},\qquad
 \operatorname{diam}_{d_{g_t}}\phi_t(E_{t,A})
 \le C_{P,A}+O(m_t^{-1})
\]
が成り立つ．異なる normal moment に属する点も同時に含み，
正規化した ambient diameter は $O_{P,A}(R_t^{-1})$ である．
\end{corollary}

\begin{proof}
$m_t\ge m_P$ により，前命題の片側 moment 評価に必要な固定最小深度を超える．
また仮定の上限から $m_t\le\vartheta M_t$ および $A/M_t\le c/(2m_t)$ となるので，
中央帯の moment 値は同じ共通 interval に入る．
$s\le\vartheta M_t$ での $\mu_t\ge c_0/s$ と
$|\mu_t|\le A/M_t$ を合わせれば $s\ge c_0M_t/A$ を得る．
右側では符号を逆にした評価を用いる．従って
\[
 \kappa_A=\min\{\vartheta,c_0/A,1/4\}>0
\]
と置けば，中央帯の全点で
$\kappa_A M_t\le s\le(1-\kappa_A)M_t$ である．
特に，異なる moment に属する任意の二点の深度 $s_1,s_2$ に対しても
\[
 \left|\int_{s_1}^{s_2}\frac{ds}{\min(s,M_t-s)}\right|
 \le 2\log\frac{1}{2\kappa_A}
\]
となる．前命題の固定深度での水平 lift，終点の残余点での動径弧，
最後の phase 弧を用いれば，この二点を一様な長さで結べる．
経路は深部 neck 内に留まるが，中央帯や moment level 内に留まる必要はない．
両端に長さ $O(m_t^{-1})$ の Moser trajectories を加えれば像の評価も従う．
\end{proof}

以上の命題と系は，固定した source-free 基底上で異なる残余点を結ぶ portal 費用を扱う．
Source の pinched fibre を含む全共通 interval の同一 moment 集合は，
系\ref{ke:source-inclusive-same-moment-collapse}で評価する．
三重点へ動く残余 collars は固定段階の exhaustion として扱い，
命題\ref{ke:collapsed-tree-common-moment}の共通 moment parameter 化と合わせて，
命題\ref{ke:collapsed-tree-full-fiber-collapse}で潰される木上の全 fiber 合併径を得る．
任意に異なる moments に属する neck 全体の径が消えるという主張ではない．
全曲面の no-shortcut と距離上界は，定理\ref{ke:global-ambient-lower}，
\ref{ke:global-ambient-upper}で別に証明する．

\begin{corollary}[Central source を含む深部での短い円]
\label{ke:source-short-circle}
Source の固定正則座標 $xy=tz$ で，十分小さい固定 $r>0$ を取る．
$\rho_a(x,y,z)=(e^{2\pi\sqrt{-1}a}x,e^{-2\pi\sqrt{-1}a}y,z)$ の
周期 $1$ の生成場を $V_{\rm src}$ とする．
$|x|,|y|,|z|\le r$ の範囲で，$L_t=-\log|t|^2$ として
\[
 |V_{\rm src}|_{g_t}
 \le \frac{C}{\min\{-\log|x|^2,-\log|y|^2,L_t\}},
\]
ただし零の座標の深度は $+\infty$ とする．特に
$|x|,|y|\le c|t|^{1/2}$ の central source で短さは $O(L_t^{-1})$ である．
\end{corollary}

\begin{proof}
小さい固定 $c>0$ について，$x$-oriented 部分
$|x|\ge c|t|^{1/2}$ を考える．$(x,z)$ を用いる annulus--disc extension
$y=(t/x)z$ に対し，関連 KE 主稿
\texttt{coarse:lem:pair-arm-density} の capped density 比較と
\texttt{coarse:lem:complete-volume-upper} により $\varphi_t\le C$ を得る．
実際，$r$ は extension の人工境界から固定幅を取って選び，
$s_x=-\log|x|^2+O(1)\le L_t/2+O(1)$ なので，
$\alpha_y\asymp L_t$，$\alpha_x\asymp s_x$ であり，
capped density と product envelope は比較可能である．
これには gate $|x|=c|t|^{1/2}$ 上も含まれる．
補題\ref{ke:capped-projective-lower}を使うと
$g_t\ge c' g_{\mathrm{FS}}$ が成立する．固定した ambient coordinate $z$ の
微分は射影計量に対して有界なので，残余 rank-one 下界
$p_z^*h_D\le C\omega_t$ が成り立つ．
体積上界と \texttt{neck:lem:rank-one-circle} を使い，
$x$ の角速度を持つ $V_{\rm src}$ の長さを $C/s_x$ 以下に抑える．
$y$-oriented 部分も対称である．

残る領域は $(x,y)$ で表した rescaled bidisc
$\{|x|,|y|<c|t|^{1/2}\}$ である．$c^2<r$ とすれば
その閉包は $|z|\le c^2$ で source chart 内に含まれる．
境界では少なくとも一方の oriented 評価により
$|V_{\rm src}|_{g_t}=O(L_t^{-1})$ である．
$V_{\rm src}$ は正則ベクトル場
$W=x\partial_x-y\partial_y$ の虚方向の実生成場であり，
$|V_{\rm src}|_{g_t}^2$ は $|W|_{g_t}^2$ の固定倍である．
正則ベクトル場の Bochner 恒等式
$\Delta_{\omega_t}|W|^2=|\nabla^{1,0}W|^2-\operatorname{Ric}(W,\bar W)\ge0$
と最大値原理で，同じ上界が bidisc 内部にも成り立つ．
\end{proof}

この系は source の円作用そのものを評価し，固定点での作用の退化を許す．
移動する normal 境界までの actual buffer と閉じた pinched torus の構成は，
次の系\ref{ke:source-diagonal-moser}および定理\ref{ke:projected-source-neck}で扱う．
Source 中央帯の全点径は定理\ref{ke:source-central-diameter}，
全共通 interval の同一 moment 退出は定理\ref{ke:source-all-moment-exit}で別に評価する．

\begin{lemma}[Source の零 moment 商と閉じた pinched fibre]
\label{ke:source-reduction-pinched-fiber}
$t\ne0$ を固定し，境界付き compact 残余曲面 $P$ の近傍に
\[
 X=\{xy=t h(p)\},\qquad
 \rho_a(x,y,p)=(e^{2\pi\sqrt{-1}a}x,e^{-2\pi\sqrt{-1}a}y,p)
\]
を取る．$h$ の零点は $P$ の内点にある単純零点とし，
$\omega=d\lambda$ は $X$ 上で滑らかな K\"ahler form とする．
$\eta=\mathcal A\omega$，$\bar\lambda=\mathcal A\lambda$，
$\mu=\bar\lambda(V)$ と置く．固定 $R>0$ と
$0\in\operatorname{int}I$ なる閉区間 $I$ を選び，
$|t|\sup_P|h|<R^2$，および全 $p\in P$ で
\[
 \left.\mu\right|_{|y|=R}<\inf I<\sup I
       <\left.\mu\right|_{|x|=R}
\]
を仮定する．$K=\{p\in P,\ |x|,|y|\le R,\ \mu\in I\}$ の
飽和開近傍で centered normal Moser flow が全時刻存在するとする．
このとき次が成り立つ．
\begin{enumerate}
 \item 各 source $o=(0,0,q)$ で $\mu(o)=0$ であり，
 centered Moser flow は $o$ を固定する滑らかな微分同相である．
 Actual circle の symplectic isotropy representation は weights $(1,-1)$ を保ち，actual moment は
 $\widehat\mu=\mu\circ\phi^{-1}$ である．その値は actual 軌道の
 $\lambda$-周期に等しく，source ごとに加法定数を選び直す必要はない．
 \item $(\mu,p)$ は $K/S^1\cong I\times P$ を与える．
 $(0,q)$ 以外ではこれは滑らかな orbit space の座標で，
 $K\setminus\{o:h(q)=0\}$ 上の作用は自由である．
 \item $P$ の marked Morse foliation で，各 source が
 saddle を含まない相異なる閉 regular leaves 上にあるとする．
 $\mu$ を固定して葉を持ち上げた閉 fibre は，通常 $T^2$，
 単純 saddle leaf では $S^1\times\figeight$，
 $\mu=0$ で source を一つ含む葉では一重 pinched torus となる．
 二次元 smooth strata は Lagrangian で，各 source は
 nondegenerate focus--focus point である．近傍の regular torus 上の
 normal circle は primitive な vanishing cycle を与える．
 これらの結論は $\phi$ によって actual $\omega$ へ移される．
\end{enumerate}
ここでの滑らかさは各固定 $t\ne0$ に対する主張であり，
$t\to0$ での flow の導関数や正規形の一様性は含めない．
\end{lemma}
\begin{proof}
Source の近くでは $h'(q)\ne0$ により $(x,y)$ が $X$ の正則座標である．
平均 $\eta,\bar\lambda$ と
$\beta=\int_0^1(1/2-a)\rho_a^*(\iota_V\omega)\,da$ はここでも滑らかである．
$V(o)=0$ なので $\mu(o)=\beta(o)=0$．
$\omega_s=(1-s)\eta+s\omega$ は全時刻で K\"ahler form であり，
$\iota_{Y_s}\omega_s=\beta$ は滑らかな場 $Y_s$ を与え，$Y_s(o)=0$ となる．
補題\ref{ke:circle-exact-transport}の証明は自由性を使わないので，
固定点を含めてその exact identity と周期の結論が成立する．
共役によって固定点での円表現も変わらない．

$h(p)\ne0$ の fibre では，$x\mapsto e^r x$，$y\mapsto e^{-r}y$ の
実動径場は $-JV/(2\pi)$ である．従って
\[
 d\mu\bigl(-JV/(2\pi)\bigr)=|V|_\eta^2/(2\pi)>0.
\]
各 moment level はただ一つの円となる．$h(p)=0$ では二軸が原点を共有し，
$x$ 軸で $\mu$ は零から正へ，$y$ 軸で零から負へ狭義単調に変化する．
Cut の仮定により各 $(\mu,p)\in I\times P$ の逆像は一軌道である．
$K/S^1$ の compact 性から連続全単射は同相となる．
固定点の外では $dp$ は complex rank $1$ を持ち，
$dp(JV)=0$，$d\mu(JV)=-|V|_\eta^2\ne0$ であるから，
$(\mu,p)$ は滑らかな submersion である．Weights が $(1,-1)$ なので
固定点以外に有限 stabilizer もない．

持ち上げた葉の smooth 二次元接空間は $V$ と
$\ker d\mu$ 内の一方向で張られ，$\iota_V\eta=-d\mu$ より isotropic である．
固定点を含まない円及び saddle graph 上の主円束は，
基底の $H^2$ が零なので自明であり，所記の fibre 型を得る．
Source $q$ を一つ含む円葉 $\gamma$ では，
$\gamma\setminus\{q\}$ 上の主円束は $S^1\times(0,1)$ である．
閉 fibre は compact Hausdorff で，その補集合はただ一つの固定点である．
従ってこれは cylinder の一点 compactification，すなわち
$S^2$ の二点を同一視した一重 pinched torus である．

残余葉の局所 submersion を $a(p)$ とする．
これは円不変なので $\mu$ と Poisson 可換である．
Source で $\eta(o)$ は正の Hermitian 形式であり，円不変性から
$(x,y)$ の二つの異なる weights に関して対角である．
複素線形 rescaling 後，moment の二次項は
$\pi(|u|^2-|v|^2)$ となる．
$da(q)\ne0$ と $p-q=xy/(t h'(q))+O((xy)^2)$ により，
$a(p)-a(q)$ の二次項は $\Re(cuv)$，$c\ne0$ である．
Phase 変更と実線形な基底値の再規格化により，この組を
$F=(|u|^2-|v|^2)/2$，$H=\Re(uv)$ に直す．
$u=\alpha+\sqrt{-1}\beta$，$v=\gamma+\sqrt{-1}\delta$ と書くと，
\[
 (Q_1,P_1,Q_2,P_2)
 =2^{-1/2}(\alpha+\delta,\,\beta-\gamma,\,\beta+\gamma,\,\delta-\alpha)
\]
は $d\alpha\wedge d\beta+d\gamma\wedge d\delta=\sum_{j=1}^2dQ_j\wedge dP_j$ を満たし，
\[
 F=-(Q_1P_2-Q_2P_1),\qquad H=-(Q_1P_1+Q_2P_2)
\]
となる．従ってこれは focus--focus の標準二次型
\cite[Section~2]{ZungFocus}であり，局所的な node は非退化で，
零次元 stratum は固定点そのものである．
Regular torus は円上の主円束なので，normal 軌道はその
$H_1(T^2;\mathbb Z)$ の primitive class を表す．
$p=q$ 上で $\mu\to0$ とすると軸上のこの円が固定点へ縮むため，
それが vanishing cycle である．重複被覆による multiple cycle は生じない．
最後に $\phi^*\omega=\eta$ を用いて全ての strata と circle を移す．
\end{proof}

\begin{remark}[零 reduction の標準模型と ambient の滑らかさ]
標準 product
$\omega_0=(\sqrt{-1}/2)(du\wedge d\bar u+dv\wedge d\bar v)$，
$\mu=\pi(|u|^2-|v|^2)$，$w=uv$ では，$w\ne0$ 上の reduced form は
\[
 \omega_{\mathrm{red},\mu}
 =\frac{\sqrt{-1}\,dw\wedge d\bar w}
        {2\sqrt{(\mu/\pi)^2+4|w|^2}}.
\]
実際，$m=\mu/(2\pi)$，$\rho=|w|$ とすれば
$(|u|^2+|v|^2)/2=\sqrt{m^2+\rho^2}$ であり，
reduction は $\frac12 d\sqrt{m^2+\rho^2}\wedge d\arg w$ となる．
$\mu=0$ の商は位相的には $w$-平面だが，その symplectic density は
$w=0$ で $|w|^{-1}$ の特異性を持つ．この商を通常の滑らかな
symplectic 曲面とみなす必要はなく，前補題は ambient の
滑らかな形式と接空間を用いて Lagrangian 性を証明している．
この明式は標準模型の計算であり，任意の平均された $\eta$ に対して
Darboux 座標が元の residual 座標 $p$ 又は $w=t h(p)$ を保つとは主張しない．
Focus--focus の局所型と monodromy は既知の理論であり，ここで行ったことは
固定点を含む actual KE neck の exact transport と閉 fibre への適用である．
\end{remark}

\begin{corollary}[Moment に依存する残余関数の focus 判定]
\label{ke:source-focus-transversality}
前補題の固定 source $o=(0,0,q)$ で，$A(\mu,p)$ を通常の実座標
$(\mu,\Re p,\Im p)$ の滑らかな関数とする．
可換な写像 $F=(\mu,A(\mu,p))$ 自身が $o$ で非退化 focus--focus 型であるための
必要十分条件は $d_pA(0,q)\ne0$ である．
\end{corollary}
\begin{proof}
前補題の正規化を用いると，二次項は
\[
 Q_\mu=\pi(|u|^2-|v|^2),\qquad
 Q_A=A_\mu(0,q)Q_\mu+\Re(cuv)
\]
であり，$c\ne0$ と $d_pA(0,q)\ne0$ は同値である．
このとき $Q_\mu$ の倍を引いて前補題の focus 二次型を得る．
$c=0$ なら二次形式の span は高々一次元なので非退化ではない．
\end{proof}

この系は指定した momentum map の判定である．例えば残余 submersion $a$ に対し
$A=(a-a(q))^3$ と置くと同じ葉の分割を与えるが，写像の二次項は退化する．
本稿の source 近傍では $a$ 自身を用い，その germ を後続の相対補間で固定する．

\begin{corollary}[Source を含む遅い exhaustion 上の actual 円作用]
\label{ke:source-diagonal-moser}
上の source chart で固定した十分小さい $b>0$ を取る．
$r_t\to0$ を十分遅く選べば，$xy=tz$ の飽和集合
$K_{t,r_t}=\{|x|,|y|\le r_t^2,\ |z|\le b/2\}$ は
normal Moser map の source と像の双方に含まれる．変位と円軌道像の
ambient 直径は $O(1/|\log r_t|)$ で，source の固定点は移らない．
この近傍に actual $\omega_t$ に対する focus--focus germ が得られる．
\end{corollary}

\begin{proof}
まず $r$ を固定して $U_{t,r}=\{|x|,|y|<r,\ |z|<b\}$ とする．
$K_{t,r}$ からの最初の退出が $|z|=b$ で起こる場合，経路は
例えば $3b/4\le |z|\le b$ の固定 shell を横切る．これは source を避ける
固定残余 compact に属し，十分小さい $r$ と $t$ で
定理\ref{ke:fixed-base-neck-moser}の証明中の射影下界（又は局所積比較）が
$z$ の変位に一定の actual 費用を与える．この shell では元の $z$ 自体が
source-free pair の残余座標となる．

$|x|=r$ で退出する場合は $r^2\le|x|\le r$ を横切る．
補題\ref{ke:source-inclusive-fixed-shell}の shell 評価を，固定 source 座標
$xy=tz$，$|z|\le b$ に適用する．各 $r$ を先に固定してから全ての
十分小さい複素 $t\ne0$ を取れば，この横断費用は $r$ に独立に正である．
$z=0$ も同補題の固定成分 compact に含まれる．$y$ も同様である．

よって各固定 $r$ について $|t|$ を十分小さくすれば，$K_{t,r}$ と
$U_{t,r}$ の境界の間の actual buffer は $r$ に独立な正の定数以上である．
補題\ref{ke:short-circle-actual-buffer}と系\ref{ke:short-circle-image-buffer}は
作用の自由性を使わないので，$U_{t,r}$ 上の
$|V_{\rm src}|=O(1/|\log r|)$ に適用できる．
必要な比較的 compact な中間領域は，境界を少し内側に動かすか，
$K_{t,r}$ の十分小さい actual 距離近傍を用いて取れる．
$r_n\downarrow0$ と，各 $r_n$ に対する閾値 $t_n\downarrow0$ を選び，
$t_{n+1}<|t|\le t_n$ では $r_t=r_n$ とする．
これで被覆と変位上界を得る．$(x,y)$ は $X_t$ 上の正則座標であり，
$dx\wedge dy$ に対する対数体積は滑らかな $dd^c$-primitive を与える．
補題\ref{ke:source-reduction-pinched-fiber}の固定点での transport と
局所二次型の議論を $h(p)=p=z$，残余 submersion $\Re z$ に適用すれば，
固定点の保存と actual focus--focus germ を得る．
\end{proof}

この diagonal 選択は所定の $t$ 依存の速さを保証しない．
外側の foliation と接続する閉じた pinched torus は
定理\ref{ke:projected-source-neck}で構成し，
同一 moment の source-free 境界への接続費用は定理\ref{ke:source-all-moment-exit}で与える．

\subsection{Source を含む一つの射影的 neck}

\begin{lemma}[射影を保つ source-inclusive normal form]
\label{ke:projected-source-normal-form}
固定二重線 $D=D_{ij}$ 上で，三重点を除く固定 compact 領域の近傍を取る．
関連 KE 主稿 \texttt{pair:prop:bundle-conic} と同じ bundle projection を用い，
$h=-F_d/\prod_{k\ne i,j}L_k$ とする．固定幅の零切断近傍で
射影を保つ正則座標変換があり，退化族の式は
\[
 x'y'=t h_t(p),\qquad h_t=h|_D+O(t)
\]
となる．ここで $x',y'$ は元の二つの直線束に値を持ち，座標変換は
恒等写像から全固定階で $O(t)$ ずれる．$h_t$ の零点は旧 source の近傍に
一つずつある単純零点である．従って全基底上で一つの正則円作用
$(x',y')\mapsto(e^{2\pi\sqrt{-1}a}x',e^{-2\pi\sqrt{-1}a}y')$ を得る．

より正確には，基底と Hermitian norms を先に固定し，
$E_a=\{\|x\|<a,\ \|y\|<a\}$ と置くと，正数 $r,\epsilon$ を選べる．
座標変換 $\Phi_t$ は $|t|<\epsilon$ を含めて正則に依存し，
$E_{2r}$ 上で定義され，その像は $E_r$ を含む．$E_r$ 上の逆写像と共に，
各固定階 $q$ に対して
\begin{equation}\label{ke:eq:source-normal-form-common-domain}
 \|\Phi_t-\id\|_{C^q(E_{2r})}
 +\|\Phi_t^{-1}-\id\|_{C^q(E_r)}\le C_q|t|
\end{equation}
を満たす．Norm は固定した有限個の束 charts で測り，必要なら外側の
tube を最初に少し縮める．円作用の定義域は
$\Phi_t^{-1}(E_r)\cap X_t$ と取れば，その全 orbit が同じ定義域に入る．
各 source は $t=0$ を含む円板上の一価な正則断面として追跡でき，
$t\ne0$ ではこれらの断面が円作用の固定点集合である．ここでは基底を三重点へ動かした
ときの $r,\epsilon,C_q$ の一様性は主張しない．
\end{lemma}

\begin{proof}
束の値を持つ関数 $f_t(p,x,y)=xy-th(p,x,y)$ の fibre derivative は
$t=0$ の零切断で非退化な Hessian を持つ．陰関数定理から一意な
正則臨界断面 $a_t=O(t)$ を得る．臨界点の一意性により線形 frame の間で貼り合う．
$\xi=(x,y)-a_t$ とし，
$f_t-f_t(p,a_t)=\xi^{\mathsf T} Q_t(p,\xi)\xi$ と Taylor 展開する．
$Q_t$ は Hessian の線分積分で取れる対称な行列であり，
$Q_t=Q_0+O(t)$，$\xi^{\mathsf T}Q_0\xi=\xi_1\xi_2$ である．
$M=Q_0^{-1}Q_t$ は $Q_0$ に関し自己共役である．恒等行列の近傍で
二項級数により $S=M^{1/2}$ を取ると $S^{\mathsf T}Q_0 S=Q_t$．
Frame 変更で $M,S$ は共役変換するので，$\xi'=S\xi$ は束全体で定義される．
固定幅を縮めれば $\xi\mapsto\xi'$ は一様に恒等写像に近い微分を持つ正則同型である．
$t h_t=-f_t(p,a_t)$ とすれば所要の式を得る．
$a_t=O(t)$ により $h_t=h|_D+O(t)$ となり，零点は陰関数定理と
零点近傍の外の一様な非消滅から従う．

量化と定義域を補足する．基底 compact を含む少し広い固定近傍と有限個の
束 charts を先に選ぶ．その上で Hessian の非退化性，臨界断面の陰関数定理，
および $M$ の恒等行列からの距離を一様に取れる．平方根は
$\|M-I\|<1/2$ における同じ収束二項級数で定義するので，
$t$ の argument による枝の変更はない．各 frame で
$\Phi_t(z)=z+B_t(p,z)$ と書けば，外側の固定 tube 上で
$B_t=O(t)$ が全固定階で成立する．$|t|$ を小さくした後は
$\|D_zB_t\|<1/2$ であり，内側 tube の任意の $w$ に対する
$z=w-B_t(p,z)$ は外側 tube 内の固定小球で一意に解ける．
これが共通像と逆写像を与え，微分した同じ方程式から
\eqref{ke:eq:source-normal-form-common-domain}を得る．
Frame での一意な逆写像は重なりで一致する．

$f_t(p,a_t)$ は $t=0$ で恒等的に零なので，$h_t$ は $t=0$ にも正則に延びる．
各単純零点 $q$ に対する $h_t(p)=0$ の陰関数定理は，
一つの円板 $|t|<\epsilon$ 上で $q(t)=q+O(t)$ を与える．
円板を共通に縮めればこれらは互いに交わらず，他に零点はない．
標準回転の固定点は $x'=y'=0$ であり，$t\ne0$ では族の式から $p=q(t)$ に限られる．
元の座標では $(p,x,y)=(q(t),a_t(q(t)))$ である．
最後に $E_r$ は二つの束 norms により定義したため標準回転で不変であり，
$\Phi_t^{-1}$ による移送は定義域全体での円作用を与える．
従って部分 orbit から平均を定義する必要はない．
\end{proof}

\begin{lemma}[Source を横切る固定 end と shell]
\label{ke:source-inclusive-fixed-shell}
固定次数 $d\ge5$ の本稿の一般位置全平面退化を考える．三重点を避けた
二重線上の固定 compact $P\Subset\operatorname{int}P^+$ を取り，
補題\ref{ke:projected-source-normal-form}の一つの affine frame で
$x'y'=t h_t(p)$ と書く．$P^+$ は source を含んでよい．固定 source 座標
$xy=tz$ における同じ結論も含める．$V$ は $x'$ を正方向，$y'$ を負方向に
回す周期 $1$ の生成場，$\mathcal A$ はその円平均とする．
次の一つの体積形式と primitive を用いる：
\begin{equation}\label{ke:eq:source-inclusive-fixed-frame}
\begin{aligned}
 \Omega_t=d\log x'\wedge dp=-d\log y'\wedge dp,
 \\
 \lambda_t=d^c\log\frac{\omega_t^2}{|\Omega_t|^2},
 \qquad \mu_t=(\mathcal A\lambda_t)(V).
\end{aligned}
\end{equation}
最初の二つの表示はそれぞれ $x'\ne0$，$y'\ne0$ 上で読む．
この $\Omega_t$ は source にも非零正則に延び，$d\lambda_t=\omega_t$ である．

ある $\delta>0$ が存在し，十分小さい各固定 $r>0$ と $s_r=-\log r^2$ に対し
\begin{align}
 \mu_t|_{\{p\in P,\ |x'|=r\}}
   &=\frac{4\pi}{s_r}+O_{P^+}(s_r^{-1-\delta})+o_{t\to0;r,P^+}(1),
       \label{ke:eq:source-inclusive-positive-end}\\
 \mu_t|_{\{p\in P,\ |y'|=r\}}
   &=-\frac{4\pi}{s_r}+O_{P^+}(s_r^{-1-\delta})+o_{t\to0;r,P^+}(1)
       \label{ke:eq:source-inclusive-negative-end}
\end{align}
が全 base points と全 phase に一様に成り立つ．$O_{P^+}$ の定数は小さい $r$ に
共通であり，収束誤差を抑える $t$ の閾値は $r$ に依存してよい．
さらに $c_{P^+}>0$ を $r$ に独立に選べる．各固定 $r$ に対して十分小さい全 $t$ では，
\[
 \Sigma^x_{t,r}=\{p\in P^+,\ r^2\le|x'|\le r,\ y'=t h_t(p)/x'\}
\]
内で $|x'|=r^2$ と $|x'|=r$ を結ぶ全経路の actual 長さは
$c_{P^+}\log2$ 以上である．$x',y'$ を交換した shell でも同じ結論が成り立つ．
\end{lemma}

\begin{proof}
まず固定成分への収束の適用条件を確認する．全空間は平坦な射影 hypersurface family で，
relative canonical bundle は adjunction により
$K_{\mathcal X/\Delta}=\mathcal O_{\mathcal X}(d-4)$ であり，Cartier かつ relatively ample である．
中央 fibre は一般位置平面の reduced normal-crossing union なので semi-log canonical である．
全空間の特異点は source の $xy=tz$ のみで，三次元 ordinary double points である．
原点の blowup の discrepancy は $3-2=1$ なので全空間は normal かつ canonical である．
従って Song~\cite[Definition~2.3, Lemmas~4.4--4.5]{SongNegativeKE}の stable degeneration
に対する固定 compact の $C^\infty$ 評価と極限の一意性を適用できる．
関連 KE 主稿\cite{KELocalCurrent}の \texttt{pair:thm:fixed-end-moment} の証明と同じく，
各 component 上の極限は $K_{\PP^2}+D=\mathcal O_{\PP^2}(d-4)>0$ に対応する
complete cusp KE 計量である．この同定には Song~\cite[Lemma~3.7]{SongNegativeKE}の
一意性と，cusp potential が任意の正係数の logarithmic pole より緩やかなことを用いる．

$r$ を固定すると
\[
 (p,x')\longmapsto(p,x',t h_t(p)/x')
\]
は $P^+\times\{r^2\le|x'|\le r\}$ の少し広い近傍を，$t=0$ まで滑らかに
同定する．極限の $x'$ は零でないから，その像は中央成分 $y'=0$ の smooth locus の
固定 compact に入る．特に $h_0(p)=0$ でも中央 fibre の二重 locus を踏まない．
射影 Morse 変換は全固定階で $O(t)$ なのでこの固定 gauge の収束を保つ．
Song の compact 評価から任意の $t_\nu\to0$ に収束部分列があり，極限の一意性から
全複素 parameter $t\to0$ で同じ $C^\infty$ 収束を得る．
従ってこの領域上の $\log(\omega_t^2/|\Omega_t|^2)$ とその全固定階微分も，
$d\log x'\wedge dp$ で測った component の対数体積 $F_0$ へ収束する．

この成分上の唯一の局所 boundary divisor は $x'=0$ であり，source を指定する
$h_0(p)=0$ は追加 divisor ではない．$q=x'$，$s=-\log|q|^2$ とすると，
Rochon--Zhang~\cite[\S5, (5.5)--(5.6), Theorem~5.1]{RochonZhang}の表示は
\[
 F_0=-2\log(s+b(p,q))+B(p,q)+\upsilon(p,q),\qquad
 \upsilon-\Xi\upsilon_D\in s^{-\nu}C^\infty_{\rm fc},\quad \nu>0
\]
となる．$b,B$ は $q=0$ を跨いで滑らかである．固定 $P^+$ は他の boundary divisors
から離れているので，$\Xi\upsilon_D$ の radial derivative は
$O_{P^+}(s^{-2})+O_{P^+}(|q|)$ である．また $|q|\partial_{|q|}$ の cusp norm は
$O(s^{-1})$ だから，remainder の一階評価を合わせて
\[
 |q|\partial_{|q|}F_0=\frac4s+O_{P^+}(s^{-1-\delta}),\qquad
 \delta=\min\{\nu,1\}>0
\]
を得る．有限個の fixed patches と二つの components で定数と指数を共通に取る．
周期 $1$ の正の回転 $V_q$ について
$(d^cF_0)(V_q)=\pi |q|\partial_{|q|}F_0$ であり，円平均もこの一様評価を保つ．
$x'$-end では $V=V_{x'}$，$y'$-end では $V=-V_{y'}$ なので，固定 gauge の
収束誤差を加えて両 end の式を得る．

同じ component cusp 比較から，固定 $P^+$ 上の十分小さい全 $|x'|>0$ で
$g_0\ge c\,(d\log s)^2$ となる $c>0$ を取れる．固定 shell 上での smooth convergence により，
各 $r$ に対して $t$ を十分小さくすれば $g_t\ge(c/2)(d\log s)^2$ である．
Shell の両端では $s=2s_r,s_r$ なので，任意の横断経路の長さは
$\sqrt{c/2}\,|\log(2s_r)-\log s_r|=\sqrt{c/2}\log2$ 以上である．
これは shell 内の部分経路の評価であり，任意の全曲面経路が shell を横切る場合にも
その部分へ適用できる．Shell 外を許す二点間の距離を直接評価する主張ではない．

最後に primitive の規格化を確認する．$x'y'=t h_t(p)$ の微分を $dp$ と wedge すると，
二つの logarithmic 表示が重なりで厳密に一致する．Source 近傍では単純零点の仮定より
\[
 \Omega_t=\frac{dx'\wedge dy'}{t h'_t(p)}
\]
と非零正則に延びる．よって一つの $\lambda_t$ が全域に定義され，負 KE 方程式から
$d\lambda_t=\omega_t$ を得る．Source に追加の pole や留数補正は生じない．
この主張は今固定した同一の $\Omega_t$ に関するものである．別の非零正則 frame
$\widehat\Omega_t=b\Omega_t$ に変えれば primitive の差は
$-d^c\log|b|^2$ という閉形式となる．特に $b=b(p)$ なら normal moment への寄与は零だが，
residual periods が零とは限らない．異なる正規形の座標写像から定めた primitive の
無条件な一致も主張しない．
\end{proof}

\begin{theorem}[固定基底上の共通 actual neck と閉じた特異ファイバー]
\label{ke:projected-source-neck}
$P$ を三重点を避けた固定 compact な境界付き残余曲面とし，
source は内点にあり，境界近傍にはないとする．
$P$ 上に閉じた regular circle leaves と単純 saddle leaves を持つ
marked Morse foliation を与える．各 saddle leaf は一つの saddle を含み，
各 source は saddle のない相異なる regular leaves に置くとする
（動く source に合わせた小さい isotopy を許す）．
このとき，十分小さい $|t|$ について，ある共通 moment interval $I_t$ 上で
同じ円作用を用いる actual $\omega_t$-Lagrangian neck fibration が存在する．
Regular fibre は $T^2$，saddle fibre は $S^1\times\figeight$，
各 source の critical fibre は一重 pinched torus である．
\end{theorem}

\begin{proof}
$P$ を少し広げた固定基底 $P^+$ に補題\ref{ke:projected-source-normal-form}を適用する．
三重点の一つを無限遠とする affine 座標 $p$ を取り，二つの直線束を
自明化する．以下の絶対値はこの固定 frame による．
$r>0$ を固定し，外側を $p\in\operatorname{int}P^+$，$|x'|,|y'|<r$，
内側を $p\in P$，$|x'|,|y'|\le r^2$ とする．
Source の近傍では $z=h_t(p)$ が正則残余座標で
$x'y'=tz$ となる．この $t$ 依存座標でも系\ref{ke:source-short-circle}の
oriented arm の体積比較は一様である．実際，oriented 側
$|x'|\ge c|t|^{1/2}$ では元の平面座標 $x=x'+O(t)$ と $x'$ が比較可能で，
$s_x=-\log|x'|^2+O(1)$ である．もう一方は $y'=tz/x'$ と $y=y'+O(t)$
から $|y|=O(|t|^{1/2})$ である．Canonical residue は
$d\log x'\wedge dz$ の一様に有界な非零倍である．これは
$f_t$ と元の方程式の比が非零 unit で，座標変換の Jacobian と
$h'_t(p)$ が一様に非退化であることによる．従って capped density は
arm の annulus--disc envelope と同じ率を持つ．反対側も同様で，
central bidisc には Bochner 最大値原理を同じ作用に適用する．
Source を避ける固定部分では $h_t$ は一様に非消滅であり，
$x'y'=th_t(p)$ は $y'/h_t(p)$ によって固定 pair の積座標になる．
各 cap は元の座標と新座標の絶対値を $|t|^{1/2}$ で切ったものに関して
一様比較可能なので，ordinary pair の体積比較と
定理\ref{ke:fixed-base-neck-moser}の証明が適用できる．
有限被覆で外側領域全体の生成場長は $O(1/|\log r|)$ となる．

残余方向の最初の退出には $P^+\setminus P$ 内の source-free 固定幅を
横切る必要があり，そこでの射影下界によって一定の費用がある．
法線方向の退出については $r^2\le|x'|\le r$ 又は
$r^2\le|y'|\le r$ の shell を横切る．それぞれ $r$ を固定して
$t\to0$ とすれば，$P^+$ 全体で成分の smooth locus に固定 gauge で
収束する．座標変換は $O(t)$ であり，補題\ref{ke:source-inclusive-fixed-shell}の
smooth convergence と cusp 下界により，この横断費用も $r$ に一様に正である．
$ r=r_t\to0$ を対角的に選び，系\ref{ke:short-circle-image-buffer}を
適用して内側を source と像の両方に含む一つの Moser map を得る．
固定点を許すことは系\ref{ke:source-diagonal-moser}と同様である．

正則体積 frame $\Omega=d\log x'\wedge dp$ は $y'\ne0$ で
$-d\log y'\wedge dp$ と一致し，source では
$dx'\wedge dy'/(t h'_t(p))$ として非零に延びる．
従って $\lambda_t=d^c\log(\omega_t^2/|\Omega|^2)$ を全体で用いられる．
$\eta=\mathcal A\omega_t$ とし，$\mu=(\mathcal A\lambda_t)(V)$ を取る．
ここで normal Moser map には centered kernel $1/2-s$ を用いる．
同じ選択を三重点側にも用い，補題\ref{ke:normal-flow-naturality}で
共通の作用に関する flow の自然性を適用する．
内側 cutoff の $|x'|=r_t^2$ と $|y'|=r_t^2$ における $\mu$ は逆符号で，
絶対値は $c/|\log r_t|$ 以上ある．補題\ref{ke:source-inclusive-fixed-shell}に
半径 $r^2$ を代入し，まず各 $r$ で固定成分の remainder を主項より小さくしてから，
全複素 $t$ に対する収束誤差も同時に小さくする対角選択を用いる．
これには $h_t=0$ 上の cutoffs と全 phase が含まれ，primitive は上記の
一つの $\Omega$ から定めたもののままである．
さらに対角閾値を縮めて $|t|\sup_{P^+}|h_t|<r_t^4$ とし，
$I_t$ は両 cutoff の値の間にある $0$ を含む閉区間に取る．

これで補題\ref{ke:source-reduction-pinched-fiber}の仮定が
$R=r_t^2$，$h=h_t$，$I=I_t$ に対して満たされる．
同補題の商 $I_t\times P$ 上で残余 foliation の葉を取り，
exact Moser map で移せば，閉 fibre 全体の Lagrangian strata，
連結性と三種類の位相型を得る．各 source の vanishing cycle は
共通 normal circle の primitive class であり，その critical moment は正確に零である．
\end{proof}

ここで構成した neck は三重点を除く固定残余基底に沿う．
三重点 collar における同一の fiber partition は
系\ref{ke:neck-corner-saturated-union}，
特異 fibers を含む全 fiber 合併の正規化 ambient 直径評価は
命題\ref{ke:collapsed-tree-full-fiber-collapse}で別に与える．

\subsection{閉じた focus fibre の period と近傍全体の移送}

Focus 点の局所正規形と，compact singular fibre の飽和近傍全体の
symplectic germ を区別する．後者の Taylor 不変量と，primitive を指定した場合の
追加の定数項は V\~u Ng\d{o}c\cite[\S\S3, 7]{VuNgoc}の既知の不変量である．
以下では本稿の exact transport に必要な period 条件を直接記す．

\begin{lemma}[一重 focus 近傍での一つの残余 period]
\label{ke:focus-neighborhood-period}
$F:N\to D$ を円板上の proper な Lagrangian torus fibration とし，
唯一の特異 fibre $\Lambda=F^{-1}(0)$ は一つの非退化 focus 点 $o$ を持つ
compact connected pinched torus とする．$D$ を十分小さく取れば
\[
 H_1(N;\ZZ)\cong\ZZ,\qquad H_2(N;\ZZ)\cong\ZZ.
\]
また $\omega|_N$ は exact である．
正則 torus $T$ 上で，primitive な vanishing cycle を $\alpha$，
それを含む整数基底を $(\alpha,\beta)$ とすると，包含写像は
\[
 H_1(T;\ZZ)/\ZZ\alpha\ \xrightarrow{\ \sim\ }\ H_1(N;\ZZ)
\]
を誘導する．従って $N$ 上の滑らかな closed one-form $\zeta$ に対し
\begin{equation}\label{ke:eq:focus-residual-exactness}
 \zeta\text{ が exact}\quad\Longleftrightarrow\quad
 \int_\beta\zeta=0.
\end{equation}
Normal period $\int_\alpha\zeta$ は自動的に零である．
特に一つの外側の正則 fibre $T=F^{-1}(c)$，$c\ne0$ の近傍で $\zeta=0$ なら，
$\zeta$ は $N$ 全体で exact である．
\end{lemma}
\begin{proof}
Eliasson の局所正規形に適合する複素表示で
$q_1+\sqrt{-1}q_2=\bar z_1z_2$ とする
\cite[\S2]{ZungFocus}．局所の base 座標変更を用い，
$|z_1|,|z_2|<r$，$|\bar z_1z_2|<\delta\ll r^2$ で定める node block $B$ を取る．
同時の radial contraction が $B$ を原点へ縮めるので $B$ は可縮である．
二つの boundary faces はそれぞれ円板上の circle bundle であり，
corner $|z_1|=|z_2|=r$ は $|q|<\delta$ に現れない．
これらに沿って少し重ねた外側の regular block を $C$ とする．
$\Lambda$ の $B$ 外の部分は compact annulus である．
Base をさらに縮め，閉じた annulus block 上の boundary を含む proper submersion に
Ehresmann の議論を適用する．その collar を少し厚くした開集合を $C$ とすれば，
円板上の自明な annulus bundle と同じ homotopy type を持つ．従って
\[
 B\simeq\{*\},\qquad C\simeq S^1,
 \qquad B\cap C\simeq S^1\sqcup S^1.
\]
Mayer--Vietoris の写像 $H_1(B\cap C)\to H_1(C)$ は，
向きを揃えれば $(m,n)\mapsto m-n$ という全射である．
その kernel と，$H_0(B\cap C)\to H_0(B)\oplus H_0(C)$ の kernel はともに $\ZZ$ なので，
表示した $H_2$ と $H_1$ を得る．
$H_2$ の generator は，$B$ 内の二つの局所 Lagrangian discs と
$C$ 内の annulus からなる pinched fibre の二次 cycle である．
その上の $\omega$ の積分は零なので，de Rham の定理より $\omega|_N$ は exact である．
Normal circle は $B$ 内で潰れ，$\beta$ は二つの overlap components を
$B$ と $C$ を通って一周するので後者の kernel の generator へ写る．
これで包含写像の主張と，de Rham の定理による\eqref{ke:eq:focus-residual-exactness}が従う．
外側の任意の regular torus 上でも同じ $\beta$ が generator である．
\end{proof}

\begin{example}[Node の primitive germ を固定しても残る定数]
\label{ke:focus-local-primitive-shift}
前補題の regular block を $C\cong S^1\times(0,1)\times D$ と書き，
$(0,1)$ の座標を $s$ とする．両端近傍で零で $\int_0^1f(s)\,ds=1$ となる
滑らかな関数を取り，さらに $B\cap C$ の二つの collars 上で $f=0$ として，
$\zeta=f(s)\,ds$ を $B$ 側へ零延長する．
これは node 近傍で零の closed form で，$\int_\beta\zeta=1$ である．
従って $\lambda_c=\lambda+c\zeta$ は同じ $d\lambda$ と node 上の同じ primitive germ を持つが，
恒等写像の primitive 差は $c\ne0$ なら exact ではない．
この変更は fibration 自体と constant-zero Taylor 不変量を保ち，残余作用の原点だけを $c$ 変える．
非零の $c$ について，この closed 補正を base の外周 collar でも零にはできない．
そこには $H_1(N)$ の generator を持つ regular torus があり，前補題に反するからである．
\end{example}

\begin{proposition}[全飽和近傍の移送と focus 不変量]
\label{ke:focus-whole-neighborhood-transport}
前補題の $N$ 上で $\omega=d\lambda$ とし，
$\Phi:N\to N'$ が $\Phi^*\omega'=\omega$ を満たす微分同相であるとする．
$N'$ の葉の分割を $F'=F\circ\Phi^{-1}$ と定める．
全 compact fibre，その Lagrangian strata，primitive vanishing cycle と monodromy は
$\Phi$ によって移される．向き，Eliasson 座標，base cut と return cycle も一緒に運べば，
正則化作用の smooth germ $S$ が保存される．
さらに $\omega'=d\lambda'$ なら，
\[
 \Phi^*\lambda'-\lambda\text{ が exact}
 \quad\Longleftrightarrow\quad
 \int_{\Phi\beta}\lambda'=\int_\beta\lambda.
\]
この場合，全 closed cycles の作用積分と残余作用の定数項も保存される．
\end{proposition}
\begin{proof}
各 fibre は厳密に $\Phi(F^{-1}(c))$ なので，その位相型・連結性・strata が移る．
適合する Hamiltonians とその flows も共役されるため，return periods が一致する．
規格化を明示すると，normal の向きを固定し，周期 $1$ の moment を $\mu(o)=0$ と正規化する．
Eliasson 基底座標 $(c_1,c_2)$ を $c_2=\mu/(2\pi)$ と取り，
$c=c_1+\sqrt{-1}c_2$，$\alpha$ を normal cycle とする．
$c_1$ は一般には元の residual submersion そのものではない．
同じ slit disk と logarithm の枝，return cycle $\beta_c$ を選び，
$A(c)=\int_{\beta_c}\lambda$ と置く．V\~u Ng\d{o}c の規約では
\begin{align}
 dA&=\tau_1\,dc_1+\tau_2\,dc_2,\notag\\
 dS&=(\tau_1+\log|c|)\,dc_1+(\tau_2-\arg c)\,dc_2,\notag\\
 S(c)&=A(c)-A(0)+\Re(c\operatorname{Log}c-c),\qquad S(0)=0.
 \label{ke:eq:focus-regularized-action}
\end{align}
$dS$ は原点を越えて滑らかな閉形式となる\cite[Proposition 3.1, Definition 3.1]{VuNgoc}．
Return periods と正則化の選択が一致するので $S'=S$ である．
独立の向きを揃えた Eliasson 座標を使う際の flat な自由度については，
同文献の Lemma~4.1 と \S4 により Taylor series の一致を得る．

$\zeta=\Phi^*\lambda'-\lambda$ は closed なので，前補題が exactness の同値を与える．
一般には $A'(c)-A(c)=\int_\beta\zeta$ は定数であり，
$A'(0)-A(0)$ も同じ値なので $S$ では相殺される．
Exact の場合は任意の closed cycle 上の積分差が零となり，$A(0)$ も保存される．
これは\cite[\S7, ``Exact version'']{VuNgoc}の追加定数項である．
\end{proof}

Return cycle の変更 $\beta\mapsto\beta+k\alpha$ は $A$ を $2\pi k c_2$，
logarithm の枝の変更 $\operatorname{Log}c\mapsto\operatorname{Log}c+2\pi\sqrt{-1}\ell$ は
正則化項を $-2\pi\ell c_2$ だけ変える．両方を共通に運び，
$\partial_{c_2}S(0)\in[0,2\pi)$ の規格化を揃える．
作用を $J=A/(2\pi)$ で測る場合には，正則化も $S/(2\pi)$ とする．
上の命題は既に構成した葉の分割の push-forward を扱い，
独立に先指定した別の focus fibration との同値を与えるものではない．
実際，\cite[Theorem 2.1, \S6]{VuNgoc}により，規格化した Taylor series が
$0$ と $X^2$ の模型を実現できる．両者は同じ局所 focus 型と一重 monodromy を持つが，
その半大域不変量は異なる．この既知の実現原理の具体例と，
例\ref{ke:focus-local-primitive-shift}の primitive 原点だけの変更は別の現象である．

定理\ref{ke:projected-source-neck}への適用では，source を含む閉 residual leaf の
annular 近傍に submersion $a(p)$ を取り，$F=(\mu,a(p)-a(q))$ とする．
Source が相異なる regular leaves 上にあることから，他の source と saddle を避ける
小さい base disk と，その全逆像 $N$ を取れる．
Annulus の両境界と normal cuts の像を避けて base disk を縮めると，
任意の compact base subset の逆像は元の compact $K$ 内で閉じるので $F|_N$ は proper である．
共通 Moser domain はこの $N$ の全開近傍を含み，
補題\ref{ke:circle-exact-transport}は
$(\eta,\mathcal A\lambda_t)$ から $(\omega_t,\lambda_t)$ への
$\phi^*\lambda_t=\mathcal A\lambda_t+dQ$ を与えるので，
命題\ref{ke:focus-whole-neighborhood-transport}を適用できる．
別の canonical volume frame で primitive を作り直す場合は，
系\ref{ke:transition-residual-period}の残余 winding を別に確認する．
以上は各固定 $t$ と固定 stage の移送に関する保存則であり，
異なる $t$ や stage の focus 不変量が一定であるとは主張しない．

\subsection{残る円の平均と同一 normal moment}

今度は $\rho_{1,s}$ を neck の領域 $U$ 全体に与えた normal 円作用とする．
End の開領域 $E\subset U$ は $\rho_1$ に飽和しており，
$E$ 上には $\rho_1$ と可換な自由な正則円作用 $\rho_2$ があり，
両者で自由な $T^2$ 作用を生成すると仮定する．
各平均と homotopy を $\mathcal A_j,K_j$ と書き，
\[
 \eta=\mathcal A_1\omega,\qquad
 \kappa=\mathcal A_2\eta=\mathcal A_1\mathcal A_2\omega,\qquad
 \lambda_1=\mathcal A_1\lambda,\qquad
 \lambda_{12}=\mathcal A_2\lambda_1
\]
と置く．$\rho_2$ は $E$ のみに必要である．

第一の Moser flow $\phi_s$ は補題\ref{ke:circle-exact-transport}を
normal action に適用して $U$ 内に構成する．
第二の flow $\psi_s$ は $E$ 内で
\begin{equation}\label{ke:eq:residual-moser-path}
 \eta_s=(1-s)\kappa+s\eta,\qquad
 \beta_2=K_2\eta=\mathcal A_1K_2\omega,\qquad
 \iota_{Z_s}\eta_s=\beta_2,\quad \psi_0=\mathrm{id}
\end{equation}
によって構成する．以下では $\psi_s$ の定義域を $T^2$ に飽和した内側領域に取り，
$\psi_1$ の像が $\phi_s$ の共通定義域に含まれる buffer を課す．
$\phi=\phi_1,\ \psi=\psi_1,\ \Theta=\phi\circ\psi$ と書く．

\begin{theorem}[共通 normal circle と全作用周期を保つ二段階平均]
\label{ke:staged-neck-averaging}
上記の仮定の下で，$\psi_s$ は $\rho_1$ と可換であり，
\[
 \Theta^*\omega=\kappa,\qquad
 \Theta^*\lambda=\lambda_{12}+dq
\]
となる一価関数 $q$ がある．元の各 $T^2$ 軌道 $\mathcal O$ に対して
$\Theta(\mathcal O)$ は actual $\omega$-Lagrangian torus であり，
neck 全体の同じ actual 円作用
$\widehat\rho_1=\phi\rho_1\phi^{-1}$ の軌道で飽和している．

Normal moment を
$\mu=\lambda_1(V_1)$，
$\widehat\mu=\mu\circ\phi^{-1}$ と規格化すると，
\begin{equation}\label{ke:eq:staged-moment-identity}
 \mu\circ\psi=\mathcal A_2\mu,\qquad
 \widehat\mu\circ\Theta=\mathcal A_2\mu.
\end{equation}
より一般に，整数 cycle $b=(b_1,b_2)\in\ZZ^2$，
$V_b=b_1V_1+b_2V_2$ に対し，
元の軌道上のその一周を $\gamma_{b,x}$ とすると
\begin{equation}\label{ke:eq:staged-period-identity}
 \int_{\Theta(\gamma_{b,x})}\lambda
 =\int_{\gamma_{b,x}}\lambda_{12}
 =\lambda_{12}(V_b)(x).
\end{equation}
従って各 end で action の加法定数を独立に取り直す必要はない．
\end{theorem}
\begin{proof}
$\eta,\kappa,\beta_2$ は $\rho_1$ に不変である．
非退化形式に関する vector field の一意性から $Z_s$ は $\rho_1$-invariant であり，
flow の一意性により $\psi_s$ は $\rho_1$ と可換である．
補題\ref{ke:circle-exact-transport}を各段に用いると
\[
 \phi^*\omega=\eta,\quad \psi^*\eta=\kappa,\qquad
 \phi^*\lambda=\lambda_1+dq_1,\quad
 \psi^*\lambda_1=\lambda_{12}+dq_2.
\]
したがって $q=q_2+q_1\circ\psi$ とすれば所要の exact identity となる．
$\lambda_{12}$ は全 $T^2$ に不変なので，可換な生成場に対して
\[
 \kappa(V_1,V_2)
 =V_1\lambda_{12}(V_2)-V_2\lambda_{12}(V_1)
     -\lambda_{12}([V_1,V_2])=0.
\]
実次元二の自由な軌道は $\kappa$-Lagrangian であり，
$\Theta^*\omega=\kappa$ によりその像も Lagrangian である．
さらに
\[
 \widehat\rho_{1,s}(\Theta x)
 =\phi\rho_{1,s}\psi x
 =\phi\psi\rho_{1,s}x
 =\Theta\rho_{1,s}x
\]
なので，同じ actual normal circle が各 torus に含まれる．

Moment の等式は加法定数まで直接計算する．
$\iota_{V_1}\eta=-d\mu$ と可換性から
\[
 \beta_2(V_1)=K_2(d\mu)=\mathcal A_2\mu-\mu.
\]
$\mu_s=(1-s)\mathcal A_2\mu+s\mu$ と置けば
$\iota_{V_1}\eta_s=-d\mu_s$ である．従って
\[
 \frac d{ds}(\mu_s\circ\psi_s)
 =\bigl(\mu-\mathcal A_2\mu+d\mu_s(Z_s)\bigr)\circ\psi_s
 =\bigl(\mu-\mathcal A_2\mu+\beta_2(V_1)\bigr)\circ\psi_s=0.
\]
初期値は $\mathcal A_2\mu$ なので
\eqref{ke:eq:staged-moment-identity}を得る．
最後に exact differential の閉 cycle 上の積分は零であり，
$\lambda_{12}(V_b)$ は軌道上で一定なので
\eqref{ke:eq:staged-period-identity}が従う．
\end{proof}

\begin{corollary}[同じ normal action を持つ開 overlap 上の一致]
\label{ke:staged-normal-overlap}
Neck と triple chart に同じ $\omega$ と同じ normal 円作用 $\rho_1$ を与え，
同じ homotopy kernel を用いるとする．
内側 overlap からの normal Moser trajectories と，
residual stage の像からの trajectories が
$\rho_1$ に飽和した共通の開 overlap に留まると仮定する．
このとき両側の normal flow $\phi_s$ と actual normal action は一致する．
Triple 側に定理\ref{ke:staged-neck-averaging}の tori を用いれば，
それらは neck 側の同じ actual normal circle を含む．
また primitive $\lambda$ と cycle markings も同じものを使えば，
moment と作用周期は\eqref{ke:eq:staged-moment-identity}，
\eqref{ke:eq:staged-period-identity}によって共通に規格化される．
独立に固定した neck foliation の第二の周期や fiber 自体の一致には，
さらに reduced leaves を同じものに選ぶ必要がある．
\end{corollary}
\begin{proof}
共通 overlap では $\mathcal A_1\omega$ と $K_1\omega$ が一致する．
従って normal Moser vector field も一致し，
trajectory がその overlap に留まるという仮定の下で ODE の一意性を使える．
残りは定理\ref{ke:staged-neck-averaging}による．
\end{proof}

\begin{lemma}[短い作用に対する二段階構成の変位]
\label{ke:staged-displacement}
Normal action に不変な Hermitian metric $h$ を $U$ に取り，
$E$ 上では全 $T^2$ に不変とする．
$ch\le g_\omega\le Ch$，
$|V_1|_h\le e_1$ on $U$，
$|V_2|_h\le e_2$ on $E$ と仮定する．
上の二段階 flow を収める内在的 $h$-buffer があるなら，
\begin{equation}\label{ke:eq:staged-displacement}
 d_h(x,\Theta x)\le\frac{C}{2c}(e_1+e_2).
\end{equation}
各元の torus が内側領域に含まれる場合，その像の ambient diameter は
\begin{equation}\label{ke:eq:staged-ambient-diameter}
 \operatorname{diam}_{g_\omega}\Theta(\mathcal O)
 \le\sqrt C(e_1+e_2)\left(\frac Cc+\frac12\right).
\end{equation}
Centered homotopies を用いれば各段の変位定数 $C/(2c)$ は $C/(4c)$ に改善できる．
\end{lemma}
\begin{proof}
平均と凸結合は同じ二側比較を満たす．
$|K_1\omega|_h\le Ce_1/2$，
$|K_2\eta|_h\le Ce_2/2$ なので，各段の vector field の長さは
$Ce_j/(2c)$ 以下である．二つの trajectory を順に結べば
\eqref{ke:eq:staged-displacement}となる．
元の同じ torus の二点は，各 circle を高々半周ずつ進むことで
$h$-長さ $(e_1+e_2)/2$ 以下の経路で結べる．
その経路の両端に Moser trajectories を付け，
$g_\omega\le Ch$ を用いると\eqref{ke:eq:staged-ambient-diameter}を得る．
Centered kernel の絶対値の積分は $1/4$ なので最後の主張も従う．
\end{proof}

\subsection{一段階平均との Hamiltonian 比較}

二段階の endpoint map は，一段階の Moser map と一般に同一ではない．
ただし exact primitive を共通に固定した場合には，次の比較ができる．

\begin{proposition}[平均の順序と整数基底に関する Hamiltonian 比較]
\label{ke:staged-hamiltonian-comparison}
$E$ 上の同じ $T^2$ 作用と $\omega=d\lambda$ を用いる．
一段階の homotopy $H=K_1+\mathcal A_1K_2$ による map を $\Phi$，
上記の二段階 map を $\Theta$ とする．
以下の path homotopy の全 flow と inverse compositions を収める共通 buffer があるなら，
共通のより小さい開領域上で $\Phi^{-1}\Theta$ は
$\kappa$ に関する Hamiltonian isotopy の endpoint である．
従って $\Theta\Phi^{-1}$ は $\omega$ に関する Hamiltonian map である．

同じ結論は，$T^2$ の任意の二つの固定した整数基底で作る一段階または
二段階の平均の間にも成立する．
Centered homotopies を用い，$\omega$ が全 $T^2$ に不変な飽和 collar がある場合，
比較 isotopy はその collar を点ごとに固定できる．
\end{proposition}
\begin{proof}
三角形 $0\le v\le u\le1$ 上で
\[
 \Omega_{u,v}=(1-u)\kappa+(u-v)\eta+v\omega,\qquad
 \Lambda_{u,v}=(1-u)\lambda_{12}+(u-v)\lambda_1+v\lambda
\]
を定める．係数が非負で和が $1$ なので $\Omega_{u,v}$ は K\"ahler であり，
$d\Lambda_{u,v}=\Omega_{u,v}$ である．
$(0,0)$ から $(1,1)$ への対角 path は一段階構成，
$(0,0)\to(1,0)\to(1,1)$ の二辺は二段階構成に対応する．
二辺は各区間の端で stationary になるよう滑らかに parameter 化する．
この path と対角 path の凸結合によって，両座標が単調な滑らかな path family
$(u_r(s),v_r(s))$ を得る．

$\beta_1=K_1\omega$，$\beta_2=K_2\eta$ とし，
各 path 上の Moser primitive を
\[
 B_{r,s}=\dot u_r(s)\beta_2+\dot v_r(s)\beta_1
\]
とする．Homotopy identity より
\begin{align*}
 B_{r,s}
 &=-\partial_s\Lambda_{u_r,v_r}
       -d\bigl(\dot u_r K_2\lambda_1+\dot v_rK_1\lambda\bigr),\\
 dB_{r,s}&=-\partial_s\Omega_{u_r,v_r}.
\end{align*}
補題\ref{ke:circle-exact-transport}と同じ微分計算により，
各 endpoint map $F_r$ は
\begin{equation}\label{ke:eq:staged-path-exactness}
 F_r^*\omega=\kappa,\qquad F_r^*\lambda=\lambda_{12}+dQ_r
\end{equation}
を満たし，$Q_r$ は $r$ に滑らかである．
両座標の総変化は各々 $1$ なので，補題\ref{ke:staged-displacement}の
総変位の上界を全 $r$ に共通に使える．
さらに $F_r(D_{\mathrm{small}})\subset F_0(D_{\mathrm{large}})$ を保証する
image buffer を取ることが，相対写像の定義に必要である．

$W_r=(\partial_rF_r)\circ F_r^{-1}$ と置く．
\eqref{ke:eq:staged-path-exactness}を $r$ で微分すると
\[
 F_r^*(\iota_{W_r}\omega)
 =d\bigl(\partial_rQ_r-F_r^*(\lambda(W_r))\bigr).
\]
従ってこの deformation は Hamiltonian であり，
$F_0^{-1}F_r$ は $\kappa$-Hamiltonian isotopy となる．
端点は $F_0=\Phi$，$F_1=\Theta$ である．

整数基底を変えても全 Haar average $\mathcal A=\mathcal A_1\mathcal A_2$ は同じである．
二つの homotopies $H,H'$ について
\[
 H\omega=\mathcal A\lambda-\lambda-d(H\lambda),\qquad
 H'\omega=\mathcal A\lambda-\lambda-d(H'\lambda)
\]
なので，その差は exact である．
同じ線形 path $(1-s)\kappa+s\omega$ に対して primitive を
$(1-r)H\omega+rH'\omega$ と補間すれば，
再び endpoint の exact identity と上の Hamiltonian 計算が成立する．
二段階の場合は既に証明した比較を前後に組み合わせればよい．

最後に，centered homotopy は invariant form を零に写す．
指定 collar 上では $\beta_1=\beta_2=H\omega=H'\omega=0$ なので，
すべての vector fields と endpoint maps はそこで零および恒等となる．
従って比較 isotopy も同 collar を点ごとに固定する．
\end{proof}

\subsection{指定した normal subgroup を持つ actual closed corner}

\begin{corollary}[二段階 closed-corner bundle と同じ GH 誤差]
\label{ke:staged-closed-corner-gh}
系\ref{ke:actual-closed-corner}の固定次数 $d\ge5$ の族と exact triple chart を用いる．
その標準 $T^2$ の任意の固定した primitive normal subgroup $\alpha$ を選び，
$\alpha,\beta$ が $\ZZ^2$ の整数基底になるよう $\beta$ を補う．
$L=-\log|t|^2$，$R=\log L$ とし，
第一段を $\alpha$，第二段を $\beta$ とする上記の二段階 map $\Theta_t$ を取る．
十分小さい $t\ne0$ についてこの map は既存の softmax domain
$\mathcal H_t=\mathcal U_t(H_{L,R})$ の開近傍に定義できる．
\[
 \mathcal W_{t,J}^{\mathrm{st}}=\Theta_t(\mathcal H_t),\qquad
 \pi_{t,J}^{\mathrm{st}}=P_t\circ\Theta_t^{-1}
\]
とすると，その全 fibers は actual Lagrangian tori であり，
同じ actual normal action
$\widehat\rho_\alpha=\phi_t^\alpha\rho_\alpha(\phi_t^\alpha)^{-1}$ を含む．

模型 $\gamma_t=g_L^{\mathrm{sf}}$ に対し
\begin{equation}\label{ke:eq:staged-actual-metric}
 (1-C/R)\gamma_t\le\Theta_t^*g_t\le(1+C/R)\gamma_t.
\end{equation}
$R^{-2}g_t$ による $\mathcal W_{t,J}^{\mathrm{st}}$ の内在距離，
および $\mathcal V_t=\{\min\ell_a>R/2\}$ 内の全経路を許す距離のどちらについても
\begin{equation}\label{ke:eq:staged-closed-distortion}
 \sup_{x,y\in\mathcal W_{t,J}^{\mathrm{st}}}
 \left|d_t(x,y)-d_J(\pi_{t,J}^{\mathrm{st}}x,\pi_{t,J}^{\mathrm{st}}y)\right|
 \le C\frac{\log R}{R}.
\end{equation}
Fiber 自身の非正規化 intrinsic diameter は $O(R^{-1})$，
正規化後は $O(R^{-2})$ である．
Raw map と retraction に関する
\eqref{ke:eq:actual-closed-corner-target}と同じ誤差評価も成立する．
定数には固定した整数基底への依存を許す．

この bundle は前節の bundle と同一の写像としては主張しない．
両者は三重点の共通の開近傍上で actual $\omega_t$-Hamiltonian map により比較でき，
その map は同じ softmax base label を持つ fibers を対応させる．
\end{corollary}
\begin{proof}
補題\ref{ke:negative-exact}により三重点の actual form は exact である．
$h=h_t$ を constrained cusp metric とする．
外側領域 $\min\ell_a>R/2$ での全深部比較と模型の一様幾何から
\[
 ch\le\gamma_t\le Ch,\qquad
 |\nabla^h\gamma_t|_h\le C,\qquad |g_t-\gamma_t|_h\le C/R,
\]
\[
 |\omega_t|_h+|\nabla^h\omega_t|_h\le C,\qquad
 |V_\alpha|_h+|\nabla^hV_\alpha|_h
 +|V_\beta|_h+|\nabla^hV_\beta|_h
 \le C/R.
\]
生成場の評価は\eqref{ke:eq:full-corner-generator-c1}の固定整数線形結合である．
$h$ は全 torus に不変なので，$\eta_t,\kappa_t$ と両 convex paths も
$h$ に関する同じ二側比較および一様 $C^1$ bounds を持つ．
Isometry の pullback は $\nabla^h$ と可換であり，
chain contraction の微分には
$|\nabla^hV|_h|\omega|_h+|V|_h|\nabla^h\omega|_h$ だけが現れる．
従って各段の primitive とその一階微分は $O(R^{-1})$ であり，
非退化な Moser 方程式とその共変微分から
\[
 |Y_{t,s}|_h+|\nabla^hY_{t,s}|_h
 \le C/R
\]
を得る．これは補題\ref{ke:open-cell-moser-c1}および
\eqref{ke:eq:full-corner-moser-c1}と同じ計算である．

Constrained cusp norm では各 $d\log\ell_a$ が一様有界である．
従って $\min\ell_a>R$ からの二つの trajectories は，
$R$ が大きければ $\min\ell_a>3R/4$ に留まり，
$\min\ell_a>R/2$ との間の固定幅の buffer により全時刻で存在する．
必要な inverse compositions と path homotopy にも，少し大きい内側領域を用いて
同じ buffer を与えられる．

各段で，$\nabla^h\gamma_t$ の上界と上記の $C^1$ 評価により
$|\mathcal L_Y\gamma_t|_{\gamma_t}\le C/R$ であるから，
Gronwall の評価と二つの flow の合成により
$\Theta_t^*\gamma_t=(1+O(R^{-1}))\gamma_t$ を得る．
Actual/model 比較を最後に引き戻せば
\eqref{ke:eq:staged-actual-metric}となる．
模型の全閉 corner の距離評価
\eqref{ke:eq:closed-corner-model-input}と一様な正規化 diameter を使えば，
内在距離について\eqref{ke:eq:staged-closed-distortion}が従う．

保護領域内の距離には，系\ref{ke:closed-corner-protected-gh}と同じ
clipped projection $\Psi$ を用いる．二つの trajectories の模型長さの和は $C/R$ 以下であり，
\eqref{ke:eq:protected-pruning-length}から
\[
 d_J(\Psi(\Theta_tz),\Psi(z))\le C/R^2
\]
となる．保護領域全体の actual/model 下比較と，
既に得た内在距離の上界を合わせれば同じ distortion bound となる．
Fiber intrinsic diameter は，模型の phase torus の直径 $C/R$ と
\eqref{ke:eq:staged-actual-metric}から従う．
Raw map の評価は補題\ref{ke:raw-values}と
補題\ref{ke:raw-lipschitz}を二つの短い trajectories に用いればよい．

最後に，一段階・二段階および整数基底の変更に対して
命題\ref{ke:staged-hamiltonian-comparison}を適用する．
共通のソース形式は全 Haar average $\mathcal A\omega_t$ であり，
同じ softmax torus の二つの像を対応させる actual Hamiltonian map を得る．
\end{proof}

\subsection{Reduced foliations による延長と actual 適用の条件}

\begin{corollary}[共通 moment strip 上の条件付き neck 延長]
\label{ke:staged-reduced-extension}
定理\ref{ke:staged-neck-averaging}の normal 平均後の
$(U,\eta,\rho_1,\mu)$ において，共通の compact moment interval $I$ を取り，
reduced surfaces が境界付き曲面 $P$ の滑らかな束として同定されているとする．
さらに $P$ に，境界 coorientations と内部の embedded marking を固定した
基準 marked Morse foliation $\mathcal F_0$ があり，
境界近傍では radial，内部では simple saddles のみを持ち，
一つの critical leaf に一つの saddle だけを含むと仮定する．
各境界成分の互いに素な開 collar に，二段階構成から得た reduced circle foliations があり，
基準の radial circles をそれらへ写す collar 同定を $I$ に滑らかに選べるとする．
この同定は基準の collar 標識と coorientations に適合し，
互いに素な少し広い collars に台を持つ $\Diff_0(P)$ の滑らかな族へ延長できると仮定する．
このとき，それらの開 collar germs と一致する marked simple-saddle foliation family を
$P$ 上へ延ばせる．その円持ち上げを $\phi$ で移すと，
actual $\omega$-Lagrangian neck fibration が得られ，end では上記の staged tori と一致する．
その regular fiber は $T^2$，simple-saddle fiber は $S^1\times\figeight$ である．
\end{corollary}
\begin{proof}
仮定した $\Diff_0(P)$ の族により $\mathcal F_0$ を移す．
得られる族は指定した開 end germs と一致し，
内部の marked decomposition と一 critical-leaf 一 saddle の性質を保つ．
補題\ref{neck:marked-interpolation}と同様に，parameter germs の追加の補間が必要なら
collar 標識を固定した isotopies を用いる．
各固定 $\mu$ の reduced surface では一次元葉は isotropic であり，
その normal circle lift は Lagrangian である
（補題\ref{neck:lift}の reduction の証明）．
最後に $\phi^*\omega=\eta$ を用いて actual form へ移す．
\end{proof}

\begin{remark}[Actual neck への適用条件]
\label{ke:staged-neck-scope}
正則 normal action と K\"ahler 性から，
moment は各 complexified normal orbit に沿って狭義単調になる．
しかし，それだけでは全 base points に共通な moment interval や，
退化 parameter に一様な flow buffer は得られない．
系\ref{ke:staged-reduced-extension}では，これらの geometric identifications を仮定している．
本研究の actual KE 族では，source を除いた固定 compact base 上の正則 normal model に対し，
単一の global canonical primitive と共通 moment interval は
関連 KE 主稿\cite[\texttt{pair:cor:global-primitive}，
\texttt{pair:cor:actual-common-moment}]{KELocalCurrent}で検証される．
前者は base の nowhere-vanishing holomorphic $1$-form $\tau$ に対する
$\Omega=d\log v\wedge\tau$ と
$\lambda_t=d^c\log(\omega_t^2/(\Omega\wedge\overline\Omega))$ を用い，
$(\mathcal A_1\lambda_t)(V_1)$ が canonical normal moment と一致する．
後者は固定した end radius の後で $t\to0$ とする評価である．
Actual 適用には，さらに neck 全体の normal Moser flow の生存と buffer，
triple の各 end における同一作用・reduced foliation の同時一致を用いる．
これらは補題\ref{ke:fixed-stage-source-neck}と
系\ref{ke:neck-corner-saturated-union}で検証する．
Critical actions の直線性を保つには，さらに節\ref{ng8:section}の
relative area と marking の条件が必要である．

系\ref{ke:staged-closed-corner-gh}は，三重点の標準 torus 内の任意の
指定 primitive subgroup に適合できることを示す．
別々に構成した neck action との開 collar 上の一致を，
一つの三重点の複数の end で同時に満たすことは，同系だけの結論ではない．
節\ref{ke:relative-corner-section}で平均順序を相対化し，
系\ref{ke:neck-corner-saturated-union}で全交わりを同じ葉に揃える．
その後の source・branch・core を含む全域接着は
定理\ref{ke:actual-global-fibration}による．
非定数 holomorphic unit による orbit partition の変更も，
単なる整数基底変更の Hamiltonian 比較では処理されない．
全曲面の距離と singular fibers の一様な ambient 径は，
定理\ref{ke:global-ambient-lower}，\ref{ke:global-ambient-upper}および
命題\ref{ke:collapsed-tree-full-fiber-collapse}で，別の解析入力を組み合わせて証明する．
\end{remark}

\input{sections/paper/ke_neck_corner_transition.tex}
\input{sections/paper/ke_relative_corner_patch.tex}
% !TeX root = ../../surface-lagrangian-torus-fibration.tex
% 2026-11-21 / v1.3: 全 stage に共通な reduced marking と source-inclusive な飽和被覆．
\section{Source を含む neck と三重点 end の共通 exhaustion}
\label{ke:neck-corner-exhaustion-section}

固定残余基底の定理\ref{ke:projected-source-neck}と，
moving arm の命題\ref{ke:arm-common-moment-strip}を接続する．
選択の順序は，残余 compact 領域，normal cutoff，moment interval，
最後に退化 parameter の閾値である．この順序によって，固定 compact の定理を
未証明の moving-base 一様評価として用いることを避ける．
最初に source と saddle を含む neck の end germs を三重点 bundle と合わせ，
さらに cut の位置を離して選ぶことで，closed corner と切詰めた neck の
全ての実際の交わりで fiber partition を一致させる．
本節の結論は部分 fibration の接着である．Core への延長と全域被覆は
定理\ref{ke:actual-global-fibration}，全曲面を通る経路の距離比較は
定理\ref{ke:global-ambient-lower}，\ref{ke:global-ambient-upper}で別に扱う．

\subsection{残余基底と end 座標の区別}

二重線 $D=D_{ij}$ と，節\ref{ke:neck-corner-transition-section}で固定した
許容射影を用いる．三重点の集合を $Q\subset D$ とする．
射影 Morse 変換の後の方程式は，束の局所 frame において
\[
 x'y'=t h_t(p),\qquad p\in D\setminus Q
\]
である．Neck の残余基底は元の $p\in D$ である．
$h_t$ は束の section を frame で表示した関数であり，
全基底上の一対一座標ではない．単純 source 近傍で $h_t'(p)\ne0$ のときに限り，
$z=h_t(p)$ を局所座標として $x'y'=tz$ を使う．

一つの三重点 $q_k$ の近傍では local base coordinate を $p_k$ とし，
\eqref{ke:eq:arm-gauged-coordinates}の第三座標を
\[
 Z_{k,t}(p_k)=\Xi_{kk}(p_k)=C_D(p_k,t)^2/h_{k,t}(p_k)
\]
と書く．添字 $k$ は $h_t$ の当該局所 frame での表示も表す．
\[
 Z_{k,0}(p_k)=p_k/C_D(p_k,0)
\]
は $q_k$ で単純零点を持つ正則座標である．
各固定 punctured annulus 上では $Z_{k,t}\to Z_{k,0}$ が全固定階で成り立つ．
従って $t$ を十分小さくすれば，同じ annulus の少し小さい部分で
$Z_{k,t}$ も一対一の正則座標である．これは $h_t$ を global coordinate とする主張ではない．

Global neck の frame と各 end の frame の変換は，残余基底だけに依存する．
補題\ref{ke:projected-source-normal-form}の臨界断面と恒等近傍の平方根は
frame 変更に関して同変なので，normal circle は両表示で厳密に同じである．
一つの三重点を無限遠とする base coordinate $p_g$ と束の frame を取れば，
全ての compact exhaustion 上で $\Omega_N=d\log x'_g\wedge dp_g$ を使える．
Local frame における同じ体積形式との差は base-only unit である．
Canonical primitive の差は normal generator を消す閉形式なので，
normal moment の加法定数は変わらない．
第二 period については同じ結論を主張しない．

\begin{lemma}[固定残余 compact 上で先に選べる normal cutoff]
\label{ke:fixed-stage-source-neck}
$P\Subset\operatorname{int}P^+\subset D\setminus Q$ を固定し，
source は $P$ の内部に，$P^+\setminus\operatorname{int}P$ は source の外に置く．
固定 frame で $r=e^{-n/2}$ とし，
\[
 U_{t,n}=\{p\in\operatorname{int}P^+,\ |x'_g|,|y'_g|<r\},
 \qquad
 K_{t,n}=\{p\in P,\ |x'_g|,|y'_g|\le r^2\}
\]
とする．$P,P^+$ に依存する定数 $C_P,\delta_P,c_P>0$ と $N_P$ があり，
任意の固定した $n\ge N_P$ に対して，$|t|$ を十分小さくすれば次が成立する．
正則 circle generator $V$ は outer の作用不変領域 $U_{t,n}$ 全体で
\begin{equation}\label{ke:eq:source-fixed-stage-generator-bound}
 \sup_{U_{t,n}}|V|_{g_t}\le C_P/n
\end{equation}
を満たす．従って各 orbit の長さも $C_P/n$ 以下である．
さらに $K_{t,n}$ と外側の境界の actual buffer は
$\delta_P$ 以上であり，内側の二つの normal cut の canonical moment は逆符号で
絶対値が $c_P/n$ 以上である．
従って $n$ をさらに大きくして $C_P/n\ll\delta_P$ とすれば，
centered normal Moser map $\phi_{N,t}$ の source と像は共に $K_{t,n}$ を含む．
その actual 変位は $C_P/n$ 以下である．
十分小さい対称 interval を $(-c_P/n,c_P/n)$ 内に選ぶと，
円作用の商は $(\mu_N,p)$ によりその interval と $P$ の積に位相同定される．
\end{lemma}

\begin{proof}
定理\ref{ke:projected-source-neck}の証明を，$r$ を先に固定して適用する．
固定された source-free 部分には定理\ref{ke:fixed-base-neck-moser}の比較を使い，
source 近傍には系\ref{ke:source-short-circle}を使う．有限被覆により
生成場の\emph{点ごとの}上界の定数を一つの $C_P$ にできる．
実際，source-free patch では
定理\ref{ke:fixed-base-neck-moser}の比較から
$|V|_{g_t}\le C/\min(s,M_t-s)$ であり，
source patch では系\ref{ke:source-short-circle}から
$|V|_{g_t}\le C/\min(s_x,s_y,L_t)$ を得る．
$U_{t,n}$ 上ではいずれの分母も $n-O_P(1)$ 以上で，
固定 frame の有界なずれは $n$ を大きくして吸収できる．
従って\eqref{ke:eq:source-fixed-stage-generator-bound}を得る．
これは単に一周の積分長から点ごとの速度を推論したものではない．
残余方向では $P^+\setminus P$ の固定幅に射影下界を使う．
Normal 方向では $r^2\le|x'_g|\le r$ または対称な shell の成分極限と
cusp 下界を使う．この shell の深度比は $2$ なので，$n$ を大きくした後にも
固定された $P^+$ 上で正の下界 $\delta_P$ を選べる．
同じ成分の fixed-shell canonical moment 評価は normal cuts 上で
$c_P/n$ の符号付き下界を与える．
これらは各固定 $n$ の後に $t$ を小さくする順序の評価である．

系\ref{ke:short-circle-image-buffer}により，共通 source と像を持つ centered flow を得る．
Source の固定点では生成場と chain primitive が零であり，flow は固定点を保つ．
Moment の radial 単調性と source の二軸の記述は
定理\ref{ke:projected-source-neck}と同じであり，商の位相同定を与える．
この段階では $P$ を動かした定数の一様性を主張していない．
\end{proof}

\subsection{全 stage に共通な reduced marking}

\begin{lemma}[一つの marked pants 木系の共通選択]
\label{ke:fixed-marked-tree-exhaustion}
固定した退化族の各二重線 $D=D_{ab}$ に対し，一つの embedded marked
simple-saddle foliation $\mathcal F_D^0$ を $D\setminus Q$ 上に選べる．
これは固定 compact $K_D$ の外で $|Z_{k,0}|=\mathrm{constant}$ の円を葉とし，
各 source は互いに異なる full regular circle leaves に属する．
全 critical leaves の小さい飽和 pants neighborhoods と，これらの source leaves の
小さい飽和 annuli は $K_D$ の内部に入る．

本節の全ての十分大きい exhaustion stages と，それぞれの閾値以下の全ての
複素 $t\ne0$ について，この一つの葉層を基準として neck を構成できる．
各 end の全 regular circle で切った compact reduction の Reeb tree は，
葉の labels $k\notin\{a,b\}$，saddle labels，および $\mu=0$ の source marks を保って，
一つの固定した $T_{ab}=T_{ba}$ に同定される．後続の
命題\ref{ke:link-relative-collar-gluing}の相対補間と
補題\ref{ke:whole-fiber-core-trim}の切詰めを行っても，この木系の組合せ型は変わらない．
これは一つの許容された木系の存在であり，任意に指定した木系又は作用周期への
適合を主張しない．
\end{lemma}

\begin{proof}
\emph{固定した pants と full leaves．}
三重点の互いに素な小円板を除いた球面を取り，一つの labelled boundary を root とする．
三つ穴球面を木に沿って貼り，各 pants に saddle 一個の標準 Morse function を置く．
Regular annuli 上で単調に接続し，saddle の値を互いに異なるよう選ぶと，
内部の critical points が $d-4$ 個の index-one points だけである first integral
$h_D^0$ を得る．各 end へ radial に延長する．
有限個の source 点の近傍で，boundary-relative な小さい diffeomorphism により
葉層を動かし，全 source 値を互いに異なる regular values にできる．
実際，各 source の互いに素な小円板内に，既に選んだ有限個の値と全 critical values を
避ける regular point を一つずつ取り，それを source 点へ移す bump isotopy を使えばよい．
その後の first integral を改めて $h_D^0$，その葉層を $\mathcal F_D^0$ とする．

Source 点 $q_r(0)$ を含む full regular circle を $L_r$ と書く．
各 $L_r$ の小さい飽和 annulus $A_r$ と，各 full critical leaf の飽和 pants
$Q_v$ を互いに素に取れる．これら全体と，その間の有限個の regular annuli を
一つの固定 compact $K_D$ の内部へ入れる．$K_D$ の外では end の radial 葉を使う．
End collars を長くしても新しい vertex は加わらず，その外端で切った Reeb tree は，
leaf edges の長さを変えるだけで同じ labelled $T_{ab}$ になる．
$D_{ab}=D_{ba}$ は同じ実際の曲線なので，二方向に別の pants 分解を選ばない．

\emph{Moving sources の一価な標識．}
補題\ref{ke:projected-source-normal-form}の $h_t$ は，$t=0$ を含む円板上で
正則である．同補題の臨界断面は陰関数定理による一意な断面であり，行列の平方根は
恒等行列の近傍の級数で定められている．従って $t^{1/2}$ の枝を選ぶ操作はない．
$h_0$ の単純零点 $q_r(0)$ ごとに，$\overline U_r\subset\operatorname{int}A_r$ となる
互いに素な座標円板 $U_r$ を固定する．正則陰関数定理と $\partial U_r$ 上の非消滅により，
ある一つの円板 $|t|<\epsilon_D$ 上で，各 $U_r$ 内の唯一の零点
$q_r(t)$ が正則かつ一価に定まる．特に source labels の置換は起こらない．
この選択は $K_D$ 上で先に行うので，後の exhaustion stage に依存しない．

$z_r(q_r(0))=0$ とし，$\chi_r$ を $0$ の近傍で $1$，$U_r$ の境界近傍で
$0$ となる固定 bump function とする．$a_r(t)=z_r(q_r(t))$ と置き，
\[
 z_r(S_t(p))=z_r(p)+\chi_r(z_r(p))a_r(t)\quad(p\in U_r),
 \qquad S_t(p)=p\quad(p\notin\textstyle\bigcup_r U_r)
\]
と定める．$|t|$ を小さくすれば微分は恒等写像に十分近く，境界近くで恒等なので
これは diffeomorphism である．$a_r(t)$ を $u a_r(t)$，$0\le u\le1$ と置き換えると
同じ支持を持つ isotopy を得る．従って $S_t\in\operatorname{Diff}_0(D\setminus Q)$，
$S_0=\mathrm{id}$ であり，実二変数 $(\operatorname{Re}t,\operatorname{Im}t)$ に滑らかに依存する．
全 $Q_v$ と boundary/end collars は点ごとに固定する．
また $S_t(A_r)=A_r$ であり，$S_t(L_r)$ は $q_r(t)$ を含み，
他の source を含まない full regular circle である．

\emph{全 stage の end 補正．}
$b_j$ を十分大きいところから始め，$K_D\subset\operatorname{int}P_j$ とし，
本節で使う end の補正 annuli を全て $K_D$ の外に置く．
命題\ref{ke:single-arm-reduced-graph}と
補題\ref{ke:reduced-collar-graph}で使う normalized coordinates では，
指定された reduced 葉は $a=Q_{j,\mu,t}(c,\theta)$ という小さい circle graph である．
同補題の明示的な写像
\[
 D_{j,\mu,t}(c,\theta)
 =\bigl(c+\chi(c)(Q_{j,\mu,t}(c,\theta)-c),\theta\bigr)
\]
を使い，外へ恒等写像で延長する．第一成分の $c$ 微分は正であり，
$\chi$ を $u\chi$ に変える isotopy も同じ支持を持つ．
Graph の小ささ，共通 moment interval，全 trajectories の余裕は，
次の定理\ref{ke:neck-corner-diagonal-matching}で選ぶ既存の stage 条件そのものである．
ここで新しい moving-base 一様評価は必要ない．

従って neck の reduced 葉族を
\[
 \mathcal F_{j,\mu,t}
   =(D_{j,\mu,t}\circ S_t)_*\mathcal F_D^0
\]
と選べる．二つの補正の支持は互いに素である．全 full critical leaves は
同じ $Q_v$ にあり，source leaves は全 $\mu$ で $S_t(L_r)$ のままである．
End の cut は $D_{j,\mu,t}$ による radial circle の像で行うので，
内側の葉を切断せず，その Reeb tree は同じ $T_{ab}$ である．
特に $\mu=0$ で一つの source だけを含む閉じた円葉を丸ごと持ち上げられる．
これは点の近傍の marking だけを保つ主張ではない．

\emph{Core collars と critical values．}
後続の相対接着では $q_a=\delta_j=4\pi/N_j$ の葉を anchor に取り，
Liouville flow による葉族との比較を
定理\ref{ke:saddle-moving-relative-interpolation}の
boundary-relative な diffeomorphisms によって行う．
各 moment の全 partition が一つの diffeomorphism の像なので，
full critical leaf の連結性と一 saddle の性質，regular annuli の接続順序が保たれる．
支持は $\mu=0$ を避けるため，source の full pinched fibers とその marking は変わらない．
Source 点 $q_r(t)$ を非零 moment の reduced 曲面上でも marked points として指定する．
Anchor では全て相異なる regular leaves 上にあるので，小さい飽和 annuli 内で
regularity と分離を保つよう band を縮め，同定理の fixed-marked-point 版を用いる．
全 fibers ごとの trim と cone filling は，この同じ labelled end を corner annuli で
接続するため，$N_j$ の変更によって別の木を選ぶ操作にはならない．

初期の $h_D^0$ では source 値と critical values を分離して選んだ．
その後も葉の partition は上の diffeomorphisms によって同定されるので，
異なる source leaves が合体したり，source が saddle の full leaf に入ったりしない．
別々の connected leaves に与える数値が，first integral の再表示により一致しても，
それだけでは quotient の点は同一視されない．
同じ connected leaf に二つの saddle が生じる変形は用いていない．
ここで固定した補助 first integral の値は，第二 action の値を指定する条件ではない．

\emph{Parameter を一周するときの範囲．}
$S_t$ は $t=0$ を含む円板上の一価な公式であり，
end では各固定 annulus 上の $Z_{k,t}/Z_{k,0}$ が恒等に近い正則 unit なので，
その対数も $t=0$ で $0$ となる枝で取れる．
Collar graph の補正は circle 座標 $\theta$ に関して周期的で，
labelled end を入れ替えない．各 $t$ でこの基準から marking を移すため，
$\arg t$ を一周させても source labels と抽象的な pants 木系を同じものに選べる．
これは $X_t$ 全体の parameter monodromy の自明化ではない．
K2 で共通にする周期 $1$ の normal circle とその向きは保つが，
torus の全 $H_1$ 基底を punctured parameter 円板全体で固定するとは主張しない．
第二 cycle は各固定 $t$ の共通 torus 上で選ぶ従来の範囲に留める．
相対 core 補間も各 $t$ で同じ anchor marking から行えばよく，
stage 間又は $t$ 全体で滑らかな fibration 族を作る必要はない．
有限個の二重線についてこれらの選択を同時に行い，木系 $T=(T_{ab})$ を得る．
\end{proof}

この補題は，既知の Morse 葉層の isotopy 一般論を新たな分類定理として主張するものではない．
Morse mapping space の path components の分類
\cite[Main Theorem]{Maksymenko}だけからは，指定した full critical leaves 又は
embedded pants marking の不変性は従わない．
ここでは小円板と互いに素な end annuli に支持を持つ具体的な isotopies を用いた．

\subsection{三つの尺度を分ける diagonal 選択}

\begin{theorem}[Source-inclusive neck と全 end の開 collar の一致]
\label{ke:neck-corner-diagonal-matching}
固定二重線 $D$ 上に，source を相異なる regular leaves に置く
marked simple-saddle foliation を取り，三重点の punctured ends では
$|Z_{k,0}|=\mathrm{const}$ の円を葉とする．
この foliation は，source と saddle を含む固定 compact 部分の外では radial とする．
任意に遅い段階的な選択として，列
\[
 b_j\longrightarrow\infty,\qquad
 n_j\longrightarrow\infty,\qquad
 \iota_j>0,\qquad t_j\downarrow0
\]
を次の性質を満たすよう選べる．
$t_{j+1}<|t|\le t_j$ のとき，$I_j=[-\iota_j,\iota_j]$ 上に
source と simple saddles を含む actual $\omega_t$-Lagrangian neck があり，
各 end の開 collar で三重点側の指定 normal-first torus bundle と
fiber partition が厳密に一致する．Normal Moser map と canonical normal moment も一致する．
Source の一重 pinched torus，saddle の $S^1\times\figeight$，
および内部の marking は保持される．

この一致 collar は，triple の射影座標で
\[
 b=-\log|Z_{k,t}(p_k)|^2\asymp b_j,
 \qquad a,c\ge4n_j
\]
となる moment strip から構成する．$n_j/b_j$ は十分大きく選び，
$b\asymp b_j$ で $a$ または $c$ が $b_j$ 程度しかない領域まで
neck が覆うとは主張しない．
法線円作用の Moser 変位は neck 全体で $1/j$ 以下にできる．
また $m_t=b_j/64$ と取った三重点の cut は
\[
 m_t\to\infty,\qquad \log m_t=o(R_t)
\]
を満たすように選べるので，定理\ref{ke:slow-cut-closed-corner-gh}と両立する．
この結論は時間 parameter に関する微分可能性を要求しない．
\end{theorem}

\begin{proof}
\emph{残余領域を最初に固定する．}
初期葉層と動く source の標識には，補題\ref{ke:fixed-marked-tree-exhaustion}の
一つの共通選択を用いる．
十分大きい $b_j\uparrow\infty$ を取り，各 end の
$s_{k,0}=-\log|Z_{k,0}|^2$ を用いて
\[
 P_j=D\setminus\bigcup_k\{s_{k,0}>8b_j\},\qquad
 P_j^+=D\setminus\bigcup_k\{s_{k,0}>16b_j\}
\]
とする．大きい $j$ では，source と saddle は全て両領域の固定内部にあり，
残余 annuli $b_j/2<s_{k,0}<4b_j$ は互いに素である．
これらは各段階では固定 compact 領域であり，
補題\ref{ke:fixed-stage-source-neck}の定数を
$C_j,\delta_j,c_j,N_{P,j}$ と書く．
$N_{P,j}$ はこの固定残余領域に対する normal cut の下限であり，
後で core を作るために選ぶ component cut $N_j$ とは別の量である．

各 end の local frame，$C_D$，global neck の frame の間に現れる
対数絶対値の差を，上の少し広い固定 annuli 上で $M_j$ 以下に抑える．
$t$ に依存するものも，固定領域で $t=0$ の近傍へ制限して同じ上界に入れられる．
従って local normal 深度と global frame の深度の差は $M_j$ 以下である．
無限遠に置いた end の frame 変換にもこの bound を適用する．

\emph{Normal cutoff を次に選ぶ．}
$n_j$ は増加列とし，
\[
 n_j\ge N_{P,j},\qquad n_j>4M_j,\qquad
 n_j>100(1+2\sigma)b_j,\qquad
 \frac{C_j}{n_j}<\min\left(\frac1j,\frac{\delta_j}{100}\right)
\]
を満たすよう選ぶ．必要なら最初の有限個の項は除く．
$r_j=e^{-n_j/2}$ とし，global inner tube の深度は $2n_j$，
outer tube の深度は $n_j$ となる．

End では，命題\ref{ke:arm-common-moment-strip}と同じ canonical volume の計算を，
固定比 $a=Kb$ の代わりに $a=4n_j$ で行う．
Dominant-product 比較は $a/b$ に定数が依存しないので，
$b_j/2<b<4b_j$ において
\begin{equation}\label{ke:eq:exhaustion-normal-cut-moment}
 \mu,\nu\big|_{a=4n_j}
 =\frac{\pi}{n_j}
  +O\left(\frac1{n_jb_j}
         +\frac1{n_j}\left(\frac{n_j+b_j}{L_t}\right)^2\right)
\end{equation}
となり，$c=4n_j$ では主項の符号が逆になる．
座標接続の指数誤差も $C/(n_jb_j)$ に吸収される．
この計算では小さい triple primitive を使い，
系\ref{ke:transition-normal-primitive}によって normal moment を global neck の
primitive と厳密に同定する．

\emph{Moment interval を選ぶ．}
各 $j$ について
\[
 0<\iota_j<\min\left\{\frac{c_j}{16n_j},\frac{\pi}{64n_j}\right\}
\]
とする．構成は目標 interval $I_j$ より広い $|\mu|<2\iota_j$ の上で行う．
固定段階の global normal cuts と
\eqref{ke:eq:exhaustion-normal-cut-moment}の四つの境界値の間に
この広い interval が入るよう，後で $t$ の閾値を取る．
Normal 単調性により，三重点側の $\nu$-level と $\mu$-level は全て
$a,c>4n_j$ にある．Global neck 側では，source を含む全 $P_j$ 上で
同じ interval の円，または $\mu=0$ の source 固定点を持つ．

\emph{最後に $t_j$ を選ぶ．}
各段階で有限個の要求を同時に満たす閾値を選ぶ．具体的には，固定 $P_j^+$ と
$r_j$ に対する補題\ref{ke:fixed-stage-source-neck}の全評価，
source 零点の内部への保持，frame と Morse 座標の共通性，
$|t|\sup_{P_j^+}|h_t|<r_j^8$，および
\eqref{ke:eq:exhaustion-normal-cut-moment}の誤差を $\pi/(4n_j)$ 以下にする条件を課す．
さらに各固定 annulus 上で $Z_{k,t}/Z_{k,0}$ の正則対数が存在し，
\[
 (q_t,\theta_t)=
 \left(\log\frac{-\log|Z_{k,t}|^2}{b_j},\arg Z_{k,t}\right)
\]
とその $t=0$ 版の $C^3$ 差が $b_j^{-2}$ 以下になるようにする．
この微分は固定 annulus の座標での微分であり，各 $j$ を固定した後には
$t\to0$ により実現できる．

$t_j<t_{j-1}/2$ とし，さらに
\[
 R_t\ge j\bigl(1+\log(n_j+1)\bigr),\qquad
 L_t>100n_j
 \quad\text{for }0<|t|\le t_j
\]
を課す．固定段階の定数や $n_j$ に大きな増大があっても，
これらは全て $t_j$ を小さくすることで満たせる．
各評価は $0<|t|\le t_j$ 全体で成り立つので，
$t_{j+1}<|t|\le t_j$ で第 $j$ 段階を使える．
この選択から $\log m_t/R_t\le1/j$ が従う．

\emph{全 flow の共通開領域を確保する．}
内側の残余 annulus $b_j<b<2b_j$ を取り，外側を
$b_j/2<b<4b_j$ とする．上の座標比較により，この外側 annulus は
全て $P_j$ の内部にある．Normal 深度 $a,c>3n_j$ の部分は，
global frame に移しても深度 $>3n_j-M_j>2n_j$ なので $K_{t,n_j}$ 内にある．
また $n_j\gg b_j$ により，外側 annulus 全体で固定比の正則 plateau arm 内にある．

$\nu$-level の全 torus は $a,c>4n_j$ にあり，第二段の $\psi_s$ の
モデル長さと $C^1$ 誤差は $Cb_j^{-2}$ 以下である．
同じ領域で normal generator の長さは $C/n_j$，模型との誤差は $C/b_j$ なので，
centered normal Moser 場のモデル長さは $C/(n_jb_j)$ 以下である．
正規化 log 深度と残余 $q$ の微分は一様に有界であるから，
大きい $j$ では，第二段の全画像と，そこからの第一段の全 trajectories は
$a,c>3n_j$ と外側残余 annulus の中に留まる．
逆向きの flow と平均順序比較の小さい trajectories も，少し広い同じ margin に収められる．
比較的 compact な飽和開近傍を取れば，
全円 orbit，residual-stage の全画像，および normal-stage の全 trajectories を含む
共通の開 plateau が得られる．

Global neck と三重点側で同じ centered kernel を用いる．
両側の bare normal circle は正則座標同定 $G$ により厳密に共役されるので，
補題\ref{ke:normal-flow-naturality}から
\[
 \phi_{N,t,s}\circ G=G\circ\phi_{\alpha,t,s}
\]
がこの共通領域で成り立つ．Primitive の近似だけから flow の一致を推論していない．
共通 canonical moment は上で確認したものである．

\emph{Reduced 葉を合わせる．}
命題\ref{ke:single-arm-reduced-graph}をより狭い $I_j$ に制限する．
三重点 torus の第二段の像は，neck の同じ $\mu_N$-level の reduction 上で
$(q_t,\theta_t)$ に関して $C^1$ 誤差 $O(b_j^{-2})$ の円 graph を与える．
固定座標 $(q_0,\theta_0)$ への変更も上で $O(b_j^{-2})$ に抑えたので，
これらの円は基準の radial 葉の小さい graph である．
補題\ref{ke:reduced-collar-graph}を互いに素な広い annuli で用い，
各 moment に滑らかに依存し，内部の source と saddle 近傍では恒等な
foliation の変形を作る．無重みの moment 微分に一様な小ささは課さない．
命題\ref{ke:moment-dependent-neck-foliation}によって，
この葉層族を全 neck へ持ち上げられる．
Source の $\mu=0$ を含む全商では $(\mu_N,p)$ の位相同定を用い，
滑らかな reduction 座標を使うのは source-free な end のみである．

系\ref{ke:corner-leaves-in-neck-reduction}により，選んだ end の torus は集合として
\[
 G\phi_\alpha\psi(T)=\phi_N G\psi(T)
\]
であり，変形した neck 葉の持ち上げと完全に一致する．
三重点の三つの端に別々の normal-first 順序が必要な場合は，
定理\ref{ke:relative-corner-patch}を使う．支持は互いに素な arm にあり，
本証明で確保した共通 source/image の余裕の中に選べる．
これで一つの corner bundle の指定された各 end strip と一致させられる．

最後に，neck の終端を一致 annulus 内の一つの regular 葉で切り詰める．
例えば基準座標で $q_0=\log(3/2)$ の円を用い，その moment-dependent な像を
境界として取る．これは foliation に飽和した切断であり，
その両側に一致する開 collar がある．内部の sources と saddles は変形の支持外にあるので，
定理\ref{ke:projected-source-neck}の特異 fiber と marking はそのまま保持される．
この切詰めにより，離れた部分で二つの bundle の偶然の交差があることを
そのまま接着と呼ぶ必要もない．
\end{proof}

この議論は有限個の二重線について同時に行える．
残余 cut を共通にし，必要な $n_j$ の最大値，moment 幅と $t_j$ の最小値を取ればよい．
各三重点での三 arm の座標は既に一つの共通 chart に合わせている．
ただし，これで得るのは指定された end strips の相対的一致である．
浅い normal 領域と component core を含む全曲面の被覆は，
補題\ref{ke:fixed-stage-stratum-coverage}と定理\ref{ke:actual-global-fibration}で確認する．
Source を含む全 fiber 合併の径は命題\ref{ke:collapsed-tree-full-fiber-collapse}，
全曲面の no-shortcut は定理\ref{ke:global-ambient-lower}で評価する．
これらは局所 flow の小さい変位から直接従うものではない．

\subsection{Closed corner と切詰めた neck の全交わり}

Germ の一致から二つの実際の領域の和集合上の fibration を得るには，
指定した collar 以外の交点も除外する必要がある．そのため normal-first の相対化を，
最終的に使う moment strip より先に，少し浅い normal 領域まで施す．
以下の証明では，交点の moment を同定した後に初めて深い normal 下界を使う．

\begin{corollary}[全交わりが共通の葉で飽和した部分 fibration]
\label{ke:neck-corner-saturated-union}
定理\ref{ke:neck-corner-diagonal-matching}の有限個の二重線と三重点について，
diagonal 選択と cut の位置を取り直せば，次を同時に実現できる．
三重点側では
\[
 m_t=6b_j,\qquad t_{j+1}<|t|\le t_j
\]
とした全 softmax domain $\mathcal H_{t,m_t}$ を使い，
三つの arm の normal-first 構成を相対化した actual corner bundle を
$\mathcal C_{t,J}$ とする．Neck の各 end は，対応する $\kappa$-torus の残余深度
$b=7b_j$ によって指定される moment-dependent reduced 葉で切り詰める．
その actual 像を $\mathcal N_{t,D}$ とする．

各 $\mathcal C_{t,J}\cap\mathcal N_{t,D}$ は双方の全 torus fibers の和であり，
そこで二つの fiber partition と向き付き normal circle は厳密に一致する．
End の第二 cycle はこの共通 torus 同定で移送して標識する．独立に先指定した
全整数基底同士の一致や第二 period の一致を主張するものではない．
各 end には空でない開 overlap がある．小さい開近傍を用いて同じ構成を行えるので，
これらの和集合上に一つの actual $\omega_t$-Lagrangian 部分 fibration が得られる．
Neck 内部の sources と simple saddles の局所モデルは変わらない．

この選択も $m_t\to\infty$，$\log m_t=o(R_t)$ を満たす．
従って各 closed corner について，定理\ref{ke:slow-cut-closed-corner-gh}の
内在・protected GH 評価は保持される．和集合が曲面全体を覆うこと，
component core への延長，およびこの和集合の外を通る shortcut の評価は，
本系の結論には含まれない．
\end{corollary}

\begin{proof}
\emph{定数と cut の位置を選ぶ．}
前定理の残余領域を
\[
 P_j=D\setminus\bigcup_k\{s_{k,0}>12b_j\},\qquad
 P_j^+=D\setminus\bigcup_k\{s_{k,0}>16b_j\}
\]
に替える．Frame 差 $M_j$ は少なくとも $3b_j<s_{k,0}<12b_j$ を含む固定 annuli で取り，
$n_j>4M_j$，$n_j>1000(1+2\sigma)b_j$ とする．
固定段階の全 buffer 条件，$C_j/n_j<1/j$，および有限個の outer neck tubes の
分離も，必要なら $n_j$ をさらに大きくして満たす．
実際，異なる二重線上の固定 compact $P_j^+$ は三重点を除いているので互いに素であり，
まずその十分細い束近傍を互いに素に選び，最後に $t$ を小さくすればよい．
各三重点 $q_J$ の固定 chart の閉包は，他の三重点と，
$D_{ab}$，$\{a,b\}\not\subset J$ なる非 incident な二重線からも離して選ぶ．
一般位置条件により $q_J$ を通る二重線は $J$ 内の三本だけである．
固定段階の outer neck tubes を細くし，最後に $t_j$ を小さくすれば，
非 incident な $\mathcal C_{t,J}$ と $\mathcal N_{t,D_{ab}}$ は交わらない．
従って以下では対応する end の交点だけを扱えばよい．
Moment 幅は前定理と同じく
\[
 0<\iota_j<\min\{c_j/(16n_j),\pi/(64n_j)\}
\]
とし，$|\mu|<2\iota_j$ 上で構成する．
式\eqref{ke:eq:exhaustion-normal-cut-moment}は残余 annulus を
$4b_j<b<10b_j$ に替えても成立する．
$t_j$ を最後に小さくして，この annulus 上の全 $\mu,\nu$-levels が
$a,c>4n_j$ にあること，および $P_j^+$ の各 end 上で
$Z_{k,t}/Z_{k,0}$ の必要な $C^3$ 比較を確保する．
ここで $a,c$ は $\Xi_k$ の normal 深度，$b=-\log|Z_{k,t}|^2$ である．

\emph{Moment の制限より広い領域で先に平均順序を合わせる．}
新しい triple 座標の torus 飽和領域として
\[
 E_k^-=\{5b_j\le b\le9b_j,\quad a,c\ge n_j/2\}
\]
を取り，その少し広い $E_k^0$ と，さらに広い
\[
 E_k^+\Subset\{4b_j<b<10b_j,\quad a,c>n_j/4\}
\]
を取る．最初の不等式には必要に応じて小さい開近傍を付ける．
例えば残余の中間 cut を $4.5b_j,9.5b_j$，normal の中間 cut を $2n_j/5$
に取れば，入れ子領域の buffer は正の定数で下から抑えられる．
$\log(b/b_j),\log(a/n_j),\log(c/n_j)$ の滑らかな関数で作る
cutoffs の固定階の $h$-norms も一様である．
三つの $E_k^+$ は互いに素で，全て最小深度 $m_t/2=3b_j$ より深い．

座標接続の逆写像を $G$ とし，$\alpha=G^*\omega_t$ に対して，
一段階 map $\Phi$ も各 normal-first map $\phi_{\alpha,k}\psi_k$ も
centered kernel を用いて選ぶ．
定理\ref{ke:relative-corner-patch}を上の領域へ適用し，
\[
 C_t=G\Upsilon\Phi,\qquad
 C_t(z)=G\phi_{\alpha,k}\psi_k(z)\quad(z\in E_k^-)
\]
を得る．$C_t(\mathcal H_{t,m_t})=\mathcal C_{t,J}$ と置く．
$\Upsilon\Phi$ の変位と一次微分の変化は $O(b_j^{-2})$，
$G$ は指数的小摂動である．
従って，source と像のどちらでも各正規化 log 深度は $O(b_j^{-2})$ しか変わらない．
この相対化は moment を $I_j$ に制限する前に行っている．

\emph{Neck を共通の reduced 葉で切る．}
$4b_j<b<10b_j$ と $|\mu|<2\iota_j$ の共通 strip に
命題\ref{ke:single-arm-reduced-graph}を適用する．
$B$ を source の $\kappa$-torus の残余深度とし，
その第二段の像の reduced 葉を $\gamma_{\mu,B}$ と書く．
$B$ は torus 全体で一定であり，像の点ごとの残余深度 $b$ と区別する．
Graph 評価から
\[
 b=B\bigl(1+O(b_j^{-2})\bigr)
 \quad\text{on }\gamma_{\mu,B},
\]
またこれらの円は $B$ の順序を保つ．
補題\ref{ke:reduced-collar-graph}による foliation の移送を，
$5b_j<B<9b_j$ 上ではこの全 family と一致し，
source・saddle の近傍では恒等となるように取る．
Neck source を $\gamma_{\mu,7b_j}$ で切り，残余内部を残す．
その内側では，end 座標が定義される範囲で
$b\le7b_j+O(b_j^{-1})$ である．
実際の neck はこの領域を global normal Moser map $\phi_N$ で移したものとし，
その moment 座標は
\[
 \widehat\mu_N=\mu_N\circ\phi_N^{-1}
\]
である．平均前の点における $\mu_N$ と混同しない．

\emph{Neck の全 normal path の残余方向を閉じ込める．}
Global normal flow の全 trajectories は outer tube 内にあり，
global frame での二つの normal 深度は $n_j$ より大きい．
$4b_j<b<10b_j$ にいる間は frame 差から local normal 深度は
$n_j-M_j>3n_j/4$ である．
Projected 深度から旧座標の arm 条件へ戻す粗評価も確認する．
補題\ref{ke:arm-coordinate-ratio}の scaled Morse 写像の逆写像は，
$|x'|+|y'|<c|p|$，$|t|<c|p|^3$ 上で
\[
 |x|+|y|\le C\bigl(|x'|+|y'|+|t|/|p|^2\bigr)
\]
を満たす．これは臨界断面 $a_t=O(t/p^2)$，行列平方根
$S-I=O(t/p^3)$ と固定幅の逆写像定理から従う．
現在の固定 annulus を含む $A_j\Subset D\setminus Q$ 上で，
最後の $t_j$ の選択に
\[
 |t|\sup_{A_j}|p|^{-2}\le e^{-n_j/2},\qquad
 |t|\sup_{A_j}|p|^{-3}<c
\]
を追加する．固定 unit の上下界と $n_j\gg b_j$ により逆写像の領域条件が満たされ，
$|x'|,|y'|\le C e^{-3n_j/8}$ から旧 normal 深度は
$3n_j/4-O(1)$ 以上となる．$q_k=p+\alpha x+\beta y$ と
$Z_{k,t}/Z_{k,0}\to1$ から旧 residual 深度は $b+O(1)$ である．
従って旧最小深度は $3b_j$ より大きく，旧 normal/residual 深度比は
$1+2\sigma$ より大きい．この評価は projected normal circle 全体で保たれるので，
座標補間の cutoff は全 orbit 上で $\chi_k=1$ となる．
よってその全 normal circle は正則な $\Xi_k$ plateau に含まれる．
補題\ref{ke:normal-flow-naturality}をここでは vector field の段階で用いると，
$q=\log(b/b_j)$ に対して
\begin{equation}\label{ke:eq:neck-residual-first-exit}
 |Y_N|_h\le\frac{C}{n_jb_j},\qquad
 |dq(Y_N)|\le\frac{C}{n_jb_j}
\end{equation}
を得る．この評価には global と local の endpoint maps の一致は使わない．
Centered homotopy が不変模型を消すこと，normal generator の長さが $C/n_j$ 以下で
actual/model 誤差が $C/b_j$ 以下であることから直接従う．

Source の residual cut と $8b_j$ の間には固定の $q$-幅がある．
従って最初に $b=8b_j$ に到達する時刻があると仮定すると，
その直前の annulus を通る $q$-変位が
\eqref{ke:eq:neck-residual-first-exit}の時間 $1$ の積分より大きくなり矛盾する．
よって neck の全 image と全途中像は $b<8b_j$ にある．
Source が end の外にある場合も，同じ annulus を横切る必要があるので同じ議論になる．

\emph{実際の交点を先に広い plateau に入れる．}
$x\in\mathcal C_{t,J}\cap\mathcal N_{t,D}$ とする．
Corner の source は全深度 $\ell_a>6b_j$ を持つので，
小さい変位により $x$ の旧 triple 深度は
$\ell_a(x)\ge6b_j-O(b_j^{-1})$ である．
元の射影座標 $p$ は固定 affine triple 座標 $q_i,q_j,q_k$ の線形式であり，
それらと $\zeta_a$ の比は固定 unit である．従って
$|p(x)|\le C e^{-3b_j}$，すなわち
\[
 b(x)\ge6b_j-C>\tfrac{11}{2}b_j
\]
となる．ここでは $F=\Xi_k$ をまだ用いず，
固定 end 上の $Z_{k,t}/Z_{k,0}$ の比較だけを用いている．
既に得た neck の上界と合わせて
$11b_j/2<b(x)<8b_j$ を得る．

$x$ から global neck の逆 normal trajectory を時刻 $0$ まで追う．
全軌道は outer tube 内にあり，帯内では local normal 深度が $3n_j/4$ より大きい．
\eqref{ke:eq:neck-residual-first-exit}と固定幅の margin によって，
逆軌道も全時刻で $5b_j<b<9b_j$ に留まる．
従って全軌道と全 normal orbits 上で $F=\Xi_k$ が成立する．
特に $F(x)$ の残余深度は実際の $b(x)$ と同じである．

$x=C_t(z)$ と書く．$\Upsilon\Phi$ の小変位から，
$z$ の残余深度は $(5b_j,9b_j)$ 内，normal 深度は $n_j/2$ より大きい．
従って $z\in E_k^-$ であり，
$x=G\phi_{\alpha,k}\psi_k(z)$ という staged 表示を使える．
Global neck の逆軌道を $F$ で移すと local normal の逆方程式を満たすので，
ODE の一意性から
\begin{equation}\label{ke:eq:intersection-inverse-normal}
 \phi_N^{-1}(x)=G\psi_k(z)
\end{equation}
を得る．この時点までは，$z$ の normal 深度が $4n_j$ より大きいことを
仮定していない．

\emph{Moment と全 fiber を同定する．}
式\eqref{ke:eq:intersection-inverse-normal}，
系\ref{ke:transition-normal-primitive}，および第二段の moment の保存から
\[
 \widehat\mu_N(x)=\mu_N(G\psi_k(z))
                =\mu_\alpha(\psi_k(z))=\nu_k(z)
\]
が厳密に成立する．
$x$ が neck の目標 moment interval にあるので $\nu_k(z)\in I_j$ であり，
ここで初めて normal 単調性を使って $a(z),c(z)>4n_j$ と結論する．

$z$ を含む全 $\kappa$-torus を $T_z$ とし，その残余深度を $B$ とする．
Reduced graph の同定から，$p(G\psi_k(T_z))$ は
$\gamma_{\widehat\mu_N(x),B}$ そのものである．
Neck は $B=7b_j$ の葉で切られており，graph family は $B$ の順序を保つから，
交点が存在すれば $B\le7b_j$ である．
その全 torus 上で第二段と第一段の trajectories は
$a,c>3n_j$ と $4b_j<b<10b_j$ にあり，global inner tube 内に含まれる．
従って全 torus に自然性を適用でき，
\[
 C_t(T_z)=G\phi_{\alpha,k}\psi_k(T_z)
          =\phi_NG\psi_k(T_z)
\]
は neck の一つの全 fiber と一致する．
$\mathcal H_{t,m_t}$ は全 phase tori の和なので，
この torus の一部だけが corner に含まれることもない．
これで交わりが双方で飽和していることと fiber partition の一致が従う．

$|\mu|<\iota_j$ と，例えば $6.4b_j<B<6.6b_j$ の共通 levels は空でない．
Normal 深度は $4n_j$ より大きく，
$\min_a(\ell_a-m_t)\ge0.4b_j>1$ なので softmax の全比条件を厳密に満たす．
従ってこれは両側の内部に含まれる開 overlap である．
Closed corner の source を，全最小深度が $23b_j/4$ より大きい
小さい torus 飽和開近傍へ広げても，上の交点の下界と first-exit の余裕は保たれる．
Neck 側でも少し広い moment interval を使える．
こうして開領域の fiber partitions が接着する．
その overlap の基底座標は $(\nu_k,B)=(\widehat\mu_N,B)$ であり，
normal 単調性と reduced graph の順序保存により座標同定は滑らかである．
まず閉領域の基底を用いると，各基底は compact Hausdorff であり，
飽和した交わりの像は閉である．同じ全 fiber による同定は，proper quotient を介して
これらの閉集合の間の同相写像を与える．同定の graph は閉であり，
異なる corner 同士と異なる neck 同士は交わらないので，同値類の元は高々二つである．
従って有限 compact 空間上の同値関係は閉で，貼合せ基底は Hausdorff となる．
上の小さい開 collar を同じ $(\nu_k,B)$ 座標で延長すれば，
開領域の部分 fibration もこの基底の近傍上で得られる．

向き付き normal cycle は作用の同変性により一致する．第二 cycle は共通 torus
同定で移送する．Base-only frame unit の winding により，別に選んだ residual cycle
とは normal cycle の整数倍だけ異なる場合があるので，その二つの基底を同一視しない．
Source・saddle の内部標識は end の補正の支持外にあるため保たれる．

異なる三重点 chart は互いに素に選ぶ．異なる neck images は，最初に互いに素にした
outer tubes の中に留まる．従って有限個を同時に貼り合わせられる．
内部の singular fibers と marking は end の変形の支持外にある．

最後に $t_j$ の選択に
$R_t\ge j(1+\log(6b_j)+\log(n_j+1))$ を追加すれば，
$m_t=6b_j$ に対しても $\log m_t/R_t\le1/j$ である．
相対化の変位と $C^1$ 誤差は $O(b_j^{-2})$ なので，
節\ref{ke:relative-corner-section}の pullback metric の評価により
同じ closed-corner GH distortion bound が成立する．
この距離評価を neck，core，または保護領域の外まで延長してはいない．
\end{proof}

\input{sections/paper/ke_source_connectors.tex}
\input{sections/paper/ke_source_all_moment_connectors.tex}
\input{sections/paper/ke_component_core_geometry.tex}
\input{sections/paper/ke_component_core_transfer.tex}
\input{sections/paper/ke_component_link_cover.tex}
% !TeX root = ../../surface-lagrangian-torus-fibration.tex
% 2026-11-23 / v1.2: Adapted primitiveとpositive residue gaugeの大域exact差を明示．
\section{固定成分 end における actual invariant primitive の評価}
\label{ke:core-invariant-primitive-section}

本節では actual circle generator の一階微分を評価し，
primitive の相対不変化が core 外に新しい trapping を作らないことを証明する．
対象は，固定段階の後に選ぶ \emph{成分側の固定 compact collar} である．
Source の critical moment を通る moving central neck 全体について
同じ一階評価を主張するものではない．
閉じた component link の被覆と全 fiber による cut は
節\ref{ke:component-link-cover-section}，Liouville 葉族の全域接着は
定理\ref{ke:actual-global-fibration}で扱い，本節ではその primitive の条件を検証する．

\subsection{重み付き cusp bounds と順序を固定した選択}

正の residues を持つ potential は，成分の affine volume frame を用いた
式\eqref{ke:eq:positive-residue-ke-potential}の明示形
$\Psi_w=F+\sum_{a<m}(1-w_a)\log|\ell_a|^2$ を選ぶ．
End では $\lambda_0=d^c\Psi_w$，$d\lambda_0=\omega_\infty$ とする．
Core の近傍で必要な compactly supported exact 補正と対応する Lyapunov
function の変更は，系\ref{ke:actual-limit-weinstein-core}で既に行ったものとする．
すなわち全体では $\lambda_0=d^c\Psi_w+dQ_{\rm core}$，
$Q_{\rm core}\in C_c^\infty(W_i)$ であり，Lyapunov function を $\Psi$ と書く．
その外側では $\Psi=\Psi_w$ とする．
$s_a=-\log\|s_{D_a}\|^2>1$，$H=1+\sum_a s_a$ とし，
標準 cusp background を $h$ と書く．十分高い end では
\begin{equation}\label{ke:eq:weighted-core-input}
 \begin{split}
 cH\le |d\Psi|_h\le CH,&\qquad
 |\lambda_0|_h+|\nabla^h\lambda_0|_h\le CH,\\
 |d\log H|_h\le C,&\qquad Z_0\Psi\ge cH^2,
 \qquad \iota_{Z_0}\omega_\infty=\lambda_0.
 \end{split}
\end{equation}
また $\omega_\infty$，その逆，およびそれらの全固定階の cusp derivatives
は一様に有界である．これらは
Rochon--Zhang\cite[\S5, (5.3)--(5.6) と直後の段落]{RochonZhang}の
bounded $C^\infty_{\mathrm{fc}}$ potential と cusp metric の二側比較から従う．
同論文の Theorem~5.1 の decay を用いてもよいが，以下には bounded derivatives
で十分であり，smooth divisor に対する \S7 の polyhomogeneous expansion は不要である．
各 SNC chart で $|\nabla^r s_a|_h\le C_r s_a$ を用いると上界が従う．
最大の $s_a$ に対する normalized radial direction では
$d\Psi$ の主項が $w_a s_a$，残りが $O(1)$ であるから下界を得る．
固定された正の weights の最小値を使えば，深度比には依存しない．

選択の順序は次のとおりである．まず K2 の残余段階 $b=b_j$ と
$P_j\Subset P_j^+$，次に $n=n_j$ と共通 moment interval
$[-\iota_j,\iota_j]$ を固定する．さらに
$N_0\gg n_j+\iota_j^{-1}$ を選び，その後に必要な上限深度も有限に固定する．
最後に $|t|$ の閾値を小さくする．
深さ $N_0$ の選択後に取るすべての領域，円軌道，Moser trajectories，
および cutoff supports は一つの $V_j\Subset W_i$ に入れる．
この順序では $t\to0$ の fixed-gauge smooth convergence を任意の固定階で使える．

\subsection{短い軌道からではなく，正規化座標から一階評価を得る}

\begin{lemma}[固定残余 compact 上の actual normal generator]
\label{ke:component-normal-generator-c1}
$P\Subset P^+$ を固定する．成分極限の normal end で
centered normal averaging により得る symplectic circle action の
周期 $1$ の生成場を $V_N$ とする．
全円軌道と全 averaging trajectories が固定した外側 buffer に含まれる領域では，
任意の固定 $r$ に対して
\[
 \|V_N\|_{C^r(h)}\le C_{P,r}/N_0\qquad(s\ge N_0)
\]
となる．特に $r=1$ が成り立つ．
深度の有限上限を固定した後には，既存の actual neck の同じ作用について
\[
 \|V_{N,t}\|_{C^1(h)}\le 2C_{P,1}/N_0
\]
が十分小さい $|t|$ に対して成り立つ．
ここでは $h$ を graph gauge で用いる．
\end{lemma}

\begin{proof}
局所的な normal coordinate を $v$，$s=-\log|v|^2$ とする．
$s_0=s(v_0)$ の点で，角度の universal cover 上に
\[
 \xi=\frac{\log v-\log v_0}{s_0}
\]
を取り，$P^+$ の固定座標と組み合わせる．固定半径の小さい $\xi$-chart 上で
$s/s_0$ は二側比較可能であり，$h,h^{-1}$ と全固定階の係数微分は有界である．
Bare rotation $v\mapsto e^{2\pi\sqrt{-1}\tau}v$ の生成場は
\[
 V_{\mathrm{bare}}=\frac{2\pi}{s_0}\partial_{\operatorname{Im}\xi}.
\]
従ってその全固定階の covariant derivatives は $O_P(s_0^{-1})$ である．
この結論は軌道長から微分を推測したものではなく，座標表示と connection bounds
から得ている．一周の移動も $O(s_0^{-1})$ なので，全 orbit は外側 chart に留まる．

Bare circle の平均を $\mathcal A_N$，centered homotopy を $K_N$ とすると，
$\beta_N=K_N\omega_\infty$ の $C^q(h)$ norm は
$C_{P,q}/s_0$ 以下である．
ここでは $\omega_\infty$ の全固定階 bounds と，回転の明示的な translation
表示を使う．$\Omega_\tau=(1-\tau)\mathcal A_N\omega_\infty+
\tau\omega_\infty$ は一様に非退化である．
Bare action が正則な normal bundle chart では K\"ahler forms の凸結合であり，
同値に $\mathcal A_N\omega_\infty-\omega_\infty=d\beta_N=O_{C^q}(s_0^{-1})$
という小摂動だけでも非退化性を示せる．
\[
 \iota_{Y_\tau}\Omega_\tau=\beta_N
\]
で定める Moser field は全固定階で $O_P(s_0^{-1})$ である．
Flow 方程式を二回微分すると，averaging map $A_N$ とその逆は
$C^2$ で一様に有界であり，実際に $\id$ との差は $O_P(s_0^{-1})$ である．
高階にも同じ議論を適用できる．

Actual generator は $(A_N)_*V_{\mathrm{bare}}$ である．座標微分は
\[
 D\bigl((A_N)_*V_{\mathrm{bare}}\bigr)
 =D^2A_N\,V_{\mathrm{bare}}\,DA_N^{-1}
   +DA_N\,DV_{\mathrm{bare}}\,DA_N^{-1}
\]
の形なので，$A_N$ の $C^2$ bound と bare generator の $C^1$ smallness から
所要の $C^1$ smallness を得る．$C^1$ の conjugating map だけでは，
この式の第一項を評価できない．

有限個の固定残余 charts と外側 buffers を使えば $P$ 全体で同じ結論となる．
Source の残余点を $P$ に含めても，今扱う成分側では $v\ne0$ であり，
$x'y'=t h(p)$ の反対側の座標を $x'=t h(p)/y'$ と解ける．
従って projected source coordinates は固定 compact 上で $C^\infty$ に収束する．
Metric，averaging forms，Moser fields と flows も同じ kernel の下で収束するので，
最後の actual $t$ の評価を得る．Central fixed point を含む moving rescaled
source 領域にこの fixed-compact 議論を適用してはいない．
\end{proof}

\begin{lemma}[Corner の相対補正に必要な一階高い評価]
\label{ke:relative-corner-patch-c2}
命題\ref{ke:quantitative-order-hamiltonian}の入力を $C^5$ に，
定理\ref{ke:relative-corner-patch}の cutoff bound を $C^4$ に上げれば，
比較 Hamiltonian と相対化 flow に
\[
 \|H_r\|_{C^3(h)}+\|W_r\|_{C^2(h)}\le Ce\delta,
 \qquad \|\Upsilon-\id\|_{C^2(h)}\le C_Be\delta
\]
が成り立つ．従って K2 の actual corner torus action の各固定整数生成場は
\[
 \|V_{C,a}\|_{C^1(h)}\le C/b
\]
を満たす．定数は $n/b$ 又は component 側の二つの深度比には依存しない．
\end{lemma}

\begin{proof}
既存証明の $B_{r,s},C_{r,s},Y_{r,s}$ を同じ式で用いる．
入力を一階増やせば $K_j\omega,K_j\lambda$ は $C^4$ で $O(e\delta)$，
従って $Y_{r,s},\partial_rY_{r,s}$ も $C^4$ で同じ bound を持つ．
Flow は $C^4$ で制御し，その $r$ 変分を三回まで空間微分して Gronwall を使う．
変分方程式の三階空間微分には $D^4Y$ が現れるので，
$Y$ の $C^3$ bound だけでこの parameter derivative を処理してはならない．
明示的な primitive
\[
 Q_r=\int_0^1F_{r,s}^*
       \bigl(\Lambda_{u_r,v_r}(Y_{r,s})-C_{r,s}\bigr)\,ds,
 \qquad
 H_r=(\partial_rQ_r)\circ F_r^{-1}-\lambda(W_r)
\]
を三回微分すると，両項とも $O_{C^3}(e\delta)$ である．
従って $d(\chi H_r)$ は $C^2$ で $O_B(e\delta)$，
その Hamiltonian field と flow は $C^2$ で同じ評価を満たす．
基点からの積分は使わず，領域の直径に依存する定数は生じない．

Actual triple の canonical primitive は full-triple tensor estimate の
determinant と一回の $d^c$ 微分から全固定階で模型に $O(b^{-1})$-近い．
必要な $C^5$ 入力は，既存の \emph{全固定階}の結果から供給される．
$G$ による座標補間も補題\ref{ke:arm-log-derivatives}の全固定階の評価を用い，
ここでは $C^6$ まで取れば十分である．
Cutoffs は $\log(b/b_j),\log(a/n_j),\log(c/n_j)$ の固定滑らかな関数なので，
固定整数基底に対する normalized charts で $C^4$ bound が一様である．

Bare torus generator は各 logarithmic chart で深度の逆数を係数とする定数場で，
$C^r$ norm は $C_r/b$ 以下である．
K2 の actual parametrization $G\Upsilon\Phi$ とその逆は，上の補強により
$C^2$ で一様に有界である．前補題の pushforward の式を用いれば結論を得る．
これは既存の $C^1$ map estimate だけを引用して得た結論ではない．
\end{proof}

以上より，固定成分側の必要な有限個の patches で
\begin{equation}\label{ke:eq:actual-generator-delta}
 |V|_h+|\nabla^hV|_h\le\delta_j,
 \qquad \delta_j=C_*/b_j+C_j/N_0
\end{equation}
を得る．$C_*$ は corner の一様定数，$C_j$ は固定残余 compact の定数である．
$b_j$ を選んだ後に $N_0$ を大きくするため，$C_j$ の増大は障害にならない．
同じ normal-first plateau 上の二つの normal actions は
補題\ref{ke:normal-flow-naturality}により \emph{同じ parameterized action} である．
同じ無向 orbit を持つというだけの意味ではない．

\subsection{Actual symplectic action に対する exact 平均}

\begin{lemma}[Centered scalar primitive と重み付き一階評価]
\label{ke:actual-centered-primitive-weighted}
$d\lambda=\omega$ とし，$R_s$ を周期 $1$ の actual symplectic circle action，
$V$ をその generator とする．全 orbit を含む領域上で
\[
 \mathcal A\eta=\int_0^1R_s^*\eta\,ds,
 \qquad K\eta=\int_0^1(1/2-s)R_s^*(\iota_V\eta)\,ds
\]
と置く．このとき
\begin{equation}\label{ke:eq:actual-centered-exact}
 \mathcal A\lambda-\lambda=dF,
 \qquad F=K\lambda
   =\int_0^1(1/2-s)\lambda(V)\circ R_s\,ds.
\end{equation}
$|\lambda|+|\nabla\lambda|\le CH$，$|d\log H|\le C$，
$|V|+|\nabla V|\le\delta\le1$ なら
\begin{equation}\label{ke:eq:centered-F-weighted-bound}
 |F|+|dF|\le C'\delta H.
\end{equation}
可換な二つの円では
$F=K_1\lambda+\mathcal A_1(K_2\lambda)$ とすれば，
$\mathcal A_1\mathcal A_2\lambda-\lambda=dF$ と同じ評価を得る．
\end{lemma}

\begin{proof}
Cartan の公式を部分積分すると
$dK+Kd=\mathcal A-\id$ となる．
作用が symplectic なので $R_s^*(\iota_V\omega)=\iota_V\omega$ であり，
centered kernel の積分が零だから $K\omega=0$ である．
これが \eqref{ke:eq:actual-centered-exact} を与える．
特に flux の消滅を仮定して局所 primitive を選んでいるのではない．

$|d(\lambda(V))|\le|\nabla\lambda||V|+|\lambda||\nabla V|
\le C\delta H$ である．
Flow の変分方程式から $\|DR_s\|\le\exp(C\delta)$，
$|d\log H(V)|\le C\delta$ から
$H\circ R_s\le e^{C\delta}H$ を得る．
明示積分を一回微分すれば所要の評価となる．
Rank two でも $K_1\lambda$ と $K_2\lambda$ の scalar $C^1$ bounds，
および scalar の $\mathcal A_1$ 平均の $C^1$ bound だけで十分である．
逐次平均した one-form の一階微分を新たに評価して $C^2F$ を要求する必要はない．
この選択は積分定数も固定しており，領域の直径や深度比に由来する定数を持ち込まない．
\end{proof}

\subsection{共通 action 上の貼合せと nontrapping}

次の補題は解析的な貼合せ結果であり，入力の end atlas が閉じた link の全点を
覆うことを自動的に主張しない．K2 の normal-first plateau を使う際の
着手条件を明示しておく．

\begin{proposition}[深い star cover 上の adapted primitive]
\label{ke:star-cover-adapted-primitive}
固定成分の compact end 領域を，normal patches $\mathcal N_a$ と
互いに素な corner patches $\mathcal C_{ab}$ で扱う．
Normal patch は一つの global normal circle action を持ち，
corner patch は actual $T^2$ action を持つとする．
必要な全 overlaps と全 orbits を含む外側 buffers を取り，次を仮定する．
\begin{enumerate}
\item 各 normal--corner overlap は全 torus で飽和し，normal action は
corner の指定 primitive subgroup と厳密に同じである．
同じ corner に至る二つの interpolation bands は互いに素であり，
二つの normal patches が交わる箇所は corner gauge の plateau に含まれる．
\item Normal だけの領域から full corner gauge へ移る cutoffs は common normal
action に不変で，$|d\chi|_h\le C$ を満たす．
保持したい全 corner collars では補間が既に終わり，corner gauge がそのまま成立する．
\item この union 上で組み立てる scalar function を，その外へ零で延長するための
cutoff $\eta$ があり，$|d\eta|_h\le C$，
$\eta=1$ on the retained collars，かつ支持は core の外の固定 compact に含まれる．
$\eta$ の支持近傍全体が上の union に含まれる．
\end{enumerate}
これらの cutoffs の定数と patches の数は段階によらず有界とする．
\eqref{ke:eq:weighted-core-input} と \eqref{ke:eq:actual-generator-delta} が成り立ち，
$\delta_j$ が十分小さければ，$F\in C_c^\infty(W_i)$ を取って
\[
 \lambda^{\mathrm{ad}}=\lambda_0+dF
\]
を retained normal collars 上で normal-invariant，
retained corner collars 上で $T^2$-invariant にできる．
全 cutoff 支持上で $|F|+|dF|\le C\delta_jH$ であり，
\begin{equation}\label{ke:eq:adapted-nontrapping}
 Z^{\mathrm{ad}}\Psi\ge(c-C\delta_j)H^2>0,
 \qquad \iota_{Z^{\mathrm{ad}}}\omega_\infty=\lambda^{\mathrm{ad}}.
\end{equation}
特に元の core の近傍を変えず，その外に新しい trapped compact dynamics を作らない．
\end{proposition}

\begin{proof}
同じ一つの $\lambda_0$ を，各 normal action と各 corner action に沿って平均する．
前補題の明示積分で得る scalar primitives をそれぞれ $F_N,F_C$ とする．
共通 normal generator $V$ に対し，
両 one-forms $\lambda_0+dF_N$ と $\lambda_0+dF_C$ は不変である．
従って $d(V(F_C-F_N))=0$．各連結成分で $V(F_C-F_N)$ は定数で，
周期 orbit に沿う積分が零なのでその定数は零である．
すなわち $F_C-F_N$ 自体が common normal action に不変である．

共通 band 上で
\[
 F_{NC}=F_N+\chi(F_C-F_N)
\]
と置く．$\lambda_0+dF_{NC}$ は normal-invariant であり，
$\chi=1$ の plateau 上では full corner average と厳密に一致する．
各 corner への二つの bands は互いに素で，他の重なりでは既に $F_C$ だから，
この処方は全 union 上に一つの smooth scalar $F_*$ を定める．
局所 primitive を別々に基点正規化していないため，
$|F_*|+|dF_*|\le C\delta_jH$ がそのまま保たれる．
最後に $F=\eta F_*$ とし，零で延長する．
Normal/Torus invariance を使う retained collars 上では $\eta=1$ である．
他の cutoff bands には不変性を要求せず，重み付き smallness を使う．

$\omega_\infty^{-1}$ の一様 bound により
$|Z^{\mathrm{ad}}-Z_0|_h\le C|dF|_h$ である．
\eqref{ke:eq:weighted-core-input} と組み合わせれば
\eqref{ke:eq:adapted-nontrapping} を得る．
支持の外では元の field と同じであり，支持を core より高い $\Psi$ の領域に取ったので，
新しい forward-trapped orbit が高い end から低い critical region へ入ることもない．
ここで end における正時間 flow の完備性は必要としていない．
\end{proof}

\paragraph{K2 で使える具体的な cut の位置．}
定理\ref{ke:neck-corner-saturated-union}では retained corner の深度 cut は $6b_j$，
全 actual neck--corner intersections は $11b_j/2<b<8b_j$ に入り，
normal-first plateau は $5b_j\le b\le9b_j$ に広げてある．
Corner の構成写像自体は retained domain より外側の
$\min\ell_a>3b_j$ 上に存在する．
その source torus の残余深度を $B$ とし，
$B/b_j\in[5.2,5.4]$ に補間支持を取ればよい．
$b=B(1+O(b_j^{-2}))$ なので，十分大きい $b_j$ では，
全 actual intersections に入る前に full corner gauge へ切り替わる．
この band は normal-first plateau の内部にあり，
normal 深度を $N_0\gg n_j$ に取れば両側の全 trajectories にも余裕がある．
$B$ は \emph{corner source torus の label} であり，
actual image point の見かけの残余深度で代用しない．
また $B$ の cutoff は corner action に不変なので common normal action にも不変である．

一つの corner の二つの bands はそれぞれ
$(s_a,s_b)\sim(N_0,5.3b_j)$ と $(5.3b_j,N_0)$ に位置するため，
$N_0\gg b_j$ なら互いに素である．
深い double overlap では両方とも corner gauge へ切替済みとなる．
Cutoffs を $\log(B/b_j)$ の固定関数として選べば，
その cusp 一階微分は一様に有界である．
このように global normal patch ごとに一つの平均を使い，
任意の chartwise partition of unity は導入しない．
そのため残余 compact の大きな定数 $C_j$ を corner 誤差 $b_j^{-1}$ に掛ける必要がない．

最後の零延長 cutoff の一案は，高い固定 shell を先に atlas で覆い，
$\eta=\eta_0(\Psi/N)$ を使うことである．
保持する link collar の $\Psi/N$ の範囲を $\eta_0=1$ の内部に取り，
その外側の固定範囲で零にする．$\Psi\asymp H$，$|d\Psi|\le CH$ なので
その cutoff の微分は一様に有界である．
ただし \emph{支持する shell 全体が actual atlas の source と像の共通領域に
入ること}は別の被覆評価であり，次の系で確認する．
閉じた link の存在から，この buffer を省略してよいわけではない．
Retained closed corner の境界 $6b_j$ だけを構成領域の境界とみなすと，
上の補間支持は取れないので，既存の外側 construction domain を明示して使う．

\begin{corollary}[固定 component shell 上の cutoff supports]
\label{ke:component-shell-primitive-support}
K2 の広い construction domains を用い，固定段階の後で $N$ を十分大きくする．
補題\ref{ke:closed-component-link}の link とその小さい collar には，
命題\ref{ke:star-cover-adapted-primitive}の全仮定を満たす star cover と
零延長 cutoff を取れる．従ってこの適用では，cutoff support の全点被覆を
追加の未証明仮定として残す必要はない．
\end{corollary}

\begin{proof}
Link の normal depths は $[N/2,2N]$，corner rounding の二つの depths はともに
$N$ と比較可能であり，残余 directions の depths は $3N$ 以下でよい．
正の weights と $\Psi\asymp H$ から，$N$ に依存しない $0<c_S<C_S$ があって
link と小さい collar 上で
$c_SN\le\Psi\le C_SN$ となる．
$\eta_0=1$ on $[c_S/2,2C_S]$，
$\operatorname{supp}\eta_0\Subset(c_S/4,4C_S)$ とし，
$\eta=\eta_0(\Psi/N)$ を選ぶ．
支持を含む少し広い shell $c_-N<\Psi<c_+N$ を先に取る．
ここで $c_-,c_+$ は固定した正の定数である．

この shell の全点は，$N\to\infty$ で component divisor に近づく．
有限個の SNC charts では，一つ以上の active depth が $c'N$ 以上であり，
全 depths は $C'N$ 以下である．$N$ を固定段階の frame constants より十分大きく
取れば，それらは projected normal coordinate でも同じ二側比較を満たす．
At most two divisors が交わることを用い，深い方を normal direction とする．
他方の残余 depth が $5.3b_j$ より浅い領域は，
固定 residual compact $P_j^+$ の内部に十分な余裕を持って入り，
normal depth は $c'N$ 以上である．
前補題の component normal Moser map は normalized $C^2$ で
$O(C_j/N)$-近いので，その source と像は，少し小さい shell と residual region
をともに含む．同じことは actual $t$ でも最後の閾値選択により成立する．
また $N\iota_j\gg1$ により，この shell の canonical normal moments は
既存の共通 interval の内部にある．

残余 depth が $5.3b_j$ より深い領域では，二つの component depths がともに
$5.3b_j-O(1)$ より大きく，従属 triple depth は固定 $N$ の後で $t\to0$ により
無限大へ発散する．$G\Upsilon\Phi$ とその逆の normalized $C^2$ smallness から，
source の全 depths は $5.1b_j$ より大きいとできる．
従って点とその全 torus orbit，および必要な forward/backward trajectories は
corner の広い construction domain $\min\ell_a>3b_j$ に含まれる．
これは retained softmax domain の内部に入る，という主張ではなく，
primitive の平均と補間に使う \emph{同じ作用の外側の定義域} に入る，という主張である．

境界値 $5.3b_j$ の両側には固定比の余裕を取れる．
その共通領域では source residual label $B$ を使い，
$B/b_j\in[5.2,5.4]$ の band で前述の切替えを行う．
Normal depth は $c'N\gg n_j$ なので全 band が既存 normal-first plateau に入り，
補題\ref{ke:normal-flow-naturality}の全 orbit と trajectory 条件も満たす．
同じ corner の反対側の band は互いに素であり，
深い double overlap は pure corner plateau に入る．
この処方で上の少し広い shell 全体を覆う．

すべての上限 depths は固定した $N$ の定数倍以下であり，
有限個の full orbits と flow buffers も有限上限を増やすだけで含められる．
従って一つの fixed component compact に入る．
$|d\eta|\le C$ は既述の重み付き評価による．
Normal/corner 切替えの微分も，source log-depths の固定幅の関数と
normalized $C^2$ bounds から一様である．
ここで $N$ をさらに大きくすれば，normal generator の評価の
$N_0$ は $c'N/2$ と取れ，$C_j/N_0$ を所要の小ささにできる．
先に取った compact exact Moser による $\Psi$ の $C^1$-小変形を用いる場合も，
同じ strict margins と最後の $t$ 閾値選択により全主張が保たれる．
\end{proof}

\subsection{Graph gauge の exact Moser を先に行う}

Graph gauge の $\sigma_t=E_t^*\omega_t$ と
式\eqref{ke:eq:component-positive-gauge}の positive-residue primitive $\lambda_t^w$
を取る．Closed logarithmic correction は $W_i$ 上の同じ形式として選ぶ．
Core の compact exact 補正を両側に共通に加え，
$\lambda_t=\lambda_t^w+dQ_{\rm core}$，
$\lambda_0=d^c\Psi_w+dQ_{\rm core}$ として比較する．
一般に $d(\lambda_t-\lambda_0)=\sigma_t-\omega_\infty$ であり，
その差を直接 exact differential とみなして cutoff することはできない．
固定 compact 移送\ref{ke:fixed-component-transfer}を先に用いて
\[
 \varphi_t^*\sigma_t=\omega_\infty,
 \qquad \varphi_t^*\lambda_t=\lambda_0+dQ_t
\]
とする．必要な全 flow buffers を含めて $V_j$ を先に広く取る．
\[
 \lambda_{\mathrm{int},t}=(\varphi_t^{-1})^*\lambda_0,
 \qquad \Psi_t=\Psi\circ\varphi_t^{-1}
\]
と置けば，$d\lambda_{\mathrm{int},t}=\sigma_t$ であり，
\begin{equation}\label{ke:eq:transported-core-primitive-periods}
 \lambda_t-\lambda_{\mathrm{int},t}
   =d(Q_t\circ\varphi_t^{-1}).
\end{equation}
従って元の period gauge についても
$\lambda_t^w-\lambda_{\mathrm{int},t}
=d(Q_t\circ\varphi_t^{-1}-Q_{\rm core})$ である．
従って $\lambda_{\mathrm{int},t}$ に対し，既存の \emph{actual} circle/torus actions
で前節の平均を直接行えばよい．
その結果
\[
 \lambda^{\mathrm{ad}}_t=\lambda_{\mathrm{int},t}+dF_t
\]
を得る．ここで共通の固定 component 領域上で
$\lambda_t=\lambda_t^w+dQ_{\rm core}$ を代入すれば，
\begin{equation}\label{ke:eq:adapted-positive-period-exact}
 \lambda^{\mathrm{ad}}_t-\lambda_t^w
 =d\left(F_t-Q_t\circ\varphi_t^{-1}+Q_{\rm core}\right)
\end{equation}
と\emph{一つの大域的な一価函数の微分}として表示できる．
従ってこの共通領域内の全閉曲線 $\gamma$ について
$\int_\gamma\lambda^{\mathrm{ad}}_t=\int_\gamma\lambda_t^w$ である．
全 meridian と corner torus の周期，およびそれらの整数結合も保存される．
内側と外側に別の primitive を用意する追加の gluing は不要である．$\varphi_t\to\id$ の全固定階収束を使えば，
式\eqref{ke:eq:weighted-core-input}の transported 版と
\eqref{ke:eq:adapted-nontrapping}の actual 版は，最後の $t$ 閾値を小さくして成立する．
支持外の transported core は変更されない．

Retained invariant collar 上で，周期 $1$ の正の meridian を $V_a$，
$M_a=\lambda^{\mathrm{ad}}_t(V_a)$ とすると
\[
 dM_a=-\iota_{V_a}\sigma_t,
 \qquad [Z^{\mathrm{ad}}_t,V_a]=0,
 \qquad Z^{\mathrm{ad}}_tM_a=M_a.
\]
また $M_a$ は orbit period そのものだから，
補題\ref{ke:component-meridian-normalization}の厳密な規格化により
$M_a=-\rho_a+q_a$，$\rho_a=2\pi w_a>0$ となる．
Component 側の固定深い cut では $q_a\to0$ なので $M_a<0$ である．
精密な $q_a\asymp s_a^{-1}$ 展開を用いる場合の $s_a$ は averaging source gauge
の深度であり，別の image-depth coordinate へ無条件に移してはならない．

\paragraph{相対 primitive の結論と全域構成への適用．}
Actual generator の $C^1$ smallness は，fixed component collar に限定すれば
上の normal lemma と corner の高階化により，追加の解析仮定ではなくなる．
Centered scalar average の exactness，weighted $C^1$ bound，
適合 star cover 上の全 overlap の primitive gluing と nontrapping も証明される．
補題\ref{ke:closed-component-link}の link については，
系\ref{ke:component-shell-primitive-support}が零延長 cutoff の支持全体の
source/image buffers も与える．閉じた link と新しい cuts の全 fiber 飽和性には，
補題\ref{ke:closed-component-link}，\ref{ke:whole-fiber-core-trim}を用いる．
相対 filling と全 $X_t$ の被覆は定理\ref{ke:actual-global-fibration}で得る．
本節の定性的 fixed-stage 評価だけから大域距離歪みの消失を結論せず，
そのための解析は定理\ref{ke:global-ambient-lower}，\ref{ke:global-ambient-upper}で行う．

% Primary source checked: https://arxiv.org/pdf/1106.0873v2,
% pp.38--39, (5.3)--(5.6) and the paragraph after (5.6), Theorem 5.1.

\input{sections/paper/ke_relative_reduced_foliations.tex}
% !TeX root = ../../surface-lagrangian-torus-fibration.tex
% 2026-11-23 / v1.6: Source/Core同時buffer、平行linkのLiouville時間上限とactual O(1)移動費用．
\section{Actual KE 形式上の全域 core 接着}
\label{ke:actual-core-gluing-section}

本節では，閉じた component link，その全 collar 上の invariant primitive，
および saddle を動かす相対葉族を組み合わせる．
有限型 Weinstein core の存在だけから既存葉との一致を推論せず，
actual link の既存葉を出発点に Liouville flow で内部を充填する．
固定 collar を点ごとに保存する symplectic germ は指定しない．
従って命題\ref{ke:component-relative-flux}の相対 flux 条件を
新たな未証明仮定として加える必要はない．

\subsection{固定した Liouville domain の cone filling}

次の補題は Evans--Mauri の cone extension
\cite[Lemma~2.16 and Remark~2.18]{EM}の compact-domain 版を，
境界の contact form を変更しない形で記すものである．
正時間の flow を finite-volume end の無限遠まで延ばすことは使わない．

\begin{lemma}[実際の境界葉を保つ compact cone filling]
\label{ke:actual-link-cone-filling}
$(K,\omega=d\lambda)$ を compact connected Liouville domain とし，
$Z$ を $\iota_Z\omega=\lambda$ によって定める．
$Z$ は $S=\partial K$ に厳密に外向きであり，内側の固定 Weinstein subdomain
$K_0\Subset\operatorname{int}K$ の外では proper Lyapunov function を厳密に増加させるとする．
内側では有限型 Weinstein structure と一致するものとし，その skeleton を
$\Sigma\subset\operatorname{int}K_0$ と書く．内側の structure は
補題\ref{ke:compact-relative-morse-smale}の閉 isotropic filtration を持つよう選ぶ．

$S$ 上に compact graph $\Gamma$ への connected-fiber continuous surjection
$\pi_S:S\to\Gamma$ があり，graph の開辺上では smooth submersion とする．
Regular fibers は $\omega$-Lagrangian $T^2$，
三価頂点上の fibers は smooth strata が isotropic な $S^1\times\figeight$ とする．
このとき $\pi_S$ を境界上でそのまま保つ proper continuous surjection
\[
 \Pi:K\longrightarrow\operatorname{Cone}_{[0,1]}(\Gamma)
\]
があり，cone point の逆像は $\Sigma$，その他の fibers は上記の regular または
branch fibers の Liouville images である．全 fibers は連結であり，
全 smooth strata は isotropic である．
\end{lemma}

\begin{proof}
正時間に $S$ へ達しない点の集合は $\Sigma$ である．
実際，$K_0$ の外で新しい critical point 又は trapping はなく，
外側の任意の compact shell 上で Lyapunov function の増加率は正の下界を持つ．
また $\partial K_0$ は外向きであるので，その外から正時間に $K_0$ へ戻れない．
従って $K\setminus\Sigma$ の各点は有限正時間で $S$ に到達する．

負時間の flow は $K$ の中に留まり，compactness により全負時間で存在する．
境界での横断性と ODE の一意性により，
\[
 (-\infty,0]\times S\longrightarrow K\setminus\Sigma,
 \qquad(s,y)\longmapsto\varphi_s(y)
\]
は diffeomorphism である．$\mathcal L_Z\lambda=\lambda$ より
\[
 \varphi_s^*\lambda=e^s\lambda|_S,\qquad
 \varphi^*\omega=d(e^s\lambda|_S).
\]
従って
\[
 \Pi(\varphi_s(y))=[e^s,\pi_S(y)],\qquad \Pi(\Sigma)=\{0\}
\]
と定義する．各正の radial level の fiber は $\pi_S$ の一つの全 fiber の像であり，
flow の conformally symplectic 性から Lagrangian 性と isotropic strata が保たれる．

連続性を確認する．任意の有限 $T>0$ について，
$\{\varphi_s(y):-T\le s\le0,\ y\in S\}$ は $\Sigma$ と交わらない compact set である．
従って $\Sigma$ へ近づく点列の退出時間は無限大に発散し，$e^s\to0$ となる．
逆に compactness と上の表示から，radial coordinate が零に近づく点列の集積点は
$\Sigma$ に含まれる．これで cone point における連続性が従う．
Skeleton の連結性は flow から直接従う．実際，
\[
 K_T=\varphi_{-T}(K)\quad(T\ge0),\qquad
 \Sigma=\bigcap_{T\ge0}K_T
\]
である．後者の等式は，$K_T$ が正時間 $T$ まで $K$ 内に留まる点全体であることと，
上の有限退出の結論による．$K_T$ は非空な compact connected sets の減少族なので，
その共通部分 $\Sigma$ も非空で連結である．
最後に compact から Hausdorff 空間への連続全射なので proper である．
\end{proof}

\begin{remark}[Contact form の変更と指定された境界葉]
\label{ke:contact-rescaling-collar-scope}
前補題の Liouville flow による延長は，Evans--Mauri の既知の cone-extension
argument を，指定された境界へ相対化したものである．同論文の
Lemma~2.16 と Remark~2.18\cite[pp.~543--544]{EM}は symplectic completion
$\overline W$ を扱い，link の元の形式 $\beta=\lambda|_{\mathrm{Link}(W)}$ と異なる
$\alpha=e^h\beta$ に関する fibration も許す．このとき二つの symplectizations は，
$d(e^r\alpha)$ 側から $d(e^s\beta)$ 側への写像
$(r,x)\mapsto(r+h(x),x)$ で同定される．一般には元の固定 link は保たれない．

また，元の $d\beta$-Lagrangian fiber $F$ に対して
\[
 d(e^h\beta)|_F=e^h\,dh|_F\wedge\beta|_F
\]
だから，同じ $F$ が変更後も Lagrangian であるためには右辺の消滅が必要である．
例えば $\TT^3=(\RR/2\pi\ZZ)^3$ の座標を $(x,y,z)$ とし，
\[
 \beta=\cos z\,dx+\sin z\,dy,\qquad h=\sin x
\]
と置く．$\beta$ は contact form で，各 $z$-level torus は $d\beta$-Lagrangian である．
しかし $F=\{z=\pi/2\}$ では
\[
 d(e^h\beta)|_F=e^{\sin x}\cos x\,dx\wedge dy\not\equiv0 .
\]
これは contact 構造を変えない rescaling でも，固定された葉を保つとは限らない
という比較例である．本稿では $\beta=\lambda|_S$ 自体に関する既存の全境界葉を
入力とし，その形式を変更せず compact domain 内へ延長する．
Contactomorphism 型の一致だけから，この相対的な葉の保存を結論していない．
\end{remark}

\subsection{外向き collar と既存 neck 葉の整合}

\begin{proposition}[Actual link を起点とする相対 collar 接着]
\label{ke:link-relative-collar-gluing}
閉 link $S_{i,j,t}$ と，その近傍で既存 normal circle および corner $T^2$ action に
不変な primitive $\lambda^{\rm ad}_{i,j,t}$ を取る．
$S_{i,j,t}$ は既存部分 fibration の領域内にあり，その全 fibers で飽和しているとする．
Normal face は corrected moment
\[
 M_a=q_a-\rho_a,\qquad q_a=\delta_j>0,\qquad
 0<\delta_j<\min(\iota_j/8,\rho_a/8)
\]
で指定し，$q_a=\pm\widehat\mu_N$ とする．
Corner の丸めは同じ moment functions の smooth monotone minimum を用いる．
内部に新しい trapping がなく，内側では固定 Weinstein core の actual 移送と一致するとする．

このとき既存部分 fibration の葉族を source critical moment から離れた小さい
normal bands 内だけで変形し，前補題の core filling と全開 collar 上で一致させられる．
保持する collar 内の全 corner torus fibers は保ち，source の全 pinched tori とその近傍，
normal marking，saddle fibers の型および Reeb marking を保つ．
Saddle の位置，第二 action および critical values の数値は固定しない．
\end{proposition}

\begin{proof}
$Z$ を $\lambda^{\rm ad}_{i,j,t}$ の Liouville field とする．
各 invariant generator $V_a$ について
\[
 [Z,V_a]=0,\qquad ZM_a=M_a
\]
である．前者は $\mathcal L_{V_a}\lambda^{\rm ad}=0$ と
$\mathcal L_{V_a}\omega_t=0$，後者は $\mathcal L_Z\lambda^{\rm ad}=\lambda^{\rm ad}$ による．
全 $M_a$ は link とその小さい collar 上で負である．
従って normal face で Liouville time と $q_a$ は
\[
 q_a(s)=\rho_a+e^s(\delta_j-\rho_a)
\]
によって互いに滑らかに対応する．

各 normal face の既存 reduced 葉を anchor の葉として取り，
Liouville flow により $q_a=\delta_j$ の両側の小さい band へ運ぶ．
Compactness と transverse moment の非消滅から共通 flow buffer がある．
これは simple-saddle foliation の smooth family であり，全葉は連結なままである．
Corner では $Z$ が全 $T^2$ action と可換なので，flow は torus orbits を torus orbits へ写す．
従って end collars 上の新旧二族は同じ full torus partition である．

Reduction の座標を共通 corner partition に適合させて取る．
そのために固定した一つの moment level から regular annular end charts を続ければよく，
parameter ごとに first integral の表示値を揃える必要はない．
End collars はこの band 内の全 neck--corner intersections を含むように選ぶ．
実際，節\ref{ke:core-invariant-primitive-section}の切替は source label
$B\in[5.2b_j,5.4b_j]$ で終わり，全交わりでは
$b_{\rm res}>11b_j/2$，$B=b_{\rm res}(1+O(b_j^{-2}))$ なので，
恒等域の内側境界を切替帯と全交わりの間に取れる．
定理\ref{ke:saddle-moving-relative-interpolation}と
系\ref{ke:actual-saddle-moving-neck-interpolation}を用い，
$q_a=\delta_j$ の近くでは Liouville 葉族，band の外側では旧葉族となるよう補間する．
Band は例えば $|q_a-\delta_j|<\delta_j/4$ の中へさらに縮める．
Source critical moment は $\widehat\mu_N=0$ なので，その全 fiber と近傍は支持外である．
同じ double line の反対側 component には opposite moment sign の band を使うため，
二つの補間の支持は互いに素である．

Corner rounding 部分の collar では torus orbits をそのまま使える．
各 normal band の補間は end collars で恒等なので，これらの局所変更は全 overlaps 上で一致する．
全て actual $\omega_t$ の Hamiltonian reduction の中で行うため，
補間後の fibers は actual Lagrangian である．
内側 cone filling とは一つの Liouville 葉族として開 collar 全体で一致する．
\end{proof}

\begin{lemma}[補間後の全葉を保つ平行 Liouville link]
\label{ke:post-interpolation-parallel-link}
補題\ref{ke:closed-component-link}と
命題\ref{ke:link-relative-collar-gluing}の固定 stage の設定を取る．
正の residue constants を $\rho_a>0$ とし，
$\rho_-=\min_a\rho_a$，$\rho_+=\max_a\rho_a$ と置く．
$\delta=4\pi/N_j<\rho_-/8$ とし，corner rounding に用いる
$0<\varepsilon<\delta/20$ を固定する．

相対補間後に境界葉が Liouville 葉族と一致する共通開 collar の中で，
十分小さい $s_0>0$ と平行 link
\[
 S^-_{i,j,t}:=\varphi_{-s_0}^Z(S_{i,j,t})
       \subset\operatorname{int}K_{i,j,t}
\]
を取れる．この link とその小さい両側 collar は，補間\emph{後}の
全 regular torus，branch fiber のいずれについても飽和しており，
core 側の Liouville cone partition と外側の補間後の partition は
その開 collar の全 fibers として厳密に一致する．
さらに normal face の anchor $q_a=\delta$ から出る Liouville trajectory は
\begin{equation}\label{ke:eq:post-interpolation-parallel-moment}
 q_a(\varphi_s^Zx)
  =\rho_a+e^s(\delta-\rho_a)
\end{equation}
を満たす．$s_0$ は，各固定 $j,t$ について
\[
 0<s_0<\frac{\delta}{128\rho_+}
\]
となるよう選べる．この不等式は共通 collar の選択であり，
$s_0$ の $t$ 一様な正の下界や新たな距離比較を主張しない．
\end{lemma}

\begin{proof}
$\lambda^{\rm ad}$ は指定した normal action と corner $T^2$ action に
不変で，$M_a=\lambda^{\rm ad}(V_a)=q_a-\rho_a$ である．
$\iota_Z\omega_t=\lambda^{\rm ad}$ と
$\mathcal L_Z\lambda^{\rm ad}=\lambda^{\rm ad}$ から
$ZM_a=M_a$ となる．よって Liouville flow の定義された共通
collar では
\[
 q_a(\varphi_s^Zx)=\rho_a+e^s(q_a(x)-\rho_a).
\]
この式は，固定点を含まない normal face の全相と全 reduced 葉に
同じく成立し，anchor 上では表示された式になる．

Corner rounding で用いる
$F_{ac}=(q_a+q_c-2\delta-\chi_\varepsilon(q_a-q_c))/2$
を考える．$\chi_\varepsilon=|\cdot|$ outside
$[-\varepsilon,\varepsilon]$ だから，丸めの変化している
$S_{i,j,t}\cap\{|q_a-q_c|\le\varepsilon\}$ 上では
\[
 q_a+q_c=2\delta+\chi_\varepsilon(q_a-q_c),\quad
 \chi_\varepsilon(q_a-q_c)\le\varepsilon,
 \quad\Longrightarrow\quad
 \delta\le q_a,q_c\le\delta+\varepsilon .
\]
したがって，この rounding 上の両 moment は
$\rho_-$ より十分小さい．$|s|\le\delta/(128\rho_+)$
で $\lvert e^s-1\rvert\le2|s|$ を使えば，
normal face 上の cut moment と rounding 上の
両 moment の変位はいずれも $\delta/64$ 以下である．
特に，normal face は anchor の
$|q_a-\delta|<\delta/8$ の内部に，rounding は同じ
corner $T^2$ の全 orbit が定義される固定の相対開 collar
の内部に保てる．この評価は $N_j$ より前に選んだ
residue constants 以外の moving depth の微分を用いない．

相対葉族補間の cutoff は，normal face の anchor のさらに小さい
開近傍では Liouville 葉族そのものと一致するよう選ばれている．
Corner では $[Z,V_a]=[Z,V_c]=0$ なので Liouville flow は
各\emph{全} torus orbit を全 torus orbit へ運ぶ．
上の数式だけで共通 flow buffer の存在を結論している
わけではない．その存在は補題\ref{ke:closed-component-link}の
compact で横断的な link，および
命題\ref{ke:link-relative-collar-gluing}で用意した
flow・補間両者の開 collar から得る．
それらの有限個の collar 幅より十分小さい $s_0$ を選び，
同時に上の $\delta/(128\rho_+)$ の上限を課す．

$Z$ は $S_{i,j,t}$ 上で外向きだから，
負時間の像 $S^-=\varphi_{-s_0}^Z(S)$ は core 内部にある．
その全点は補間が Liouville 葉族と一致する開域内にあり，
$\varphi_{-s_0}^Z$ は境界の各 whole fiber を
whole fiber へ写す．従って，$S^-$ およびその小さい
両側 collar は，補間\emph{後}の全 fiber の和である．
この同じ collar の二つの partition はもともと
境界 $S$ からの同一 Liouville flow によって定義されたため，
部分葉や異なる reduced 葉の暗黙の同一視なしに厳密に一致する．
最後の主張は定性的な saturation であり，
Liouville trajectories の actual 長さの一様上界や
任意の $t$ に対する固定正幅を必要としない．
\end{proof}

\begin{lemma}[正の normal moment を保つ Liouville 正時間の上限]
\label{ke:positive-normal-moment-time-ceiling}
補題\ref{ke:post-interpolation-parallel-link}の固定stageの normal face
\(q_a(0)=\delta=4\pi/N_j\) を取る．\(0<\delta<\rho_a\) とし，
正時間 \(0\le s\le\tau\) の Liouville trajectory が存在して，
全区間で同じ正の normal moment chart \(q_a(s)>0\) に留まるなら
\[
 0\le\tau<\log\frac{\rho_a}{\rho_a-\delta}.
\]
全 normal faces を同じ正時間だけ正の moment charts 内に保つなら，
\[
 \tau<\min_a\log\frac{\rho_a}{\rho_a-\delta}
 =\log\frac{\rho_+}{\rho_+-\delta}
 =\frac{4\pi}{\rho_+N_j}+O(N_j^{-2}).
\]
したがって，この正の moment chart を保持するという条件の下では，
\(j\to\infty\) に一様な正の Liouville 時間幅は存在しない．
これは同じ chart 内の正時間に関する必要条件であり，
負時間の平行linkの存在や，chart を離れた全flowの生存・距離評価には
何も主張しない．
\end{lemma}
\begin{proof}
不変 primitive の恒等式 \(ZM_a=M_a\) と
\(M_a=q_a-\rho_a\) を積分すると
\[
 q_a(s)=\rho_a+e^s(\delta-\rho_a)
\]
である．\(q_a(\tau)>0\) は
\(e^\tau<\rho_a/(\rho_a-\delta)\) と同値である．
右辺の対数は \(\rho_a>\delta\) 上で \(\rho_a\) の減少函数であり，
共通の上限は \(\rho_+=\max_a\rho_a\) で決まる．
最後は \(-\log(1-\delta/\rho_+)=\delta/\rho_++O(\delta^2)\)
による．
\end{proof}

\begin{lemma}[平行 Liouville link の actual 計量移動費用]
\label{ke:parallel-liouville-metric-cost}
命題\ref{ke:joint-source-component-buffer}の順序で固定した十分大きい
stage \(j\) と component cut \(N=N_j\) を取る．
補題\ref{ke:closed-component-link}の閉 link を \(S_{i,j,t}\)，
その近傍の adapted Liouville 場を \(Z_{i,j,t}^{\rm ad}\) と書く．
Core の compact exact Moser によって移送した正 residue
exhaustion を \(\Psi_{i,j,t}\) とする．

固定次数・退化族・正 residue weights のみに依存する
\(0<c_0<C_0\)，\(A>0\)，\(\theta>0\) があり，
各固定 \(j,N\) の後に \(\lvert t\rvert\) を十分小さくすると，
全ての \(x\in S_{i,j,t}\) について
\(\varphi_s^{Z^{\rm ad}}(x)\) は
\(\lvert s\rvert\le\theta/N\) で存在し，
\begin{equation}\label{ke:eq:parallel-liouville-shell-and-cost}
 \frac{c_0N}{2}
 \le\Psi_{i,j,t}(\varphi_s^{Z^{\rm ad}}x)\le2C_0N,
 \qquad
 \operatorname{Length}_{g_t}
 \bigl(\varphi_{[0,s]}^{Z^{\rm ad}}x\bigr)
 \le2AC_0N\lvert s\rvert
\end{equation}
を満たす．移送した component graph 上で計量を読み，
軌跡は全曲面の actual 計量により測る．

特に補題\ref{ke:post-interpolation-parallel-link}の負時間
\(s_0=s_0(j,t)>0\) はさらに
\(s_0<\theta/N\) を課して選べる．その平行 link \(S^-\) に対し
\[
 \sup_{x\in S_{i,j,t}}
 d_{g_t}\bigl(x,\varphi_{-s_0}^{Z^{\rm ad}}x\bigr)
 \le2AC_0\theta,\qquad
 R_t^{-1}\sup_xd_{g_t}
 \bigl(x,\varphi_{-s_0}^{Z^{\rm ad}}x\bigr)
 \le\frac{2AC_0\theta}{R_t}.
\]
これは全ての固定 stage に同じ定数で成り立つが，
補間後の葉族と一致する開 collar の \emph{時間幅}に
\(j,t\) 一様な正の下界を与えるものではない．
任意の core 内 Liouville trajectory 全体の長さも評価しない．
\end{lemma}

\begin{proof}
式\eqref{ke:eq:positive-residue-ke-potential}と正の固定 weights から，
十分深い component end では \(\Psi\asymp H=1+\sum_as_a\) である．
系\ref{ke:component-shell-primitive-support}の link 深度評価から，
固定した二定数 \(c_0,C_0\) で
\(c_0N\le\Psi_{i,j,t}\le C_0N\) on \(S_{i,j,t}\) となる．
元の成分極限上で
\(\{c_0N/3\le\Psi\le3C_0N\}\) を含む
\(t\) に依存しない compact \(V_j\Subset W_i\) とその外側 buffer を，
\(N\) を固定した後，\(t\) の閾値より先に取る．
固定 compact graph 収束と輸送の \(C^1\) 小ささにより，
十分小さい \(t\) では，輸送した shell
\[
 \mathcal S_{j,t}
 =\left\{\frac{c_0N}{2}\le\Psi_{i,j,t}\le2C_0N\right\}
\]
はこの graph image と buffer の内部に入る．

式\eqref{ke:eq:weighted-core-input}は
\(\lvert\lambda_0\rvert_h+\lvert d\Psi\rvert_h\le CH\) を与える．
命題\ref{ke:star-cover-adapted-primitive}の
\(\lvert dF\rvert_h\le C\delta_jH\) と
式\eqref{ke:eq:adapted-positive-period-exact}の
\emph{固定 compact 上の} exact transport を合わせる．
各固定 \(j,N\) の後に
式\eqref{ke:eq:component-fixed-convergence}と
命題\ref{ke:fixed-component-transfer}の
\(C^1\) 小ささを必要なだけ課せば，
Rochon--Zhang の cusp metric 比較を含めて，
\(\mathcal S_{j,t}\) 上で
\[
 \lvert\lambda^{\rm ad}_{i,j,t}\rvert_{g_t}
 +\lvert d\Psi_{i,j,t}\rvert_{g_t}
 \le A_0\Psi_{i,j,t}
\]
となる．ここで \(A_0\) は \(j,N,t\) に依存しない．
\(t\) の閾値は各固定 stage の広い compact に依存してよい．
Kähler 計量の同一視
\(\lvert Z^{\rm ad}\rvert_{g_t}
 =\lvert\lambda^{\rm ad}\rvert_{g_t}\) より，
定数を一度増やして
\[
 \lvert Z^{\rm ad}\rvert_{g_t}\le A\Psi_{i,j,t},
 \qquad
 \lvert Z^{\rm ad}\Psi_{i,j,t}\rvert
 \le A\Psi_{i,j,t}^{\,2}
\]
を得る．ここでは source 中央の moving graph を使わず，
各固定した component shell 上の比較だけを使う．

Liouville time を \(s\) とし，軌道上で
\(u(s)=1/\Psi_{i,j,t}(\varphi_s^{Z^{\rm ad}}x)\) と置くと
\(\lvert u'(s)\rvert\le A\) である．初期値は
\(1/(C_0N)\le u(0)\le1/(c_0N)\) なので，
\[
 0<\theta<
 \min\left\{\frac{1}{4AC_0},
                \frac{4\pi}{128\rho_+}\right\}
\]
を選べば，\(\lvert s\rvert\le\theta/N\) で
\(c_0N/2\le\Psi_{i,j,t}\le2C_0N\) が保たれる．
この shell は固定した compact の内部にあるから，
常微分方程式の延長定理によりその全時間で flow が存在する．
最後に
\[
 \operatorname{Length}_{g_t}
 \bigl(\varphi_{[0,s]}^{Z^{\rm ad}}x\bigr)
 \le\int_0^{\lvert s\rvert}A(2C_0N)\,dr
 =2AC_0N\lvert s\rvert .
\]
既存の開 collar との葉族一致には
補題\ref{ke:post-interpolation-parallel-link}をそのまま使い，
その \(s_0\) を必要ならさらに小さくする．
このため \(j,t\) 一様な collar 幅は不要である．
\end{proof}

\begin{remark}[時間幅の消失と距離費用の違い]
\label{ke:parallel-liouville-cusp-sharpness}
単純な cusp radial 模型
\(g=(du^2+d\vartheta^2)/u^2\)，\(q=4\pi/u\) で
\(Zq=q-\rho\)，\(q(0)=4\pi/N\) とすると，
\(\tau=\theta/N\)，\(0<\rho\theta<4\pi\) に対して
\[
 \lim_{N\to\infty}
 \operatorname{Length}_{g}(\varphi_{[0,\tau]}^Zx)
 =\log\frac{4\pi}{4\pi-\rho\theta}>0.
\]
従って \emph{時間が零へ縮む}という事実だけでは，
正規化前の軌跡長が \(o(1)\) になるとは限らない．
上の \(O(1)\) はこの cusp 模型の尺度と整合する．
これは固定モデルでの例であり，actual KE の
parallel link の長さの最適下界を主張しない．
\end{remark}

\subsection{全曲面上の actual fibration と固定段階の費用}

\begin{proposition}[Source と core の buffer・primitive の同時有限段階選択]
\label{ke:joint-source-component-buffer}
固定した次数・一般位置族と十分大きい K2 residual stage $j$ を取る．
Source を含む有限個の $P\Subset P^+$，全 marked source/end data を
normal cut より先に固定する．共通の normal circle と centered homotopy を
source から K2 の end まで用いる．
このとき $b_j,n_j,\iota_j,N_j$，最後に $t$ の閾値という順序で，
以下を同時に満たすように選べる：
\begin{enumerate}
\item Source の全 $K_{t,n_j}$ に対して normal Moser の
forward/backward flow が存在し，source と image は共に
$K_{t,n_j}$ の開近傍を含む．Fixed points は保たれ，
actual trajectory 長は $C_P/n_j$ 以下である．
\item Component 側の全 link とその averaging orbits，
cutoff supports，flow buffers は固定 component compact に入り，
retained collars で adapted primitive が normal/corner actions に不変となる．
\[
 \sigma_j=\frac{C_*}{b_j}+\frac{C_j^{\rm comp}}{N_j},\qquad
 Z^{\rm ad}\Psi\ge(c-C\sigma_j)H^2\ge\tfrac c2H^2>0
\]
とできる．
\item 同じ固定 compact 上で
\eqref{ke:eq:adapted-positive-period-exact} により adapted primitive と
positive-residue primitive の全閉周期が厳密に一致する．
Normal face の corrected moment は
$M_a=q_a-\rho_a$，$\rho_a=2\pi w_a>0$ であり，
neck--corner overlap に別の moment 原点を導入しない．
\end{enumerate}
これは各\emph{固定}stage の同時実現であり，$j$ 全体に一様な
source buffer や Liouville 時間幅，あるいは大域 GH 速度を主張しない．
\end{proposition}
\begin{proof}
最初に $b_j$ を十分大きい stage とし，
式\eqref{ke:eq:actual-generator-delta} の $C_*/b_j$ を
$c/(4C)$ より小さく選ぶ．この選択で固定した residual
data を以下の全段階で保持する．
各固定 $P\Subset P^+$ について
補題\ref{ke:fixed-stage-source-neck}は，$n$ に依存しない
$\delta_P,C_P,c_P>0$ を与える．有限個の最大・最小を取り，
\[
 (2c_w+\tfrac12)C_P/n_j<\delta_P/2,\qquad
 c_w=\pi/(4\sqrt3)
\]
となるまで $n_j$ を大きく選ぶ．
系\ref{ke:short-circle-image-buffer}の仮定が満たされ，
全 source の正逆Moser flowは共通の
source/image領域を持つ．Source fixed point では生成場も
centered homotopy primitiveも零なので flow は固定点を保つ．
その後，$0<\iota_j<c_P/(2n_j)$ を K2 の他の moment 幅条件も
満たすように固定する．$n_j,\iota_j$ の選択後に
$t$ を小さくすることはまだ行わない．

Source stage と既存K2の
全 residual data を固定した後，
補題\ref{ke:component-normal-generator-c1} と
補題\ref{ke:relative-corner-patch-c2} の有限な
$C_j^{\rm comp}$ を取る．$N_j$ をさらに十分大きくして
\[
 N_j\iota_j\gg1,\qquad
 \frac{4\pi}{N_j}<\min\{\iota_j/8,\rho_-/8\},\qquad
 C_j^{\rm comp}/N_j<c/(4C),\quad \rho_-=\min_a\rho_a
\]
を課せる．系\ref{ke:component-shell-primitive-support}の
star cover が必要な full orbits，supports と trajectories を
固定 compact 内に含める．命題\ref{ke:star-cover-adapted-primitive}は
retained invariant primitive と上の nontrapping bound を与える．
ここで source の $C^0$ trajectory 長から component generator の
$C^1$ 小ささを導いたのではない．

固定 component compact では
命題\ref{ke:fixed-component-transfer}の exact Moser を先に行い，
その後に同じ actual circle/torus actions で平均する．
式\eqref{ke:eq:adapted-positive-period-exact}によって全閉周期が保たれ，
補題\ref{ke:component-meridian-normalization}の winding 計算により
neck と corner の normal moment は定数のずれなく一致する．
最後に固定した全 source/corner/core domains に必要な smooth
convergence・正逆flow・高階比較の有限個の $t$ 閾値の最小を取る．
各段階の全要求を満たすが，その閾値や定数に $j$ 一様性は不要である．
\end{proof}

\begin{theorem}[全平面退化の actual KE--Lagrangian 全域構成]
\label{ke:actual-global-fibration}
節\ref{ke:global-realization-section}の設定で，任意の固定 $d\ge5$ と退化族を取る．
十分大きい各固定 exhaustion stage $j$ について，$|t|$ を十分小さくすれば，
一つの admissible marked tree system $T$ と connected-fiber continuous surjection
\[
 \pi_{j,t}:(X_t,\omega_t)\longrightarrow B_d(T)
\]
を構成できる．Regular fibers は actual $\omega_t$-Lagrangian $T^2$，
source fibers は一重 pinched tori，branch fibers は $S^1\times\figeight$ であり，
$d$ 個の core fibers は成分極限 KE 形式の finite-type Weinstein skeletons の
symplectic images である．全 smooth singular strata は isotropic である．

構成は系\ref{ke:neck-corner-saturated-union}の部分 fibration を，
component link より内側で全 fibers ごと切り詰めて core と接着する．
各 corner の保持された全 torus fibers と source の近傍は変えない．
Normal bands 内の saddle positions と基底の parameterization は $t$ とともに動かしてよい．
木系の組合せ型と component labels は固定できる．

また，core と今回変更した collar の全体は，各 $j$ に対して固定した component compacts の
actual images に含められる．各 component ごとの領域の actual ambient diameter と，
その少し広い固定された連結 compact domain の actual image 内で選ぶ接続経路，
および小さい Moser 移送の費用は
$C_j=C(1+\log(N_j+1))$ で抑えられる．ここで $C$ は固定した次数・族・weights
のみに依存し，stage と $t$ に依存しない．
従って十分遅い対角選択 $j=j(t)\to\infty$ により，それらの
$R_t^{-1}$-正規化費用は零へ収束する．
全曲面の uniform distance distortion は本定理の結論ではない．
異なる components の core domains を全部合わせた集合の diameter が
$O(C_j)$ であるとは主張していない．
\end{theorem}

\begin{proof}
命題\ref{ke:joint-source-component-buffer}の同時選択を行った上で，
K2 の fixed stage を取り，補題\ref{ke:fixed-stage-stratum-coverage}を使う．
既存部分 fibration の外部は，各成分の互いに素な固定 compact graph images に含まれる．
それらの compact set と全 finite-type cores を内部に含むほど深く normal cuts を選び，
補題\ref{ke:closed-component-link}の閉じた飽和 link $S_{i,j,t}$ を作る．
選択の順序は $b_j$，$n_j$，$\iota_j$，その後に component 深度 $N_j$，最後に $t$ の閾値である．
特に $N_j\iota_j\gg1$ を課す．

各成分の固定 compact neighborhood 上で先に exact Moser 移送を行い，
finite-type Weinstein primitive を actual form $\omega_t$ の primitive として移す．
これを actual circle actions で centered 平均し，normal-invariant cutoffs を用いて貼り合わせる．
Primitive 補正は exact であり，component core の近傍では零である．
生成場の $C^1$ 評価と正 residue potential の coercive bound により，
新しい trapping を作らず，全 link 上で指定された invariance を得る．
必要な cutoff support 全体の被覆と flow buffers は
系\ref{ke:component-shell-primitive-support}による．
補題\ref{ke:component-meridian-normalization}の厳密な period 規格化の下で，
negative corrected moments と monotone rounding によって $Z$ は各 link に外向きである．

Link の quotient も具体的に同定できる．各 normal face の quotient は $T_{ij}$ であり，
corner rounding の quotient は対応する labelled ends を結ぶ interval である．
後者には新しい branch point はない．従って link quotient $\Gamma_i$ は，
各 $T_{ij}$ の対応する葉を交点ごとに貼り，二価点を除いた graph であり，
補題\ref{topo:counts}の $G_i$ と同じ labelled graph である．

Link 自身の葉は既存部分 fibration の全 fibers なので，
補題\ref{ke:actual-link-cone-filling}を用いて内部を充填する．
命題\ref{ke:link-relative-collar-gluing}により，外向き collar で旧葉族へ戻す．
この補間後の部分葉族を $\widetilde f_{j,t}$ と書く．
切詰めは，補間前の partition ではなく，この $\widetilde f_{j,t}$ に対して行う．
両葉族が Liouville 葉族に一致する band 内で，$S_{i,j,t}$ の少し内側に
平行 moment link $S^-_{i,j,t}=\partial K^-_{i,j,t}$ を取る．
Normal face は同じ moment の level，corner face は同じ全 $T^2$ orbits の和なので，
$S^-_{i,j,t}$ は補間後も全 fibers で飽和している．
平行linkの全葉飽和性と共通開collar上の一致は
補題\ref{ke:post-interpolation-parallel-link}で確認した．
この有限移動の actual 長さは補題\ref{ke:parallel-liouville-metric-cost}により
stage 一様に $O(1)$ であり，既存の対数接続予算に吸収できる．
従って補題\ref{ke:whole-fiber-core-trim}の連結 fiber による議論をそのまま適用し，
$\widetilde f_{j,t}$ を
$U_{j,t}\setminus\bigcup_i\operatorname{int}K^-_{i,j,t}$ へ制限する．
Core 側は元の $K_{i,j,t}$ のままとすれば，その交わりは両 partition が厳密に一致する
共通 Liouville collar となる．$X_t\setminus U_{j,t}$ を覆う各指定 component compact は
対応する $\operatorname{int}K_{i,j,t}$ に入っているので，平行 link を十分近く選べば
$\operatorname{int}K^-_{i,j,t}$ にも入り，全 $X_t$ の被覆が保たれる．
共通 collar 内での基底の分割は，補題\ref{ke:anchor-trimmed-neck-domain}のとおり
元の anchor $S_{i,j,t}$ へ取り直せる．

全 pieces の fibers は連結で，全 overlaps 上で同じ全 fiber として一致する．
Core 側は基底の compact cone 全体，保持 neck は閉 moment interval と compact Reeb domain，
corner は閉 moment domain の全逆像として切る．全域被覆と全 fiber の保持を
上で確認したので，これらは有限個の compact な飽和閉片を与える．
各片の base は compact Hausdorff であり，
補題\ref{glue:closed-whole-fiber-quotient}が適用できる．
従って実際の fiber partition の商 $Q_{j,t}$ は compact Hausdorff で，
$X_t\to Q_{j,t}$ は proper な連結 fiber の連続全射となる．
各 base piece は $Q_{j,t}$ へ閉埋込みされ，global fiber は一つの local fiber と一致する．
この段階では $B_d(T)$ との同定も，後の計量に適合する parameter 化も用いない．
基底の combinatorics は，component に cone over $\Gamma_i$，double line に
interval times $T_{ij}$，triple point に regular two-dimensional block を置いたものとなる．
Saddle tracks は対応する二つの core vertices を結び，
$T_{ij}$ の内部辺は bigons，labelled ends は三重点の triangle に接続する．
保持 corner の quotient が disk のままであり，core cut の境界が proper arcs で
あることは補題\ref{ke:corner-cap-disk}による．同補題の証明は moment cuts と
深度評価のみを用い，ここでの $B_d(T)$ との同定には依存しない．
相対葉族補間は continuation marking を変えないので，この基底は
定義\ref{topo:definition}の cellular $B_d(T)$ と同相である．
Source の nodal values は，その二次元 cells の内部の標識点として残る．
実数値 actions を合わせる追加条件又は大域 matroidal atlas は使わない．

固定 $j$ で用いた component 領域と全補間 supports は，
一つの固定 $V_{i,j}\Subset W_i$ の中にある．
補題\ref{ke:core-depth-upper}と命題\ref{ke:component-logarithmic-connectors}により，
各 component image 内の二点は，少し広い固定 compact image 内の経路で
長さ $C(1+\log(N_j+1))$ 以下で結べる．
従って定理の $C_j$ はこの対数評価で取れる．
ここで component 全域の二点を結ぶために用いるのは cusp の動径経路であり，
任意の Liouville trajectory の全長ではない．
ただし閉 link から近い平行 link への特定の Liouville 移動は
補題\ref{ke:parallel-liouville-metric-cost}で $O(1)$ に抑えている．Reduced diffeomorphisms の moment 微分について
$j$ 一様の bound を仮定する必要もない．

最後に $t_j\downarrow0$ を取り，既存 K1--K2 の全条件とともに
\[
 R_t\ge jC\bigl(1+\log(N_j+1)\bigr)
 \qquad(0<|t|\le t_j)
\]
を課す．$t_{j+1}<|t|\le t_j$ では第 $j$ 段階を使えば，
core と補間領域の正規化費用は $1/j$ 以下となる．
この選択は各固定次数・固定族について行い，次数一様性は主張しない．
\end{proof}

本定理は actual KE 形式上の全域 fibration と，各 component の固定段階の接続費用を与える．
同一 moment の source 接続は定理\ref{ke:source-all-moment-exit}，
基底値の整合は命題\ref{ke:global-base-parameterization}で扱う．
これらを用いて，$d\ge6$ で潰される木上の全 fiber 合併径を
命題\ref{ke:collapsed-tree-full-fiber-collapse}で評価し，
全曲面の距離比較を定理\ref{ke:global-ambient-lower}，
\ref{ke:global-ambient-upper}で証明する．
閉 corner から core 側へ切り取った部分の intrinsic GH 評価を，
新しい全曲面上の距離評価と同一視しない．

\subsection{Core の深度と対数接続費用，除いた corner の旧 target 像}
\label{ke:core-target-caps-section}

以下では十分大きい段階 $j$ を固定し，$m_j=6b_j$，$N=N_j$，
$\delta_j=4\pi/N_j$，$L=L_t$，$R=\log L$ とする．
相対化済みの旧 closed corner を
\[
 W_{t,J}=C_{t,J}(\mathcal H_{t,m_j}),\qquad
 P^{\rm old}_{t,J}=P_{t,m_j}\circ C_{t,J}^{-1}
\]
と書く．$K_{i,j,t}$ は補題\ref{ke:closed-component-link}と
定理\ref{ke:actual-global-fibration}で実際に作った compact core domain とする．
小さい共通 collars を確保するため平行 moment levels へ動かした場合も含める．
それらの levels は $q_a\in[\delta_j/2,2\delta_j]$ の範囲に取る．

\begin{lemma}[固定 compact core の深度の上界]
\label{ke:core-depth-upper}
$N_j$ の下限を必要に応じて大きくすることで，固定した族・chart・正 weights のみに
依存する $C>0$ に対し次を実現できる．固定 $j$ の後に $|t|$ を十分小さくすると，
\[
 \max_{a\ne i}s_a\le C N_j
 \quad\text{on }E_{i,t}^{-1}(K_{i,j,t}),
\]
が成立する．ここで $s_a=-\log\|s_{D_{ia}}\|^2>1$ は成分 $W_i$ 上で固定した
global depth functions である．従って $i\in J=\{i,a,c\}$ と
$x\in K_{i,j,t}\cap W_{t,J}$ に対し，旧 corner source
$z=C_{t,J}^{-1}(x)$ の二つの独立深度は
\begin{equation}\label{ke:eq:removed-corner-depth-upper}
 \ell_a(z),\ell_c(z)\le C N_j
\end{equation}
を満たす．$i\notin J$ ならば，十分小さい $t$ に対して
$K_{i,j,t}\cap W_{t,J}=\varnothing$ である．
\end{lemma}

\begin{proof}
まず link 上を調べる．Normal face の source level は fixed-end moment 評価と
単調性により深度 $N_j$ と固定比の間にある．残余基底の compact は，三重点近傍の
end 深度 $O(b_j)$ までを含むように固定されている．従って normal face 上では，
切っている divisor の深度は $O(N_j)$，他の divisor の深度は $O(b_j)$ である．
Corner face と丸め部分では，二つの独立深度がともに $O(N_j)$ 以下である．
これは\eqref{ke:eq:component-two-action-estimate}と
補題\ref{ke:closed-component-link}の moment-level graphs による．

この評価を actual link へ移す．Normal faces では，同じ centered normal flow を，
normal 深度が $N_j$ と固定比の間にある source 上へ制限する．
補題\ref{ke:fixed-stage-source-neck}の短円評価と ODE の一意性により，
各 trajectory の actual 長さは $C_{P_j}/N_j$ 以下である．
初期点の全 component depths は，固定 frame 差と残余深度の上界を $N_j$ に吸収して，
固定した $C_0$ に対し $C_0N_j$ 以下にできる．固定 component の cusp 比較から
全 $W_i$ 上で $|d\log s_a|_{g_{\infty,i}}\le K$ と取れる．
拡大 compact $\{\max s_a\le2C_0N_j\}$ の少し広い graph gauge を用意し，
最後の $t$ の閾値で $E_{i,t}^*g_t\ge g_{\infty,i}/2$ を課す．先に
\[
 \sqrt2 K C_{P_j}/N_j<\log2
\]
としておけば，いずれかの深度が初めて $2C_0N_j$ に達するまでの
$\log s_a$ の変化は $\log2$ 未満となり矛盾する．Graph gauge は拡大 compact の
近傍に定義してあるので，その像からの最初の退出にも同じ議論を適用できる．
従って normal faces の全 image で深度は $2C_0N_j$ 以下である．

Corner faces と rounding には，この短円長評価を流用せず，corner map とその
逆写像の既存の log-depth 変位 $O(b_j^{-2})$ を使う．Source の二つの独立深度は
$O(N_j)$ 以下なので，actual physical depths も固定倍の $N_j$ 以下である．
Graph gauge と global divisor norms の unit 差を固定段階で吸収すれば，
component depths にも同じ上界が成立する．この議論は
$q_a\in[\delta_j/2,2\delta_j]$ の平行 levels と小さい collars にも適用できる．
全ての $N_j$ の下限は固定 residual compact の定数を決めた後に課し，その後で
拡大 compact の graph convergence に必要な $t$ の閾値を取っている．

境界上の評価を domain 全体へ移すため，成分極限の固定 strictly plurisubharmonic
potential を用いる：
\[
 \Psi_i=\sum_{a\ne i}w_as_a-2\sum_{a\ne i}\log s_a+h+v .
\]
$h,v$ は有界であり，$w_a>0$ なので，定数 $c,C_0,C_1>0$ に対して
\[
 c\sum_{a\ne i}s_a-C_0\le\Psi_i
 \le C_1\sum_{a\ne i}s_a+C_0
\]
が成立する．$E_{i,t}^{-1}(K_{i,j,t})$ は $W_i$ 内の compact domain である．
Strictly plurisubharmonic function はその内部で最大値を取れないから，
\[
 \sup_{E_{i,t}^{-1}(K_{i,j,t})}\Psi_i
 \le\sup_{E_{i,t}^{-1}(\partial K_{i,j,t})}\Psi_i
 \le C_2N_j .
\]
上の下界から domain 全体で $s_a\le C N_j$ を得る．
これは Liouville trajectories の全長や，primitive 補正後の potential の存在を
要求しない最大原理の議論である．

Triple chart の non-dominant defining functions と global divisor norms の間の
unit は成分極限の固定した chart 上で有界である．さらに，今得た固定 compact
$D_{i,N}=\{\Psi_i\le C_2N_j\}$ を少し広げて graph gauge を取る．
$t$ の最後の閾値を小さくすれば，この compact 上で
$E_{i,t}^*\zeta_a/\zeta_{a,0}$ と $E_{i,t}^*\zeta_c/\zeta_{c,0}$ の絶対値は
例えば $[1/2,2]$ に入る．これは固定 compact 上の graph convergence と，
その上では defining functions が消えないことによる．従って actual physical depths と
component depths の差は固定 chart の差を含めて $O(1)$ となり，
$\ell_a(x),\ell_c(x)$ にも同じ $C N_j$ 上界がある．
さらに $x=C_{t,J}(z)$ に対して
\[
 |\log\ell_q(x)-\log\ell_q(z)|\le Cb_j^{-2}
 \quad(q=a,c)
\]
だから，$\ell_q(z)\le e^{Cb_j^{-2}}\ell_q(x)\le C'N_j$ となり，
\eqref{ke:eq:removed-corner-depth-upper}が従う．
最後に，$i\notin J$ の component は当該三重点を通らない．Triple chart を先に
十分小さく固定すればその閉包は $H_i$ と交わらず，固定 compact の graph image
$K_{i,j,t}$ も，小さい $t$ に対して当該 corner と交わらない．
\end{proof}

\paragraph{この証明が使わない主張．}
切り取る部分は，実際に構成した $K_{i,j,t}\cap W_{t,J}$ である．
それを full corner 全体での不等式 $q_a,q_c\ge\delta_j$ の全解集合と同一視して
いない．その不等式の別成分の不存在を示す必要はなく，dominant-product 評価を
その適用範囲外に延長してもいない．固定 compact への包含だけでも，ある段階依存の
$A_j<\infty$ による深度上界は得られる．上の最大原理はそれを $C N_j$ に整理する．

\begin{proposition}[Component の対数接続費用]
\label{ke:component-logarithmic-connectors}
一つの component $W_i$ とその固定 KE 計量 $g_{\infty,i}$ を取る．
固定した $C_0>0$ に対し，$N\ge2$ と
\[
 D_{i,N}=\{p\in W_i:\max_{a\ne i}s_a(p)\le C_0N\}
\]
の任意の二点を，長さ $C(1+\log N)$ 以下の経路で結べる．
全経路は $\{\Psi_i\le AN\}$ に選べ，$C,A$ は固定した族・charts・weights と
$C_0$ のみに依存する．従って $N=N_j$ を先に固定してから $t$ を小さくすれば，
actual graph image 内でも同じ評価が成立する．特に，実際に構成した一つの
component の core と変更 collar の全体の ambient 径と接続費用は
\begin{equation}\label{ke:eq:component-logarithmic-budget}
 C\bigl(1+\log(N_j+1)\bigr)
\end{equation}
以下である．経路を core 自身又は一つの fiber の内部に制限する主張ではない．
\end{proposition}

\begin{proof}
固定した有限個の nested SNC coordinate polydiscs で divisor を覆う．
一つの chart では，divisor の局所方程式を $z_1\cdots z_k=0$，$k=1,2$ とし，
$\sigma_a=-\log|z_a|^2$ と置く．標準 cusp 比較
\cite[\S5, (5.3)--(5.6), Theorem~5.1]{RochonZhang}により，固定した定数で
\[
 g_{\infty,i}\le C\left(
   \sum_{a=1}^k\frac{|dz_a|^2}{|z_a|^2\sigma_a^2}
   +\sum_{a=k+1}^2|dz_a|^2\right).
\]
Chart を予め小さく取れば，局所深度と対応する global depth の差は有界であり，
他の divisors の defining functions は units である．
内側の polydisc にある点から，各 normal coordinate の角度と接方向を固定し，
その絶対値を外側 polydisc 内の固定した小半径まで一つずつ増やす．
この経路の長さは
\[
 C\sum_{a=1}^k\int_{\sigma_*}^{\max\{\sigma_a,\sigma_*\}}\frac{du}{u}
 \le C'(1+\log N)
\]
であり，途中の全 global depths も $C'N$ 以下に留まる．
各 chart の出口集合は divisor から離れた固定 compact である．
$W_i$ は連結なので，有限個の出口集合と charts 外の固定 compact を含む
一つの連結な compact domain $K_*\Subset W_i$ を取り，その中で固定基点へ
一様に有界な長さで結べる．二点をこの基点経由で結べば主張の長さ評価を得る．
深度と $\Psi_i$ の比較から，固定した $A$ で全経路は
$\{\Psi_i\le AN\}\Subset\{\Psi_i<2AN\}$ に含まれる．

$N=N_j$ を固定した後，この拡大 compact 上の fixed-gauge smooth convergence により
$E_{i,t}^*g_t\le2g_{\infty,i}$ を課す．小さい exact Moser 移送の変位も
最後の $t$ の閾値で $1$ 以下にできる．
補題\ref{ke:core-depth-upper}は actual core と平行 moment collars に適用され，
葉族の補間はその collar の点集合を変えない．従ってそれらの全点に上の経路を適用できる．
有限個の components について定数の最大を取れば，$C$ は stage に依存しない．
\end{proof}

\begin{lemma}[除いた corner 部分は旧 target の成分頂点へ近づく]
\label{ke:removed-corner-target-cap}
補題\ref{ke:core-depth-upper}の設定で $L\ge C' N_j$ とする．
定数を大きく取って
\[
 \eta_{j,t}=\frac{\sqrt2\log(C' N_j)}{R_t}
\]
と置けば，全ての $i\in J$ について
\[
 P^{\rm old}_{t,J}(K_{i,j,t}\cap W_{t,J})
 \subset B_{d_\square}(e_i,\eta_{j,t}).
\]
特に，既存の slow diagonal で $\log N_{j(t)}=o(R_t)$ としたとき，
\[
 \max_{i\in J}\operatorname{diam}_{d_\square}
 P^{\rm old}_{t,J}(K_{i,j,t}\cap W_{t,J})
 \le2\eta_{j,t}\longrightarrow0 .
\]
従って各成分 cap の旧 target 径が一様に零へ収束する．異なる成分 caps の
和集合の径が零へ収束するという主張ではない．
\end{lemma}

\begin{proof}
$z=C_{t,J}^{-1}(x)$，$J=\{i,a,c\}$，$r_q=\ell_q(z)-m_j$ とする．
\eqref{ke:eq:removed-corner-depth-upper}と $\sum\ell_q=L$ から，$L/N_j$ を大きく
すれば $r_i=\max_qr_q\ge L/2$ である．Softmax の厳密な逆式は
\[
 w_q=1+\frac{\log(r_q/r_i)}R .
\]
従って
\[
 w_i=1,\qquad
 0\le w_a,w_c\le\frac{\log(2C N_j)}R .
\]
これらの点は $e_i$ と同じ unit square に入る．その square の線分の長さを使えば
$d_\square(w,e_i)\le\sqrt{w_a^2+w_c^2}\le\eta_{j,t}$ となる．
\end{proof}

\subsubsection{保持した corner で維持される距離評価}

旧 corner の保持部分を
\[
 A_{t,J}=W_{t,J}\setminus\bigcup_{i\in J}\operatorname{int}K_{i,j,t}
\]
とし，旧 projection の restriction を $P_A$ と書く．
旧 closed-corner theorem の誤差を
\[
 E_{j,t}=C\left(\frac1{m_j}
       +\frac{\log(m_j+1)+\log R_t}{R_t}\right)
\]
と置く．$d_W^{\rm int}$ は\emph{旧 full corner} 内の距離，
$d^{\rm out}$ は旧 protected domain 内の距離であり，ともに $R_t^{-2}g_t$ で測る．
これらの幾何学的な領域や距離は，その一部分の fiber partition を core filling へ
置き換えても存在し続ける．したがって $x,y\in A_{t,J}$ に対して
\begin{equation}\label{ke:eq:retained-protected-estimates}
 d_\square(P_Ax,P_Ay)-E_{j,t}
 \le d^{\rm out}(x,y)
 \le d_W^{\rm int}(x,y)
 \le d_\square(P_Ax,P_Ay)+E_{j,t}
\end{equation}
がそのまま成立する．ここでは $d_J=d_\square|_{\mathcal C_J}$ を用いた．
また，actual 全曲面の距離には
\[
 R_t^{-1}d_{g_t}(x,y)
 \le d_\square(P_Ax,P_Ay)+E_{j,t}
\]
という上界が維持される．旧 full-corner lift が除いた cap を通っても，それは
依然として $X_t$ 内の許された経路だからである．

$A_{t,J}$ \emph{自身}の新しい内在距離については，下界
\[
 d_A^{\rm int}(x,y)\ge d_\square(P_Ax,P_Ay)-E_{j,t}
\]
だけが，経路集合の包含から直ちに従う．その上界には，除いた cap を横切る旧 lift を
保持境界上の短い経路へ置き換える追加の detour 評価が必要である．
固定 compact の中を通る接続路の評価だけでは，その経路が $A_{t,J}$ の中に留まると
まではいえない．ここではこの新しい内在上界を主張しない．
後続の定理\ref{ke:global-ambient-upper}は旧 full corner 内の経路を使うため，
この保持領域の内在 detour 評価を必要としない．
一方，$R_t^{-1}d_{g_t}$ の下界には protected domain の外を通る全経路の制御が必要であり，
定理\ref{ke:global-ambient-lower}で別に証明する．

Projection の全射性の損失も小さい．補題\ref{ke:removed-corner-target-cap}より，
$P_A(A_{t,J})$ に含まれない点は $e_i$ の $\eta_{j,t}$ 近傍のいずれかに属する．
各 $i\in J$ の link の corner rounding は，二つの moments が
$\delta_j+O(\varepsilon)$ となる点を含み，その source は旧 full softmax domain の
内部にある．従って
$x_i\in S_{i,j,t}\cap W_{t,J}\subset A_{t,J}$ を選べ，
$P_A(x_i)$ も $e_i$ の同じ近傍に属する．失われた任意の target 点は，対応する
$P_A(x_i)$ から $2\eta_{j,t}$ 以下にある．
従って $P_A(A_{t,J})$ は $\mathcal C_J$ の $2\eta_{j,t}$-net である．
特に\eqref{ke:eq:retained-protected-estimates}から，
\[
 d_{\rm GH}\bigl((A_{t,J},d^{\rm out}|_{A_{t,J}}),
                  (\mathcal C_J,d_\square)\bigr)
 \le \tfrac12E_{j,t}+2\eta_{j,t}
\]
が得られる．実際，$d_\square(P_A(x),w)\le2\eta_{j,t}$ を満たす組
$(x,w)$ を全て取った correspondence の distortion は
$E_{j,t}+4\eta_{j,t}$ 以下である．左辺の source metric は protected distance の制限であり，
$A_{t,J}$ の新しい内在距離でも，全曲面の距離でもない．

\paragraph{旧 softmax 値と全域基底値の比較．}
以上は旧 softmax 値 $P^{\rm old}_{t,J}$ に対する主張である．
全域 quotient の位相が $B_d(T)$ であることと，保持 corner の全 torus fibers が同じで
あることだけから，特定の parameterization に対して
$p_T\circ\pi_{j,t}=P_A$，または両者の差が $o(1)$ であるとは結論できない．
Core と変更 collar の actual 径が小さいことも，任意に再 parameter 化した target
でのその像の径を小さくすることは保証しない．

命題\ref{ke:global-base-parameterization}では，各 labelled cell の parameterization
を選び，core 側の cone sectors を上の小さい vertex caps へ対応させ，neck の共通
collars でも値を連続に一致させる．保持 corner の旧値との一致は $O(\eta_{j,t})$ の
近似であり，旧 softmax 辺の regular tori を branch locus に移す厳密一致は要求しない．
木を潰す $p_T$ との整合には，命題\ref{ke:collapsed-tree-common-moment}の
全枝共通 moment を用いる．本補題の cap の位置はこれらの構成への入力であり，
全域の基底値比較や定理\ref{ke:global-ambient-lower}の no-shortcut 評価を
単独で与えるものではない．

% !TeX root = ../../surface-lagrangian-torus-fibration.tex
% 2026-11-22 / v1.2: 共通 normal moment に相対的な fiber partition の変更と基底値の安定性．
\section{旧 softmax 値に適合する全域基底の parameter 化}
\label{ke:global-base-compatibility-section}

定理\ref{ke:actual-global-fibration}で得た actual fiber partition の quotient を
$Q_{j,t}$ とし，まだ $B_d(T)$ との同相写像を選ばずに
$\pi^0_{j,t}:X_t\to Q_{j,t}$ と書く．
以下の操作はこの fiber partition を変えず，quotient の parameter 化だけを選ぶ．
十分大きい固定 $j$ に対して
\[
 m=6b_j,\qquad N=N_j,\qquad \delta=4\pi/N,
 \qquad L=-\log|t|^2,\quad R=\log L
\]
とする．各 component の normal face の anchor level を同じ $q_a=\delta$ に取り，
相対葉族の補間支持を $q_a\in[\delta/2,2\delta]$ の中に取る．
Core filling を少し外まで延ばした場合にも，基底の分割にはこの anchor level を使える．
以下の $K_{i,j,t}$ と $A_{t,J}$ はこの共通 anchor で切った core と保持集合を指す．
実際の接着に平行 link を使った場合は，その間の共通 collar で旧 corner tori が
保たれるので，基底の分割をこの anchor へ取り直す．変更領域は同じ固定 component
compact 内にあり，既存の深度上界と費用評価は定数を増やすだけで保たれる．
補題\ref{ke:anchor-trimmed-neck-domain}により，この分割後の保持 neck の
moment 範囲は正確に $[-\delta,\delta]$ である．その外側の旧 neck 尾部は
対応する component core に入り，平行 link との差は変更 collar に含める．
既存の選択順序の中で $N$ を大きくし，最後に $t$ を小さくすることを許す．

三重点 $J$ の旧 corner と旧 projection を $W_{t,J}$，$P^{\rm old}_{t,J}$ とし，
その保持部分を $A_{t,J}$ とする．その quotient の旧 target 像を
\[
 D_{t,J}:=P^{\rm old}_{t,J}(A_{t,J})\subset\mathcal C_J
\]
と書く．補題\ref{ke:removed-corner-target-cap}の誤差より大きい
\begin{equation}\label{ke:eq:base-compatibility-small-scale}
 \eta=\eta_{j,t}=A\frac{1+\log(N+1)}R<\frac1{100}
\end{equation}
を取り，$A$ は固定した次数・族・chart のみに依存する十分大きい定数とする．

\subsection{全 ends に共通な normal longitude の定量的構成}

Target の二 label 部分を長さ $2$ の区間として
\[
 E_{ij}(u)=
 \begin{cases}
  e_i+u e_j,&0\le u\le1,\\
  (2-u)e_i+e_j,&1\le u\le2
 \end{cases}
\]
と表す．Normal moment の向きは，component $i$ の face が $\mu=-\delta$，
component $j$ の face が $\mu=\delta$ となるように一度だけ固定する．
この規約では generator は triple lattice の $(1,-1,0)$ に対応する．

\begin{lemma}[旧 corner 辺での共通 model profile]
\label{ke:common-edge-moment-profile}
同じ pair $ij$ に接する各三重点 $J=\{i,j,k\}$ について，旧 softmax 辺
$w_k=0$ のうち $\mu\in[-\delta,\delta]$ の部分は，$\mu$ で一対一に
parameter 化される．その旧辺 longitude を $u_{J,t}(\mu)$ と書く．
Complete simplex model のみによる一つの strictly increasing profile
$u_{0,t}(\mu)$ があり，全 ends に一様に
\begin{equation}\label{ke:eq:common-edge-profile-error}
 \sup_{|\mu|\le\delta}|u_{J,t}(\mu)-u_{0,t}(\mu)|
 \le\frac C{mR}.
\end{equation}
その endpoints は
\[
 u_{J,t}(-\delta)=\frac{\log N+O(1)}R,
 \qquad 2-u_{J,t}(\delta)=\frac{\log N+O(1)}R
\]
を満たし，model profile についても同様である．
特に $u_{0,t}(0)=1$ である．
\end{lemma}

\begin{proof}
$r_a=\ell_a-m$，$S=L-3m$ と置く．旧 softmax 辺は
$r_k=\max(r_i,r_j)/L$ である．その全体を
$v=\log(r_j/r_i)\in[-R,R]$ により
\begin{equation}\label{ke:eq:hexagon-edge-profile}
 r_i=\frac{S}{1+e^v+\max(1,e^v)/L},\quad
 r_j=e^v r_i,\quad r_k=\frac{\max(r_i,r_j)}L
\end{equation}
と明示できる．旧辺 longitude は正確に $u=1+v/R$ である．
$B=\ell_k=m+O(1)$ と書くと，$v=0$ の両側で
\[
 \frac{dB}{d\ell_i}=
 \begin{cases}1/L,&r_i\ge r_j,\\-1/(L+1),&r_j\ge r_i\end{cases}
\]
が成立する．以下の微分は $v=0$ で片側微分として読み，そこで連続につなぐ．

Normal generator を $V$，$\mathfrak a=|V|_h$ とすると，normal 深度
$a=\ell_i,c=\ell_j\gg m$ 上で
\[
 \mathfrak a\asymp a^{-1}+c^{-1},\qquad
 |\nabla_h^rV|_h\le C_r\mathfrak a
\]
である．$\partial_v$ の normal 成分の長さは上下から正定数で抑えられる．
その residual 成分の長さは $C/m$ 以下である．実際，
$|da/dv|\asymp\mathfrak a^{-1}$ と上の $dB/da$ から
\[
 \frac{|dB/dv|}B\le\frac C{\mathfrak a BL}\le\frac C m .
\]
Model moment を $\nu_0$，実際の全平均 moment を $\nu_J$ とする．
Tame 下界と
$d\nu(V')=-\kappa(V,V')$ により，向きを上のように取ると
\begin{equation}\label{ke:eq:edge-moment-weighted-derivative}
 c\mathfrak a\le\frac{d\nu_0}{dv}\le C\mathfrak a,
 \qquad
 c\mathfrak a\le\frac{d\nu_J}{dv}\le C\mathfrak a
\end{equation}
となる．Residual 成分の寄与は normal 成分の正の主項に対し $O(m^{-1})$ なので，
大きい $j$ で吸収できる．$\log\mathfrak a$ の $v$ 微分も有界である．

ここで加法定数を無視してはいけない．
補題\ref{ke:component-meridian-normalization}および
系\ref{ke:transition-normal-primitive}によって，この $\nu_J$ は全 ends で
同じ canonical normal moment $\mu$ である．
命題\ref{ke:relative-corner-patch}の証明で用いた full-triple comparison と
相対座標接続の primitive 評価，および系\ref{ke:relative-corner-patch-c2}から
\[
 \|G^*\lambda_{\log,t}-\lambda_{0,L}\|_{C^2(h)}\le C/m.
\]
Model primitive は全平均で変わらないので，generator の重みを保持して評価すると
\begin{equation}\label{ke:eq:edge-moment-weighted-error}
 |\nu_J-\nu_0|+
 \left|\frac d{dv}(\nu_J-\nu_0)\right|
 \le\frac C m\,\mathfrak a .
\end{equation}
全三重点で同じ $L,m$ と同じ complete simplex model を使うため，$\nu_0$ は
$k$ に依存しない．非正則な補間帯でも source primitive は $G^*\lambda$ である．

Dominant-product moment 評価により，$\mu=\pm\delta$ の model levels は
片方の normal 深度が $N$ と固定比の間，他方が $L-O(N+m)$ の領域にある．
$N\gg n_j\gg m$ とし，その後 $L/N$ を大きく選べば，両端とその固定相対幅の
buffer は全て上の estimates の範囲に入る．単調性により，その間の
$\mu\in[-\delta,\delta]$ の levels も同じ normal 下界を持つ．

$\nu_0(v_0)=\mu$ とする．$v_0\pm K/m$ において
\eqref{ke:eq:edge-moment-weighted-derivative}を積分し，
$\mathfrak a$ の比の有界性と\eqref{ke:eq:edge-moment-weighted-error}を用いる．
$K$ を固定して十分大きくすると，$\nu_J(v_0-K/m)<\mu<\nu_J(v_0+K/m)$ となる．
従って唯一の actual solution $v_J(\mu)$ は $|v_J-v_0|\le C/m$ を満たす．
これを $R$ で割れば\eqref{ke:eq:common-edge-profile-error}である．
端では小さい normal 深度が $N(1+O(m^{-1})+o_{t\to0}(1))$，大きい方が
$L-O(N+m)$ なので，endpoint の式も従う．
Model の $i,j$ 交換対称性により $\nu_0(0)=0$ であり，最後の主張を得る．
\end{proof}

\begin{corollary}[共通 endpoints を持つ全枝共通の profile]
\label{ke:common-neck-longitude}
$\kappa=\eta$ とする．前補題の profiles の endpoints を，target の
$4\eta$-vertex neighborhoods 内だけの increasing piecewise smooth rescaling に
よって揃えられる．その結果，
\[
 \sigma_t:[-\delta,\delta]\longrightarrow[\kappa,2-\kappa],\qquad
 \widetilde u_{J,t}:[-\delta,\delta]\longrightarrow[\kappa,2-\kappa]
\]
という increasing homeomorphisms を取り，
\[
 \sigma_t(0)=1,\qquad
 \sup_\mu|\widetilde u_{J,t}(\mu)-\sigma_t(\mu)|\le C/(mR)
\]
とできる．さらに $\mu\in[-\delta,-\delta/2]$ では両 profile は $C\eta$ 以下，
$\mu\in[\delta/2,\delta]$ では $2$ との差が $C\eta$ 以下である．
\end{corollary}

\begin{proof}
$\alpha_-=u(-\delta)$，$\alpha_+=2-u(\delta)$ と置く．
$[0,2]$ の increasing piecewise linear homeomorphism $\psi$ を，
\[
 0\mapsto0,\quad\alpha_-\mapsto\kappa,\quad
 4\eta\mapsto4\eta,\quad2-4\eta\mapsto2-4\eta,
 \quad2-\alpha_+\mapsto2-\kappa,\quad2\mapsto2
\]
で定め，中央区間では恒等とする．$\alpha_\pm\asymp\log N/R$，
$\eta=A(1+\log(N+1))/R$ だから，slopes と endpoints に関するこの式の
Lipschitz constants は一様に抑えられる．
Model と各 end にこの規則を適用し，$\sigma_t=\psi_0\circ u_{0,t}$，
$\widetilde u_{J,t}=\psi_J\circ u_{J,t}$ とする．
前補題の endpoint 誤差を含めて評価すれば表示された誤差が従う．
$\mu\in[-\delta,-\delta/2]$ の独立深度は $O(N)$ なので，元の longitude は
$O((1+\log N)/R)$ である．Rescaling の支持も同じ尺度だから最後の主張を得る．
\end{proof}

\subsection{実際の cap 境界と保持 corner の disk 構造}

\begin{lemma}[Moment cut の境界は一本の proper arc である]
\label{ke:corner-cap-disk}
上の $N$ と $t$ の選択をさらに強くすれば，各 $J=\{i,a,c\}$ について
core $i$ の quotient 境界と旧 corner の交わりは，$e_i$ に接する二つの旧境界辺を
結ぶ一本の proper embedded arc である．除かれる部分はその arc と $e_i$ 側の
二つの境界辺で囲まれた disk であり，$e_i$ を含む．
異なる三つの core の caps は互いに素である．従って $D_{t,J}$ は閉 disk であり，
その境界は三本の cap arcs と三本の旧境界 arcs が交互に並ぶ．

さらに $D_{t,J}$ から $\mathcal C_J$ への homeomorphism $r_{t,J}$ を，
旧 target からの変位 $C\eta$ 以下となるように取れる．
この写像は各 vertex の $[0,4\eta]^2$ の外で恒等であり，cap arc を
その vertex に接する二つの境界辺の長さ $\kappa$ の部分へ写す．
旧境界 arc での longitude は前系の $\psi_J$ による rescaling と一致する．
\end{lemma}

\begin{proof}
小さい target cap に含まれることだけでは，前半の位相的主張は従わない．
ここでは補題\ref{ke:closed-component-link}の境界の明示式を使う．
Core $i$ 側の source で独立深度を $a,c$ とすると，boundary の二つの normal
faces は $q_a=\delta$，$q_c=\delta$ であり，両者は monotone smooth minimum
$F(q_a-\delta,q_c-\delta)=0$ によって接続される．
Joint moment map の derivative の対称部分は負定値であるため，
$\partial_aq_a<0$，$\det Dq>0$ である．
従って $q_a=\delta$ 上では
\[
 \frac{dq_c}{dc}\bigg|_{q_a=\delta}
 =\frac{\det Dq}{\partial_aq_a}<0.
\]
Normal level は $a\asymp N$ の唯一の graph であり，$c$ が旧 boundary の
$m+O(1)$ から $2N$ へ増える間に，$q_c$ は $\delta$ より十分大きい値から
小さい値へ移る．従って二つの faces は唯一の joint-moment 近傍で接続する．
Smooth minimum の zero set はその moment coordinates で一本の滑らかな arc であり，
余分な component は生じない．

旧 boundary が $6.2b_j$ より浅い部分では，共通 plateau 上の neck face を用いる．
そこで同じ normal moment cut と厳密に一致するので，上の arc の続きを与える．
旧 softmax 境界との交点の一意性は，補題\ref{ke:common-edge-moment-profile}の
strict monotonicity による．その交点はそれぞれ $\mu=\pm\delta$ の指定された端で
ある．Fixed compact core の深度上界により，当該 corner の交点は component $i$ の
dominant chart に含まれる．この chart では他の boundary divisor の defining
functions は非零 unit であり，その $N$-level faces は入ってこない．
従って当該 chart 内での boundary は，いま列挙した二つの normal faces と一つの
rounding で尽くされる．これが余分な boundary component を排除する理由である．
旧 $e_i$ の source では二つの独立深度は $m+O(1)\ll N$ であり，これは
link の内側にある．従って除かれる側は $e_i$ を含む disk である．
三つの caps は\ref{ke:removed-corner-target-cap}の互いに素な vertex neighborhoods
に含まれるので，残る $D_{t,J}$ は表示された六つの boundary arcs を持つ disk である．

変位を制御した homeomorphism を具体的に作る．
$e_i$ を原点とする square の二座標を用い，$U_i=[0,4\eta]^2$ とする．
$A$ を大きく選んだので cap は $[0,\eta/2]^2$ の中にある．
$D_{t,J}\cap U_i$ は一つの corner を上の arc で切った閉 disk である．
その境界上で，外側の二辺は固定し，旧 coordinate-axis の二部分は
前系の $\psi_J$ によって parameter 化する．残る cap arc を，
$(\kappa,0)$ から原点を通り $(0,\kappa)$ へ至る二本の axis segments へ
orientation を保って写す．これで boundary circles の homeomorphism を得る．

ここで disk extension を仮定として残す必要はない．両方の curvilinear disks を
有限個の boundary arcs に適合する閉 disk 座標で表し，boundary circle の写像を
$e^{i\theta}\mapsto e^{i\varphi(\theta)}$ と lift する．
$\varphi(\theta+2\pi)=\varphi(\theta)+2\pi$ で狭義単調である．
\[
 (r,e^{i\theta})\longmapsto
 \left(r,e^{i[\theta+\chi(r)(\varphi(\theta)-\theta)]}\right)
\]
を用いれば，$\chi=0$ near $0$，$\chi=1$ near $1$ として boundary map を disk
全体へ延長できる．各 $r$ で angular map は狭義単調だから homeomorphism である．
ここで用いる disk 構造は上で確認した一本の proper arc による切断から得られるもの
であり，未知の global extension property ではない．有限の piecewise smooth
refinement を使えば regular strata 上の smoothness も保てる．

この extension の source と image はともに $U_i$ 内にあるので，変位は
$4\sqrt2\eta$ 以下である．三つの $U_i$ は互いに素で，外側の二辺を固定したため
その外の恒等写像と貼り合わさる．これが $r_{t,J}$ である．
\end{proof}

\subsection{全域 parameter 化の構成と旧 corner 値の近似}

\begin{proposition}[旧 softmax 値に近い全域基底値]
\label{ke:global-base-parameterization}
固定した marked tree system $T$ と map $p_T:B_d(T)\to\mathcal C_d$ を用いる．
上の $j,N,t$ の選択の下で，labels と core vertices を保つ homeomorphism
\[
 h_{j,t}:Q_{j,t}\longrightarrow B_d(T)
\]
を取り，$\pi_{j,t}=h_{j,t}\pi^0_{j,t}$，$f_{j,t}=p_T\pi_{j,t}$ とすると，次を
同時に満たすようにできる．
\begin{enumerate}
\item 全ての保持 corner について
\begin{equation}\label{ke:eq:global-base-old-corner-error}
 \sup_{x\in A_{t,J}}d_\square
 \bigl(f_{j,t}(x),P^{\rm old}_{t,J}(x)\bigr)\le C\eta_{j,t}.
\end{equation}
\item Component $i$ の core および相対葉族を変更した collar の像は
\[
 f_{j,t}(K_{i,j,t}\cup\text{変更 collar})
 \subset B_{d_\square}(e_i,C\eta_{j,t})
\]
に含まれる．
\item 三重点に向かう各 neck leaf arm の外側部分は，その end の旧 softmax
辺と共通の normal moment を用いて parameter 化される．その longitudinal
coordinate と，全枝共通の $\sigma_t(\mu)$ の差は $C\eta_{j,t}$ 以下である．
\end{enumerate}
Fiber partition，actual Lagrangian 性，連結性，source・branch・core の fiber 型は
変わらない．Parameter 化を有限の piecewise smooth strata に適合させられる．
\end{proposition}

\begin{proof}
\emph{Target triangle の明示的な inset と panels．}
各 $J$ で $\rho_J(w)=w/\sum_{a\in J}w_a$ とし，$0<\varepsilon=\eta<1/100$ に対して
\[
 H_s(w)=\rho_J^{-1}\bigl((1-3s\varepsilon)\rho_J(w)
                   +s\varepsilon(1,1,1)\bigr),\qquad0\le s\le1
\]
と置く．$H_1$ は $\mathcal C_J$ をその内部の閉 disk へ写す embedding であり，
$\sup_wd_J(H_s(w),w)\le C\varepsilon$ である．
この評価は表示された path の piecewise smooth length を積分すれば従う．

外周と inset disk の間の closed annulus は，正確に
$H_s(w)$，$w\in\partial\mathcal C_J$，$0\le s\le1$ で parameter 化される．
実際，simplex 座標 $\lambda$ では
$s=\min_a\lambda_a/\varepsilon$ であり，元の boundary point は
$(\lambda-s\varepsilon\mathbf1)/(1-3s\varepsilon)$ で回復できる．
従って annulus の存在やその trivialization を追加仮定していない．

Outer boundary を六つの arcs に分ける．Vertex $e_i$ の arc $I_i$ は，隣接する
二つの辺上で $e_i$ からの長さが $\kappa$ 以下の部分の和とする．
残る三つの edge arcs は $I_{ij}=E_{ij}([\kappa,2-\kappa])$ である．
これらの arcs と $s\in[0,1]$ の積の像を target panels とする．
Vertex panel $H([0,1]\times I_i)$ は
$B_{d_J}(e_i,C(\kappa+\varepsilon))$ に含まれる．

\emph{保持 corner 全体を中央 disk として使う．}
$Q_{j,t}$ 内の $A_{t,J}$ の quotient は $D_{t,J}$ と同一視できる．
この\emph{全体}に対して
\[
 g_{J,t}=H_1\circ r_{t,J}
\]
を用いる．補題\ref{ke:corner-cap-disk}により，これは inset disk への
homeomorphism であり，旧 $P^{\rm old}$ との差は $C\eta$ 以下である．
保持 corner をさらに狭い $B=6.5b_j$ の帯などへ切り直してはいないので，
\eqref{ke:eq:global-base-old-corner-error}は全保持点で成立する．

\emph{Actual quotient の残りの pieces．}
一つの triangle cell $\tau_J$ から $D_{t,J}$ の内部を除いた領域は，三つの core
sectors と三つの neck leaf-arm rectangles に分かれる．
Core sector は core cone の link graph における該当 edge の cone である．
その境界 link edge は，二つの branch vertices からの neck tails と，中央の cap arc
からなる．Leaf-arm rectangle は，pair $ij$ の marked tree の葉 $k$ に向かう arm
を，branch vertex と旧 softmax 辺の間で切ったものに対応する．

この分割を actual quotient 内で行える理由を確認する．
旧辺 $w_k=0$ は $B=m+O(1)$ の graph で，normal 深度は
$\mu\in[-\delta,\delta]$ に一様に $cN$ 以上である．
$N\gg n_j$ なので，この graph の全 fibers は K2 の共通 plateau overlap 内にある．
そこで normal moment は同じ $\mu$ で，reduced 葉は順序を保った同じ全 tori である．
この marked Reeb continuation は，既存の diffeomorphism 族から得られる．
補題\ref{ke:fixed-marked-tree-exhaustion}の $D_{j,\mu,t}\circ S_t$ と，
定理\ref{ke:saddle-moving-relative-interpolation}の $h_\mu^{\chi(\mu)}$
（各補間 band の外では恒等）を合成し，固定した基準葉族
$\mathcal F^0$ を構成後の $\mathcal F_\mu$ へ送る族 $\Xi_{\mu,t}$ を取る．
系\ref{ke:actual-saddle-moving-neck-interpolation}はこの葉族の全逆像を使っている．
従って商上の写像
\[
 (\mu,[p]_{\mathcal F_\mu})
 \longmapsto(\mu,[\Xi_{\mu,t}^{-1}(p)]_{\mathcal F^0})
\]
は $\mu$ と labels を保つ graph の自明化である．
各 leaf arm 上で branch vertex から boundary graph までの正の長さを
再規格化すれば，arm は $[-\delta,\delta]\times[0,1]$ と parameter 化できる．
最後の座標 $r$ は branch で $0$，旧 corner boundary で $1$ となる．
各 $\mu$ で端の長さを再 parameter 化するだけであり，foliation の変更ではない．
Source がこの rectangle 内にあっても quotient は同じ product であり，対応する
一つの base value の fiber が pinched torus になるだけである．

\emph{Leaf-arm rectangles の具体的な写像．}
該当する rectangle に対して
\begin{equation}\label{ke:eq:leaf-arm-base-interpolation}
 U(\mu,r)=(1-r)\sigma_t(\mu)+r\widetilde u_{J,t}(\mu),\qquad
 F_J(\mu,r)=H_r(E_{ij}(U(\mu,r)))
\end{equation}
と定める．各 $r$ で $U$ は $[-\delta,\delta]$ から
$[\kappa,2-\kappa]$ への increasing homeomorphism である．
従ってこれは rectangle から該当 edge panel への homeomorphism である．
$r=1$ では $g_{J,t}$ と一致し，$r=0$ では $E_{ij}(\sigma_t(\mu))$ となる．
$\mu=\pm\delta$ では $U$ がそれぞれ $\kappa,2-\kappa$ に等しいため，
core sector との二本の共通 sides も既に指定される．
全 ends の longitude の共通選択と endpoint の一致は前系が保証する．

\emph{Core sectors の extension．}
Component $i$ の core cone の radial coordinate を $r_i\in[0,1]$ とする．
各 branch ray 上では，$i$ から出発する target edge の longitude を
$\kappa r_i$ と指定する．Core sector の残りの境界では，二つの neck tails に
\eqref{ke:eq:leaf-arm-base-interpolation}の $\mu=\pm\delta$ での値を，cap arc に
$g_{J,t}$ の値を用いる．これらは core sector の boundary circle を vertex panel
の boundary circle へ orientation を保って一対一に写す．
補題\ref{ke:corner-cap-disk}の証明に記した angular lift の式を用いて，この境界写像
を disk 全体へ延長する．従って core sector 全体の像は該当 vertex panel に入り，
その全点は $e_i$ から $C\eta$ 以下である．

\emph{内部木と bigons．}
Pair $ij$ に属する全 branch edges と bigons の和集合は
\[
 \bigl(T_{ij}^{\rm int}\times[0,1]\bigr)/
 \bigl(T_{ij}^{\rm int}\times\{0\}\sim v_i,
       T_{ij}^{\rm int}\times\{1\}\sim v_j\bigr)
\]
である．Actual quotient にも，中央の $\mu\in[-\delta,\delta]$ の product と，
両端の core cones による同じ表示がある．
中央の全内部木で longitude を共通の $\sigma_t(\mu)$ とし，両 core cones では
それぞれ $\kappa r_i$，$2-\kappa r_j$ とする．内部木の座標は変えない．
$p_T$ の定義で使う simplex-edge coordinate は
\[
 \Theta(u)=\bigl(\rho_d(E_{ij}(u))\bigr)_j
 =\begin{cases}u/(1+u),&u\le1,\\1/(3-u),&u\ge1\end{cases}
\]
だから，target の上の表示の第二因子には $\Theta(u)$ を用いる．
これにより全 bigons と branch edges の homeomorphism を得る．
その branch 上の値は，隣接する全 triangle panels で指定した値と同じである．

\emph{全域への貼合せと $p_T$ との関係．}
各 triangle では，上で構成した写像は全 $\mathcal C_J$ への homeomorphism である．
その後に $(p_T|_{\tau_J})^{-1}$ を合成し，bigons と branch edges では直前の
product-coordinate map を用いる．全共通 boundaries で一致し，core vertices を
正しく同定しているので，まず有限個の飽和閉片上で
$X_t\to B_d(T)$ を連続に貼り合わせることができる．
これは全 fiber 上で一定なので，quotient の普遍性により
連続写像 $h_{j,t}:Q_{j,t}\to B_d(T)$ を誘導する．
各 labelled cell の内部で一対一かつ全射であり，boundary の identifications も同じ
だから，この写像は一対一の連続全射である．Compact から Hausdorff への写像なので
homeomorphism である．これで $f_{j,t}=p_Th_{j,t}\pi^0_{j,t}$ が，構成中の値と
厳密に同じであることが確認できた．
系\ref{topo:cubical-composite-connected}により，この同じ $f_{j,t}$ の全 fiber は連結である．
この結論は後の距離評価又は GH 収束を入力にしていない．

Regular strata の座標には有限の piecewise smooth refinement を用いる．
これは quotient の座標だけの変更であり，全 actual fibers とその isotropic strata
を保存する．必要なら angular extension の cutoff を中心近傍で零として，
各 open two-dimensional piece 上で smooth な座標を選べる．

最後に quantitative assertions を確認する．保持 corner の誤差は
$r_{t,J}$ と $H_1$ の変位を加えれば $C\eta$ 以下である．
Core の triangle sectors は vertex panels に，bigon sectors と branch rays は
$e_i$ から長さ $\kappa$ 以下の edge parts に入る．
変更 collar の外側部分では，normal moment が $\delta$ と固定比の間にあるため，
前系の最後の評価から $\sigma_t$ と $\widetilde u_{J,t}$ は該当する endpoint の
$C\eta$ 近傍にある．\eqref{ke:eq:leaf-arm-base-interpolation}もその近傍を保ち，
$H_r$ による追加変位は $C\eta$ 以下である．Corner の変更 collar は既存の target
cap estimate と保持 corner の比較で扱える．これで全ての主張を得る．
\end{proof}

\subsection{内部木の共通 moment と neck 全体の値}

\begin{proposition}[潰される内部木上では共通 normal moment を保てる]
\label{ke:collapsed-tree-common-moment}
命題\ref{ke:global-base-parameterization}の parameter 化を用いる．
$u\in(\kappa,2-\kappa)$ とし，$\mu=\sigma_t^{-1}(u)$ と置く．
このとき
\[
 f_{j,t}^{-1}(E_{ij}(u))
\]
は，pair $ij$ の neck において，同じ normal moment $\mu$ 上で
$T_{ij}^{\rm int}$ 全体に属する構成後の neck fibers を合併した集合に等しい．
相対補間の支持外では，これらは旧 K2 fibers そのものである．
Core に近い $u\in[0,\kappa]$ 又は $u\in[2-\kappa,2]$ の部分は，それぞれ
component $i$ 又は $j$ の固定 compact core と変更 collar に含まれる．
前者の厳密な同一 moment の主張は，corner 値の $O(\eta)$ 近似とは別の性質である．
\end{proposition}

\begin{proof}
$p_T$ の $E_{ij}(u)$ 上の逆像は，定義通り一つの $T_{ij}^{\rm int}$ である．
Triangle の内部は三 label の support を持つ target 内部へ homeomorphically 写るので，
そこから追加の fiber は来ない．中央の book の全内部木では longitude を
$\sigma_t(\mu)$ と共通に選んだため，その逆像は同じ $\mu$ の slice に厳密に一致する．
両端の longitude ranges には各 component の core cones を対応させたので，残りの
主張も従う．変更 collar の境界を別の平行 level に取った場合は，例外をその
component の固定 compact collar に拡大すれば同じ結論を使える．
\end{proof}

\begin{corollary}[旧 neck と source critical values の target control]
\label{ke:neck-base-value-control}
保持された pair $ij$ の neck 上で，normal moment を $\mu$ とすると，
\[
 d_\square\bigl(f_{j,t}(x),E_{ij}(\sigma_t(\mu))\bigr)\le C\eta_{j,t}.
\]
Core collars の外にある内部木の部分では左辺は零である．
特に全 source critical values の像は
$B_{d_\square}(e_i+e_j,C\eta_{j,t})$ に含まれる．
\end{corollary}

\begin{proof}
内部木では前命題を使う．旧 corner の外の leaf-arm rectangle では
\eqref{ke:eq:leaf-arm-base-interpolation}と
$|\widetilde u_{J,t}-\sigma_t|\le C/(mR)$ を使えばよい．
残る old neck--corner overlap では $B\in[m+O(1),7b_j]$ である．
固定 $\mu$ の source level を $a=a(B)$ と書くと，weighted normal estimates より
\[
 |\partial_B\nu|\le C\mathfrak a/B,
 \qquad -\partial_a\nu\ge c\mathfrak a^2.
\]
従って $|da/d\log B|\le C/\mathfrak a$ であり，
\[
 \left|\frac{d}{d\log B}\log\frac{r_j}{r_i}\right|\le C.
\]
この $B$-interval の ratio は有界なので，旧 boundary $w_k=0$ と同じ $\mu$ で比較
した longitude の差は $C/R$ 以下である．また当該 overlap の第三座標は
\[
 0\le w_k=1+R^{-1}\log(r_k/r_{\max})
 \le C(1+\log b_j)/R\le C\eta_{j,t}.
\]
第三座標を零にする線分と，辺に沿う longitude の差を合わせ，保持 corner 上の
\eqref{ke:eq:global-base-old-corner-error}を用いれば求める評価を得る．
Source の critical moment は $\mu=0$ であり，$\sigma_t(0)=1$，
$E_{ij}(1)=e_i+e_j$ だから最後の主張が従う．
\end{proof}

\paragraph{厳密一致の障害と，得られた結論の範囲．}
旧 softmax 辺 $w_k=0$ の torus は actual quotient では triangle の regular
interior にあり，対応する cellular boundary は branch locus である．
$p_T$ は triangle interior を target triangle interior へ写すので，保持 corner
全体で $p_T\pi=P^{\rm old}$ とすることは一般に不可能である．
上の inset はこの障害を $O(\eta)$ の誤差で解消する．

既存の slow diagonal によって $\log N_{j(t)}=o(R_t)$ とすれば
$\eta_{j(t),t}\to0$ である．従って保持 corner の旧 target 値，component の小さい
target caps，および潰される内部木の共通 moment は，全域の一つの
$f_t=p_T\pi_t$ に対して同時に実現される．
本節は target の parameter 化を構成したものであり，normal levels 又はその全木の
actual ambient 径，全曲面経路に対する no-shortcut 下界を証明したものではない．

\subsection{Normal moment を保つ相対変更の安定性}
\label{ke:relative-replacement-section}

作用 atlas のために reduced 葉や core filling を選び直す場合，
それらを結ぶ diffeomorphism の微分を actual KE 計量で小さく抑える必要はない．
以下は，変更後の局所構成とその全 fiber 接着が別途得られた場合に，
同じ coarse target 上の値を制御する命題である．
この命題自体は，critical actions を指定する局所構成の存在を与えない．

\begin{proposition}[共通 moment に相対的な変更の基底値比較]
\label{ke:moment-relative-replacement}
命題\ref{ke:global-base-parameterization}の許容される固定 $j,t$ と，
同じ anchor cuts，保持 corner $A_{t,J}$，component core と変更 collar，
保持 neck の有限閉被覆を用いる．この被覆上に別の局所 actual Lagrangian
fibration データを与え，次を仮定する．
\begin{enumerate}
\item 各 $A_{t,J}$ と全 K2 交わりを含む開 end collars では，
元の全 fiber partition を保つ．Source の compact nodal neighborhoods も保つ．
\item Pair $ij$ の保持 neck では，同じ actual normal circle，
向きと加法定数を固定した同じ canonical moment $\mu$，
および同じ anchor levels $\mu=\pm\delta$ を用いる．
Reduced 葉の商には，$d-2$ 個の labelled ends を持つ許容される三価の
固定 pants tree $\widehat T_{ij}$ による
実際の marked product 自明化があり，$\mu$ と全 end labels を保つ．
各 leaf arm は branch から旧 corner boundary まで
$[-\delta,\delta]\times[0,1]$ と parameter 化され，
source values を含めて連続で，正則 strata 上で piecewise smooth とする．
木の抽象的な同型だけをこの自明化の代わりには用いない．
\item 各 component の変更は同じ core と変更 collar の内部に限り，
局所商は $\widehat T=(\widehat T_{ij})$ の labelled link の cone である．
その boundary の neck tails，branch rays，および保持 corner の cap arcs は
上の marked product と一致する．特に cap arcs と anchor cuts は変更しない．
\item 各局所商は compact Hausdorff で，全交わりは双方の全 fibers の和であり，
その partition は一致する．局所 fibers は連結で，正則部分の Lagrangian 性と
特異部分の等方的 strata は各局所構成で既に確認されているものとする．
\end{enumerate}
すると，これらのデータの閉商 $\widehat Q_{j,t}$ と，
labels を保つ homeomorphism
$\widehat h_{j,t}:\widehat Q_{j,t}\to B_d(\widehat T)$ が得られる．
対応する全域 fibration を $\widehat\pi_{j,t}$，
$\widehat f_{j,t}=p_{\widehat T}\widehat\pi_{j,t}$ とすると，
parameter 化を
\begin{equation}\label{ke:eq:relative-replacement-value}
 \widehat f_{j,t}=f_{j,t}\quad\hbox{on every }A_{t,J},\qquad
 \sup_{x\in X_t}d_\square
       \bigl(\widehat f_{j,t}(x),f_{j,t}(x)\bigr)
 \le C\eta_{j,t}
\end{equation}
となるように選べる．$C$ は固定した次数・族と target の有限幾何に依存してよいが，
reduced 葉を変更する自明化や diffeomorphism の微分には依存しない．
\end{proposition}

\begin{proof}
まず補題\ref{glue:closed-whole-fiber-quotient}を適用する．
交わる local fibers は同じ全 ambient subset なので，local fiber relations の
有限和は既に推移的で，閉同値関係になる．従って compact Hausdorff な
$\widehat Q_{j,t}$ を，後の $B_d(\widehat T)$ との同定に依存せず得る．
各 global fiber は local fiber そのものであり，連結性と actual Lagrangian 性，
特異 strata の条件も保持される．

保持 corner の商には，命題\ref{ke:global-base-parameterization}で用いた
\emph{同じ} $g_{J,t}=H_1r_{t,J}$ を用いる．Cap arcs を変えていないので，
domain $D_{t,J}$ とこの写像はそのまま使える．
新しい leaf-arm rectangle の座標を $(\mu,\widehat r)$ とすると，
\[
 \widehat U(\mu,\widehat r)
 =(1-\widehat r)\sigma_t(\mu)+\widehat r\widetilde u_{J,t}(\mu),
 \qquad
 \widehat F_J(\mu,\widehat r)
 =H_{\widehat r}\bigl(E_{ij}(\widehat U(\mu,\widehat r))\bigr)
\]
と置く．各 $\widehat r$ で $\widehat U$ は同じ interval の increasing
homeomorphism である．従ってこれは同じ target panel への homeomorphism で，
$\widehat r=1$ では保持 corner，$\widehat r=0$ では
$E_{ij}(\sigma_t(\mu))$ と一致する．
ここで新しい $\widehat r$ と古い $r$ が近いことは要求しない．

内部木全体では，同じ $\sigma_t(\mu)$ と simplex-edge coordinate $\Theta$
を使う．Core cones の branch rays は $\kappa r_i$ または $2-\kappa r_i$
へ送り，triangle sectors は既に指定した boundary map を vertex panels へ
延長する．Cap arcs と全 end labels の整合性により，
命題\ref{ke:global-base-parameterization}の disk extension をそのまま使える．
新しい内部木の組合せ型が元のものと異なっても，対応する bigons と branch
edges に同じ product-coordinate construction を用いればよい．

これらの写像は有限閉 pieces 上で一致し，その全 fibers 上で一定なので，
$X_t$ 上の連続写像を先に貼り合わせ，閉商の普遍性で
$\widehat h_{j,t}$ を得る．各 labelled cell の内部で全単射であり，
boundary の identifications も一致するため，
compact-to-Hausdorff の判定により同相である．
新たな atlas が与えられている場合には，この同相でその atlas を移すだけである．
Parameter 化から critical-action の直線性を作り出す操作ではない．

値の比較は全被覆の三種類の領域ごとに行う．保持 corner では二つの写像は同じである．
一つの component の core と変更 collar では，両像とも同じ $e_i$ の
$C\eta_{j,t}$ 近傍にある．残る pair $ij$ の neck では，元の写像に
系\ref{ke:neck-base-value-control}を使う．新しい写像にも，表示した
$\widehat F_J$ と $|\widetilde u_{J,t}-\sigma_t|\le C/(mR)$ から
\[
 d_\square\bigl(\widehat f_{j,t}(x),
                    E_{ij}(\sigma_t(\mu(x)))\bigr)\le C\eta_{j,t}
\]
が成り立つ．内部木ではこの距離は零である．共通の同じ moment を使うため，
三角不等式で\eqref{ke:eq:relative-replacement-value}を得る．
この最後の評価には新旧の graph 座標の差も，その微分も入らない．
\end{proof}

\begin{corollary}[相対変更後の同じ射影による GH 実現]
\label{ke:relative-replacement-gh}
定理\ref{ke:global-metric-realization}の遅い対角選択に沿って，
前命題の局所データが一つの固定 $\widehat T$ に対して得られたとする．
すると同じ actual KE 距離について
\[
 \operatorname{dis}(\widehat f_t)
 \le\operatorname{dis}(f_t)+C\eta_{j(t),t}\longrightarrow0.
\]
各 $\widehat f_t$-fiber，すなわち一つの $p_{\widehat T}$-fiber 上の
全 $\widehat\pi_t$-fibers の合併について，
\[
 \sup_{z\in\mathcal C_d}
 \operatorname{diam}_{R_t^{-1}d_{g_t}}\widehat f_t^{-1}(z)
 \le\operatorname{dis}(\widehat f_t)\longrightarrow0
\]
である．従って個々の $\widehat\pi_t$-fibers の ambient 直径も一様に零へ向かう．
特に $d=5,6$ で前命題の局所データが作用適合 atlas を持つよう構成できれば，
その atlas と GH 実現は同じ $\widehat\pi_t$ によって達成される．
\end{corollary}

\begin{proof}
各二点の target 距離の差は\eqref{ke:eq:relative-replacement-value}の二倍以下である．
これと $\eta_{j(t),t}\to0$ を使い，定数 $C$ を増やす．
$\widehat f_t$ は全射であり，同じ target 値を持つ二点の正規化距離は
$\operatorname{dis}(\widehat f_t)$ 以下だから，fiber の主張も従う．
Atlas の存在は前命題の入力であって，この距離比較からの結論ではない．
\end{proof}

\begin{remark}[Normal moment の絶対誤差だけでは足りない]
\label{ke:raw-moment-error-insufficient}
共通の canonical moment を保つ条件を，単に
$|\widehat\mu-\mu|=o(1)$ へ弱めることはできない．
実際，固定 stage と $0<\alpha<\beta<1$ に対し，
$\mu_t=-L_t^{-\alpha}$，$\nu_t=-L_t^{-\beta}$ と取る．
十分小さい $t$ で両者は保持 moment interval に入り，
補題\ref{ke:common-profile-log-depth}より
\[
 |\mu_t-\nu_t|\longrightarrow0,\qquad
 \sigma_t(\mu_t)\longrightarrow\alpha,\qquad
 \sigma_t(\nu_t)\longrightarrow\beta.
\]
従って対応する辺上の距離は $\beta-\alpha>0$ へ収束する．
近似的な matching を使うなら，必要なのは
$|\widehat\sigma_t(\widehat\mu)-\sigma_t(\mu)|$ の一様な制御である．
この例は profile の誤差尺度を区別するものであり，
actual atlas fibration の不存在を示す反例ではない．
\end{remark}

\input{sections/paper/ke_sharp_normal.tex}
% !TeX root = ../../surface-lagrangian-torus-fibration.tex
% 2026-11-22 / v1.3: 固定 stage の有限 co-Gram 評価と最後の対角極限の分離．
\section{切り詰めた対数深度による全曲面の距離下界}
\label{ke:global-scalar-lower-section}

本節では $L=-\log|t|^2$，$R=\log L$，$\bar g_t=R^{-2}g_t$ とする．
Actual 全域 fibration の存在と component cap の評価は
定理\ref{ke:actual-global-fibration}および
補題\ref{ke:removed-corner-target-cap}を入力とする．
以下では clipped map を構成し，triple 外の単一 scalar 評価へ還元する．
その scalar 評価は節\ref{ke:sharp-normal-section}で証明した
normal Schwarz barrier によって供給される．さらに actual fibration の
基底値との全域比較を行い，同じ写像に対する全曲面の距離下界を得る．
外部 KE 研究の $\mathrm{(CP)}_d$，$\mathrm{(PBE)}_d$ を仮定しない．
以下の定数 $C_d$ は固定した次数と退化族に依存してよく，
$C_j$ はさらに固定 stage に依存してよい．

\subsection{全曲面上の clipped logarithmic map}

$\mathcal O_{\mathbf P^3}(1)$ に固定 Hermitian metric を取り，
各 $L_i$ を定数倍して
\[
 \ell_i=-\log\|L_i\|^2>1
\]
とする．$L_i=0$ では $\ell_i=+\infty$ と解釈する．
定数倍の積は defining equation の $F_d$ に吸収する．
$b_j$ を K2 の residual exhaustion scale，$m_{\rm old}=6b_j$ を旧 corner cut とする．
これとは独立に，固定 stage の全定数と core depth $N_j$ を決めた後で
lower clipping depth $m=\mathfrak m_j\gg m_{\rm old}+N_j+n_j+\mathfrak n_j$ と
plateau 定数 $K=K_j>2d$ を選ぶ．
いずれも stage とともに無限大へ発散させ，最後の slow diagonal に
$\log\mathfrak m_j+\log K_j=o(R_t)$ を課す．以下では $m=\mathfrak m_j$，$K=K_j$ と略記し，
\[
 M=\frac{L}{K},\qquad S=\log(M/m)=R-\log K-\log m,
 \qquad
 \Theta_i=\left[\frac{\log(\ell_i/m)}S\right]_0^1
\]
と置く．$[q]_0^1=\max(0,\min(q,1))$ である．

\begin{lemma}[Source 軸も含む大域的な cubical map]
\label{ke:global-clipped-depth-map}
十分小さい $|t|$ に対し
\[
 \Theta_t=(\Theta_1,\ldots,\Theta_d):X_t\longrightarrow\mathcal C_d
\]
は well-defined な continuous map で，固定 $t$ ごとに locally Lipschitz である．
少なくとも一つの coordinate は $1$ に等しく，非零 coordinate は高々三つである．
従って微分が非零となり得る active coordinates
\[
 A_t(x)=\{i:m<\ell_i(x)<M\}
\]
は高々二つである．Source 軸上でも $L_i=0$ の coordinate は plateau $1$ に入り，
その対数の発散による不連続性はない．
\end{lemma}

\begin{proof}
族の積方程式と誘導 Hermitian metric により，零でない点では
\[
 \sum_i\ell_i=L-\log\|F_d\|^2\ge L-C.
\]
従って $\max_i\ell_i\ge(L-C)/d>M$ である．
いずれかの $L_i$ が零ならこの結論は自明である．
一方，一般位置の任意の四平面には共通点がない．
各四元集合 $I$ について $\sum_{i\in I}\|L_i\|^2$ は
$\mathbf P^3$ 上で正の最小値を持つので，四番目に大きい $\ell_i$ は一様に有界である．
有限個の $I$ に共通の定数を越えて $m$ を取れば，$\ell_i>m$ は高々三つとなる．
Clip 関数は局所 Lipschitz であり，$L_i=0$ の開近傍では $\Theta_i=1$ である．
これで全主張が従う．
\end{proof}

\begin{lemma}[Co-Gram による全経路の長さ評価]
\label{ke:global-clipped-cogram-lower}
一つの $t$ と clipping parameters $m,K$ を固定し，$0<S\le R$ とする．
ある $\epsilon\ge0$ に対し，各非空集合 $I\subset\{1,\ldots,d\}$，
$|I|\le2$ の開集合
\[
 U_I=\{x\in X_t:m<\ell_i(x)<M\ \text{for all }i\in I\}
\]
上で，該当する covectors の co-Gram matrix がほとんど至る所で
\begin{equation}\label{ke:eq:global-active-upper}
 \left(\langle d\log\ell_i,d\log\ell_j\rangle_{g_t^{-1}}
       \right)_{i,j\in I}\le(1+\epsilon)\mathrm{Id}
\end{equation}
を満たすとする．$A_t(x)\ne\varnothing$ なら $I=A_t(x)$ と取ることで
全 active covectors を制御する．
このとき $X_t$ の任意の rectifiable curve $\gamma$ について
\[
 L_\square(\Theta_t\circ\gamma)
 \le\frac{R\sqrt{1+\epsilon}}S\,L_{\bar g_t}(\gamma).
\]
特に全点対で
\begin{equation}\label{ke:eq:global-raw-lower}
 R^{-1}d_{g_t}(x,y)
 \ge d_\square(\Theta_t x,\Theta_t y)
 -C_d\left(\epsilon+\frac{\log m+\log K}{R}\right).
\end{equation}
さらに $f_t=p_T\pi_t$ と
$\sup_xd_\square(f_t(x),\Theta_t(x))\le a_t$ なら，
右辺を $d_\square(f_t(x),f_t(y))$ に置き換えて誤差 $2a_t$ を加えられる．
これらは固定 parameter の不等式であり，$\epsilon$ の上限，
$\epsilon\to0$，又は $S/R\to1$ は仮定しない．
\end{lemma}

\begin{proof}
各 $U_I$ で co-Gram は smooth なので，ほとんど至る所の不等式は
連続性によりその全点へ延びる．他の coordinates が clipping level にある点も
$U_I$ から除いていないため，strict に active な残りの covectors に同じ評価が成立する．
従って ambient measure zero の集合に曲線が含まれる場合を除外していない．
Active region では $d\Theta_i=S^{-1}d\log\ell_i$ である．
Clip levels に沿って動く区間ではその coordinate の微分は零であり，
一般の absolutely continuous curve でもその level set 上でほとんど至る所同じことが成り立つ．
従って\eqref{ke:eq:global-active-upper}は
\[
 |(\Theta_t\circ\gamma)'|_{\ell^2}
 \le S^{-1}\sqrt{1+\epsilon}|\gamma'|_{g_t}
\]
を与える．有限 Euclidean cubical complex の curve の長さは，各 face の
Euclidean speed の積分である．これは relative interiors の逆像で計算し，
seams 上ではその法線 coordinate の微分がほとんど至る所零になることを用いればよい．
従って square seams や clipping levels の往復を含めて同じ積分評価が成り立つ．
全経路の長さについて下限を取ると，正確な有限評価
\[
 R^{-1}d_{g_t}(x,y)\ge q\,d_\square(\Theta_t x,\Theta_t y),
 \qquad q=\frac{S}{R\sqrt{1+\epsilon}}\in(0,1]
\]
を得る．$\eta=(\log m+\log K)/R=1-S/R\in[0,1)$ と置けば，
$\epsilon\ge0$ に対する $1-(1+\epsilon)^{-1/2}\le\epsilon/2$ より
\[
 1-q
 =\eta+(1-\eta)\{1-(1+\epsilon)^{-1/2}\}
 \le\eta+\epsilon/2.
\]
$D_d=\operatorname{diam}(\mathcal C_d,d_\square)<\infty$ を用いて
$(1-q)d_\square\le D_d(\eta+\epsilon/2)$ とすれば
\eqref{ke:eq:global-raw-lower}が従う．
最後に端点の三角不等式から
$d_\square(\Theta_t x,\Theta_t y)\ge d_\square(f_t(x),f_t(y))-2a_t$ であり，
$q\le1$ と合わせれば同じ誤差に $2a_t$ を加えるだけでよい．
領域間の往復回数を数えず，加法誤差も経路の各部分へ配分していない．
\end{proof}

\subsection{二つの active directions は既存 triple 比較で制御できる}

\begin{proposition}[大域下界の rank-one 評価への還元]
\label{ke:global-lower-rank-one-reduction}
$\mathcal T_{j,t}$ を，三つの異なる global depths が $b_j/2$ 以上となる領域とする．
この領域上では，既存 full-triple comparison と rank-one Schwarz estimate により
\eqref{ke:eq:global-active-upper}が
\[
 \epsilon_{j,t}=C(b_j^{-1}+K_j^{-2})
\]
で成立する．$X_t\setminus\mathcal T_{j,t}$ では active coordinate は高々一つである．
従って大域長さ下界のために新たに十分な解析入力は，次だけである：
\begin{equation}\label{ke:eq:single-active-normal-input}
 \sup_{\substack{x\notin\mathcal T_{j,t}\\\mathfrak m_j<\ell_i(x)<L/K_j}}
       |d\log\ell_i|_{g_t}(x)\le1+\varepsilon_t,
 \qquad \varepsilon_t\longrightarrow0.
\end{equation}
この入力は source を含む pair 領域の \emph{一方向の上 co-gradient bound} であり，
全 tensor の二側比較，angular derivatives，あるいは portal 上界を要求していない．
\end{proposition}

\begin{proof}
三つの $L_i$ が十分小さい点は，対応する一つの三重点の固定近傍に入る．
そこで exact triple coordinates $\zeta_i\zeta_j\zeta_k=t$ を取り，
$s_a=-\log|\zeta_a|^2$ とする．固定 smooth units と Hermitian factors の差により
\[
 \ell_a=s_a+h_a,
 \qquad |h_a|+|dh_a|_{g_t}\le C.
\]
この微分評価には，当該 triple 内での full-triple tensor 比較と
complete simplex model の $g_L^{\rm sf}\ge c g_{\rm FS}$ を使う．
ここで全曲面上の一様 projective lower bound を新たな入力にしていない．
従って active directions について
\[
 d\log\ell_a=d\log s_a+O_{g_t}(b_j^{-1}).
\]
ここで $|d\log s_a|_{g_t}\le1+O(b_j^{-1})$ は，
関連 KE 主稿 \texttt{simplex:lem:rank-one-schwarz} と
\texttt{coarse:thm:full-triple-tensor}から得る．

二つの active depths は $s_a,s_b\le L/K_j+O(1)$，
残る一つは $L-s_a-s_b$ である．
同稿 \texttt{simplex:prop:dominant-product} 又は
\texttt{simplex:cor:active-cogram} により，二方向の co-Gram の上界は
$\mathrm{Id}+O(b_j^{-1}+K_j^{-2})$ となる．
一方向だけの場合は上述の rank-one upper bound で十分で，
co-Gram が $1$ に近いという両側の結論まで要求する必要はない．
Smooth unit の誤差を加えて所要の上界を得る．

二つの active coordinates があれば，それらと plateau $1$ を持つ別の coordinate
の三つが全て $\mathfrak m_j\gg b_j$ より深い．従って必ず $\mathcal T_{j,t}$ に属する．
その補集合では一方向だけが残り，式\eqref{ke:eq:single-active-normal-input}で十分となる．
\end{proof}

\begin{corollary}[全域 co-Gram と actual 距離に関する Lipschitz 評価]
\label{ke:global-depth-lipschitz}
定理\ref{ke:sharp-global-one-active}の順序でパラメータを選ぶ．
このとき全 $X_t$ で，補題\ref{ke:global-clipped-cogram-lower}の仮定は
\[
 \epsilon_{j,t}\le C\bigl(\varepsilon_j+b_j^{-1}+K_j^{-2}\bigr)
\]
で成立する．特に任意の二点について
\begin{equation}\label{ke:eq:global-depth-lipschitz}
 d_\square(\Theta_t(x),\Theta_t(y))
 \le\frac{\sqrt{1+\epsilon_{j,t}}}{S}\,d_{g_t}(x,y).
\end{equation}
この主張は以下の fibration の値との比較を入力にしていない．
\end{corollary}

\begin{proof}
$\mathcal T_{j,t}$ 内では前命題を用い，外では active coordinate が高々一つ
なので定理\ref{ke:sharp-global-one-active}を用いる．境界では連続性と
principal submatrix によって同じ不等式を得る．
補題\ref{ke:global-clipped-cogram-lower}の全曲線評価で下限を取ればよい．
\end{proof}

\subsection{同じ大域 fibration の基底値との比較}

本項では命題\ref{ke:global-base-parameterization}で選んだ parameter 化を用い，
$f_{j,t}=p_T\pi_{j,t}$ とする．辺の向きを同節に合わせて
\[
 E_{ik}(u)=
 \begin{cases}e_i+u e_k,&0\le u\le1,\\
 (2-u)e_i+e_k,&1\le u\le2
 \end{cases}
\]
とする．Normal generator は $(1,-1)$ であり，component $i$ の側は
$\mu<0$，component $k$ の側は $\mu>0$ である．

\begin{lemma}[共通 profile と reciprocal moment の対数]
\label{ke:common-profile-log-depth}
系\ref{ke:common-neck-longitude}の共通 profile $\sigma_t$ について，
各固定 stage に対して十分大きい有限な $A_j$ を取ると，
$|\mu|\le\delta_j$ で
\[
\begin{array}{ll}
 \displaystyle
 \sigma_t(\mu)=\frac{\log(1/|\mu|)}R
       +O(\eta_{j,t}+R^{-1}),
       & \mu<-A_j/L,\\[4pt]
 \displaystyle
 \sigma_t(\mu)=2-\frac{\log(1/\mu)}R
       +O(\eta_{j,t}+R^{-1}),
       & \mu>A_j/L,\\[4pt]
 \displaystyle
 \sigma_t(\mu)=1+O_j(R^{-1}),
       &|\mu|\le A_j/L,
\end{array}
\]
となる．ここで $\eta_{j,t}=C(1+\log(N_j+1))/R$ であり，
最初の二式の定数は固定族と degree に依存してよい．
\end{lemma}

\begin{proof}
補題\ref{ke:common-edge-moment-profile}の $v$ を用いる．
$\mu\in[-\delta_j,\delta_j]$ に対応する model levels では
両 normal depths が $N_j$ と固定比以上に深く，
$N_j\gg m_{\rm old}$ である．
式\eqref{ke:eq:hexagon-edge-profile}と
式\eqref{ke:eq:edge-moment-weighted-derivative}から
\[
 \frac{d\nu_0}{dv}\asymp
 \frac1a+\frac1c\asymp\frac{1+e^{|v|}}L,
 \qquad \nu_0(0)=0.
\]
従って $|v|\ge1$ では $|\nu_0(v)|\asymp e^{|v|}/L$ であり，
$|\nu_0(v)|\le A_j/L$ では $|v|\le C_{A_j}$ である．
$u_{0,t}=1+v/R$ を用いれば表示式を得る．
Endpoints の rescaling は $O(\eta_{j,t})$ しか値を動かさず，
中央の固定区間では恒等なので最後の式も保たれる．
\end{proof}

この補題は共通 profile の独立な定量的説明である．以下の全点比較には，
normal 深度による中央帯の再分類に代えて，同一 moment の actual 接続を用いる．

\begin{lemma}[同一 moment を持つ旧 corner 辺への参照点]
\label{ke:same-moment-corner-reference}
固定 stage $j$ の各 pair $ik$ について，元の K2 neck の全残余領域，source discs，
および必要な全 end annuli を含む連結な固定 residual compact
$P_{ik,j}\Subset P^+_{ik,j}$ を，component cut $N$ より先に選べる．
Source discs を除いた部分も連結とする．ある有限な $N_{0,j}$ と
$C_j^{\rm neck}$ があり，$N\ge N_{0,j}$，$\delta=4\pi/N$ とし，その後 $t$ を
十分小さくすると次が成立する．

Pair $ik$ の構成後の neck の点 $x$ が actual normal moment
$\widehat\mu(x)=\mu\in[-\delta,\delta]$ を持つとする．任意の incident triple
$J=\{i,k,a\}$ について，旧 softmax 辺 $w_a=0$ 上に同じ moment を持つ全 torus
$\Lambda_{ik,J,t}(\mu)\subset A_{t,J}$ が存在し，その任意の点 $y$ に対して
\begin{equation}\label{ke:eq:same-moment-corner-reference}
 d_{g_t}(x,y)\le C_j^{\rm neck},\qquad
 P^{\rm old}_{t,J}(y)=E_{ik}(u_{J,t}(\mu))
\end{equation}
が成立する．定数 $C_j^{\rm neck}$ は $t,\mu,N$ に依存しない．
全 phase と $\mu=\pm\delta$ を含み，$x$ は source 又は saddle の特異 fiber 上に
あってよい．ここで接続経路は一つの fiber 又は moment level の中に留まる必要はない．
\end{lemma}

\begin{proof}
旧 cut を $m_{\rm old}=6b_j$ とする．
式\eqref{ke:eq:hexagon-edge-profile}より，旧辺の全 source parameter 範囲で
その residual 深度は
\[
 B=m_{\rm old}+r_a,\qquad
 0\le r_a=\max(r_i,r_k)/L<1
\]
である．従って，この全範囲を含む固定 end annuli と，marked tree 及び source
discs を含む元の neck の全残余領域を先に合わせて $P_{ik,j}$ を選べる．
End の moment-dependent な境界 graph の像も，K2 の固定 buffer によって
あらかじめ用意した広い residual domain の内部に入っている．少し広い
$P^+_{ik,j}$ も K2 で用意した residual domain 内に取れる．必要な annuli は
旧 trim $B=7b_j$ より内側にあり，$Z_{a,t}/Z_{a,0}\to1$ と固定 frame の差を
吸収する余裕を持たせる．この選択に $N$ 又は $[-\delta,\delta]$ は使わない．
ここで全旧辺を projected neck chart と同一視しているのではなく，その source
residual 範囲から annuli を先に固定している．

補題\ref{ke:common-edge-moment-profile}によれば，$N$ が十分大きいと
$|\mu|\le\delta$ に対応する旧辺の部分では両 normal 深度が $cN$ 以上となる．
$N\gg n_j$ と固定 frame の全定数を $N_{0,j}$ に含め，最後に $t$ を小さくすれば，
この必要な部分の全 tori は K2 の共通 plateau overlap に入る．そこでは
source residual coordinate が上で選んだ annulus を与える．
同補題の strict monotonicity が $\Lambda_{ik,J,t}(\mu)$ の存在と一意性を与え，
補題\ref{ke:corner-cap-disk}の proper arc の端点が $\mu=\pm\delta$ であることから，
両端を含めてこの torus は保持集合に属する．

系\ref{ke:neck-corner-saturated-union}と定理\ref{ke:staged-neck-averaging}により，
この全 torus 上の actual normal action と canonical moment は neck のものと
厳密に同じである．すなわち同じ normal Moser map $\Phi^N_t$ の前では
$\mu_N$，後では $\widehat\mu=\mu_N\circ(\Phi^N_t)^{-1}$ であり，
corner source では $\mu_\alpha=\mu_N\circ G$ と書けば，
式\eqref{ke:eq:staged-moment-identity}と K2 の flow 自然性により
\[
 \mu_\alpha\circ\psi_t=\mathcal A_2\mu_\alpha=\nu_J,
 \qquad \widehat\mu\circ C_{t,J}=\nu_J
\]
となる．第二段階の前後で同じ関数 $\mu_\alpha$ が点ごとに不変という意味ではなく，
この exact transport によって全平均 moment と actual moment を同定する．
新たな加法定数又は phase の
選択はない．相対葉族の補間も actual normal moment を保つので，構成後の $x$ と
$\Lambda_{ik,J,t}(\mu)$ は同じ actual moment 集合に入る．

固定基底に対する系\ref{ke:source-inclusive-same-moment-collapse}を適用する．
そこで用いる定理\ref{ke:fixed-base-neck-moser}と
命題\ref{ke:fixed-base-moment-connectors}は，固定基底の有限な下限以上の
normal cut $m$ と $m\le\vartheta M_t$ に対し，$m,t$ に独立な定数を持つ．
前者は固定 actual buffer と $C_P/m$ の短円評価，後者は固定した凸性の割線定数
$c_0,C_0,\vartheta,m_0$ による．従って固定 stage でも適用でき，この使用に
同じ stage 内で $m\to\infty$ とする仮定は要らない．必要な normal cut の下限は
K2 の $n_j$ の選択に含め，元の共通 moment interval を固定した後に
$[-\delta,\delta]$ をその内部へ入れる．Source からの全 moment 退出にも
この interval 全体に共通な定数がある．

以上の系は，$P_{ik,j}$ 上の同じ moment を持つ全点の合併に対して
$C+C_{P_{ik,j}}$ の ambient 直径上界を与える．同じ Moser map の両端の短い
trajectories の費用を加えれば，actual 像にも一つの $C_j^{\rm neck}$ を使える．
これらの定数は $N$ より先に固定され，後から moment interval を狭めても変わらない．
全 phase，異なる reduced 葉，source 固定点，saddle 葉も同じ集合評価に含まれる．
旧辺の projection 値は定義通り表示された $E_{ik}(u_{J,t}(\mu))$ である．
\end{proof}

以下では，この有限個の pair に対する参照点の定数と，
系\ref{ke:source-inclusive-same-moment-collapse}の同じ固定基底上の接続定数の
最大を $C_j^{\rm neck}$ とする．これを $N_j$ より先に固定し，component cut の
既存の下限に
\begin{equation}\label{ke:eq:neck-constant-depth-choice}
 C_j^{\rm neck}\le1+\log(N_j+1)
\end{equation}
を加える．これは許容される cut の選択条件であり，$C_j^{\rm neck}$ の
exhaustion に関する増大率を評価した結果ではない．後節の距離上界でも，
この同じ固定基底，定数，および cut 条件を用いる．

\begin{proposition}[Clipped depth と actual 基底値の全域比較]
\label{ke:global-depth-base-comparison}
式\eqref{ke:eq:neck-constant-depth-choice}を満たすよう $N_j$ を選ぶ．
Sharp normal 評価の下側 cut $\mathfrak m_j$ を，
固定 stage の core と変更 collar の非成分深度の全上界，
および保持 neck の残余 compact 上の非 pair 深度の全上界より
大きく取る．
この選択は定理\ref{ke:sharp-global-one-active}のパラメータ順序と両立する．
固定した次数・退化族と予め固定した chart・weights にのみ依存し，
$j,t$ に依存しない定数 $C$ に対して
\begin{equation}\label{ke:eq:global-depth-base-comparison}
 \sup_{x\in X_t}d_\square(\Theta_t(x),f_{j,t}(x))
 \le \frac{C\{1+\log\mathfrak m_j+\log K_j+\log(N_j+1)\}}R
\end{equation}
となる．$\Theta_t$ 自体が Lagrangian fibers 上で一定である必要はない．
\end{proposition}

\begin{proof}
\emph{Corner．}
旧 corner の source softmax は
\[
 s_a=m_{\rm old}+
     (L-3m_{\rm old})\frac{L^{w_a}}{\sum_bL^{w_b}}.
\]
$0\le w_a\le1$，$\max w_a=1$ より，その全域で
$|\log s_a-Rw_a|\le C\log(m_{\rm old}+1)$ である．
Global section norm との差は固定 smooth units なので有界である．
Dominant coordinate は旧写像と $\Theta_t$ の双方で $1$ となり，
他の coordinates は同じ三つの labels に限られる．
従って同じ square 内の線分を用いて，Moser--Hamiltonian 補正前の点 $y$ では
\[
 d_\square(\Theta_t(y),w)
 \le C\frac{\log\mathfrak m_j+\log K_j}R.
\]
ここまでの定数は有限個の固定 triple charts の unit bounds によるもので，
stage に依存しない．旧 corner map $C_{t,J}=G\Upsilon\Phi$ の変位にも
同じ一様性がある．実際，系\ref{ke:neck-corner-saturated-union}の証明では
入れ子領域の cutoffs を
$\log(b/b_j),\log(a/n_j),\log(c/n_j)$ の固定した滑らかな関数で作り，
その全固定階の模型 norm は一様である．相対化は moment interval に制限する前に
行われ，$\iota_j^{-1}$ による微分損失は生じない．
命題\ref{ke:quantitative-order-hamiltonian}と
定理\ref{ke:relative-corner-patch}の小摂動量を $\delta_{\rm av}$ と書くと，
$e=O(b_j^{-1})$，$\delta_{\rm av}=O(b_j^{-1})$ であり，
相対 Hamiltonian flow の軌跡長は模型計量で $O(b_j^{-2})$ となる．
全平均の Moser flow には定理\ref{ke:slow-cut-closed-corner-gh}の
生成場評価 $C/b_j$ を使えば十分である．また
命題\ref{ke:three-arm-coordinate-interpolation}の $G_s$ の軌跡長は
$O(e^{-cb_j}/b_j)$ である．全軌跡は最小深度が $c'b_j$ 以上の三重点域に留まり，
full-triple 比較はそこで actual 計量と模型計量の一様な二側比較を与える．従って
定理\ref{ke:slow-cut-closed-corner-gh}と同じ全 phase の領域で
\begin{equation}\label{ke:eq:corner-uniform-actual-displacement}
 \sup_{z\in\mathcal H_{t,m_{\rm old}}}
 d_{g_t}(z,C_{t,J}(z))\le C/b_j
\end{equation}
と評価でき，この $C$ も $j,t$ に依存しない．
補正前後の点を式\eqref{ke:eq:global-depth-lipschitz}で比較すると，
その追加誤差は $C/(b_jR)$ 以下である．この順序により，
補正の像で raw physical depth の微分を別に評価する必要はない．
命題\ref{ke:global-base-parameterization}を用いれば，
全保持 corner で所要の値比較が従う．

\emph{Core と変更 collar．}
補題\ref{ke:core-depth-upper}と固定 compact の graph convergence により，
component $i$ のこの領域では
\[
 \ell_a\le D_j\quad(a\ne i),\qquad \ell_i\ge L-D_j
\]
となる有限な $D_j$ がある．$\mathfrak m_j>D_j$，
$K_j>2$，その後 $t$ を十分小さく取れば $\Theta_t=e_i$ が厳密に成立する．
同じ基底 parameter 化の cap 評価より
$d_\square(f_{j,t},e_i)\le C\eta_{j,t}$ である．

\emph{同一 moment の参照点による neck の比較．}
共通 anchor interval 内の保持 neck の点 $x$ を取り，
$\mu=\widehat\mu(x)\in[-\delta_j,\delta_j]$ とする．
補題\ref{ke:same-moment-corner-reference}により，任意の incident triple $J$ の
旧 corner 辺の torus から参照点 $y$ を選べる．既証の
式\eqref{ke:eq:global-depth-lipschitz}と
$S=R-\log\mathfrak m_j-\log K_j\ge R/2$，および
式\eqref{ke:eq:neck-constant-depth-choice}により
\[
 d_\square(\Theta_t(x),\Theta_t(y))
 \le\frac{\sqrt{1+\epsilon_{j,t}}}{S}C_j^{\rm neck}
 \le\frac{C\{1+\log(N_j+1)\}}R .
\]
上の corner の議論は $y$ において
\[
 d_\square(\Theta_t(y),P^{\rm old}_{t,J}(y))
 \le C\frac{1+\log\mathfrak m_j+\log K_j}R
\]
を与える．さらに $P^{\rm old}_{t,J}(y)=E_{ik}(u_{J,t}(\mu))$ であり，
系\ref{ke:common-neck-longitude}の endpoint rescaling を含めた比較から
\[
 d_\square(P^{\rm old}_{t,J}(y),E_{ik}(\sigma_t(\mu)))\le C\eta_{j,t}.
\]
系\ref{ke:neck-base-value-control}を同じ actual moment $\mu$ に適用すると，
$d_\square(f_{j,t}(x),E_{ik}(\sigma_t(\mu)))\le C\eta_{j,t}$ である．
この四つの不等式を加えれば，neck 上で所要の比較が従う．Source 固定点，
pinched fiber の全点，branch fiber も同一 moment 集合の評価で扱っており，
それらの physical normal 深度を別に分類していない．

\emph{全被覆と依存の順序．}
命題\ref{ke:global-base-parameterization}と同じ共通 anchor で分割し，
core と変更 collar，保持 corner，anchor interval 内の保持 neck を用いる．
補題\ref{ke:anchor-trimmed-neck-domain}により，旧 neck の
$|\widehat\mu|>\delta_j$ の部分は対応する component の core 側にあり，
平行 link との差と補間支持は変更 collar に含められる．従って
$\sigma_t$ をその定義区間の外へ延長して用いる必要はない．
Corner との全交わりでは K2 の normal action と moment が同じである．
定理\ref{ke:actual-global-fibration}の全域被覆により，これらの領域で全 $X_t$ が
尽くされる．同一 moment の接続は
系\ref{ke:source-inclusive-same-moment-collapse}から直接用いており，
後節の全曲面距離上界又はその connector 予算は入力にしていない．
また $\Theta_t$ の Lipschitz 評価は本命題に先立って証明済みなので，
基底値比較と全域下界の間にも循環はない．
Core の深度上界 $D_j$ は lower cut と最後の $t$ 閾値にだけ用い，値誤差には
残らない．Neck の接続定数は既に $1+\log(N_j+1)$ に吸収したので，
有限個の固定 chart の定数を一つの $C$ にまとめれば結論を得る．
\end{proof}

\begin{theorem}[Actual fibration に対する全曲面の距離下界]
\label{ke:global-ambient-lower}
任意の固定 $d\ge5$ と本稿の一般性条件を満たす族を考える．
定理\ref{ke:actual-global-fibration}と
命題\ref{ke:global-base-parameterization}の同じ写像
$f_{j,t}=p_T\pi_{j,t}$ を用いる．
式\eqref{ke:eq:neck-constant-depth-choice}を含め，
固定 stage のパラメータを上述の順序で選べば，
全 $x,y\in X_t$ に対して
\begin{equation}\label{ke:eq:global-ambient-lower}
\begin{split}
 R^{-1}d_{g_t}(x,y)\ge {}&
 d_\square(f_{j,t}(x),f_{j,t}(y))
 -C_d\bigl(\varepsilon_j+b_j^{-1}+K_j^{-2}\bigr)\\
 &-\frac{C\{1+\log\mathfrak m_j+\log K_j+\log(N_j+1)\}}R .
\end{split}
\end{equation}
ここで $C_d,C$ は固定した次数・族と構成 data にのみ依存し，$j,t$ に依存しない．
従って $j=j(t)\to\infty$ を十分遅く選ぶと右辺の全誤差が消える．
左辺は $X_t$ 全体の任意の経路を許す ambient 距離であり，
protected-corner 距離ではない．
\end{theorem}

\begin{proof}
系\ref{ke:global-depth-lipschitz}で得た co-Gram 上界を
補題\ref{ke:global-clipped-cogram-lower}へ代入し，
命題\ref{ke:global-depth-base-comparison}で両端点を比較する．
全域の長さ評価を先に証明しているので，neck，corner，core 間の
経路の出入りの回数による加法誤差は生じない．
各固定 stage の全定数を決めた後，
\[
 R_t\ge j\{1+\log\mathfrak m_j+\log K_j+\log(N_j+1)\}
\]
を課すまで $t$ 閾値を小さくしてよい．これとは別に，各固定 compact の比較や
構成に必要な既存の全ての閾値も満たすようにする．
$\varepsilon_j\to0$，$b_j,K_j\to\infty$ と合わせれば最後の主張を得る．
この下界から fibration の全射性や距離上界を逆に推論してはいない．
表示された評価は許容された cut の下での誤差予算であり，
$\mathfrak m_j,K_j,N_j$ の有効な増大率，又は固定 compact の収束閾値は評価していない．
従って最適な収束速度を結論するものではない．
\end{proof}

\subsection{任意回の出入りを扱う protected-buffer 原理}

\begin{lemma}[加法誤差を持つ局所下界の正しい大域化]
\label{ke:buffered-local-lower-globalization}
$(X,d_X)$ を length space，$(Y,d_Y)$ を直径 $D$ の metric space とし，
$f:X\to Y$ を continuous map とする．$U_\alpha\subset V_\alpha$ を
$X=\bigcup_\alpha U_\alpha$ なる cover とし，$V_\alpha$ は開とする．
全 $V_\alpha$ 内の曲線 $c$ に対し
\[
 d_Y(f(c(0)),f(c(1)))\le\lambda L_X(c)+\epsilon,
 \qquad \lambda\ge1,
\]
かつ $U_\alpha$ から $X\setminus V_\alpha$ への first-exit curve の長さは
$b>0$ 以上と仮定する．このとき
\[
 d_X(x,y)\ge
 \frac{d_Y(fx,fy)-\epsilon}{\lambda+\epsilon/b}
 \ge d_Y(fx,fy)-D(\lambda-1+\epsilon/b)-\epsilon.
\]
従って局所誤差 $\epsilon\to0$ に加えて，
$\lambda\to1$ と $\epsilon/b\to0$ が十分である．
\end{lemma}

\begin{proof}
有限長の任意の曲線 $\gamma$ の出発点を含む $U_\alpha$ を選び，
最初に $V_\alpha$ を出る時刻までを一つの部分曲線とする．
その時刻の点を含む次の $U_\beta$ を選んで繰り返す．
完了した部分曲線はそれぞれ長さ $b$ 以上なので，その個数は
$L_X(\gamma)/b$ 以下である．無限回の first exits が有限長の中で起きることも
この下界により排除される．最後の一片を含めて局所不等式を加えると
\[
 d_Y(fx,fy)\le(\lambda+\epsilon/b)L_X(\gamma)+\epsilon.
\]
全経路で下限を取ればよい．元の曲線が二領域間を実際に何回往復したかは制限していない．
また $f(U_\alpha)$ と $f(\partial V_\alpha)$ の target separation が $b_0>\epsilon$
なら，局所不等式から $b=(b_0-\epsilon)/\lambda$ を得られる．
\end{proof}

既存の corner protection は，inner cut $m_{\rm old}$ と
outer cut $m_{\rm old}/2$ の差だけなら
正規化した幅が $O(R^{-1})$ である．一方，局所 distortion には
$(\log m_{\rm old}+\log R)/R$ が含まれる．従ってこのままでは
$\epsilon/b\to0$ は導けない．本節では，上記の scalar differential によって
全曲線に対する乗法的な長さ評価を先に得ることで，この問題を避けた．

\subsection{有限個の飽和した全被覆でも，加法誤差だけでは足りない}

\begin{example}[全被覆上の protected 比較が大域下界を与えない例]
\label{ke:full-cover-oscillating-bypass}
$Y=[0,1]\times[-1,1]$ に Euclidean metric を入れる．
$0<c<1$，整数 $N\to\infty$，$h_N=N^{-3}$ とし
\[
 a_\pm(x)=1\pm c\sin(2\pi Nx),\qquad
 A_\pm(x)=\int_0^x a_\pm(r)\,dr
\]
とする．$|A_\pm(x)-x|\le c/(\pi N)$ である．
$X_N$ は同じ rectangle に smooth metric
$g_N=dy^2+a_N(x,y)^2dx^2$ を入れたものとする．
$y\ge h_N$ では $a_N=a_+$，$y\le-h_N$ では $a_N=a_-$，
$|y|\le h_N/2$ では $a_N=1+c$ とし，間を単調に滑らかに補間する．
従って上半分では $a_N\ge a_+$，下半分では $a_N\ge a_-$ である．
$f_N:X_N\to Y$ は恒等写像とする．

$U_\pm=\{\pm y\ge0\}$ は $f_N$-飽和した二つの閉集合による全被覆で，
共通部分は一本の連結な interval である．全 fibers は一点なので径は零である．
上半分では
\[
 g_N\ge dA_+^2+dy^2.
\]
逆に transformed coordinates $(A_+(x),y)$ の線分を $y\ge h_N$ へ押し上げ，
両端に垂直線分を付ければ，追加費用は $2h_N$ 以下である．
従って全点対 $p,q\in U_+$ で
\[
 \big|d_{U_+}^{\rm int}(p,q)-d_Y(p,q)\big|
 \le C/N+2h_N.
\]
下半分も同様である．さらに protected regions
$V_+=\{y>-h_N/4\}$，$V_-=\{y<h_N/4\}$ 内でも，
対応する lower bound は同じ calibration $(A_\pm,y)$ によって成立する．
従って各 $U_\pm$ には誤差 $O(N^{-1})$ の protected distance 比較がある．

しかし $(0,0)$ から $(1,0)$ へ，$\sin(2\pi Nx)$ の符号ごとに
$y=-h_N$ と $y=h_N$ の安い側を使い，符号が変わるたびに垂直移動すると
\[
 d_{X_N}((0,0),(1,0))
 \le\int_0^1\min(a_+,a_-)\,dx+O(Nh_N)
 =1-\frac{2c}{\pi}+o(1).
\]
Target distance は $1$ なので大域下界の誤差は消えない．
この例では full cover，飽和性，connected overlap，zero fiber diameter，
および局所 protected distortion の全てが成立する．
不足するのは誤差を支配する protection であり，$\epsilon_N/b_N\asymp N^2$ である．
Actual KE family の反例ではなく，有限個の局所加法誤差を形式的に有限回だけ
足してよいとする推論への反例である．
\end{example}

\subsection{Scalar Chern--Lu の境界制御を省略できない理由}

局所正則 coordinate $z$ に対し，$s=-\log|z|^2$ と置くと
半平面への rank-one Chern--Lu は
$\Delta\log Q\ge Q-1$ という logistic inequality を与える．
Global section norm では $dd^c\ell_i$ が固定 projective curvature である．
当該 tube 上で $g_t\ge c g_{\rm FS}$ を別途確認できれば，その trace は有界となる．
ただしこの事実だけで\eqref{ke:eq:single-active-normal-input}は従わない．
微分不等式の誤差を $O(1/\ell_i)$ に整理することと，
その領域の全境界から最大原理を適用できることは別の条件である．

この区別は固定した smooth negative KE manifold の smooth divisor 近傍でも見える．
Metric が非退化な点の正則 defining coordinate を $z$ とすると
\[
 |d\log(-\log|z|^2)|_g
 \asymp\frac{1}{|z|\,|\log|z|^2|}\longrightarrow\infty.
\]
$dd^c(-\log|z|^2)=0$ であり，局所 Schwarz の微分不等式そのものも成立するが，
punctured domain の source metric は complete ではない．
最大原理の boundary control がないことが破綻の理由である．
これは本稿の退化族に対する否定例ではない．
節\ref{ke:sharp-normal-section}では，正則 projected normal frame を用いて
$\Delta s=0$ を保ち，source を含む比較領域を
$\mathfrak n_j\le s\le\kappa L$ として，その三種類の境界を全て評価した．
外側 $s=\mathfrak n_j$ の粗い有限上界 $A_j$ と内側 $s=\kappa L$ の上界 $C_j$ は，
\[
 B(s)=1+\varepsilon_j+A_j\mathfrak n_j/s+C_js/(\kappa L)
\]
という明示的な障壁により吸収される．残余境界では
full-triple の係数一の評価を用いる．
最終 active band を
$\mathfrak m_j\gg A_j\mathfrak n_j/\varepsilon_j$，
$s\le L/K_j$，$K_j\gg C_j/(\kappa\varepsilon_j)$ に取ることで
両境界の寄与が小さくなる．$s=\kappa L$ の比較境界と，
写像 $\Theta_t$ の plateau $L/K_j$ は区別している．
この境界評価と rank-two triple 比較を合わせたものが
系\ref{ke:global-depth-lipschitz}であり，局所 logistic inequality
だけから大域 Schwarz bound を推論したのではない．

% !TeX root = ../../surface-lagrangian-torus-fibration.tex
% 2026-11-21 / v1.5: 固定 anchor tori による portal の選択と全境界の整合．
\section{全 fiber 合併の収縮と actual 全曲面距離の上界}
\label{ke:global-ambient-upper-section}

命題\ref{ke:global-base-parameterization}で構成した
$f_{j,t}=p_T\pi_{j,t}:X_t\to\mathcal C_d$ を用いる．
本節の全ての source 側の距離は，\emph{全曲面の経路を許す actual ambient 距離}
$d_{g_t}$ である．保持 corner 自身の内在距離は用いない．
十分大きい固定 stage $j$ の後に $|t|$ を小さくし，記号を
\[
 \bar d_t=R_t^{-1}d_{g_t},\qquad m=m_j=6b_j,\qquad
 \delta=4\pi/N_j,\qquad
 \eta=\eta_{j,t}=A\frac{1+\log(N_j+1)}{R_t},\qquad\kappa=\eta
\]
とする．旧 corner theorem の誤差を
\begin{equation}\label{ke:eq:global-upper-corner-error}
 E=E_{j,t}=C\left(\frac1{m_j}
       +\frac{\log(m_j+1)+\log R_t}{R_t}\right)
\end{equation}
と書く．以下では $C$ は固定した次数・退化族・有限 chart 系のみに依存し，
有限定数 $C_j$ は stage $j$ に依存してよい．
固定した $m_j\ge m_*$ に対して，$R_t\ge R_*$ と
$\log(m_j+1)+4\log R_t\le R_t/8$ を最後の $t$ 閾値に課し，
定理\ref{ke:slow-cut-closed-corner-gh}の有限 cut 評価を用いる．

\subsection{同一 moment の接続費用と全 fiber 合併}

各 pair $ab$ について，補題\ref{ke:same-moment-corner-reference}で $N_j$ より先に
選んだ同じ固定残余基底 $P_{ab,j}\Subset P^+_{ab,j}$ を用いる．これは全 marked tree，
source discs，および全 labelled ends の旧 softmax 辺の残余像を含み，source discs を
除いても連結である．旧辺の残余深度は $B=m_j+r_c$，$0\le r_c<1$ なので，
選択は $[-\delta,\delta]$ に依存しない．相対葉族補間もこの固定基底の境界 collar
を保つため，上界のために後から基底を拡大し直さない．

接続定数 $C_j^{\rm neck}$ も下界節と同じものを用い，
式\eqref{ke:eq:neck-constant-depth-choice}を課してある．これは $N_j$ より先に固定できる．
成分側に残る短い neck 尾部にも，次の深度評価がある．

\begin{lemma}[Component 側の neck 尾部の深度]
\label{ke:component-neck-tail-depth}
$N_j$ の下限を大きくした後に $t$ を小さくすれば，共通 anchor interval の外側にある
各 neck 尾部は，対応する一つの component graph 上の
$\max_as_a\le CN_j$ に含まれる．$C$ は stage に依存しない．
従って尾部と同じ component の core・変更 collar・link の任意の二点の接続費用は
$C(1+\log(N_j+1))$ 以下である．
\end{lemma}

\begin{proof}
同じ normal Moser map を $\phi_N$ とし，component に向く正の moments を
\[
 q_{\rm raw}=\pm\mu_N,\qquad
 q_{\rm act}=\pm\widehat\mu_N=q_{\rm raw}\circ\phi_N^{-1}
\]
と書く．尾部は $q_{\rm act}\ge\delta=4\pi/N_j$ の側にあり，その裸の逆像は
$q_{\rm raw}\ge\delta$ を満たす．この transport は source 中心でも滑らかで，
source 軸の点ごとの保存を要求しない．固定点は moment が零なので尾部に属さない．
補題\ref{ke:component-meridian-normalization}，\ref{ke:closed-component-link}の
fixed-end 評価を，固定した $P^+_{ab,j}$ の全 phase に適用する．
裸の normal 深度 $s=2N_j$ では
\[
 q_{\rm raw}=\frac{2\pi}{N_j}+o_j(N_j^{-1})+o_t(1)<\frac{4\pi}{N_j}
\]
となるよう，先に $N_j$，最後に $t$ を選べる．Normal 動径に沿う単調性から，
$q_{\rm raw}\ge\delta$ の component 側の裸の尾部は $s\le2N_j$ に入る．
裸の source 軸でも同じ向きの正則 normal coordinate を使う．

この評価を同じ actual normal Moser map の像へ移す．
さらに広い residual buffer を先に固定し，
補題\ref{ke:fixed-stage-source-neck}の有限下限 $N_P$ を $P^+_{ab,j}$ に適用したものを
$N_{P^+_j}$ と記す．$N_j$ より先に固定深度
$M_j\ge\max\{n_j,N_{P^+_j}\}$ を取り，固定基底の定数について
$C_{P^+_j}/M_j<\epsilon_0$ を課す．この長さは残余方向の固定 buffer の半分よりも
小さくする．$2M_j\le s\le2N_j$ では，反対側の normal depth は
$L_t-s-O_j(1)\ge L_t-2N_j-O_j(1)$ なので，最後に $t$ を小さくして
これも $2M_j$ 以上にする．Source 軸では反対側の深度は無限大である．従って
補題\ref{ke:fixed-stage-stratum-coverage}の flow 制限と ODE 一意性により，
元と同じ centered flow の軌跡長を $C_{P^+_j}/M_j$ 以下にできる．
Fixed frame 差と残余深度 $O(b_j)$ を $N_j$ の下限に吸収すると，
初期点の全 component depths は $3N_j$ 以下である．
固定 cusp 比較により，stage に依存しない $K$ で
$|d\log s_a|_{g_{\infty,i}}\le K$ である．
$\max s_a\le6N_j$ の拡大 graph compact 上で
$E_{i,t}^*g_t\ge g_{\infty,i}/2$ を課し，
$\sqrt2K\epsilon_0<\log2$ とする．
すると first-exit argument により軌跡は深度 $6N_j$ の境界へ達しない．
残余方向の退出も，予め確保した $P_{ab,j}\Subset P^+_{ab,j}$ の buffer で防げる．

浅い部分 $s\le2M_j$ の flow 像は，固定 component 上の flow convergence により
先に一つの固定 compact に入る．その全深度の有限上界 $A_j$ を決め，
$N_j\ge A_j+M_j$ を追加する．従って浅い部分にも同じ $CN_j$ 上界がある．
ここまでの右辺の定数は全て $N_j$ より先に固定されている．
最後に命題\ref{ke:component-logarithmic-connectors}を用いれば，接続費用を得る．
Core の深度上界を neck 尾部へ無条件に延長する議論ではない．
\end{proof}

以上から，全接続の予算はまず
$C(1+\log(N_j+1))+C_j^{\rm neck}$ と取れる．
下界節で選んだ同じ cut 条件\eqref{ke:eq:neck-constant-depth-choice}により
第二項を吸収する．定数 $C$ を一度増やせば
\begin{equation}\label{ke:eq:global-connector-budget}
 H=H_{j,t}:=C\bigl(1+\log(N_j+1)\bigr)
\end{equation}
とする．補題\ref{ke:parallel-liouville-metric-cost}による
core 境界の平行 link への移動は，stage 一様な $O(1)$ の actual 費用を持つ．
従ってこの費用も $H_{j,t}$ に含められ，縮む Liouville 時間幅の逆数を
新たな接続定数として掛ける必要はない．
同じ $H_{j,t}$ により次の費用を同時に抑えられる：同一 pair の同一 actual normal moment 上の
任意の二点の接続，各 component の core・変更 collar・neck 尾部内の任意の二点の接続，
および必要な actual normal Moser trajectories．
第一の場合の二点は，異なる reduced 葉又は source・saddle の特異葉に属してよい．
経路自体が moment level 又は一つの fiber 内に留まるとは主張しない．
第二の場合には一つの component ごとに，少し広い固定連結 compact image 内の
経路を使う．異なる component cores 全体の径を $H$ で抑えるわけではない．

\begin{proposition}[潰される木上の全 fiber 合併を含む一様収縮]
\label{ke:collapsed-tree-full-fiber-collapse}
定数 $C$ を十分大きく選べば，全 $z\in\mathcal C_d$ に対し
\begin{equation}\label{ke:eq:full-fiber-diameter-upper}
 \operatorname{diam}_{d_{g_t}} f_{j,t}^{-1}(z)\le H_{j,t}.
\end{equation}
特に，$d\ge6$ で $p_T$ が潰す内部木上の\emph{全ての}
$\pi_{j,t}$-fibers の合併に同じ評価が成立する．
\end{proposition}

\begin{proof}
$\operatorname{supp}z=\{a,b\}$ とし，$z=E_{ab}(u)$ と書く．
$\kappa<u<2-\kappa$ なら，
命題\ref{ke:collapsed-tree-common-moment}により
$f_{j,t}^{-1}(z)$ は共通の
$\mu=\sigma_t^{-1}(u)$ を持つ構成後の neck fibers の合併である．
相対補間後も normal moment は保たれるので，この合併は
$P_{ab,j}$ 上の同一 moment 集合に含まれる．
従って\eqref{ke:eq:global-connector-budget}で径が抑えられる．
両端の $0<u\le\kappa$ 又は $2-\kappa\le u<2$ では，逆像全体が対応する
一つの component の固定 compact core と変更 collar に含まれ，その接続費用を使う．
境界 $u=\kappa,2-\kappa$ は，core cone の radial level $1$ と neck の
$\mu=-\delta,\delta$ を貼った同じ全 fibers の合併であり，両分類はそこで一致する．
$z=e_a$ の場合も同じ component compact 内にある．

$|\operatorname{supp}z|=3$ の場合，$p_T$ の逆像は一つの triangle cell の一点で
あり，$f_{j,t}^{-1}(z)$ は一つの構成後の $\pi_{j,t}$-fiber である．
定理\ref{ke:actual-global-fibration}の閉 fiber 関係の証明のとおり，交わる局所 fibers は
同じ ambient subset である．従ってこの global fiber は
一つの保持 corner，neck，または component core・変更 collar に含まれる．
保持 corner なら旧 fiber と同じ全 torus であり，
定理\ref{ke:slow-cut-closed-corner-gh}の非正規化 intrinsic diameter $C/m_j$ を使う．
Neck なら同一 moment 集合の部分集合として上の評価を使い，
core・変更 collar なら固定 compact の費用を使う．
$H\ge C/m_j$ となるよう定数を増やせば全ての場合を得る．
\end{proof}

\subsection{共通 moment を持つ corner portals}

三元集合 $J=\{a,b,c\}$ について，旧 softmax 辺 $w_c=0$ 上で
normal moment が $\mu\in[-\delta,\delta]$ である全 torus を
\[
 \Lambda_{ab,J,t}(\mu)\subset A_{t,J}
\]
と書く．補題\ref{ke:common-edge-moment-profile}と
\ref{ke:same-moment-corner-reference}によりこれは存在し，
\[
 P^{\rm old}_{t,J}(\Lambda_{ab,J,t}(\mu))
       =E_{ab}(u_{J,t}(\mu)).
\]
系\ref{ke:common-neck-longitude}の endpoint rescaling と共通 profile の評価から
\begin{equation}\label{ke:eq:portal-profile-error}
 \sup_{|\mu|\le\delta}|u_{J,t}(\mu)-\sigma_t(\mu)|\le C\eta .
\end{equation}
この評価は全 labelled ends に同じ $\sigma_t$ を用いている．
各 phase の選択は任意でよく，その違いも同一 moment の接続費用に含める．

ここで旧辺の保持部分は，端点を含めて正確に
\[
 A_{t,J}\cap(P^{\rm old}_{t,J})^{-1}(E_{ab}([0,2]))
 =\bigcup_{|\mu|\le\delta}\Lambda_{ab,J,t}(\mu)
\]
である．実際，補題\ref{ke:corner-cap-disk}の成分 $a,b$ の proper cap arcs は旧辺と
$\mu=-\delta,\delta$ で一度ずつ交わり，除去側はそれぞれ $e_a,e_b$ を含む．
残る中央閉区間の全 tori は $N_j\gg n_j$ により K2 の共通 plateau 内にあり，
補題\ref{ke:same-moment-corner-reference}の exact transport が同じ canonical moment を与える．
両端の tori はそれぞれ $S_{a,j,t}\cap A_{t,J}$，$S_{b,j,t}\cap A_{t,J}$ に属する．
小さい平行 link で接着した場合にも，補題\ref{ke:anchor-trimmed-neck-domain}の
anchor 再分割と corner collar の全 torus 保存により，この記述は変わらない．

各 $a\in J$ について $b\in J\setminus\{a\}$ を一つ固定し，この辺の $a$ 側の
anchor torus から一点 $p_{a,J}$ を固定する．$a$ から $b$ へ向きを取る記法では
$p_{a,J}\in\Lambda_{ab,J,t}(-\delta)$ である．従って
\[
 p_{a,J}\in S_{a,j,t}\cap A_{t,J},\qquad
 d_\square(P^{\rm old}_{t,J}(p_{a,J}),e_a)\le C\eta .
\]
最後の評価は補題\ref{ke:removed-corner-target-cap}による．

\begin{lemma}[隣接 corner 間で整合する portals]
\label{ke:coherent-corner-portals}
任意の $z\in\mathcal C_J$ に対し $a_J(z)\in A_{t,J}$ を選び，
\begin{equation}\label{ke:eq:corner-portal-target-error}
 d_\square(P^{\rm old}_{t,J}(a_J(z)),z)\le C\eta
\end{equation}
とできる．さらに，$z\in\mathcal C_J\cap\mathcal C_{J'}$ なら
\begin{equation}\label{ke:eq:corner-portal-junction-cost}
 d_{g_t}(a_J(z),a_{J'}(z))\le H
\end{equation}
となるよう同時に選べる．同じ $J$ を連続して使う場合は，同じ点を使ってよい．
この選択の連続性は主張せず，必要な有限個の gallery endpoints のみに用いる．
\end{lemma}

\begin{proof}
$|\operatorname{supp}z|=3$ なら，この support を含む三元集合は $J$ 自身しかない．
$z\in P^{\rm old}(A_{t,J})$ なら，その旧 fiber から一点を選ぶ．それ以外では
$P^{\rm old}(A_{t,J})$ が $\mathcal C_J$ の $2\eta$-net であることを使って
\eqref{ke:eq:corner-portal-target-error}を満たす一点を選ぶ．
特に $z=P^{\rm old}(x)$ となる指定済みの $x\in A_{t,J}$ があれば，
その $x$ を選んでもよい．異なる corner への接続はこの場合には生じない．

$z=E_{ab}(u)$，$0<u<2$ の場合を考える．
$\kappa\le u\le2-\kappa$ なら
$\mu=\sigma_t^{-1}(u)$ とし，各 $J\supset\{a,b\}$ に対し
$a_J(z)\in\Lambda_{ab,J,t}(\mu)$ を選ぶ．
\eqref{ke:eq:portal-profile-error}が target 誤差を与え，
全ての $a_J(z)$ は同じ $P_{ab,j}$ の同一 moment 集合に属するので
\eqref{ke:eq:global-connector-budget}が接続費用を与える．
$u=1$ では $\mu=0$ であり，portal は source を含まない triple の regular torus 上にある．
同じ moment 集合には source の pinched fibers も含められ，全 moment 退出定理により
それらとの接続にも例外はない．

$0<u<\kappa$ なら，上で固定した $a_J(z)=p_{a,J}$ を用いる．
$P^{\rm old}(p_{a,J})$ は $e_a$ の $C\eta$ 近傍にある．
$d_\square(z,e_a)=u<\kappa=\eta$ なので target 誤差を得る．
各 $a_J(z)$ は同じ component compact に属するため，対数接続費用 $H$ を使える．
$2-\kappa<u<2$ では $p_{b,J}$ を，$z=e_a$ では $p_{a,J}$ を用いる．
$u=\kappa,2-\kappa$ は前段の closed moment interval に含め，それぞれ
$\mu=-\delta,\delta$ の境界 torus から選ぶ．Portal の連続性は要らないので，
端の開区間での固定点選択との不一致は問題にならない．
以上は target の全 support 型を尽くす．
\end{proof}

\begin{lemma}[任意点から gallery の最初の corner への接続]
\label{ke:point-to-corner-portal}
任意の $x\in X_t$ に対し $q_x\in\mathcal C_d$ を選び，
\begin{equation}\label{ke:eq:point-anchor-target-error}
 d_\square(f_{j,t}(x),q_x)\le C\eta
\end{equation}
とできる．さらに $q_x\in\mathcal C_J$ を満たす\emph{どの}三元集合 $J$ についても，
前補題の選択に整合する $a_J(q_x)$ を取り
\begin{equation}\label{ke:eq:point-to-portal-cost}
 d_{g_t}(x,a_J(q_x))\le H
\end{equation}
とできる．
\end{lemma}

\begin{proof}
全域構成の被覆を用いて，次の三つの場合に分ける．

まず $x\in A_{t,J_0}$ かつ
$|\operatorname{supp}P^{\rm old}_{t,J_0}(x)|=3$ なら
$q_x=P^{\rm old}_{t,J_0}(x)$ とする．保持 corner 上の基底値比較により
\eqref{ke:eq:point-anchor-target-error}を得る．
$q_x$ を含む三元集合は $J_0$ に限られ，前補題で選んだ $a_{J_0}(q_x)$ は
$x$ と同じ旧 torus に属する．その径 $C/m_j\le H$ により接続できる．
この選択により，異なる endpoints が同じ anchor を持つ場合にも一つの portal を使える．

次に $x$ が pair $ab$ の保持 neck 上にあり，
$\mu(x)\in[-\delta,\delta]$ であるときは
\[
 q_x=E_{ab}(\sigma_t(\mu(x)))
\]
とする．系\ref{ke:neck-base-value-control}が target 誤差を与える．
任意の $J\supset\{a,b\}$ で
$a_J(q_x)\in\Lambda_{ab,J,t}(\mu(x))$ と選べば，
両点は同じ $P_{ab,j}$ の同一 moment 集合に属し，接続費用は $H$ 以下である．
保持 corner 上で旧 $P^{\rm old}$ の support が二元になる点もこの場合に含まれる．
実際，その全 torus は $B=m_j+O(1)$ の旧辺上にあり，
K2 の共通 plateau overlap 内の同じ neck fiber である．
Core cut により残る旧辺の interval は正確に $\mu\in[-\delta,\delta]$ である．

残る点は一つの component $a$ の core 又は変更 collar に含まれる．
$q_x=e_a$ とすれば，命題\ref{ke:global-base-parameterization}の core-image 評価が
target 誤差を与える．各 $J\ni a$ について
$a_J(q_x)=p_{a,J}$ とすれば，$x$ と同じ固定 component compact に属するので
接続費用は $H$ 以下である．
Anchor moment interval の外側の neck 尾部は
補題\ref{ke:component-neck-tail-depth}により同じ対数予算で扱える．

複数の場合が重なる点はどちらで扱ってもよい．この分類の役割は，
gallery の最初の square の label が元の corner の label と異なるときにも，
同じ moment 又は同じ component 内で最初の portal へ到達できることにある．
\end{proof}

\subsection{有限 gallery による全域上界}

\begin{theorem}[Actual 全曲面距離の一様上界]
\label{ke:global-ambient-upper}
任意の固定 $d\ge5$ と上の actual fibration に対し
\[
 M_d=3\binom{\min\{d,6\}}3\le60
\]
と置く．固定 stage の全条件を満たす十分小さい $t$ に対して，全 $x,y\in X_t$ で
\begin{equation}\label{ke:eq:global-ambient-upper}
 \bar d_t(x,y)
 \le d_\square(f_{j,t}(x),f_{j,t}(y))+U_{j,t},
\end{equation}
ただし
\begin{equation}\label{ke:eq:global-ambient-upper-error}
 U_{j,t}=M_dE_{j,t}
       +C(M_d+1)\eta_{j,t}
       +(M_d+1)\frac{H_{j,t}}{R_t}.
\end{equation}
特に上の許容 cut 選択と $N_j\ge m_j=6b_j$ の下で
\begin{equation}\label{ke:eq:global-upper-explicit-depth-error}
 U_{j,t}\le C_d\left\{\frac1{b_j}
       +\frac{1+\log R_t+\log(N_j+1)}{R_t}\right\}.
\end{equation}
従って既存の遅い対角選択を十分小さくすれば，$j=j(t)\to\infty$ に沿って
$U_{j(t),t}\to0$ である．同じ選択で，全 $f_t$-fibers の正規化 ambient 径も
一様に零へ収束する．
\end{theorem}

\begin{proof}
前補題の anchors $q_x,q_y$ を取る．
系\ref{topo:six-label-gallery}により，両者を結ぶ最短路を
$M\le M_d$ 本の unit-square 内線分に分ける：
\[
 z_0=q_x,\ z_1,\ldots,z_M=q_y,
 \qquad[z_{r-1},z_r]\subset Q_r\subset\mathcal C_{J_r}.
\]
その Euclidean 長さを $\ell_r$ とすれば
$\sum_r\ell_r=d_\square(q_x,q_y)$ である．
両 anchors が等しい場合は，その点を含む一つの square の定値線分を使えばよい．

第 $r$ 線分の両端に対応する portals
\[
 p_r^-=a_{J_r}(z_{r-1}),\qquad p_r^+=a_{J_r}(z_r)
\]
を取る．初めと終わりでは補題\ref{ke:point-to-corner-portal}の指定点を使う．
\eqref{ke:eq:retained-protected-estimates}から得られる actual ambient 上界と
\eqref{ke:eq:corner-portal-target-error}により
\begin{align*}
 \bar d_t(p_r^-,p_r^+)
 &\le d_\square(P^{\rm old}_{t,J_r}(p_r^-),
                         P^{\rm old}_{t,J_r}(p_r^+))+E\\
 &\le \ell_r+2C\eta+E .
\end{align*}
ここで使う経路は旧 full corner $W_{t,J_r}$ の中の経路であり，
core cut で除いた cap を通っても $X_t$ 内の正当な経路である．
保持領域 $A_{t,J_r}$ の中に経路を閉じ込める必要はない．

隣接する $J_r,J_{r+1}$ が同じなら共通 portal を使える．
異なる場合，$z_r$ の support は二元以下である．
二元の中央 interval なら両 portals は同じ $\sigma_t^{-1}(u)$ を持つので，
\eqref{ke:eq:corner-portal-junction-cost}で費用 $H$ 以下の接続を得る．
Vertex に近い端 interval と一元 support では，同じ component compact の費用を使う．
最初と最後の $x,y$ への接続も各 $H$ 以下である．
従って全ての接続を足し合わせると
\[
 \bar d_t(x,y)
 \le d_\square(q_x,q_y)+M(E+2C\eta)+(M+1)H/R_t.
\]
さらに\eqref{ke:eq:point-anchor-target-error}を両端に用いれば
\[
 d_\square(q_x,q_y)
 \le d_\square(f_{j,t}(x),f_{j,t}(y))+2C\eta .
\]
$M\le M_d$ と合わせ，定数をまとめれば
\eqref{ke:eq:global-ambient-upper}--\eqref{ke:eq:global-ambient-upper-error}を得る．
距離の上界を実際の経路長で書く場合には，各距離の infimum を任意の正の誤差で
近似する有限個の経路を選び，最後にその誤差を零へ送ればよい．

最後に全ての fixed-stage 定数を先に $C_j$ に含め，既存の parameter 閾値とともに
\[
 R_t\ge j\bigl(1+C_j+\log(N_j+1)\bigr)
\]
を課す．$m_j\to\infty$，$N_j\ge m_j$ であり，
$\log R_t/R_t\to0$ も成立するので
$E_{j(t),t},\eta_{j(t),t},H_{j(t),t}/R_t\to0$ である．
これは既存の有限個の stage 条件を強めるだけの遅い対角選択である．
全 fiber 合併径については\eqref{ke:eq:full-fiber-diameter-upper}を使う．
\end{proof}

\paragraph{平方の lift 条件との区別．}
本節で実測した局所条件は，近い旧 target 値を持つ endpoints の
actual ambient 距離の上界と，junction portals の同一 moment 接続である．
命題\ref{topo:finite-portal-criterion}の \textup{(S)} は，target 線分へ
\emph{厳密に}写る rectifiable lift を要求する．Inset と core cap の再 parameter 化の
後では，旧 corner の lift がこの条件を満たすとは限らない．
上の証明はその \textup{(S)} を引用せず，必要な endpoint 版の条件を
有限 gallery に直接適用している．注意\ref{topo:finite-endpoint-comparison}の
anchor 誤差，square の端点比較誤差，junction と始終点の接続費用は，この適用では
順に $C\eta$，$E+2C\eta$，$H/R_t$ である．
\textup{(F)} に相当する actual 全 fiber 合併径は
命題\ref{ke:collapsed-tree-full-fiber-collapse}で独立に証明した．
全曲面を通る任意の経路に対する下界 \textup{(N)} は本節の結論に含めない．

% !TeX root = ../../surface-lagrangian-torus-fibration.tex
% 2026-11-22 / v1.8: 主定理の統合証明と許容中間尺度の帰結．
\section{Actual Lagrangian fibration による負 KE 極限の実現}
\label{ke:global-metric-realization-section}

\begin{theorem}[全固定次数における大域的な計量実現]
\label{ke:global-metric-realization}
定理\ref{ke:global-realization}の固定次数 $d\ge5$ と退化族に対し，
節\ref{ke:cut-selection-order}の順序で有限段階のデータを選ぶと，
同定理の二条件を同時に満たす一つの固定木系 $T$ と actual Lagrangian fibration
$\pi_t:(X_t,\omega_t)\to B_d(T)$ が得られる．特に，同じ全射
$f_t=p_T\pi_t$ に対して
\begin{equation}\label{ke:eq:global-metric-realization}
 \operatorname{dis}(f_t):=
 \sup_{x,y\in X_t}
 \left|R_t^{-1}d_{g_t}(x,y)
       -d_\square(f_t(x),f_t(y))\right|\longrightarrow0
\end{equation}
となる．距離は全曲面の ambient 距離であり，量化と規格化は
定理\ref{ke:global-realization}のものを用いる．固定 stage の明示的な評価は
式\eqref{ke:eq:global-realization-explicit-cut-budget}で与える．
\end{theorem}

\begin{proof}
\emph{構成と同じ写像の選択．}
定理\ref{ke:actual-global-fibration}により，固定 stage $j$ の後で
$|t|$ を十分小さくすれば，全 $X_t$ 上に actual fibration を得る．
全被覆，連結 fibers，actual Lagrangian 性はこの定理で証明されている．
補題\ref{ke:fixed-marked-tree-exhaustion}の共通選択により，
木系の組合せ型，labelled ends，および source markings は stage と $t$ に依存せず固定できる．
ここでは各固定二重線の一つの embedded pants foliation から構成し，
動く source を小円板内で，end germs をその外の annuli で相対的に調整している．
Normal circle の向きは共通にするが，全 torus cycles の parameter-monodromy を
自明化する主張ではない．
命題\ref{ke:global-base-parameterization}を用いて，その同じ fiber partition の
基底を $B_d(T)$ と parameter 化し，$f_{j,t}=p_T\pi_{j,t}$ とする．
これは距離比較のために別の fiber partition を採用する操作ではない．
補題\ref{ke:anchor-trimmed-neck-domain}により，core と変更 collar の外で使う
保持 neck は共通 profile の定義域 $|\widehat\mu|\le\delta_j$ に入る．
補題\ref{ke:same-moment-corner-reference}は，その同じ moment の旧 corner 辺への
接続を与える．下界の値比較と上界の portals は，この同じ固定残余基底と
canonical moment を用いている．

\emph{固定 stage の両側評価．}
定理\ref{ke:global-ambient-upper}の記号で
\[
 U_{j,t}=M_dE_{j,t}+C(M_d+1)\eta_{j,t}
          +(M_d+1)H_{j,t}/R_t,
 \qquad M_d=3\binom{\min\{d,6\}}3\le60
\]
とする．下界の誤差を
\begin{equation}\label{ke:eq:global-realization-lower-budget}
 V_{j,t}=C_d\left\{\varepsilon_j+b_j^{-1}+K_j^{-2}
 +\frac{1+\log\mathfrak m_j+\log K_j+\log(N_j+1)}{R_t}\right\}
\end{equation}
と書く．定理\ref{ke:global-ambient-lower}と
定理\ref{ke:global-ambient-upper}は，\emph{同じ} $f_{j,t}$ に対し
全 $x,y\in X_t$ で
\[
 d_\square(f_{j,t}(x),f_{j,t}(y))-V_{j,t}
 \le R_t^{-1}d_{g_t}(x,y)
 \le d_\square(f_{j,t}(x),f_{j,t}(y))+U_{j,t}
\]
を与える．上界では全内部木の共通 moment，source の全 moment 退出，
各 component compact の費用，および高々 $M_d$ 本の square gallery を使う．
下界では global clipped-depth map の長さ評価を全ての経路に直接積分する．
従って局所 protected 距離を ambient 距離に読み替える段階も，
経路の領域出入りの回数を有限と仮定する段階もない．
上界の式\eqref{ke:eq:global-upper-explicit-depth-error}と合わせると，
許容 cut の下で固定族の定数 $C_d$ により
\begin{equation}\label{ke:eq:global-realization-explicit-cut-budget}
 \operatorname{dis}(f_{j,t})\le C_d\left\{
 \varepsilon_j+b_j^{-1}+K_j^{-2}
 +\frac{1+\log R_t+\log\mathfrak m_j+\log K_j+\log(N_j+1)}{R_t}
 \right\}.
\end{equation}
ここには独立の stage-dependent な乗法係数を残していない．
Stage の定数は，cut の下限と最後の parameter 閾値へ入っている．

\emph{パラメータの選択順序と対角化．}
固定領域上で既に証明した有限個の誤差と費用を一つの遅い段階選択にまとめる原理は，
関連する局所解析原稿の補題\cite[\texttt{pair:lem:diagonal-fixed-end}]{KELocalCurrent}に明記した．
以下では本稿固有の cut の選択順序を確認する．この対角化自体は，
前段で別に証明した moving buffer，source 退出，scalar 下界を供給するものではない．
まず $b_j\to\infty$，$\varepsilon_j\downarrow0$ を選び，旧 corner の cut を
$m_j=6b_j$ とする．各 stage の residual domains とその接続定数
$C_j^{\rm neck}$ を先に固定する．次に component cut $N_j$ を，既存の下限および
\eqref{ke:eq:neck-constant-depth-choice}を満たすよう選び，core と相対 collars を作る．
この順序で $H_{j,t}=C(1+\log(N_j+1))$ となり，
上界には\eqref{ke:eq:global-upper-explicit-depth-error}が使える．
定理\ref{ke:sharp-global-one-active}と
命題\ref{ke:global-depth-base-comparison}の順序に従い，
Scalar 障壁専用の補助 cut $\mathfrak n_j$ とその境界定数を決めてから，
新しい lower cut $\mathfrak m_j$ と upper cut factor $K_j\to\infty$ を選ぶ．
$\mathfrak n_j$ は既存 neck/corner 構成の $n_j$ とは独立の記号である．
この $\mathfrak m_j$ は旧 corner cut $m_j$ と同じ量ではない．
構成・比較の有限個の条件に加えて
\begin{equation}\label{ke:eq:global-realization-diagonal}
 R_t\ge j\{1+\log\mathfrak m_j+\log K_j+\log(N_j+1)\}
\end{equation}
を満たす閾値 $r_j>0$ を取る．$r_{j+1}<r_j/2$ としてよく，
$0<|t|<r_j$ なら stage $j$ の全条件が成り立つようにする．
$j(t)=\max\{j:|t|<r_j\}$ と置けば $j(t)\to\infty$ であり，
\[
 \frac{\log(m_{j(t)}+1)+\log R_t}{R_t}\to0,
 \quad \frac{1+\log(N_{j(t)}+1)}{R_t}\to0.
\]
ここでは $N_j\ge m_j$ と，$\log R_t/R_t\to0$ を用いた．
従って $U_{j(t),t},V_{j(t),t}\to0$ となる．
$\pi_t=\pi_{j(t),t}$ と選び，両側評価から
$\operatorname{dis}(f_t)\le\max\{U_{j(t),t},V_{j(t),t}\}\to0$ を得る．
この定理は $t$ に関する fibration 族の連続性・滑らかさや，
あらかじめ指定した速さの対角化を結論に含めない．
明示的な cut budget は，$C_j^{\rm neck}$ の増大率，scalar 障壁に必要な
$\mathfrak m_j,K_j$ の成長，又は fixed-compact convergence の有効な閾値を
与えるものではない．従って $R_t$ だけによる最適な全域収束速度の証明とは区別する．
\end{proof}

\begin{lemma}[全域の歪みから同半径の閉球への移送]
\label{ke:distortion-to-pointed-balls}
Compact geodesic spaces $(X,d_X)$，$(Y,d_Y)$ と，全射 $F:X\to Y$ があり，
全点対に対する distortion が $\delta\ge0$ 以下とする．
任意の $x\in X$ と $A>0$ に対し，ambient 距離を制限した二つの閉球
\[
 (\overline B_X(x,A),x),\qquad
 (\overline B_Y(F(x),A),F(x))
\]
の間には，基点対を含む，distortion が $3\delta$ 以下の対応がある．
従って基点付き GH 距離は $3\delta/2$ 以下である．定数は $x,A$ と両空間に依存しない．
\end{lemma}
\begin{proof}
$z\in\overline B_X(x,A)$ に対し，$d_Y(F(x),F(z))\le A+\delta$ である．
必要なら $F(x)$ から $F(z)$ への最短路上で距離 $\delta$ 以下だけ戻し，
半径 $A$ の球内の点 $y_z$ を選び，対 $(z,y_z)$ を取る．基点対 $(x,F(x))$ も取る．
球内の任意の $y$ に対し，全射性から $F(w)=y$ となる $w$ を選ぶと
$d_X(x,w)\le A+\delta$ である．Source の最短路上で必要なだけ戻せば，
$d_X(z,w)\le\delta$ となる $z\in\overline B_X(x,A)$ を取れる．
この対 $(z,y)$ も加えると，両閉球へ全射な対応を得る．
この対応の各対 $(z,y)$ には，元の全域写像の点 $w$ で
\[
 d_X(z,w)+d_Y(y,F(w))\le\delta
\]
となるものが付随する．二つの対に対してこの評価と $F$ の distortion を加えれば，
対応の distortion は $3\delta$ 以下である．基点対を含む対応から通常の共通距離を
構成すると，両球の Hausdorff 距離と基点間距離は $3\delta/2$ 以下となる．
これらの選択の連続性は必要ない．
\end{proof}

\begin{corollary}[同じ射影で制御される中間尺度の全移動基点極限]
\label{ke:intermediate-moving-classification}
定理\ref{ke:global-metric-realization}の同じ写像 $f_t$ と
$\epsilon_t=\operatorname{dis}(f_t)$ を用い，正数 $a_t$ が
\[
 a_t\longrightarrow\infty,\qquad a_t/R_t\longrightarrow0,
 \qquad R_t\epsilon_t/a_t\longrightarrow0
\]
を満たすとする．任意の $t_\nu\to0$ と $x_\nu\in X_{t_\nu}$ に対し，
$(X_{t_\nu},a_{t_\nu}^{-1}d_{g_{t_\nu}},x_\nu)$ は部分列を取れば pointed GH 収束する．
その全ての可能な極限は，命題\ref{topo:cubical-moving-blowups}の四種類
\[
 (C_{2,d-1},z),\quad
 (\mathbb R\times\mathcal Y_{d-2},(0,z)),\quad
 (\mathcal K_{3\pi/2},z),\quad
 (\mathbb R^2,0)
\]
であり，この条件を満たす任意の尺度と退化列に対して，各型の任意の有限基点を
実現する $x_\nu$ を選べる．これは metric spaces の分類であり，
追加の labels 又は tensor 収束を結論に含めない．
\end{corollary}
\begin{proof}
$\lambda_t=R_t/a_t\to\infty$，$\delta_t=\lambda_t\epsilon_t\to0$ と置く．
同じ全射 $f_t$ は $(X_t,a_t^{-1}d_{g_t})$ と
$(\mathcal C_d,\lambda_t d_\square)$ の全点対距離を $\delta_t$ 以下の誤差で比較する．
両者は compact geodesic spaces なので，前補題により，任意の基点と固定半径の
閉球を $O(\delta_t)$ の誤差で比較できる．命題\ref{topo:cubical-moving-blowups}を
$q_\nu=f_{t_\nu}(x_\nu)$ に適用すれば，部分列の存在と分類が従う．
逆に同命題で各型を実現する $q_\nu$ を選び，全射性によって
$x_\nu\in f_{t_\nu}^{-1}(q_\nu)$ を取れば，その任意の選択が同じ極限を実現する．
必要なのは指定した尺度で $R_t\epsilon_t/a_t\to0$ となることであり，
この条件を除いた任意の中間尺度に関する分類ではない．
\end{proof}

\begin{corollary}[選択した中間尺度での成分錐]
\label{ke:intermediate-component-cone}
定理\ref{ke:global-metric-realization}の同じ $f_t$ に対し
$\epsilon_t=\operatorname{dis}(f_t)\to0$ と置く．成分 label $i$ を固定し，
任意に $x_{i,t}\in f_t^{-1}(e_i)$ を選ぶ．正数 $a_t$ が
\begin{equation}\label{ke:eq:intermediate-component-scale}
 a_t\longrightarrow\infty,\qquad
 \frac{a_t}{R_t}\longrightarrow0,\qquad
 \frac{R_t\epsilon_t}{a_t}\longrightarrow0
\end{equation}
を満たすなら，pointed Gromov--Hausdorff の意味で
\[
 (X_t,a_t^{-1}d_{g_t},x_{i,t})
 \longrightarrow (C_{2,d-1},d_C,0).
\]
このような尺度は常に存在し，例えば
\begin{equation}\label{ke:eq:intermediate-component-choice}
 a_t=R_t\sqrt{\epsilon_t+R_t^{-1}}
\end{equation}
と取れる．極限の錐は命題\ref{topo:component-cone-germ}のものである．

より一般に，非空集合 $A_t\subset X_t$ に対し
\[
 \eta_t=\sup_{y\in A_t}d_\square(f_t(y),e_i)\longrightarrow0
\]
とする．式\eqref{ke:eq:intermediate-component-scale}の第三条件を
\begin{equation}\label{ke:eq:intermediate-basepoint-budget}
 \frac{R_t(\epsilon_t+\eta_t)}{a_t}\longrightarrow0
\end{equation}
へ強めれば，任意の $y_t\in A_t$ を基点として同じ pointed GH 収束が成り立つ．
各固定半径の閉球に対する近似誤差は $y_t$ の選択に一様である．この場合には
\[
 a_t=R_t\sqrt{\epsilon_t+\eta_t+R_t^{-1}}
\]
と選べる．特に基点を成分頂点の特異 fiber 上へ限定する必要はない．
\end{corollary}
\begin{proof}
$\lambda_t=R_t/a_t$，$\delta_t=\lambda_t\epsilon_t$ と置くと，
$\lambda_t\to\infty$，$\delta_t\to0$ である．以下，閉球には各 ambient 距離の制限を入れる．
全射 $f_t$ は基点を $e_i$ へ送り，
$(X_t,a_t^{-1}d_{g_t})$ と $(\mathcal C_d,\lambda_t d_\square)$ の
距離の差を全点対について $\delta_t$ 以下にする．
固定した $A>0$ に対し，十分小さい $t$ では $(A+2\delta_t)/\lambda_t<1/2$ なので，
命題\ref{topo:component-cone-germ}により target のその半径の球は，
錐の半径 $A+2\delta_t$ の球と等長である．
補題\ref{ke:distortion-to-pointed-balls}により source と target の同半径 $A$ の閉球を
基点を保って GH 誤差 $3\delta_t/2$ 以下で比較できる．この等長写像と合成して
錐との比較を得るので，pointed GH 収束が従う．
式\eqref{ke:eq:intermediate-component-choice}は $a_t\ge\sqrt{R_t}$，
$a_t/R_t\to0$，$R_t\epsilon_t/a_t\le\sqrt{\epsilon_t}\to0$ を満たす．

一般の基点についても，同じ $x_{i,t}\in f_t^{-1}(e_i)$ を一つ取る．
全域 distortion 評価から，全 $y\in A_t$ に対して
\[
 a_t^{-1}d_{g_t}(y,x_{i,t})
 \le\frac{R_t}{a_t}(\eta_t+\epsilon_t)=:s_t\longrightarrow0
\]
である．同じ geodesic 空間の二つの中心が距離 $s_t$ 以下なら，その半径 $A$ の
閉球の Hausdorff 距離も $s_t$ 以下である．実際，一方の球の点を他方の中心から
最短路で結び，必要な場合だけ半径 $A$ まで戻せば，移動距離は高々 $s_t$ である．
逆方向も同じであり，中心間の距離も $s_t$ 以下だから，基点を含む近似を得る．
これを $x_{i,t}$ を中心とする上の近似と合成すれば，誤差は
$O(\delta_t+s_t)$ となり，全ての $y\in A_t$ に一様な結論が従う．
最後の尺度の選択では $a_t\ge\sqrt{R_t}$ であり，
\[
 \frac{a_t}{R_t}\to0,\qquad
 \frac{R_t(\epsilon_t+\eta_t)}{a_t}
 \le\sqrt{\epsilon_t+\eta_t}\to0
\]
となる．
\end{proof}

\paragraph{Component core と変更 collar 内の基点．}
同じ $f_t$ の構成で共通 anchor によって定めた component $i$ の core domain と
変更 collar の合併を $A_{i,t}$ とする．命題\ref{ke:global-base-parameterization}により
\[
 \sup_{y\in A_{i,t}}d_\square(f_t(y),e_i)
 \le C q_t,\qquad
 q_t=\frac{1+\log(N_{j(t)}+1)}{R_t}\longrightarrow0
\]
を満たす．従って $a_t=R_t\sqrt{\epsilon_t+q_t}$ と取れば，この全領域内の
任意の基点列について同じ成分錐への pointed GH 収束を得る．実際，
$q_t\ge R_t^{-1}$ であり，
$R_t(\epsilon_t+Cq_t)/a_t\le\max(1,C)\sqrt{\epsilon_t+q_t}\to0$ である．
この結論が使うのは全域の ambient 距離であり，core 内部の経路距離や
任意の Liouville 軌道長の一様評価は要求しない．

これは既存の大域実現定理と有限複体の局所距離からの帰結であり，
固定成分の漸近錐定理を actual 族へ直接適用する証明ではない．
特に任意に先指定した $1\ll a_t\ll R_t$ については，
頂点 fiber 上の基点では式\eqref{ke:eq:intermediate-component-scale}の誤差条件を，
一般の $A_t$ では基点のずれを含む式\eqref{ke:eq:intermediate-basepoint-budget}を
確かめる必要がある．
歪みの有効な速度は未評価なので，$t$ だけで先に決めた普遍的な中間尺度，
二つの極限の交換，又は中間尺度での tensor 収束は結論しない．

\begin{corollary}[七次以上の同じ基底における core の局所位相障害]
\label{ke:global-core-nonmatroidal}
$d\ge7$ のとき，定理\ref{ke:global-metric-realization}で選んだ $\pi_t$ の基底
$B_d(T)$ のどの core 頂点の近傍にも，二次元 matroidal local model への chart は存在しない．
従って，core 頂点を含むこの同じ基底に作用適合 matroidal atlas を入れることはできない．
これは定理\ref{trop:degree-classification}の既存の位相障害を，上の計量実現で選んだ基底に適用した帰結である．
\end{corollary}
\begin{proof}
定理\ref{ke:actual-global-fibration}の構成では，各 core 頂点 $v_i$ の近傍は，
同じ木系 $T$ から作る graph $G_i$ の cone である．命題\ref{topo:counts}より
$G_i$ は連結三価であり，
\[
 b_1(G_i)=\binom{d-2}{2}\ge10.
\]
この rank は $v_i$ における二次の局所 homology の rank でもあり，局所同相で不変である．
穿孔近傍の分岐が三価だけであることも局所位相から判定される．
一方，定理\ref{trop:cubic-classification}によれば，その分岐条件を満たす二次元
matroidal local model の link は，非連結 matroids と Euclidean factors を許しても
$b_1\le8$ を満たす．従ってそのような chart は存在しない．
ここでは作用格子や計量の収束を追加仮定せず，収縮後の $\mathcal C_d$ についての
atlas の有無も結論していない．
\end{proof}

\begin{corollary}[GH 収束と全特異 fiber 合併の収縮]
\label{ke:global-gh-and-fibers}
定理\ref{ke:global-metric-realization}の $f_t$ は全射であり，
$\epsilon_t=\operatorname{dis}(f_t)$ と置くと
\begin{align*}
 d_{\mathrm{GH}}\bigl((X_t,R_t^{-1}d_{g_t}),
                     (\mathcal C_d,d_\square)\bigr)&\le\epsilon_t/2,\\
 \sup_{c\in\mathcal C_d}
 \operatorname{diam}_{R_t^{-1}d_{g_t}}f_t^{-1}(c)&\le\epsilon_t,\\
 \sup_{b\in B_d(T)}
 \operatorname{diam}_{R_t^{-1}d_{g_t}}\pi_t^{-1}(b)&\le\epsilon_t.
\end{align*}
特に $d\ge6$ では，$p_T$ が潰す各内部木上の全ての
$\pi_t$-fibers の合併が，一様に正規化 ambient 径を零にする．
\end{corollary}

\begin{proof}
$p_T$ と $\pi_t$ は全射なので，$f_t$ の graph は両空間の correspondence である．
その distortion が $\epsilon_t$ であることから第一式を得る．
同一の $f_t$-fiber の二点を\eqref{ke:eq:global-metric-realization}に代入すると
第二式を得る．各 $\pi_t$-fiber は一つの $f_t$-fiber に含まれるので第三式も従う．
内部木の主張は命題\ref{ke:collapsed-tree-common-moment}と同じ fiber 合併であり，
命題\ref{ke:collapsed-tree-full-fiber-collapse}はその非正規化費用も与えている．
\end{proof}

\begin{corollary}[直径と径一の規格化]
\label{ke:global-diameter-limit}
$D_d=\operatorname{diam}(\mathcal C_d,d_\square)$ と置くと
\[
 R_t^{-1}\operatorname{diam}(X_t,g_t)\longrightarrow D_d,
 \qquad
 \left(X_t,\frac{d_{g_t}}{\operatorname{diam}(X_t,g_t)}\right)
 \xrightarrow{\mathrm{GH}}(\mathcal C_d,D_d^{-1}d_\square).
\]
$d\ge6$ では定理\ref{topo:diameter-six-labels}により $D_d=2\sqrt5$ である．
五次の値 $D_5=\sqrt{13}$ には，
\cite[\texttt{thm:intro-diameter-C5}]{KECubicalCurrent}の有限計量結果を用いる．
\end{corollary}

\begin{proof}
全射性と distortion 評価から
$|R_t^{-1}\operatorname{diam}(X_t,g_t)-D_d|\le\epsilon_t$ である．
$D_d>0$ なので，同じ correspondence の距離を直径で割れば第二式も従う．
\end{proof}

\paragraph{結論の範囲．}
証明は，関連 KE 原稿の未検証の大域入力
$\mathrm{(CP)}_d,\mathrm{(PBE)}_d$ 又は $\mathrm{(TC)}_5$ を仮定していない．
局所 density と全深部 triple comparison，固定 compact の収束という明記した入力から，
本稿で直接 scalar bound，source exits，基底値の比較，portal 接続を証明した．
任意に先指定した木系や critical values の実現，special Lagrangian 性，
一般次数の大域 matroidal atlas，最適収束速度は本定理の結論に含めない．

\input{sections/paper/degree_partitions.tex}
\input{sections/paper/partition_topology.tex}
\input{sections/paper/tropical_obstructions.tex}
\input{sections/paper/petersen_periods.tex}
\input{sections/paper/quadric_periods.tex}
\input{sections/paper/mixed_plane_periods.tex}
\input{sections/paper/atlas_gluing.tex}
% !TeX root = ../../surface-lagrangian-torus-fibration.tex
% 2026-11-22 / v1.2: Source 保持下の成分別面積配分・finite cuts・actual 微小実現．
\section{Critical actions を直線に保つ相対 neck 接着}
\label{ng8:section}

葉層の相対 isotopy による補間だけでは，critical action curve の直線性は保たれない．
ここでは，共通の Hamiltonian circle と三重点の toric reference circles を使って
作用の加法定数を固定し，critical regions の面積を指定して補間する．
輸送を行うのは常に正の距離を置いた切断面であり，
$\mu=0,t=0$ の固定点集合を通る parallel transport は用いない．

\subsection{作用の基準と固定 moment 上の輸送}

二重曲線 $C=C_{ij}$ から三重点の小円板を除いた曲面を $P$ とする．
節~\ref{sec:necks} の actual neck には，周期 $2\pi$ の Hamiltonian circle と
正規化された moment $\mu$ がある．$i$ 側で normal moment が $\mu>0$，
$j$ 側で $-\mu>0$ となるように向きを取る．
有限個の nodal patches より十分薄い collar を選ぶ．
各 $t\ne0$ の neck では circle action は free であるが，
以下で primitive を取る全空間の tube には，中央 fiber の固定集合 $C^\circ$ も含める．

buffer を持つ全空間の tube は，$P$ と同じ homotopy 型の open curve へ retract するので，
$H^2=0$ である．
従って $\Omega=\dd\lambda_{\rm neck}$ と書ける．
circle averaging により $\lambda_{\rm neck}$ を invariant に取ると，
Cartan の公式と $\iota_V\Omega=-\dd\mu$ から
\[
 \dd\bigl(\lambda_{\rm neck}(V)\bigr)
 =-\iota_V\Omega=\dd\mu
\]
を得る．固定集合 $C^\circ$ 上では $V=0$ および $\mu=0$ なので，
$\lambda_{\rm neck}(V)=\mu$ となる．
以下の作用は，この一つの primitive で正規化する．
これは各 puncture の toric action の加法定数を独立に零にする操作ではない．

各 puncture end の toric chart で，circle generator $\alpha$ ともう一つの
primitive generator $\beta_\ell$ を固定する．
normal moment と残る toric moment の affine cuts で
compact な reduced surface $P_{t,\mu}$ を切り出す．
その境界 circle の第二作用を $R_\ell(\mu)$ と書く．
各 end の toric moment を $\nu_\ell$ とすると，これは
$\nu_\ell+c_\ell^{\rm ref}$ の cut への制限である．
$c_\ell^{\rm ref}$ は $\lambda_{\rm neck}$ で決まる定数で，
両側の成分と全ての $t$ に共通である．
実際，$\beta_\ell$ の toric flow を $R_\theta$，generator を $V_{\beta_\ell}$ とすると，
その orbit action は
\[
 \widetilde R_\ell=
 \frac1{2\pi}\int_0^{2\pi}R_\theta^*
       \bigl(\lambda_{\rm neck}(V_{\beta_\ell})\bigr)\,\dd\theta
\]
である．Cartan の公式の $\mathcal L_{V_{\beta_\ell}}\lambda_{\rm neck}$ 項は
一周期で積分すると零になるので，
$\dd\widetilde R_\ell=-\iota_{V_{\beta_\ell}}\Omega=\dd\nu_\ell$ を得る．
従って同じ連結な全空間の toric end で
$\widetilde R_\ell=\nu_\ell+c_\ell^{\rm ref}$ である．
この計算に $\lambda_{\rm neck}$ 自体の全 $T^2$-invariance は必要ない．
従って $R_\ell(\mu)$ は実数係数の affine function であり，$t$ に依存しない．
cut 自体の勾配が整数であることは要求しない．

boundary circle から同じ regular Reeb arm 上の circle $\gamma$ までの
reduced annulus を $A$ とすれば，向きと lift を一貫して選ぶことで
\begin{equation}\label{ng8:relative-action}
 J_{\beta_\ell}(\mu,\gamma)-R_\ell(\mu)
 =\frac1{2\pi}\int_A\omega_{t,\mu}
\end{equation}
となる．ここで $J_\alpha=\mu$ と正規化する．
lift の整数変更は $J_\beta$ に整数倍の $\mu$ を加える．
critical circle への値は regular annuli の exhaustion により定める．
内部 Reeb edge の作用差も，同じ式をその annulus に適用して定める．

\begin{lemma}[切断面における作用の加法定数の保存]
\label{ng8:transport-intercept}
固定した $0<\delta$ と $\mu=\delta$ または $\mu=-\delta$ の近傍で，
成分側から nearby fiber への輸送は，
toric reference circles で正規化した作用の加法定数を保存する．
従って，成分側の critical action germ が
$J_\beta=c+p\mu$ なら，輸送後も同じ $c,p$ を持つ．
\end{lemma}

\begin{proof}
補題~\ref{glue:transport} の horizontal vector field は circle invariant で，
全空間上の $\mu$ を保存する．固定 $\mu\ne0$ で symplectic reduction を取ると，
その vector field は reduced symplectic forms を保つ horizontal field に降りる．
三重点の端では二つの toric moments を保存するので，
reference circle とその作用 $R_\ell(\mu)$ も保存される．
式~\eqref{ng8:relative-action} の annulus を運べば，その面積が保存される．
critical limits に対しても exhaustion により同じ結論が従う．

一つの全空間 primitive を使うと，作用そのものの保存も直接確認できる．
輸送 path を $\tau\mapsto\tau t$，その horizontal lift を $W_\tau$，
transport を $\mathcal P_{\tau t}$ とする．
固定した $\mu=\pm\delta$ の buffers を十分小さく取れば，全輸送軌跡は上の tube に入る．
$W_\tau$ は fiber の tangent space と $\Omega$-orthogonal なので，
Cartan の公式を fiber tangent に制限すると
\[
 \left.\mathcal L_{W_\tau}\lambda_{\rm neck}\right|_{TX_{\tau t}}
 =\dd_{X_{\tau t}}\bigl(\lambda_{\rm neck}(W_\tau)\bigr).
\]
従って，輸送の定義域上で
\begin{equation}\label{ng8:exact-transport}
 \mathcal P_t^*(\lambda_{\rm neck}|_{X_t})
   -\lambda_{\rm neck}|_{D_i}
 =\dd\!\left(\int_0^1
       \mathcal P_{\tau t}^*\bigl(\lambda_{\rm neck}(W_\tau)\bigr)
       \,\dd\tau\right).
\end{equation}
閉じた transported cycles の積分は正確に保存される．
この等式は primitive の period を任意に再正規化する操作ではない．

切断面の小さい open moment interval 上でこの議論を行うので，
作用微分だけでなく関数そのものが保存される．
定理の証明に必要なのは固定した二つの切断面とその buffers だけである．
$\delta\to0$ に一様な輸送を要求していない．
\end{proof}

\begin{lemma}[相対 reduced area の affine 性]
\label{ng8:relative-dh}
固定した十分小さい $t\ne0$ に対し，
上記の toric moment cuts で切り出した $P_{t,\mu}$ の全 reduced area は
$\mu$ の affine function である．
その値は，両端の切断面で中央成分から輸送した reduced area と一致する．
\end{lemma}

\begin{proof}
$t\ne0$ の neck に固定点はなく，moment interval 上の reduced surfaces は
境界を含めて同じ曲面 $P$ と同定できる．
principal circle bundle に connection を選ぶと，reduced cohomology class の変化は
その curvature class に比例する．これは Duistermaat--Heckman の局所的な計算である
\cite[Chapter~10]{DHJH}．境界付き曲面では絶対 cohomology class だけから総面積は
決まらないため，ここで必要な主張は次の Stokes の計算で直接証明する．

境界項も確認する．$P$ は境界を持つので circle bundle は自明化できる．
この自明化と各 toric end の framing の差は，固定した整数 winding である．
固定 $\mu$ 上で $\lambda_{\rm neck}-\mu\dd\theta$ は reduced surface の primitive に降りる．
従ってその境界 circle の積分は $R_\ell(\mu)+m_\ell\mu$，
$m_\ell\in\ZZ$，という affine function になる．
Stokes の定理により全 reduced area は，これらの境界作用の符号付き和の
$2\pi$ 倍であり，affine である．
moving cut を用いる場合の境界の寄与も $R_\ell(\mu)$ に既に含まれている．
両端の値の一致は補題~\ref{ng8:transport-intercept} の reduced transport による．
\end{proof}

\begin{lemma}[固定した fiber の free band における相対面積]
\label{ng8:fixed-fiber-relative-dh}
一つの $t\ne0$ と compact moment interval $J$ を固定する．
Actual $(E_t,\omega_t)$ 上に周期 $2\pi$ の自由な Hamiltonian circle action があり，
$\iota_V\omega_t=-\dd\mu$，$E_t/S^1\simeq J\times P$ とする．
ここで $P$ は非空の境界を持つ compact connected oriented surface である．
作用不変な primitive $\lambda_t$ が
$\dd\lambda_t=\omega_t$，$\lambda_t(V)=\mu$ を満たすと仮定する．
各 boundary component の toric framing と cut を固定し，
境界の向きに合わせた第二作用 $R_{\ell,t}(\mu)$ が affine であるとする．
このとき，$\mu$ に依存しない整数 $k_{\ell,t}$ があって
\begin{equation}\label{ng8:eq:fixed-fiber-area}
 \int_{P_{t,\mu}}\omega_{t,\mu}
 =2\pi\sum_\ell\bigl(R_{\ell,t}(\mu)+k_{\ell,t}\mu\bigr)
\end{equation}
となる．特に total reduced area は $J$ 上で affine である．
Reference constants の $t$ 依存を許し，$t=0$ を含む全空間への形式の延長，
central component からの輸送，又は $t$ に一様な評価は仮定しない．
\end{lemma}
\begin{proof}
$H^2(J\times P;\ZZ)=0$ なので principal circle bundle を自明化し，
周期 $2\pi$ の angular coordinate $\theta$ を取る．
$\lambda_t-\mu\dd\theta$ は horizontal かつ invariant であり，
固定 $\mu$ の reduction 上の primitive $\beta_{t,\mu}$ に降りる．
この自明化の boundary section と toric framing の差は整数 winding
$k_{\ell,t}$ で，連結な $J$ 上で一定である．従って
$\int_{\partial_\ell P}\beta_{t,\mu}
=2\pi(R_{\ell,t}(\mu)+k_{\ell,t}\mu)$ であり，Stokes の定理で結論を得る．
Toric moment $\nu_\ell$ を用いる場合にも，上の orbit-average 計算により
第二作用は $\nu_\ell+c_{\ell,t}^{\rm ref}$ である．
この定数は固定 $t$ の連結な end で一定であれば足り，全 $t$ で同じ値である必要はない．
\end{proof}

\begin{example}[Source 固定点を横切るときの面積の折れ]
\label{ng8:source-dh-wall}
$\CC^2$ 上で $I_k=|z_k|^2/2$，角度の周期を $2\pi$ とし，
標準形式 $\omega=\sum_{k=1}^2\dd I_k\wedge\dd\theta_k$，
重み $(1,-1)$ の circle action，$\mu=I_1-I_2$ を取る．
固定した $c>0$ に対し $I_1+I_2\le c$ を reduction すると，$|\mu|<c$ でその面積は
\begin{equation}\label{ng8:eq:source-reduced-area}
 \operatorname{Area}(P_\mu)=\pi(c-|\mu|)
\end{equation}
となる．$\mu=0$ では singular reduction の regular part 上の積分を面積とし，
これは両側からの連続極限にも一致する．
\end{example}
\begin{proof}
両座標が非零の部分で $I=I_2$，$\phi=\theta_1+\theta_2$ と置くと，
$\phi$ の周期は $2\pi$ であり，reduced form は $\dd I\wedge\dd\phi$ である．
許される interval は
\[
 \max\{0,-\mu\}\le I\le(c-\mu)/2.
\]
下端は $\mu\ne0$ では滑らかな reduction の一点，$\mu=0$ では固定点の像であり，
いずれも上記の積分に面積を加えない．従って面積は
$2\pi\{(c-\mu)/2-\max(0,-\mu)\}=\pi(c-|\mu|)$ である．
外側の toric boundary action $I_2=(c-\mu)/2$ は affine だが，
原点は $\mu=0$ の固定点であり，全 interval に対する free-band の自明化は存在しない．
実際，$\mu>0$ では $\theta_1=0$，$\mu<0$ では $\theta_2=0$ の section を使うと，
boundary winding はそれぞれ $0$ と $1$ となり，上の二つの affine pieces を与える．
\end{proof}

\begin{remark}[Actual KE の相対 atlas へ使う際の入力]
\label{ng8:fixed-fiber-area-scope}
前例は，free-band の面積式を source 固定点を横切る interval へ拡張できないことを示す．
Actual P0 全体の reduced area を計算した例ではなく，既存の atlas 存在定理や
P0 の訂正でもない．Actual KE 形式に $t=0$ までの smooth total-space extension を
仮定することなく，前補題を各固定 $t$ の free compact band に適用できる．
P0 の周期 $1$ の規約からは，生成場と canonical moment をともに $2\pi$ で割って用いる．

ただし補題\ref{ng8:parametric-moser}を続けて使うには，marked identifications と
端の open foliation germs，正の target region areas，その総量と actual reduced area の一致，
および\eqref{ng8:protected-area-margin}を別に確認する．
固定領域 $K$ には保存する全 collars と caps を含め，同補題の相対構成の条件を満たす必要がある．
同じ normal moment と end の葉を保つことだけでは，これらの面積条件は従わない．
Source の近傍を保持する場合は変更の支持をそこから離し，
critical-area の調整と core collar での一致を別々に確保する．
保持 annuli により $P\setminus K$ が分かれる場合は，
定理\ref{ng8:componentwise-critical-area}の成分別の配分条件を用いる．
系\ref{ng8:p0-small-area-flexibility}は source の全飽和近傍を保持した
微小変更を与えるが，任意の所定 target への到達を主張しない．
固定 compact と cuts を選んだ後に $t$ の閾値を取る構成でも，
この相対実現の条件そのものを対角選択で代用することはできない．
\end{remark}

\subsection{Critical regions の面積を保つ補間}

球面上の固定された marked Morse foliation の critical leaves は，
有限個の puncture basins と，内部 Reeb edges に対応する annuli に球面を分ける．
以下では，これらをまとめて critical regions と呼ぶ．
各 critical leaf 自体の面積は零である．
境界付き曲面を用いる場合は，puncture caps を補助的に付けて同じ記述を用いてよい．
caps と保存する puncture collars の和を，互いに素な固定 disks の内部に取る．
この固定領域を $K$ と書く．その上の foliation と面積形式の parameter family は
あらかじめ与えておき，以下の操作では変更しない．
面積形式が parameter に依存しないことを要求しているわけではない．

\begin{lemma}[Parameter に相対的な critical-area 指定]
\label{ng8:parametric-moser}
$u\in[u_-,u_+]$ に対して，球面上の正の面積形式 $\omega_u$ と，
固定 puncture collars を持つ marked Reeb tree を与える．
両端の open intervals には，同じ相対 isotopy class に属する
Morse foliation germs が指定されているとする．

各 critical region $D_h$ の target area $a_h(u)>0$ が滑らかで，
その和が $\int_C\omega_u$ に等しく，変更しない全領域の面積に対して
\begin{equation}\label{ng8:protected-area-margin}
 a_h(u)>\int_{D_h\cap K}\omega_u
\end{equation}
を満たすとする．$K$ には caps だけでなく，保存する collars 全体を含める．
両端では指定された germs がこの面積を持つと仮定する．
このとき，端の germs と collars を保ち，
全 $u$ で $\operatorname{area}_{\omega_u}(D_h)=a_h(u)$ を満たす
滑らかな marked foliation family を構成できる．
追加の compact parameters を用いる場合は，初期の marked interpolation と
region identifications がそれらに関しても滑らかに与えられていると仮定する．
その仮定の下で，面積指定と Moser の操作も parameter に相対的に行える．
\end{lemma}

\begin{proof}
まず補題~\ref{neck:marked-interpolation} により，marked foliations を補間する．
基準 foliation と滑らかな diffeomorphisms を使って各 region を同定すれば，
それぞれの面積は parameter に関して滑らかである．

各 region 内の，固定領域 $K$ を避けた部分に非負 bump functions を取る．
正の density を掛けて面積を指定する方法は補題~\ref{per:lobe-areas} と同じである．
\eqref{ng8:protected-area-margin} と parameter interval の compactness により，
少し大きい固定近傍でも一様な面積の余裕がある．
その近傍では density を一に保ち，それ以外でまず十分小さくした後，
各 region の内部に支持された bump を加える．
局所的には，region ごとの bump coefficient が面積の独立な座標となるので，
正値性を保つ滑らかな解を取れる．
同じ target integrals を持つ正の面積形式の集合は凸であるから，
parameter 空間上の partition of unity で局所解を貼れる．
両端では元の形式が既に条件を満たしているので，そこでは変更しない．
こうして，固定 collars 上で $\omega_u$ に一致し，
指定された region integrals を持つ正の形式 $\widehat\omega_u$ を得る．

両形式の差は $K$ の近傍で零であり，その総積分も零である．
固定 disks の補集合は連結なので，その interior の compactly supported
二次 de Rham cohomology は積分により $\RR$ と同定される．
従って差には，$K$ の近傍で零となる primitive がある．
同じことは，境界付き曲面の記述では $H^2(P,\partial P;\RR)\simeq\RR$ から従う．
固定した相対 de Rham 解作用素を用いて primitive を $u$ に滑らかに選び，
差が零である両端の intervals では primitive も零に取る．
凸結合に Moser の方法を適用して
\[
 \psi_u^*\widehat\omega_u=\omega_u
\]
を満たす diffeomorphisms を得る．
端の intervals および fixed collars では $\psi_u=\id$ にできる．
最初に補間した $u$ における foliation を $\psi_u^{-1}$ で移せば，
元の $\omega_u$ に関して所期の面積を持つ．
全ての操作は parameter に滑らかで，Morse 型と marked Reeb tree を保つ．
\end{proof}

\begin{proposition}[Straight critical actions を保つ actual neck]
\label{ng8:straight-neck}
actual Hamiltonian neck の両端の fibration germs を与える．
両 germ の reduced foliations は，同じ実際の reduced surface 上で
boundary collars を固定する $\operatorname{Diff}_0$ の同じ orbit に属するものとする．
従って同じ embedded marked Reeb tree を持ち，
共通の cycle framings と toric reference actions で書いた critical actions が
同じ affine functions
\[
 J_{\beta_\ell}^{\rm crit}=c_\ell+p_\ell\mu,\qquad p_\ell\in\ZZ
\]
の制限であると仮定する．
これらから式~\eqref{ng8:relative-action} で指定される
basin および internal-annulus areas が正であり，
\eqref{ng8:protected-area-margin}の固定領域に対する余裕を持つとする．
このとき，全 interval にわたり critical actions が上の affine functions となる
Lagrangian fibration を，両端の open germs を保って構成できる．
\end{proposition}

\begin{proof}
boundary reference actions は affine なので，
target region areas は $\mu$ の affine functions である．
その和と total reduced area は，補題~\ref{ng8:relative-dh} により共に affine で，
両端で一致する．従って interval 全体で一致する．
正の target areas には compact interval 上で一様な正の下限がある．
puncture cuts と保存する collars は最初に十分小さく選び，
各 basin の target area から固定領域の面積を引いた値に一様な正の余裕を持たせる．
内部 annulus の target area の正値性は，命題の仮定として保持する．

補題~\ref{ng8:parametric-moser} を reduced surfaces に適用する．
端で一致する marked foliation family を得て，
その circle lifts の connected-fiber quotient を取る．
補題~\ref{neck:lift} により正則 fibers は Lagrangian であり，
critical fibers は $S^1\times(S^1\vee S^1)$ となる．
式~\eqref{ng8:relative-action} により critical actions は所期の affine functions である．
三重点の端では toric collars を固定しているので，既存の triple blocks と一致する．
\end{proof}

\subsection{中央成分の面積計算から nearby fiber への判定}

\begin{corollary}[中央 fiber による十分条件]
\label{ng8:central-criterion}
両成分の reduced foliations が同じ実際の曲線上の共通の
$\operatorname{Diff}_0(P\operatorname{ rel\ collars})$-orbit に属するとする．
両成分の collar germs で critical actions が整数勾配の affine functions であり，
共通の neck framings へ移した後の勾配が一致するとする．
さらに $\mu=0$ における strict double curve 上で，
対応する critical regions の実際の面積が両側で一致するとする．
このとき，十分小さい nearby fiber で命題~\ref{ng8:straight-neck} の接着ができる．
\end{corollary}

\begin{proof}
中央成分上では，$\mu=0$ で normal circle が潰れ，
式~\eqref{ng8:relative-action} は strict double curve 上の area difference になる．
framing の整数変更による $m\mu$ はこの点で零である．
三重点の toric reference actions は共通に固定されているので，
実際の basin areas の一致は critical action functions の定数項の一致を与える．
内部 annuli の作用差も，同じ面積公式で一致する．
Reeb graph は tree なので，external reference から各 critical block をたどり，
block 内の cycle relations と内部 annulus の作用差を用いれば，
全 critical actions の定数項が共通に定まる．
この過程で現れる整数 framing changes は $\mu=0$ では寄与しない．
勾配の仮定と合わせ，両側は同じ affine functions となる．

中央 curve の各 critical region の面積は正である．
puncture cuts と固定 collars を先に十分小さく取り，その後で $\delta>0$ を小さくすれば，
対応する affine target areas は $[-\delta,\delta]$ 全体で正となり，
\eqref{ng8:protected-area-margin}も満たす．

次に固定した正の切断値 $\mu=\delta$ と負の切断値 $\mu=-\delta$ および
その open buffers だけを nearby fiber へ運ぶ．
補題~\ref{ng8:transport-intercept} が同じ affine functions を保存する．
これで有限 $t$ における加法定数も一致した．
残りの moment interval には命題~\ref{ng8:straight-neck} を適用する．

この議論は，中央 fiber の面積一致から有限 $t$ の一致を無条件に推論するものではない．
固定 moment 上の相対 symplectic transport と toric reference normalization が，
その二つの計算を接続している．
\end{proof}


\subsection{固定領域の補集合ごとの面積条件}
\label{ng8:componentwise-area}

元の補題\ref{ng8:parametric-moser}では，caps と puncture collars を含む
固定 disks の補集合が連結であることを使っている．ここでは，source の全葉の
近傍などを追加して固定し，補集合が非連結となる場合を扱う．
新しい fixed set の下で初期の marked interpolation を作れることは，
面積条件とは別の入力である．特に，抽象的な labelled Reeb tree の一致を
$\operatorname{Diff}_0(P\operatorname{ rel}K)$ による同定の代わりには用いない．

$P$ を compact oriented smooth surface とし，$K\subset P$ を smooth boundary を
持つ固定 compact codimension-zero submanifold とする．$\partial P\ne\varnothing$
の場合はその collar を $K$ に含める．
\[
 P\setminus K=\coprod_{\alpha=1}^q U_\alpha
\]
は有限個の connected open surfaces の和とする．$K$ の各成分が disk であるとは
仮定しない．$\operatorname{Diff}_0(P\operatorname{ rel}K)$ は，$K$ のある近傍で
恒等な isotopy によって恒等写像と結ばれる diffeomorphisms を表す．
従ってその各元は各 $U_\alpha$ をそれ自身へ写す．

Parameter space $B$ は compact smooth manifold with corners とする．
$B$ 上の finite positive measures $\nu_b$ と，smooth な初期 marked
foliation family $\mathcal F_b^0$ を与える．各 $U_\alpha$ 上では
$\nu_b=\omega_b$ は正の smooth area form であり，固定した自由な boundary collars
を含む $\overline U_\alpha$ 上で $b$ にも smooth と仮定する．
$K$ の内部の measure が smooth area form であること，及びその parameter smoothness
は仮定しない．$\nu_b(\partial K)=0$ とする．
Critical leaves と境界の和を $\Sigma_b$ と書き，その補集合の labelled critical
regions を $D_1(b),\ldots,D_H(b)$ とする．$\nu_b(\Sigma_b)=0$ と仮定する．
各 $b_0\in B$ の近傍で，$K$ と各 $U_\alpha$ を保ち，
$\Sigma_{b_0},D_h(b_0)$ を $\Sigma_b,D_h(b)$ へ写す smooth diffeomorphism family
が与えられていると仮定する．これは region と固定領域の incidence の
実際の smooth trivialization であって，その抽象的な個数の一致だけではない．
$K$ 上の measure 又は foliation が $b$ に依存しないことは要求しない．
変更後は，各 $b$ で与えられた $K$ 上のデータをそのまま保つ．

\[
 E=\{(h,\alpha):D_h(b)\cap U_\alpha\ne\varnothing\},\qquad
 \kappa_h(b)=\nu_b(D_h(b)\cap K),\qquad
 M_\alpha(b)=\int_{U_\alpha}\omega_b
\]
と置く．$D_h(b)$ と $U_\alpha$ は開なので，$E$ の条件は交わりの内部が
非空であることと同じである．上の trivializations により $E$ は固定され，
$M_\alpha$ と各 free cell の積分は smooth である．
$\kappa_h$ 自体の smoothness は以下で使わない．

\begin{lemma}[固定支持を持つ primitive の成分別構成]
\label{ng8:componentwise-primitive}
$\eta_b$ を $\coprod_\alpha U_\alpha$ 内の一つの compact set に支持された
smooth family of two-forms とする．$K$ の近傍で零となる smooth family of
one-forms $\beta_b$ で $d\beta_b=\eta_b$ を得るための必要十分条件は
\[
 \int_{U_\alpha}\eta_b=0\qquad(1\le\alpha\le q)
\]
である．Primitive は $\eta_b$ に線型に選べ，従って $\eta_b=0$ となる
parameter 領域では $\beta_b=0$ とできる．支持も一つの compact set に収まる．
\end{lemma}

\begin{proof}
必要性は compact support に対する Stokes の定理による．十分性は各成分内で
示せばよい．支持を覆う有限個の oriented coordinate rectangles と，それらを
一つの基準 rectangle へ結ぶ有限本の経路を選ぶ．経路も有限個の重なる rectangles
で覆い，partition of unity によって $\eta$ を局所 forms に分ける．
各局所 form から，その積分と同じ質量を持つ固定 bump two-form を引くと，
積分零の局所 form を得る．残った bump masses は，重なる rectangles 内の
積分一の bumps を順に用いて基準 rectangle へ運べる．総積分が零なので，
基準 rectangle に最後の質量は残らない．

ここで必要な局所 primitive を明記する．Rectangle 内で
$\eta=f(x,y)\,dx\wedge dy$，$\int f=0$ とし，$x$ 方向の compactly supported
function $\rho$ を $\int\rho=1$ となるよう固定する．
\[
 g(y)=\int_{\mathbb R} f(x,y)\,dx,\qquad
 F(x,y)=\int_{-\infty}^{x}\bigl(f(s,y)-\rho(s)g(y)\bigr)\,ds,
 \qquad
 G(y)=\int_{-\infty}^{y}g(s)\,ds.
\]
すると $F,G$ は compactly supported であり，
\[
 \beta=F\,dy-\rho G\,dx,
 \qquad d\beta=f\,dx\wedge dy.
\]
重なる bumps の差にも同じ式を用いる．全ての選択は固定でき，演算は線型な
積分・乗算と有限和だけである．従って smooth parameter dependence と固定支持を
同時に得る．各 $U_\alpha$ の結果を足し，$K$ の近傍へ零で延長する．
\end{proof}

\begin{theorem}[有限面積配分と parameter-relative な実現]
\label{ng8:componentwise-critical-area}
上の設定で target areas $a_h(b)>0$ を与え，
$r_h(b)=a_h(b)-\kappa_h(b)$ が smooth であると仮定する．
$a_h$ と $\kappa_h$ が個別に smooth である必要はない．
初期 family の各 slice と同じ $\operatorname{Diff}_0(P\operatorname{ rel}K)$-orbit
に属し，$K$ 上の全 foliation germs を保ち，
\[
 \nu_b(D_h^{\rm new}(b))=a_h(b)
\]
となる smooth marked foliation family が存在するための必要十分条件は，
各 $b$ に対して有限行列 $x_{h\alpha}(b)$ が存在し，
\begin{align}
 x_{h\alpha}&>0&&((h,\alpha)\in E),&
 x_{h\alpha}&=0&&((h,\alpha)\notin E),\label{ng8:transport-support}\\
 \sum_\alpha x_{h\alpha}&=a_h-\kappa_h&& (1\le h\le H),&
 \sum_h x_{h\alpha}&=M_\alpha&& (1\le\alpha\le q)
 \label{ng8:transport-margins}
\end{align}
を満たすことである．ここで存在は，初期 marked family とその smooth
trivializations が既に与えられた範囲での主張である．

より具体的には，この条件の下で smooth $\psi_b\in
\operatorname{Diff}_0(P\operatorname{ rel}K)$ と smooth positive densities
$\rho_b$ を構成でき，$\widehat\nu_b=\rho_b\nu_b$ と置くと
\[
 \psi_b^*\widehat\nu_b=\nu_b,
 \qquad \mathcal F_b^{\rm new}=\psi_b^{-1}(\mathcal F_b^0)
\]
となる．$\rho_b=1$ を $K\cup\Sigma_b$ の近傍に要求してもよい．
この最後の条件は density の条件であり，$\psi_b$ が $\Sigma_b$ を固定するという意味ではない．
Free part 上では $\widehat\omega_b=\rho_b\omega_b$ が smooth positive area forms である．

初期 family が指定された open end germs と一致し，そこで
$a_h(b)=\nu_b(D_h(b))$ であるなら，その germs を保って構成できる．
より一般に，この等式が成立する開集合 $O\subset B$ と
$O_0\Subset O$ を指定すれば，$O_0$ 上で $\psi_b=\mathrm{id}$ とできる．
追加の compact parameters についても上の smooth trivializations が与えられていれば，
同じ結論が同時に成立する．
\end{theorem}

\begin{proof}
\emph{必要性．}
ある $\psi\in\operatorname{Diff}_0(P\operatorname{ rel}K)$ が初期 marked partition
を新しいものへ移すとする．$\psi$ は各 $U_\alpha$ を保ち，$K$ 上では恒等なので，
各 $D_h^{\rm new}\cap U_\alpha$ の非空性は初期 cell の非空性と同じであり，
$D_h^{\rm new}\cap K=D_h\cap K$ である．
\[
 x_{h\alpha}=\int_{D_h^{\rm new}\cap U_\alpha}\omega
\]
と置けば，非空 open cell の面積は正であり，row と column の和は
\eqref{ng8:transport-margins}になる．
従って，この必要性は指定された relative orbit の中での主張である．
別の marked isotopy class の領域が同じ incidence を持つとは主張していない．

\emph{Pointwise feasibility から smooth な配分へ．}
非空 entry だけを並べ，row と column の和を取る固定線型写像を $A$ と書く．
右辺を $y(b)$ とする．ある $b_0$ で正の解 $x^0$ を取り，
$\operatorname{im}A$ 上の線型 right inverse $R$ を固定する．
Pointwise feasibility により $y(b)-y(b_0)\in\operatorname{im}A$ なので，
\[
 x(b)=x^0+R\bigl(y(b)-y(b_0)\bigr)
\]
は $b_0$ の近傍で smooth な正の解である．解集合は凸だから，有限 cover と
parameter partition of unity により global smooth positive solution を得る．
$O$ 上では actual cell areas $\int_{D_h(b)\cap U_\alpha}\omega_b$ が正の解なので，
これを $O_0$ の近傍で選べる．

\emph{正の density の構成．}
まず一つの小さい parameter neighborhood を取り，指定された trivialization で
全 regions と $K$ を固定する．各非空 cell に一つずつ compactly supported
nonnegative bump $\varphi_{h\alpha,b}$ を取り，
\[
 \int_{D_h(b)\cap U_\alpha}\varphi_{h\alpha,b}\,\omega_b=1
\]
と正規化する．各 bump は $K$ と critical leaves から離して取る．

$K$ の外側の十分薄い近傍と $\Sigma_b$ の十分薄い近傍の，どの cell における
面積も任意に小さくできる．前者は $U_\alpha$ 内で $K$ へ収束する部分の面積が
零へ向かうこと，後者は $\Sigma_b$ の面積が零であることによる．
Smooth trivialization と parameter compactness により，これらの選択は
局所 parameter neighborhood 上で一様に行える．
従って，$K\cup\Sigma_b$ の近傍で $1$，少し大きい近傍の外で小さい正数
$\varepsilon$ となる smooth positive function $r_b$ を取り，
\[
 m_{h\alpha}(b):=
 \int_{D_h(b)\cap U_\alpha}r_b\,\omega_b
 <x_{h\alpha}(b)
 \quad ((h,\alpha)\in E)
\]
とできる．この際 bumps の支持を $r_b=\varepsilon$ の領域に置ける．
\[
 \rho_b=r_b+
 \sum_{(h,\alpha)\in E}
 \bigl(x_{h\alpha}(b)-m_{h\alpha}(b)\bigr)\varphi_{h\alpha,b}
\]
は正で，$K\cup\Sigma_b$ の近傍では $1$，各 cell の積分は指定された
$x_{h\alpha}$ である．$\widehat\omega_b=\rho_b\omega_b$ と置く．
$K$ 上では density を $1$ に保つので，row 和と
$\kappa_h+r_h=a_h$ から $\widehat\nu_b(D_h(b))=a_h(b)$ となる．
Column 和から，各 $U_\alpha$ の総積分は $M_\alpha$ になる．

この構成を有限 parameter cover 上で行い，正の densities の凸性により貼る．
Constraints は線型なので全積分を保つ．$O$ 上には $\rho_b=1$ を局所解として
用い，partition of unity を $O_0$ の近傍でこの解だけになるよう選ぶ．
また有限個の近傍を縮めれば，$K$ の一つの共通近傍上で全ての densities が
$1$ になる．従って $\widehat\omega_b-\omega_b$ の支持は
$P\setminus K$ の一つの compact set に収まる．

\emph{Componentwise relative Moser．}
Column 和により
\[
 \int_{U_\alpha}(\widehat\omega_b-\omega_b)=0
\]
なので，前補題から smooth $\beta_b$ を
$d\beta_b=\widehat\omega_b-\omega_b$，$\beta_b=0$ near $K$ と選べる．
$O_0$ 上では $\beta_b=0$ とする．
Free part 上の $\omega_{b,s}=(1-s)\omega_b+s\widehat\omega_b$ は正なので，
\[
 \iota_{X_{b,s}}\omega_{b,s}=-\beta_b
\]
は smooth vector fields を定める．これらは $K$ の共通近傍で零であり，
compactness により $0\le s\le1$ の全 flow が存在する．
その time-one map $\psi_b$ は smooth で，
free part 上で $\psi_b^*\widehat\omega_b=\omega_b$ となる．
$K$ 上では恒等なので，measure として全 $P$ 上で
$\psi_b^*\widehat\nu_b=\nu_b$ となる．
Isotopy 全体が $K$ の近傍で恒等なので relative marking を保つ．
また $O_0$ 上では全 flow が恒等である．
従って $D_h^{\rm new}(b)=\psi_b^{-1}(D_h(b))$ は所望の面積を持つ．
Morse 型，full critical leaf の型，region labels は diffeomorphism で保存される．
\end{proof}

\begin{example}[全体の総量と protected-area margin だけでは足りない]
\label{ng8:annular-protection-obstruction}
球面上に四つの puncture basins $D_1,\ldots,D_4$ と一つの internal annulus を持つ
marked Morse foliation を固定し，正の面積形式 $\omega$ を固定する．
各 puncture の小さい cap disk を保護する．さらに $D_1$ の regular circle の
小さい飽和 annulus $A\Subset D_1$ を固定し，これらの和を $K$ とする．
$A$ と $D_1$ 内の cap の間の open annulus $U_0$ は $P\setminus K$ の
一つの連結成分であり，その全体が $D_1$ に含まれる．
\[
 \kappa_1=\int_{D_1\cap K}\omega,\qquad
 M_0=\int_{U_0}\omega>0
\]
と置く．$\operatorname{Diff}_0(P\operatorname{ rel}K)$ の中で foliation を変更しても，
$U_0$ 全体は第1 basin 内に残る．従って必ず
\[
 \operatorname{area}_{\omega}(D_1^{\rm new})\ge\kappa_1+M_0
\]
である．

ところが target $a_1=\kappa_1+M_0/2$ は個別の strict margin
$a_1>\kappa_1$ を満たす．元の第1 basin から減らした面積を第2 basin に足し，
他の target areas を元のままにすれば，どの target も正で，それぞれの
protected-area margin と全体の総量一致を満たす．しかし上の不変な下界に反する．
Parameter interval の両端では元の target，中央でこの target となるよう，
非負 bump function による凸補間を取れば，open end germs を保つ smooth data
に対しても同じ障害が生じる．
これは元の disk-complement 版の反例ではなく，追加 annulus を固定する版への
無条件な拡張の反例である．
\end{example}

\begin{remark}[Full critical neighborhoods と actual 適用の範囲]
すべての full critical leaves を含む pants neighborhoods $Q_v$ まで固定し，
各 $U_\alpha$ が一つの critical region にしか接しないなら，
column 条件はその region に元の面積を全て割り当てる．この場合 target critical
areas は元の値に強制され，面積を変更する自由はない．
既存の source-marking isotopy が $Q_v$ を固定する事実から，将来の
critical-area 調整でも $Q_v$ 全体を固定できるとは結論しない．

本定理は，protected part の外側で上記の smooth reduced data と relative marked
interpolation が実際に得られる場合の実構成である．Source 固定点を含む singular
reduction 全体に smooth area-form family を仮定していない．
固定した nodal neighborhoods の外側に適切な compact free reduction を取り，
その components，incidence，free-cell masses の smoothness を確認すればよい．
例えば $a_h(b)=\nu_b(D_h(b))+\delta a_h(b)$ とすると，
$r_h(b)=\sum_\alpha\int_{D_h(b)\cap U_\alpha}\omega_b+\delta a_h(b)$
なので，$\delta a_h$ の smoothness だけでよく，protected mass の非滑らかさは相殺される．
ただしこの事実は target の有限輸送条件を自動的に保証するものではない．
また foliation の面積を実現する定理であって，異なる既製 symplectic pieces の
primitive periods 又は core collar germs を自動的に一致させる定理ではない．
\end{remark}


\subsection{面積配分の有限 cut 判定}
\label{ng8:area-cut-section}

前定理の正の配分条件は，有限個の線型不等式で検査できる．
これは max-flow/min-cut の既知の原理
\cite[Theorem~1, p.~400]{FordFulkerson1956}に対応する有限二部問題であり，
一般の組合せ原理の新規性は主張しない．以下では必要な有向の配分問題を直接証明する．

\begin{proposition}[Strict な面積輸送の cut 条件]
\label{ng8:strict-area-cuts}
前定理の incidence graph を，region vertices $h$ と可動成分 vertices $\alpha$
を結ぶ二部 graph $G$ とする．辺は $(h,\alpha)\in E$ のとき，かつそのときに限り置く．
各非孤立 region の要求量を $r_h=a_h-\kappa_h>0$，各可動成分の面積を
$M_\alpha>0$ とする．孤立 region は $r_h=0$ を要求する．
可動成分 vertex は孤立しない．
全ての辺で正の配分が存在するための必要十分条件は，次の二つである．
\begin{enumerate}
\item 各 connected component $G_0$ の region 集合 $H_0$ と可動成分集合
  $A_0$ について，
  \[
    \sum_{h\in H_0}r_h=\sum_{\alpha\in A_0}M_\alpha .
  \]
\item 各 $G_0$ と各 $\varnothing\ne S\subsetneq H_0$ について，
  $N(S)\subset A_0$ を隣接する可動成分の集合とすると，
  \begin{equation}\label{ng8:eq:strict-area-cut}
    \sum_{h\in S}r_h<
    \sum_{\alpha\in N(S)}M_\alpha .
  \end{equation}
\end{enumerate}
従って検査する cut は有限個である．$r_h,M_\alpha$ が compact parameter
space 上で連続で，graph が固定され，上の条件が全 parameter で成り立つなら，
全ての strict inequalities と正値性には一様な正の余裕がある．
\end{proposition}

\begin{proof}
必要性を示す．Component 内の row sums と column sums を足せば最初の等式を得る．
正の配分 $x$ に対して
\[
 \sum_{\alpha\in N(S)}M_\alpha-\sum_{h\in S}r_h
 =\sum_{\substack{h\notin S,\ \alpha\in N(S)\\(h,\alpha)\in E}}x_{h\alpha}.
\]
Connected component の proper nonempty row subset には，右辺に少なくとも
一つの辺がある．従って差は正である．

十分性には，まず非負配分の判定を証明する．正の margins $r,M$ と
component-wise な総和一致を仮定し，上の全 cut inequalities を非strict に
弱める．非負 entry $y_{h\alpha}$ で row sums は $r_h$ 以下，column sums は
$M_\alpha$ 以下となるものの compact polytope 上で，総和を最大にする $y$ を取る．
最大値が要求総量より小さければ，未充足の row から始める．Row から隣接する
column へは増量方向に進め，column から row へは $y_{h\alpha}>0$ の辺だけを
減量方向に進める．未充足の column に到達できれば，有限の alternating path
に沿って十分小さい正量を増減し，総和を増やせて矛盾する．

従って到達した row 集合 $S$ の全隣接 columns $N(S)$ は充足されている．
しかも $S$ 外の row から $N(S)$ への正の entry はない．あれば逆向きの辺で
その row に到達するからである．最初の row は未充足なので，
\[
 \sum_{\alpha\in N(S)}M_\alpha
 =\sum_{h\in S,\alpha}y_{h\alpha}
 <\sum_{h\in S}r_h
\]
となり，非strict cut 条件に反する．従って非負の完全配分が存在する．

今度は全辺へ同じ十分小さい $\epsilon>0$ を先に配り，
\[
 r'_h=r_h-\epsilon\deg(h),\qquad
 M'_\alpha=M_\alpha-\epsilon\deg(\alpha)
\]
を残りの margins とする．Component-wise な総和一致は保たれる．
Cut の余裕の減少は
$\epsilon\,\#E(H_0\setminus S,N(S))$ だけなので，有限個の strict inequalities
と正値性から，残りの margins は非負配分の条件を満たすようにできる．
その非負解 $y$ に $\epsilon$ を加えれば，全辺で正の解になる．
最後の一様性は有限個の正の連続関数の compactness による．
\end{proof}

\begin{remark}[端だけの検査で足りる場合]
\label{ng8:affine-area-cuts}
一つの compact interval 上で graph が固定され，
全ての $r_h(\mu),M_\alpha(\mu)$ が affine であるなら，
component-wise な等式と strict cut inequalities は両端で検査すれば足りる．
各 cut の差も affine だからである．
Source を固定した actual 構成では，free 部の面積が smooth でも affine とは
限らず，この端点判定を無条件に用いてはならない．
\end{remark}

\begin{example}[両端の面積配分だけでは足りない]
\label{ng8:interior-area-cut-obstruction}
一つの external basin $D_1$ 内に保持 annulus を置き，その内側の可動成分
$U_0$ が $D_1$ だけに接し，外側の可動成分 $U_1$ が三つの basins
$D_1,D_2,D_3$ に接するとする．Incidence edges は
$(1,0),(1,1),(2,1),(3,1)$ である．
固定領域の面積 $\kappa_h$ を加えた target を $a_h=\kappa_h+r_h$ とし，
$\mu\in[-1,1]$ 上で
\[
 (r_1,r_2,r_3)=\left(\tfrac12,\tfrac34,\tfrac34\right),\qquad
 M_0(\mu)=\tfrac34-\tfrac12\mu^2,\qquad M_1(\mu)=2-M_0(\mu)
\]
と置く．全て正で，総和は常に一致する．両端では
$x_{10}=\tfrac14,x_{11}=\tfrac14,x_{21}=x_{31}=\tfrac34$ が正の配分である．
一方 $\mu=0$ では $x_{10}=M_0=\tfrac34>r_1=\tfrac12$ となり不可能である．
同じことは $S=\{2,3\}$ の cut
$r_2+r_3<M_1$ の不成立として読める．

この面積 family は，三 punctures の球面で $D_1$ の regular arm に固定 annulus
を置き，$U_0,U_1$ の内部で正の densities のみを変えて実現できる．
固定領域とその境界近傍の面積を先に小さく取り，その外に支持した bump を用いる．
例えば初期葉層の自由部分の cell masses を
\[
 m_{10}=\tfrac34-\tfrac12\mu^2,\quad m_{11}=\tfrac14,\quad
 m_{21}=m_{31}=\tfrac12+\tfrac14\mu^2
\]
と指定すれば，各 column の和は上の $M_\alpha$ となり，
初期葉層自身が両端で target を満たす．
さらに全ての $\mu^2$ を，両端の開近傍で $1$，中央の開近傍で $0$ となる
滑らかな $h(\mu)\in[0,1]$ に置き換えれば，open end germs を満たす模型にもなる．
これは元の actual KE 族の反例ではなく，
smooth free-band data と両端の確認だけから全 interval の面積実現を
推論できないことを示す模型である．
\end{example}


\subsection{Source の全近傍を保つ actual neck の微小面積指定}

以下では一つの許容された stage と $t\ne0$ を固定する．
Normal circle と canonical moment は既存 P0 のものを用いる．
面積と第二作用を比較する場合には周期 $2\pi$ に規格化するが，
moment の向きと加法定数を選び直さない．
本節の定数又は moment interval に $t\to0$ での一様性を要求しない．

\begin{lemma}[Source の飽和 annuli と滑らかな自由部分]
\label{ng8:p0-source-protected-reduction}
定理\ref{ke:projected-source-neck}で構成した actual neck と，
その後の end collars の相対整合を考える．
必要なら $0$ を含む compact moment interval $J$ を小さく取ると，
source を通る相異なる正則閉葉の小さい飽和 annuli
$A_1,\ldots,A_s$ を，critical leaves と保持する end collars から離して取れる．
Moment に沿う marked identification を用いた記法で，これらを固定 annuli とする．
このとき次が成り立つ．
\begin{enumerate}
\item $J\times A_a$ の actual circle orbit map による逆像は，
source の full compact pinched torus とその飽和近傍を含む．
これらの全 fiber partition を保持することは，source 点だけの固定より強い条件である．
\item 保持する end collars と全 $A_a$ を含む固定領域を $K$ とする．
$P\setminus K$ 上の actual normal circle は自由であり，
その reduced area forms は $\mu\in J$ に滑らかに依存する．
これは $\mu=0$ を跨いで成立する．
$K$ 内の reduced form，又は $K$ を含む領域の面積が
$\mu=0$ で滑らかであることは主張しない．
\end{enumerate}
\end{lemma}

\begin{proof}
補題\ref{ke:source-reduction-pinched-fiber}では orbit space は
$(\mu,p)\in I_t\times P$ で，非自由軌道は有限個の $(0,q_a)$ に限られる．
各 $q_a$ は相異なる閉 regular leaves 上にあり，saddle leaf と境界から離れている．
その小さい regular annular neighborhoods を互いに素に選ぶ．
定理\ref{ke:projected-source-neck}の基準葉層は残余変数 $p$ の葉層である．
End collars の整合は少し広い境界 collars 内に支持を持つので，
ここで選んだ内部 annuli を避けられる．
既に moment に依存する marked 葉族を用いている場合にも，
補題\ref{ke:simple-saddle-family-trivialization}で $J$ を縮め，
source markings を固定する reduced identification を選べる．
この identification を source 点上で ambient smooth map に持ち上げる必要はない．
最後の変更は source の annuli 上で恒等になるからである．

$\mu=0$ で source を通る閉葉の全逆像は pinched torus であり，
その両側の regular leaves と小さい moment interval を含む逆像は
その全 fiber の飽和近傍である．これが第一の結論である．
全 source は $K$ の内部に置かれる．従ってその外の閉じた自由部分では
moment map は正則で，$(\mu,p)$ は滑らかな submersion である．
Actual form の reduction と上の reduced identification により，
自由部分の area forms は $\mu=0$ を含めて滑らかである．
Source 上の orbit-space 座標又は reduced form の滑らかさを用いていない．
\end{proof}

\begin{corollary}[P0 の source を固定した微小 critical-area 指定]
\label{ng8:p0-small-area-flexibility}
前補題の固定 $t,J$ と marked identification を用いる．
保持領域 $K$ は end collars と source regular annuli の和とし，
補助的な puncture caps を用いる場合はそれらも固定する．
Critical leaves の全近傍を $K$ に追加することは，この系では要求しない．
Annuli と collars は，各 critical region $D_h$ に自由な開部分が残るように選ぶ．

元の actual area を $a_h^0(\mu)$ とする．
ある $\epsilon>0$ が存在し，滑らかな関数 $\delta a_h:J\to\RR$ が
\begin{equation}\label{ng8:eq:p0-small-area-increments}
 \sum_h\delta a_h(\mu)=0,\qquad
 \max_h\|\delta a_h\|_{C^0(J)}<\epsilon
\end{equation}
を満たし，$J$ の両端の指定した open intervals 上で全て零なら，
critical-region areas を
\[
 a_h(\mu)=a_h^0(\mu)+\delta a_h(\mu)
\]
とする actual Lagrangian neck fibration を構成できる．
同じ canonical moment，normal circle，marked tree，end collars と端の open germs，
全 source の上記 compact nodal neighborhoods を保つ．
正則・saddle・source fibers の型と全 fiber partition の接着条件も保つ．

$a_h^0$，$a_h$ 又は $\int_{D_h\cap K}\omega_{t,\mu}$ の滑らかさは不要である．
滑らかさを要求するのは増分 $\delta a_h$ と自由部分の area forms だけである．
\end{corollary}

\begin{proof}
$U_\alpha$ を自由部分 $P\setminus K$ の連結成分とする．
Critical leaves は面積零であり，各 $D_h$ と $U_\alpha$ の内部の交わりが
非空であるとき，二部グラフ $G$ に辺 $(h,\alpha)$ を置く．
このグラフは連結である．実際，boundary collars を除いた曲面は連結である．
二点を結ぶ path を，critical leaves と source annuli の境界に横断的に選ぶ．
Critical leaf を横切るときは同じ $U_\alpha$ に留まるので，
その両側の region vertices は $U_\alpha$ を介して結ばれる．
Source annulus を横切るときは一つの regular critical region $D_h$ 内に留まるので，
その両側の自由成分は $D_h$ を介して結ばれる．
これは $G$ 内の path を与える．Annuli を十分小さく選んだので，全 region vertex が現れる．

辺ごとの元の自由面積を
\[
 m_{h\alpha}(\mu)=
 \int_{D_h\cap U_\alpha}\omega_{t,\mu}>0
\]
と置く．Marked identification の下で積分領域は固定され，
その閉包は source を避ける自由部分に含まれるので，$m_{h\alpha}$ は滑らかである．
有限グラフと compact な $J$ より
\[
 m_*:=\min_{(h,\alpha)\in E(G),\,\mu\in J}m_{h\alpha}(\mu)>0.
\]

$G$ の spanning tree を一つ固定する．総和零の vector $v=(v_h)$ に対し，
tree の葉から順に消去することで，辺上の vector $L(v)$ を線形に選び，
\[
 \sum_\alpha L(v)_{h\alpha}=v_h,
 \qquad \sum_h L(v)_{h\alpha}=0
\]
とできる．これは region vertices の需要を $v_h$，自由成分 vertices の需要を零とする
tree flow であり，最後の頂点の式は $\sum_hv_h=0$ から従う．
$\|L(v)\|_\infty\le C_G\|v\|_\infty$ と取る．
各非空 cell の内部に，$K$ と critical leaves を避けた非負の smooth bump
$\varphi_{h\alpha,\mu}$ を選び，その cell における $\omega_{t,\mu}$ に関する積分を
$1$ と正規化する．Marked identification と $J$ の compactness により
$B:=\max\{1,\sup_{h,\alpha,\mu}\|\varphi_{h\alpha,\mu}\|_\infty\}<\infty$ であり，
異なる cell の bumps は互いに素な支持を持つ．
そこで
\[
 \epsilon=\frac{\min\{m_*,B^{-1}\}}{2\max\{1,C_G\}}
\]
と取れば
\[
 x_{h\alpha}(\mu)=m_{h\alpha}(\mu)+L(\delta a(\mu))_{h\alpha}>0
\]
である．各 column の和は元の $U_\alpha$ の面積に等しく，
各 row の和だけが $\delta a_h$ だけ変わる．
端の open intervals 上では $x=m$ である．

この小摂動の場合には density 自体を明示できる．自由部分上で
\[
 \rho_\mu=1+\sum_{(h,\alpha)\in E(G)}
 L(\delta a(\mu))_{h\alpha}\,\varphi_{h\alpha,\mu},
 \qquad \widehat\omega_\mu=\rho_\mu\omega_{t,\mu}
\]
と置くと，$\rho_\mu\ge1/2$，各 cell の積分は $x_{h\alpha}$ であり，
各 $U_\alpha$ の面積差は零である．Density は $K$ と critical leaves の近傍で $1$，
また $\delta a=0$ となる全 parameter 領域で恒等的に $1$ である．
補題\ref{ng8:componentwise-primitive}により差の primitive を
$P\setminus K$ に compact support を持つよう選び，その自由成分ごとに相対 Moser を行う．
これは定理\ref{ng8:componentwise-critical-area}の明示的な微小版である．
Source を含む $K$ 内の形式を変更又は滑らかに延長する必要はない．
得られる reduced diffeomorphisms は $K$ の開近傍および
$\delta a=0$ となる全 parameter 領域で恒等となる．
移された葉層では，$K$ 内の元の面積はそのままで，自由部分の面積だけが
$\delta a_h$ だけ変わる．これが所期の等式を与える．

最後に circle lift を取る．各 regular lifted leaf の接空間は normal generator と
$\ker d\mu$ 内の一方向で張られるから，$\iota_V\omega_t=-d\mu$ により Lagrangian である．
Saddle graph の lift は元と同じ isotropic strata を持ち，source の全近傍は恒等である．
相対 Moser isotopy は自由部分に compact support を持つので，
principal circle bundle の connection で水平に持ち上げ，$K$ 上で零に延長できる．
これにより moment と circle action を保つ ambient diffeomorphism が得られ，
葉の全 partition を移す．この diffeomorphism の symplectic 性は要求しない．
Lagrangian 性は上の reduction で直接確認済みである．
従って marked quotient，connected fibers と compact Hausdorff 性も保たれる．
全 end collars および moment の両端で全葉が変わらないので，
既存 corner・core との whole-fiber overlaps はそのまま接着する．
\end{proof}

\begin{remark}[局所自由度と残る atlas 条件]
\label{ng8:p0-area-scope}
この系は，固定した actual KE neck に対する非自明な相対変形を実際に与える．
有限個の neck に同様の変更を行っても，既存の corner と core を変更せず接着できる．
同じ normal moment と全保持領域を保つため，
命題\ref{ke:moment-relative-replacement}の基底値比較にも適合する．
ただし $\epsilon$ は固定 stage・$t$・保護領域に依存してよい．

全 critical actions を指定した affine functions へ変えるためには，
まず primitive と cycle framings を固定し，source の winding jump を含めて
それに対応する \emph{relative area increments} を求めなければならない．
その増分が上の小ささ又は一般の有限輸送条件を満たすか，
端の open germs と一致するかは別の問題である．
全 interval の absolute areas を affine と仮定することも結論することもない．
Source の固定領域に含まれる面積の非滑らかさは，元と新で同じ寄与として保持される．
従ってこの微小実現を，五次・六次の同じ射影の atlas と GH の同時実現，
任意の指定 critical actions の実現，又は退化 parameter に一様な面積予算とはしない．
\end{remark}

\subsection{球面 pants 木における一意な面積配分}
\label{ng8:spherical-tree-area}

本項では，系\ref{ng8:p0-small-area-flexibility}の保持集合を具体的に用いる．
$P$ は球面から有限個の puncture disks を除いた連結な compact surface であり，
全 critical leaves は一つずつの simple saddle を持つとする．
その connected-leaf quotient を $q:P\to T$ と書くと，$T$ は pants tree である．
Critical regions $D_h$ は $T$ の開辺 $e_h$ に対応する．
補助的な puncture caps を付ける場合も，その cap と end collar は固定する．

保持集合 $K$ は end collars と，critical leaves を避ける互いに素な閉飽和
regular annuli $A_1,\ldots,A_m$ の和とする．各 $A_s$ は一つの $D_h$ に含まれ，
$q(A_s)$ は $e_h$ 内部の閉 interval である．Annuli 同士，critical leaves，
および end collars の間には開 collar の余裕を残す．同じ辺上の保護 intervals が
重なる場合には，先にその和の各連結成分へまとめ，その後の個数を $m$ と数える．
Source の複数の保護円が球面上で入れ子になることは許す．

\begin{proposition}[Source-annulus incidence graph の木性]
\label{ng8:spherical-incidence-tree}
$P\setminus K$ の連結成分を $U_\alpha$，critical regions の個数を $H$ とし，
$D_h\cap U_\alpha\ne\varnothing$ のときだけ辺 $(h,\alpha)$ を置く二部グラフを
$G$ とする．このとき $G$ は木である．さらに，$D_h$ が $n_h$ 個の保護 annuli を
含むなら，$D_h\setminus K$ の $n_h+1$ 個の連結な free pieces は相異なる
$U_\alpha$ に属し，
\[
 \deg_G(h)=n_h+1,\qquad
 \#\{U_\alpha\}=m+1,\qquad
 |E(G)|=H+m=|V(G)|-1.
\]
したがって，一つの cell $D_h\cap U_\alpha$ 自体も連結である．
$P$ がこのような球面成分の非交和なら，成分ごとの同じ主張により forest を得る．
\end{proposition}

\begin{proof}
End collars を除く操作は tree の端を短くするだけで，その連結性を変えない．
一つの保護 annulus は，$T$ の辺内部の閉 interval を取り除く操作に対応する．
Tree には別の迂回路がないので，この操作は自由部分の連結成分をちょうど一つ増やす．
逆像の連結性は，connected fibers を持つ closed quotient $q$ に対する
補題\ref{topo:connected-preimages}から従う．ここでは open subtree の逆像にも
同補題を用いられる．従って保護 annuli を順に除くと自由成分数は $m+1$ となる．
球面上の各 regular circle が separator であることによっても同じことが分かる．

一つの開辺 $e_h$ に $n_h$ 個の保護 intervals があれば，残る部分は
$n_h+1$ 個の非空 intervals である．二つの異なる intervals が同じ自由成分内で
再会すると，それらを隔てる保護 interval を避けて $T$ 内で結ぶ path ができる．
これは tree の path の一意性に反する．各 interval の逆像は一つの regular
annulus 又はその端を切った部分であるから，必要な cell の連結性も従う．

$G$ の連結性も確認する．End collars の外の連結な曲面内で path を取り，
critical vertices を避け，critical leaves と保護 annuli の境界に横断的にする．
Critical leaf を横切るときは同じ $U_\alpha$ に留まるため，二つの region labels
はその $U_\alpha$ を介して結ばれる．保護 annulus を横切るときは一つの $D_h$ 内に
留まるため，両側の自由成分はその region vertex を介して結ばれる．
このように $G$ 内の path を得る．各 region と自由成分は非空 open set を含むので，
全 vertices がこの議論の対象となる．
最後に
\[
 |E(G)|=\sum_h(n_h+1)=H+m,\qquad |V(G)|=H+m+1
\]
である．連結グラフのこの等式から $G$ は木となる．
\end{proof}

\begin{corollary}[Edge cut による一意配分と全 target の判定]
\label{ng8:tree-edge-area-test}
前命題の保持集合を固定し，定理\ref{ng8:componentwise-critical-area}と同じ記法で
\[
 r_h=a_h-\kappa_h,\qquad M_\alpha=\nu(U_\alpha)>0
\]
と置く．$e=(h,\alpha)$ を取り除いた $G$ の二つの連結成分のうち，
region vertex $h$ を含む方を $C_e$ と書き，
\begin{equation}\label{ng8:eq:tree-edge-flow}
 X_e(a)=
 \sum_{k\in C_e\cap H}(a_k-\kappa_k)
 -\sum_{\beta\in C_e\cap U}M_\beta
\end{equation}
と定める．ここで $H,U$ は二種類の vertex 集合も表す．
所定 target を同じ marked relative orbit の中で実現できるための必要十分条件は
\begin{equation}\label{ng8:eq:tree-edge-inequalities}
 \sum_h(a_h-\kappa_h)=\sum_\alpha M_\alpha,
 \qquad X_e(a)>0\quad(e\in E(G))
\end{equation}
である．可解な場合の全 cell masses は一意であり，$x_{h\alpha}=X_{(h,\alpha)}(a)$
となる．これは mass allocation の一意性であって，density 又は foliation の一意性ではない．

Parameter family では，固定した incidence と smooth free data の下で，
$r_h,M_\alpha$ が smooth なら配分も上式により smooth である．
両端の target が元の area と一致する場合，配分もそこで元の cell masses と一致する．
従って同じ relative Moser theorem をそのまま適用できる．
$\kappa_h$ 個別の smoothness，source 内部での smooth reduced form，
又は新たな whole-fiber 仮定は必要ない．
\end{corollary}

\begin{proof}
各辺を region vertex から free-component vertex へ向け，vertex の signed balance を
$b(h)=r_h$，$b(\alpha)=-M_\alpha$ とする．
$C_e$ 内の vertex equations を足すと，内部辺は相殺し，境界に残る辺は $e$ 一つだけである．
その向きは $C_e$ の外向きなので，必ず $x_e=\sum_{v\in C_e}b(v)=X_e(a)$ となる．
これで一意性と必要性を得る．

逆に total balance を仮定して $x_e=X_e(a)$ と置く．
Region vertex $h$ を除くと，その各隣接辺の反対側の成分は全ての他の vertices を
分割する．Total balance により，それらの負の balances の和は $b(h)=r_h$ である．
従って $\sum_\alpha x_{h\alpha}=r_h$．
Free vertex $\alpha$ については，隣接する各 $C_e$ が他の vertices を分割するので，
$\sum_hx_{h\alpha}=-b(\alpha)=M_\alpha$ となる．
全 $X_e>0$ なら一般の正の有限配分条件を満たし，相対 Moser の構成が適用できる．
Parameter に関する主張は有限和の表示と同定理から従う．
\end{proof}

これは一般の strict Hall cut 判定のこの tree への特殊化である．
検査する不等式は $H+m$ 個だけでよく，一つの root を選んで subtree balances を
末端から加えることで全て求められる．Forest の場合は各 connected component に
一つずつ total balance を要求し，同じ edge formulas を用いる．
$r_h,M_\alpha$ が interval 上で affine なら端での検査で足りるが，
source 保持下でそれらが smooth であるだけの場合に端点検査へ置き換えることはできない．

\subsection{複数の保護 annuli と累積面積の制約}

元の cell masses を $m_e>0$，critical-region areas を $a_h^0$ とし，
$a_h=a_h^0+\delta a_h$，$\sum_h\delta a_h=0$ とする．前系の式は
\begin{equation}\label{ng8:eq:tree-increment-flow}
 x_e=m_e+\sum_{h\in C_e\cap H}\delta a_h
\end{equation}
となる．従って全可解増分は，零総和の hyperplane 内で，有限個の明示的な
strict linear inequalities によって切り出される開凸集合である．
この集合は零を含むが，任意の大きい増分を含むとは限らない．

同一 region $h$ に $n\ge1$ 個の annuli を辺に沿って並べる．
両端の free cells を $C_{h,-},C_{h,+}$，その元の面積を $m_{h,-},m_{h,+}$ とする．
中間の $n-1$ 個の cells は，それ自体が自由成分であり，$G$ では degree-one
free vertices になる．従ってその面積は，いかなる許容された変更でも元の値に固定される．
両端の cells は相異なる自由成分に属する．

元の pants tree から辺 $e_h$ の内部を除いて得る両側の region-edge 集合を
$H_h^-,H_h^+$ とする．これらは $h$ を含まず，
$H_h^-\coprod H_h^+=H\setminus\{h\}$ である．
$-$ 側の自由成分への incidence edge を切ったとき，$h$ 側に残る regions は
$\{h\}\cup H_h^+$ だから，\eqref{ng8:eq:tree-increment-flow}と零総和条件より
\begin{equation}\label{ng8:eq:two-end-cells}
 x_{h,-}=m_{h,-}-\sum_{k\in H_h^-}\delta a_k,
 \qquad
 x_{h,+}=m_{h,+}-\sum_{k\in H_h^+}\delta a_k.
\end{equation}
所定 target の制約は，これら二つの量の strict positivity である．
中間の cells については $x=M>0$ が自動的に成り立つ．
保護 annulus がない region では一つの free cell だけが存在し，
対応する条件は $a_h>\kappa_h$ となる．

\begin{corollary}[External basin の入れ子と internal annulus の二端条件]
\label{ng8:ordered-source-area-cuts}
External basin の保護 annuli を critical leaf 側から puncture 側へ並べる．
Critical leaf 側の最初の保護 annulus から puncture までの全固定 tail の面積を
$L_h$ とする．これは保護部分だけでなく，その向こう側に隔離された全自由成分の
面積を含む．この basin に対応する条件はちょうど
\[
 a_h>L_h
\]
である．より内側に annuli を追加しても，この同じ tail 内を細分するだけなら
region 総面積について新しい不等式は生じない．

Internal region では，二つの critical leaves に最も近い保護 annuli の外側にある
両端 cells が\eqref{ng8:eq:two-end-cells}の二つの条件を与える．
その間に annuli を追加しても，中間の面積がより細かく固定されるだけで，
region 総 target に対する条件は両端の条件のままである．
ただし異なる region 内の入れ子については，両側 subtree に含まれる
\emph{全ての} target 増分を式に入れなければならない．
個々の basin の面積下限だけで internal edge の条件を置き換えてはならない．
\end{corollary}

\begin{proof}
External edge の puncture 側には他の region edge がないので，
$H_h^+=\varnothing$ としてよい．従ってその側の cell と中間 cells の面積は全て固定され，
critical leaf 側の cell の面積だけが $\delta a_h$ だけ変わる．
その positivity は $a_h-L_h>0$ である．
Internal region では前述の二端公式を用いる．中間 cells は一つの $D_h$ だけに
属する自由成分なので，column mass 条件から不変である．
これで同じ region における複数 annuli と入れ子の主張を得る．
\end{proof}

\begin{remark}[幅を縮めても消えない面積予算]
Source を通る全 regular circle の近傍を固定する条件では，保護 annulus の幅を
零へ近づけても，その circle の向こう側に隔離された面積は一般に零にならない．
従って $\nu(K)$ を小さくすることだけでは大きい target の障害を除けない．
上の判定は，既存の actual reduced forms と同じ保持領域・marked identifications を
使うものであり，source の compact nodal neighborhoods 又は全 overlap の保持に
新しい条件を加えていない．中央の小さい moment interval に適用する場合と，
離れた anchor face に適用する場合は，それぞれの実際の $K$ と面積で別に判定する．
\end{remark}

\begin{example}[正 genus では incidence cycle が生じる]
\label{ng8:positive-genus-incidence-cycle}
二つの pair of pants を，それぞれの二つの境界を二本の annular tubes で結んで貼る．
残る二つの境界を end とする曲面は genus one，boundary components 二個である．
二つの simple saddles と radial ends を持つ葉層を取り，Reeb graph の二つの
saddle vertices の間には二本の平行な internal edges があるとする．
この二つの internal regions $D_1,D_2$ に一つずつ閉飽和 regular annulus を置き，
end collars とともに固定する．自由部分は二つの pants 側 $U_L,U_R$ に分かれる．
両 $D_1,D_2$ が両 $U_L,U_R$ に接するので，incidence graph は
\[
 h_1-U_L-h_2-U_R-h_1
\]
という四辺 cycle を含む．二つの external regions は対応する $U_L,U_R$ に
degree-one vertices として付く．この例の保護円はそれぞれ非分離的である．

例えば external regions の自由需要をそれぞれ $1/2$，internal regions の需要を
$r_1=r_2=1$，自由成分の質量を $M_L=M_R=3/2$ とする．
External entries を $1/2$ とすれば，internal entries は
\[
 \begin{pmatrix}x_{1L}&x_{1R}\\x_{2L}&x_{2R}\end{pmatrix}
 =\begin{pmatrix}s&1-s\\1-s&s\end{pmatrix},\qquad 0<s<1
\]
と自由に選べ，一意配分は成立しない．
これらの正の質量は各 open cell 内の density と小さい固定 collar masses を用いて
実現できる．従って球面又は pants tree の仮定を外して，edge-cut の一意公式を
一般の protected set へ拡張することはできない．
\end{example}

\subsection{Source を跨ぐ面積要求と作用の framing}

この項では $t\ne0$ と一つの stage を固定する．以下の規格化は
補題\ref{ke:source-reduction-pinched-fiber}および
補題\ref{ke:anchor-liouville-actions}と同じ\emph{周期 $1$}である．
Normal cycle を $\alpha$，その generator を $V$ とし，
一つの primitive $\lambda_N$ に対して
\[
 \mu=\lambda_N(V)=\int_\alpha\lambda_N,\qquad
 \iota_Vd\lambda_N=-d\mu
\]
とする．$\lambda_N$ は circle 不変に取る．
Actual canonical primitive からこの選択へは circle averaging 又は
補題\ref{ke:circle-exact-transport}の exact transport を用いられ，
全 closed-cycle periods は変わらない．Source では $V=0$ なので $\mu=0$ である．
前の周期 $2\pi$ による模型と比較する際の moment は $\mu/(2\pi)$ であって，
canonical moment の向き又は加法定数を選び直さない．

\begin{lemma}[二つの capping framing と source の面積]
\label{ng8:source-capping-periods}
標準模型 $\mathbb C^2$ において
\[
 I_j=|z_j|^2/2,\quad
 \lambda_0=I_1d\theta_1+I_2d\theta_2,\quad
 m=I_1-I_2,\quad \mu=2\pi m,
 \quad w=z_1z_2=re^{\sqrt{-1}\phi}
\]
と置く．角座標 $\theta_j,\phi$ は周期 $2\pi$ である．$w\ne0$ 上の
sections $\theta_1=0$ と $\theta_2=0$ が与える reduced primitives は
\begin{equation}\label{ng8:eq:source-two-primitives}
 b_+=\frac{\sqrt{m^2+r^2}-m}{2}\,d\phi,
 \qquad
 b_-=\frac{\sqrt{m^2+r^2}+m}{2}\,d\phi,
 \qquad b_--b_+=\frac{\mu}{2\pi}\,d\phi .
\end{equation}
$\mu>0$ では $b_+$，$\mu<0$ では $b_-$ が $w=0$ へ滑らかに延びる．
Source を避けて滑らかに動く正向き境界を持つ領域 $D_\mu$ について，
その境界の $0$ の周りの winding を $n$ とし，
$B_D(\mu)=\int_{\partial D_\mu}b_+$ と置くと
\begin{equation}\label{ng8:eq:source-area-wall}
 \operatorname{Area}(D_\mu)=B_D(\mu)+n\min\{\mu,0\},
 \qquad B_D\in C^\infty .
\end{equation}
通常の領域では $n$ は含まれる source の数であり，一つの source を
含む場合は $n=1$ である．したがって面積の右微分と左微分の差は $-n$ となる．

正向き $\phi$-loop 上の capping cycles を $\beta_+,\beta_-$ とすると
\begin{equation}\label{ng8:eq:source-capping-shear}
 \beta_-=\beta_++\alpha,
 \qquad
 \int_{\beta_-}\lambda_0=\int_{\beta_+}\lambda_0+\mu .
\end{equation}
Winding $n$ の境界では右辺の補正はそれぞれ $n\alpha,n\mu$ である．
\end{lemma}

\begin{proof}
$|w|^2=4I_1I_2$ より
$I_1=(\sqrt{m^2+r^2}+m)/2$，
$I_2=(\sqrt{m^2+r^2}-m)/2$ である．
各 section への $\lambda_0$ の引き戻しが
\eqref{ng8:eq:source-two-primitives}を与え，両者の微分は
\[
 \omega_{\mathrm{red},\mu}
 =\frac{r}{2\sqrt{m^2+r^2}}\,dr\wedge d\phi
\]
に等しい．正側では $I_2$，負側では $I_1$ が $r=0$ で零となり，
それぞれ対応する primitive が延びる．Stokes の定理により
正側の面積は $B_D$，負側の面積は $B_D+n\mu$ である．
零での密度は可積分であり，両側の連続極限が零 reduction の面積を与える．
境界は $r=0$ を避けるので $B_D$ は $\mu$ の滑らかな関数である．

$(\theta_1,\theta_2)$ の winding vector では
$\alpha=(1,-1)$，$\beta_+=(0,1)$，$\beta_-=(1,0)$ である．
したがって $\beta_-=\beta_++\alpha$ であり，period の式も従う．
ここで切り替えているのは source を\emph{埋める} section である．
Source から離れた同じ境界 torus family 上で $\beta_+$ を滑らかに運ぶことと，
正負両側で capping section を使うことを混同しない．
\end{proof}

\begin{lemma}[Actual source の面積の非滑らかな項]
\label{ng8:actual-source-area-wall}
補題\ref{ng8:p0-source-protected-reduction}の固定 actual neck を取り，
$0$ を含む moment interval $J$ を必要なら小さくする．
$D_{h,\mu}$ は marked identification で固定された critical region とし，
境界は source を避けるものとする．保護領域 $K$ は各 source の全近傍を含み，
$D_h$ の内部の source は全て $K\cap D_h$ の内部に含まれるとする．
Source を除いた境界と marked quotient は $\mu$ に滑らかに依存するものとする．
これは同補題の end collars，caps と source の飽和 annuli による選択で満たされる．

$D_h$ 内の source 数を $n_h$，元の面積と保護質量をそれぞれ
$a_h^0(\mu)$，$\kappa_h(\mu)$ と書く．零 reduction の source 点には
質量を与えず，正則部分上の reduced area を積分する．このとき
\begin{equation}\label{ng8:eq:actual-wall-smooth-remainders}
 a_h^0(\mu)=A_h(\mu)-\frac{n_h}{2}|\mu|,
 \qquad
 \kappa_h(\mu)=K_h(\mu)-\frac{n_h}{2}|\mu|,
 \qquad A_h,K_h\in C^\infty(J).
\end{equation}
特に $a_h^0-\kappa_h$ は滑らかである．定数と $J$ は固定 stage，$t$，
保護領域に依存してよい．
\end{lemma}

\begin{proof}
各 source における circle representation は weights $(1,-1)$ である．
まず，ここで必要な equivariant Darboux chart を局所的に作る．
不変 Riemannian metric の exponential chart で作用を線形化し，原点の
symplectic form を定数形式 $\omega_0$ とする．十分小さい不変 ball 上で
$\omega_s=\omega_0+s(\omega-\omega_0)$ は symplectic である．
Radial homotopy により $\omega-\omega_0$ の不変 primitive を取り，
これを用いる Moser vector field を原点で零に取る．その flow は equivariant であり，
縮小した近傍を標準 symplectic 表現へ同定する．
同定後の moment の微分は $2\pi d(I_1-I_2)$ で，原点の値も零なので
$\mu=2\pi(I_1-I_2)$ が厳密に成立する．
また $\lambda_N-\lambda_0$ はこの可縮 chart 上で exact である．
したがって局所 closed-cycle periods に追加の実数定数は生じない．
この構成は元の residual 座標 $p$ を保存するとは主張しない．

各 source の chart 内で，互いに素な不変 balls
$B_s=\{I_1+I_2\le c_s\}$ を $K\cap D_h$ の持ち上げの内部に取る．
$J$ を縮めて $|\mu|<2\pi c_s$ とすれば，その reduced area は
\[
 \operatorname{Area}((B_s)_{\mathrm{red},\mu})
 =\pi c_s-\frac12|\mu|
\]
である．これは補題\ref{ng8:source-dh-wall}の積分を現在の規格化へ直した式である．
Ball の境界はこの $J$ 上で moment の submersion となり，作用は自由である．

これらの ball の内部を除いた $D_h$ と $K\cap D_h$ の持ち上げは，
境界も含めて compact な free quotient family を与える．
元の marked quotient の境界と，Darboux ball の境界はここではともに
滑らかな free 部分にあるので，同時に trivialize できる．
その reduced forms と領域は滑らかに依存し，積分は滑らかな関数となる．
球の寄与を加えれば\eqref{ng8:eq:actual-wall-smooth-remainders}を得る．
Critical boundary が saddle graph である場合は，source から離れた
有限個の滑らかな arcs への分割，又は固定した相対 rounding で同じ積分を行う．
Source 点で marked quotient の通常の滑らかさを仮定してはいない．
\end{proof}

\begin{proposition}[Source を補正した target 面積要求]
\label{ng8:source-corrected-area-demand}
前補題の同じ $K$ と marked relative isotopy class を保つ．
目標の critical-region areas を $a_h^{\mathrm{tar}}$ とすると，
定理\ref{ng8:componentwise-critical-area}に入れる自由部分の需要
$r_h=a_h^{\mathrm{tar}}-\kappa_h$ が滑らかであるための必要十分条件は
\begin{equation}\label{ng8:eq:source-target-wall-condition}
 a_h^{\mathrm{tar}}(\mu)-n_h\min\{\mu,0\}\in C^\infty(J)
\end{equation}
である．同値に $a_h^{\mathrm{tar}}+n_h|\mu|/2$ が滑らかである．
例えば両側の目標が $c_h^\pm+p_h^\pm\mu$ である場合には
\[
 c_h^-=c_h^+,
 \qquad p_h^-=p_h^++n_h
\]
が必要十分である．$n_h>0$ の領域の\emph{面積自体}を一つの affine function に
指定することは，この滑らかな相対問題の条件を満たさない．

一方，critical actions を指定するときには，保持 boundary cycles を含む
同じ primitive と同じ transported cycle framings で Stokes の式を書く．
$\partial D_h$ の向き付き cycle 表示に現れる作用を $J_\ell$，その符号付き
係数を $e_{h\ell}\in\mathbb Z$ とすれば，ある固定整数 $l_h$ により
\begin{equation}\label{ng8:eq:source-boundary-stokes}
 a_h(\mu)=\sum_\ell e_{h\ell}J_\ell(\mu)+l_h\mu
                       +n_h\min\{\mu,0\}
\end{equation}
と書ける．Cycle の選択を変えれば $l_h$ も対応して変わるので，
その整数を省略した unmarked な式として使わない．
同じ marked relative class の元と目標を比較すると
\begin{equation}\label{ng8:eq:source-relative-actions}
 a_h^{\mathrm{tar}}-a_h^0
 =\sum_\ell e_{h\ell}(J_\ell^{\mathrm{tar}}-J_\ell^0).
\end{equation}
したがって source から離れた critical cycles において，
共通の transported frame で滑らかな作用要求を与えれば，相対面積要求も滑らかである．
特にその frame で affine な critical actions を求めることは，
面積自体を affine とすることとは異なる．
\end{proposition}

\begin{proof}
最初の必要十分性は\eqref{ng8:eq:actual-wall-smooth-remainders}から直ちに従う．
Piecewise affine の場合は，$n_h\min\{\mu,0\}$ を引いた二本の直線の
値と傾きを零で合わせればよい．

Stokes の式を確認する．各 source の小 ball を除いた領域を $W_h$ とする．
$W_h$ は境界付き曲面であり，$J$ は区間なので，その上の principal circle
bundle は自明である．一つの section を固定し，周期 $1$ の circle 座標を
$a\in\mathbb R/\mathbb Z$ と書けば $\lambda_N-\mu\,da$ は basic である．
各 level でその微分は reduced form なので，$W_h$ 上で Stokes を適用できる．
各内境界の section cycle と局所の $\beta_+$ の差は整数倍の $\alpha$ であり，
その period 差は整数倍の $\mu$ である．
外境界の section cycles と，指定された critical／保持 cycles の差も同様である．
各小 ball の面積には\eqref{ng8:eq:source-area-wall}を使う．
内境界の smooth period は相殺し，残る整数倍を $l_h\mu$ とまとめると
\eqref{ng8:eq:source-boundary-stokes}を得る．
Saddle graph の境界 cycle についても，free 部分の有限 arcs 又は近傍の
regular cycles に Stokes を適用して極限を取れば同じ式になる．

相対 isotopy は source の保持近傍と boundary markings を固定する．
したがって各整数 winding と source 数は変わらず，同じ $l_h,n_h,e_{h\ell}$ を使える．
二つの式の差が\eqref{ng8:eq:source-relative-actions}である．
ここでは $\lambda_N$ が全四次元領域で closed であるとはしていない．
Homology による cycle の比較は同じ Lagrangian torus 又はその isotropic strata 上で行い，
領域間の比較には上の Stokes の式を使っている．
\end{proof}

\begin{remark}[この条件と有限面積輸送の関係]
\eqref{ng8:eq:source-relative-actions}で得られる需要が滑らかであることだけでは，
その実現は結論しない．自由成分ごとの総量，正の有限配分と strict cuts，
保持する open end germs をさらに照合する．これらが定理\ref{ng8:componentwise-critical-area}
の仮定を満たせば相対面積構成を適用でき，小さい増分の場合は
系\ref{ng8:p0-small-area-flexibility}の明示的な構成を使える．
Critical action の全線形関係と，固定された boundary actions を反映した
\eqref{ng8:eq:source-boundary-stokes}を先に作る必要がある．
Source local normal form だけから，任意の affine targets 又は実際の P0 の
大きい所定 targets がこれらを満たすとは結論しない．
\end{remark}

\begin{proposition}[Source shear を含む二つの anchor の比較条件]
\label{ng8:source-corrected-anchor-matching}
二つの core--neck collars を $s=L,R$ とし，同じ neck primitive $\lambda_N$ と
同じ canonical normal action $\mu$ へ比較する．各 collar の記号を
\[
 q_s=\sigma_s\mu,\quad \sigma_s\in\{1,-1\},\qquad
 M_s=\delta_s-\rho_s<0
\]
とし，core primitive の anchor residual period を $A_s$ とする．
Core と neck の primitive 差の residual period を
$\kappa_s^{\mathrm{prim}}$ と書く．これは保護\emph{面積} $\kappa_h$ とは別の量である．
Core cycle に対する比較 framing を
$\beta_{N,s}=\varepsilon_s\beta_{\mathrm{core},s}+k_s\alpha_s$ とすると，
補題\ref{ke:anchor-liouville-actions}の後の比較式は
\begin{equation}\label{ng8:eq:source-anchor-local}
 p_s=k_s+\varepsilon_s\frac{A_s}{M_s},\qquad
 c_s=\varepsilon_s\left(\kappa_s^{\mathrm{prim}}
                              -\rho_s\frac{A_s}{M_s}\right).
\end{equation}
ここで $c_s+p_sq_s$ が core 側の指定 affine germ である．

指定された実際の regular torus の同定と cycle transport に関して
\[
 \beta_{N,L}=\varepsilon\beta_{N,R}+n\alpha,
 \qquad \varepsilon\in\{1,-1\},\quad n\in\mathbb Z
\]
が成り立つとする．この同定で二つの affine targets が一致するための
必要十分条件は
\begin{equation}\label{ng8:eq:source-anchor-matching}
 c_L=\varepsilon c_R,
 \qquad \sigma_Lp_L=\varepsilon\sigma_Rp_R+n .
\end{equation}
従って anchor data による条件は
\begin{align*}
 \varepsilon_L\left(\kappa_L^{\mathrm{prim}}-\rho_L A_L/M_L\right)
 &=\varepsilon\varepsilon_R
   \left(\kappa_R^{\mathrm{prim}}-\rho_R A_R/M_R\right),\\
 \sigma_L\left(k_L+\varepsilon_L A_L/M_L\right)
 &=\varepsilon\sigma_R\left(k_R+\varepsilon_R A_R/M_R\right)+n .
\end{align*}
これは period の比較条件であり，全 open overlaps 上の葉の一致又は
全域 atlas の実構成までを結論しない．
\end{proposition}

\begin{proof}
\eqref{ng8:eq:source-anchor-local}は既存の anchor 計算における
$c_s+\rho_s(p_s-k_s)=\varepsilon_s\kappa_s^{\mathrm{prim}}$
を書き直したものである．同じ torus 上では $\lambda_N$ の制限は closed なので
\[
 J_L=\varepsilon J_R+n\mu .
\]
$q_s=\sigma_s\mu$ を代入して affine functions の定数項と一次係数を比較すると，
\eqref{ng8:eq:source-anchor-matching}とその同値な二式を得る．
Local source の capping framing の交換は
\eqref{ng8:eq:source-capping-shear}によりこの整数 shear に寄与する．
Source の canonical moment は零なので，この交換だけから実数の切片補正は生じない．
\end{proof}

\begin{remark}[測るべき整数と実数 period]
上の $n$ は指定された実際の cycle transport による整数である．
一つの境界が source を $n$ 回囲み，正負の capping framings を比較する
局所模型では source の winding がその整数になる．一般の neck の両端を比較するときは，
triple 側の既定 framing と transport path も含めて計算しなければならない．
単に moment interval が零を跨ぐことから，任意の transported critical cycle の
framing が source 数だけ変わるとはしない．

二つの neck primitives が異なり，その closed 差の normal period が零なら，
その residual period $C$ を明記し，切片条件を
$c_L=\varepsilon c_R+C$ とする．一つの $\lambda_N$ を使う上の設定では $C=0$ である．
Node germ を固定するだけではこの residual constant は決まらない
（例\ref{ke:focus-local-primitive-shift}）．
Source shear は整数勾配の比較を変えるが，実数切片の不一致を修復しない．
実際の $\kappa_s^{\mathrm{prim}}$ と global framing が全て測定済みであるとは主張しない．
\end{remark}

\input{sections/paper/quintic_gluing.tex}
\input{sections/paper/sextic_gluing.tex}
\section{五次・六次全平面模型の大域作用適合 atlas}
\label{ga8:section}

局所 cone charts，実際の二重曲線上の面積，normal framing，および
有限の平滑化 parameter における作用の加法定数を，ここで一つの構成にまとめる．
必要なのは，既製の二つの foliations を単に isotopy で結ぶことではない．
同じ actual Hamiltonian neck 上で critical actions を指定された直線に保ち，
その共通 Hamiltonian が各 branch star で ambient-affine になることを用いる．

\begin{theorem}[五次・六次の作用適合 matroidal atlas]
\label{ga8:global}
$d\in\{5,6\}$ とし，$X_d\subset\PP^3$ を任意の滑らかな次数 $d$ の曲面，
$\omega$ を $[\omega]=aH$，$a>0$ を満たす任意の K\"ahler form とする．
定理~\ref{main:global} の全平面構成を選び，その fibration
\[
 \pi:(X_d,\omega)\longrightarrow B_d
\]
について，定義~\ref{atl:definition} の，focus--focus 特異点を許す
作用適合 matroidal atlas を構成できる．
全ての branch denominator は $q=1$ である．

$d=5$ では五つの core charts を標準 $K_4$ fan とし，
系~\ref{qg8:six-systems} の六つの標識付き対称有向四角形系のいずれも選べる．
$d=6$ では六つの core charts を標準 Petersen fan とし，
定理~\ref{topo:petersen} の六つの標識付き同時 Petersen systems のいずれも選べる．
それぞれの実際の二重曲線上で，対応する embedded marked pants class を先に固定できる．
foliation の代表は，以下の面積条件を満たすようにその class 内で選ぶ．

構成の個数と面内の node 配分は次の通りである．
\begin{center}
\begin{tabular}{c|ccccc}
 $d$ & cores & branch edges & triangles & bigons & focus points\\ \hline
 $5$ & $5$ & $10$ & $10$ & $0$ & $50$\\
 $6$ & $6$ & $30$ & $20$ & $15$ & $90$
\end{tabular}
\end{center}
五次では各 triangle に五つの focus points があり，
六次では各 triangle に三つ，各 bigon に二つある．
全ての branch edges の primitive affine length は
\begin{equation}\label{ga8:branch-length}
 \frac{a}{\pi(d-1)}
 =\begin{cases}a/(4\pi),&d=5,\\a/(5\pi),&d=6\end{cases}
\end{equation}
である．底空間の homotopy 型はそれぞれ
$\bigvee^4S^2$ と $\bigvee^{10}S^2$ である．

focus values 自体は通常の matroidal charts の定義域から除き，
その穿孔近傍には実際の nodal fibration の作用アフィン構造を用いる．
core fibers は元の直線補集合の logarithmic skeleton の symplectic images である．
特に Petersen fan という結論は，core の symplectic filling を
標準 $M_{0,5}$ の filling と同定する主張ではない．
\end{theorem}

証明の前に，実際の成分へ移す際に使う二点を明示する．

\begin{lemma}[有限標識点も固定する divisor-relative Moser]
\label{ga8:fixed-q}
補題~\ref{glue:pair-moser} の仮定の下で，crossings を避けた divisor 上の
有限個の相異なる点を指定する．その結論の isotopy は，これらの点も
全 parameter で固定するように選べる．
\end{lemma}

\begin{proof}
$\dd\alpha_s=\partial_s\omega_s$ を取り，crossing balls の小近傍で
$\alpha_s=0$ とする，元の証明から始める．指定点 $q\in C$ ごとに
\[
 g_s(q)=0,\qquad
 (\dd g_s)|_{T_qC}=-\alpha_s|_{T_qC}
\]
という first jet を指定し，互いに素な小近傍に支持された滑らかな関数
$g_s$ として延長する．$\alpha_s$ を $\alpha_s+\dd g_s$ で置き換えると，
crossing balls の条件とその微分は変わらず，指定点での tangent restriction が零になる．

次に，$N_s=(TC)^{\omega_s}$ に沿う元の normal-jet correction
\[
 F_s|_C=0,\qquad (\dd F_s)|_{N_s}=-\alpha_s|_{N_s}
\]
を行う．$F_s|_C=0$ なので tangent restriction は変更されない．
従って $\eta_s=\alpha_s+\dd F_s$ は各指定点で covector として零である．
$\iota_{Y_s}\omega_s=-\eta_s$ を解けば，$Y_s$ は全 divisor に接し，
指定点およびより小さい crossing balls で零になる．
その Moser flow が主張を満たす．
\end{proof}

\begin{lemma}[実成分への移植と critical action germs]
\label{ga8:affine-persistence}
命題~\ref{per:k4-realization} または定理~\ref{per:local-realization} の
core model において，nodal blow-up patches を全 critical leaves から離して取る．
局所 nodal comparison，divisor-relative Moser，および
補題~\ref{glue:circle} の normal-circle conjugacy の後も，
実際の二重曲線に十分近い critical action germs は，normal moment の
整数勾配を持つ affine functions である．
共通 neck framing へ直すときの勾配の変化は，
補題~\ref{qg8:framing} または命題~\ref{sg8:framing} に記した整数 shear である．
\end{lemma}

\begin{proof}
core model の proportional-curvature collar では，これらの作用は
normal action $\rho-w$ の定数倍であり，指定した recipes ではその係数が整数である．
flat-ball の準備，局所 blow-up，および相対 almost-toric comparison は，
critical leaves の全 germ の外に支持を取れる．従ってこの段階では
critical germs とその作用関数は変更されない\cite[Theorem~1.1]{CGblowup}．

divisor-relative Moser は critical leaves を点ごとに固定する必要はない．
この symplectomorphism により，fibration の germ 全体，その周期 cycles，
primitive，および normal circle action を共に運ぶ．
作用関数の値と moment の関係は，この同時輸送で保たれる．
続く circle conjugacy は divisor を点ごとに固定し，
normal weight $+1$ の作用を actual neck の作用へ合わせる．
moments は divisor 上で零と正規化するので，両 moment の間に加法定数は残らない．
crossing germs を固定し，余分な $U(1)$ gauge winding を入れない同定を用いる．

運ばれた primitive と，節~\ref{ng8:section} の一つの neck primitive の差は，
共通 collar 上で closed である．連続に運ばれる同じ cycle に対するその積分は定数なので，
critical actions の変化は加法定数だけである．
cycle framing を $\beta\mapsto\beta+m\alpha$ と変えたときは，
さらに $m\mu$ が加わる．従って affine 性と整数勾配が保たれ，
残る勾配の比較は actual normal-bundle clutching の計算に正確に帰着する．
\end{proof}

\begin{proof}[定理~\ref{ga8:global} の証明]
\textit{偏極と有限個の marked データを準備する．}
成分を $i=0,\ldots,d-1$ と標識し，十分小さい $\eps>0$ に対して
\[
 \beta_i=(i-d+1)\eps,\qquad
 b_i=\frac{\beta_i}{2\pi},\qquad
 c_d=\frac{a}{2\pi(d-1)},\qquad
 w_{i,j}=c_d+b_j-b_i
\]
と置く．この ray は定理~\ref{pol:cone} の K\"ahler simplex の内部にあり，
例えば $\eps<a/d^2$ とすれば全 weights は正である．
実際の surgeries の capacity を確保するためには，後でさらに $\eps$ を小さくする．
各 blow-down component における直線の面積は
\[
 2\pi R_i,\qquad
 R_i=(d-1)c_d+\sum_jb_j-db_i
     =\sum_{j\ne i}w_{i,j}
\]
である．$i<j$ の二重曲線の $d$ 個の blow-up sizes は全て
$2\pi\delta_{ij}$，$\delta_{ij}=b_j-b_i$ である．

五次では六つの対称有向四角形系の一つを取り，その代表を
$f_i(j)=(i+j)/2$ on $\mathbb F_5$ としてよい．
各 $C_{ij}$ の三つの basins に，選択した一領域へ三点，残る二領域へ一点ずつ，
実際の $q_{ij,r}$ を割り当てる．
六次では選んだ同時 Petersen system に従って各 $C_{ij}$ の quartet を固定し，
四つの external basins に一点ずつ，内部 annulus に二点を割り当てる．
同じ $C_{ij}=C_{ji}$ 上で embedded marked class を先に選び，
各 $q_{ij,r}$ を相異なる regular leaves に置く．
補題~\ref{qg8:marked-allocation} と節~\ref{sg8:section} の density--Moser argument により，
crossing disks と有限個の標識点の近傍を固定してこの選択ができる．

blow-down で指定する basin areas は，五次では
$2\pi(w_{i,k}+\epsilon_{i;jk}w_{i,j})$，六次では $2\pi w_{i,k}$ である．
六次の内部 annulus の面積は $2\pi w_{i,j}$ とする．
指定点近傍の flat-ball 準備は critical leaves を避けた一領域内の exact modification なので，
その領域全体の積分を変えない．
normal-coupling の準備も divisor の制限形式を保つ．
その後の blow-up は，含まれる領域の面積から一個につき $2\pi\delta_{ij}$ を引く．
従って，実際の strict curve 上の対応する領域面積は，
定理~\ref{qg8:unique-weights} および命題~\ref{sg8:areas} によって両側で一致する．
全ての目標面積は正である．これが中央曲線における面積の条件である．

crossing germs，marked foliations，および regular patches は，
$\eps=0$ の blow-down データから小さい parameter 族として準備できる．
有限個の patch capacities は正の下限を持つので，この準備の後で $\eps$ を縮める．
実際の全空間の neck 平均化と三重点の product germs は，
節~\ref{sec:necks} の fixed-class 構成で先に準備する．

\textit{core charts と nodes を実際の成分上へ移す．}
各 blow-down component に定理~\ref{snc:main} の相対構成を施し，
上記の面積を持つ foliation を用いて，五次では命題~\ref{per:k4-realization}，
六次では定理~\ref{per:local-realization} の charts を取る．
全ての surgeries は core と critical collar germs の外の regular patches 内で行う．
局所比較定理は patch の外で恒等であるから，一個の blow-up ごとに
一個の compact focus--focus fiber が入り，core charts は保たれる．

モデルの blow-up 形式から actual $\Omega_i$ への同定には，
補題~\ref{ga8:fixed-q} を用いて各 strict curve 上の標識点も固定する．
これは crossing disks を固定する divisor-relative isotopy であるから，
両側の reduced foliations は先に選んだ同じ marked class に留まる．
全ての basin と annulus は形式と共に運ばれるので，既に確認した実面積の一致も保たれる．
両側の embedded critical leaves がこの段階で点ごとに一致する必要はない．
次の neck 補間に必要なのは，同じ marked class と対応する critical-region areas である．

次に補題~\ref{glue:circle} を適用する．
core と compact nodal fibers 全体を避けるほど細い collar で normal actions を合わせる．
補題~\ref{ga8:affine-persistence} により，実際の成分側の critical actions はなお affine である．
五次では補題~\ref{qg8:framing} と命題~\ref{qg8:slope}，
六次では命題~\ref{sg8:framing} が，両側を共通 neck framing に直した後の
整数勾配の一致を与える．この順序では，nodal monodromy と normal-bundle winding を
葉層の isotopy で消したことにはしていない．

\textit{有限の平滑化 parameter で critical actions を接続する．}
各二重曲線について，中央 strict curve 上の面積一致と整数勾配の一致が揃ったので，
系~\ref{ng8:central-criterion} を適用する．
puncture cuts と保存する collars を先に十分小さく選び，その後で正の切断値
$\delta$ を選ぶ．一つの全空間 primitive と toric reference actions を用いるため，
中央曲線上で一致した critical action intercepts は，共通の意味を持つ．
固定した $\mu=\delta$ と $\mu=-\delta$ の open buffers を nearby fiber に運ぶと，
補題~\ref{ng8:transport-intercept} はこれらの作用関数そのものを保存する．
ここでは $\mu=0,t=0$ を通る輸送も，$\delta\to0$ に一様な輸送も使わない．

残る interval では命題~\ref{ng8:straight-neck} により，
end germs と toric collars を保って measured critical regions を補間する．
total reduced area との整合は補題~\ref{ng8:relative-dh} の affine Stokes formula による．
中央曲線で正の basin および annulus areas は，$\delta$ を縮めれば全 interval で正となり，
保存する領域に対する面積の余裕も確保できる．
従って，actual neck 全体の critical actions は同じ整数勾配の affine functions となる．
三重点では固定した toric leaves を使い，節~\ref{glue:hexagon} の
飽和した affine moment cuts によって components と necks を接続する．

core と nodal fiber 全体およびその飽和近傍は，
補題~\ref{glue:transport} による固定 compact subset の輸送で保存する．
全 interface は whole fibers の和であり，open overlap 上で実際の葉層が一致する．
節~\ref{sec:gluing} の properness と飽和 quotient の議論が適用され，
追加の特異 fibers を生じずに，compact $X_t$ 上の fibration を得る．

\textit{facewise actions を一つの atlas にする．}
core の近傍には既に $K_4$ または Petersen charts がある．
各 branch の内部では，critical actions の整数勾配と
simple-saddle の transverse action の strict monotonicity により，
補題~\ref{atl:qone-chart} の chart が得られる．
その共通 Hamiltonian は，neck の normal moment $\mu$ である．
core chart 側でも，同じ作用は ambient-affine である．
$K_4$ の場合は \eqref{per:k4-lattice-map} の $\Phi$ を用い，
Petersen の場合は補題~\ref{per:branch-duality} の saturated branch-star lattice を用いる．
加法定数は core primitive と neck primitive の正規化の差として許される．
この Hamiltonian は branch tangent に沿って primitive な非零勾配を持つ．

正則 faces では実際の fibration の通常の action charts を使う．
これらと core および branch charts は同じ作用格子に適合し，
branch を横切る overlap では共通 Hamiltonian が両方で ambient-affine である．
従って命題~\ref{atl:gluing-criterion} により，transition とその逆は
ambient integer affine maps の制限となる．
五次の明示的な変換は命題~\ref{qg8:star-transition}，
六次の変換は補題~\ref{sg8:star} にも記した通りである．
cocycle identity は実際の chart transitions の合成として成り立つ．
focus--focus values の穿孔近傍には，既存の compact nodal model の作用格子を用いる．
その非自明 monodromy を通常の smooth matroidal chart で埋めることは要求していない．

\textit{個数，長さ，および指定された K\"ahler form を確認する．}
五次の一つの triangle では，三つの pair のうち選択 basin となる pair が一つなので，
node 数は $3+1+1=5$ である．
六次では各 triangle が三つの external basin nodes を持ち，
各 bigon が一つの内部 annulus の二つの nodes を持つ．
他の cell counts と homotopy 型は定理~\ref{main:global} による．
branch 上の normal action の原点は，$i$ core では $\mu=w_{i,j}$，
$j$ core では $\mu=-w_{j,i}$ にある．
全区間で primitive な勾配は一なので，その affine length は
$w_{i,j}+w_{j,i}=2c_d$ となり，\eqref{ga8:branch-length} を得る．

最後に，全空間の exact modifications を compact $X_t$ 上で Moser により戻す．
さらに，滑らかな次数 $d$ の射影超曲面の connected family に沿う symplectic transport と，
同一 class の K\"ahler forms の凸結合に対する Moser を用いて，
任意の指定された $(X_d,\omega)$ へ移す．
これらは fibration と作用格子を共に運ぶ symplectomorphisms であり，
atlas の存在，整数勾配，focus monodromy，および affine lengths を保つ．
\end{proof}

\begin{remark}[構成の量化]
本定理は，各次数について作用適合 atlas を持つ fibration の存在を述べる．
既に任意に固定された basin areas，normal holonomies，または end germs 全てが
この方法で接着できるという主張ではない．
また，K\"ahler simplex の全点で所定の容量の surgeries を実行する必要はなく，
十分小さい vertical ray を用いている．
六つの標識付き recipes または graph systems の実現は，
六つの相異なる大域 affine-isomorphism classes の分類を意味しない．
\end{remark}


\subsection{Actual link の有限 period 条件と Liouville core}

ここでは component の Liouville domain，primitive，canonical normal moments，
保持 corner と anchor cuts を固定する．比例曲率の模型への移植を入力とせず，
actual link の葉を選び直すための有限条件を記す．ただし source を含む保持領域を
避けた面積調整の存在は別問題であり，以下の period 条件だけからは結論しない．

\begin{lemma}[Anchor の period と core へ向かう critical action]
\label{ke:anchor-liouville-actions}
補題\ref{ke:actual-link-cone-filling}の $(K,d\lambda)$ と link $S$ を取り，
normal face では周期 $1$ の actual generator $V_a$ に対して
\[
 \mathcal L_{V_a}\lambda=0,\qquad
 \lambda(V_a)=M_a=q_a-\rho_a,\qquad q_a|_S=\delta_a<\rho_a
\]
とする．$q_a=\epsilon_a\widehat\mu_N$ は加法定数まででなく，
既存の canonical moment と厳密に同じものとする．
この face の一つの critical fiber 上の標識 cycle $\beta$ の period を
$A=\int_\beta\lambda$ と書く．Critical cycle の積分は，必要なら隣接する
regular annuli の極限で解釈する．

同じ link の葉を Liouville flow によって充填すると，
その critical ray 上の period は
\begin{equation}\label{ke:eq:anchor-action-ray}
 I_\alpha(r)=r(\delta_a-\rho_a),\qquad I_\beta(r)=rA,
 \qquad 0<r\le1.
\end{equation}
従って，元の actual normal circle が定義される共通 collar では
\begin{equation}\label{ke:eq:anchor-action-line}
 q_a(r)=\rho_a+r(\delta_a-\rho_a),\qquad
 I_\beta(q_a)=\frac{A}{\delta_a-\rho_a}(q_a-\rho_a).
\end{equation}
右辺の $q_a(r)$ は，collar の外では transport した local normal action の
連続延長を表す．元の canonical circle action が全 core に存在する，という
主張ではない．

同じ primitive と cycle framing で指定された neck の直線 $c+pq_a$ が
この Liouville germ と一致するための必要十分条件は
\begin{equation}\label{ke:eq:anchor-intercept}
 A=p(\delta_a-\rho_a),\qquad c=-p\rho_a
\end{equation}
である．整数勾配を要する場合は $p\in\mathbb Z$ を要求する．
\end{lemma}

\begin{proof}
$\iota_Zd\lambda=\lambda$ より $\lambda(Z)=0$，
$\mathcal L_Z\lambda=\lambda$ である．従って
$\varphi_s^*\lambda=e^s\lambda$．$r=e^s$ と置けば，全ての標識 cycles の
period が同じ倍率 $r$ を受け，\eqref{ke:eq:anchor-action-ray}を得る．
Collar では $[Z,V_a]=0$ なので，transport した normal cycle は元の actual
circle cycle と同じであり，厳密な moment 規格化から
\eqref{ke:eq:anchor-action-line}が従う．二つの affine functions の
係数を比較すると\eqref{ke:eq:anchor-intercept}を得る．
この計算は critical cycles を regular cycles で近似した場合にも連続性で成立する．
\end{proof}

\begin{remark}[Primitive と framing を変えて比較するとき]
Neck primitive $\lambda_N$ と core primitive $\lambda$ の差は共通 collar で
closed である．同じ標識 cycle $\beta$ に対し
$\kappa=\int_\beta(\lambda_N-\lambda)$ は collar 内で定数であり，
normal cycle 上の差は $\rho_a$ である．比較 cycle を
$\beta_N=\varepsilon\beta+k\alpha$，$\varepsilon\in\{1,-1\}$，
$k\in\mathbb Z$ とする場合には，neck action は
\[
 I^N_{\beta_N}(q_a)
 =\varepsilon\frac{A}{\delta_a-\rho_a}(q_a-\rho_a)
   +\varepsilon\kappa+kq_a.
\]
従って直線 $c+pq_a$ との一致条件は
\[
 p=k+\varepsilon\frac{A}{\delta_a-\rho_a},\qquad
 c+\rho_a(p-k)=\varepsilon\kappa.
\]
最後に $q_a=\epsilon_a\widehat\mu_N$ へ戻す．閉形式の残余 period を
零と仮定したり，component ごとに neck primitive の切片を独立に消してはならない．
Exact な primitive 補正 $\lambda\mapsto\lambda+dF$ は，固定した cycle の
period を変えないため，この条件の不足を自動的に修復しない．
\end{remark}

\begin{proposition}[Actual anchor periods による $K_4$／Petersen core chart]
\label{ke:actual-anchor-fan-criterion}
補題\ref{ke:actual-link-cone-filling}の actual domain と invariant primitive を
固定し，link は normal faces と保持 corner の torus-orbit caps からなるとする．
Normal face $a$ の quotient は $d-2$ 個の境界を持つ球面で，
$M_a=\delta_a-\rho_a<0$ はその face 全体で一定である．
各 normal--corner overlap では標識 cycles が既存の同じ actual torus を通じて
一致するものとする．特に二つの正の meridians は整数 torus basis であり，
保持 corner と end collars の各 meridian period は負である．

各 normal face に，次の period を持つ actual reduced Morse foliation を
与え，end collars では保持 torus partition と一致すると仮定する．
\begin{enumerate}
\item $d=5$ では四つの boundary lines の有向四角形 $f$ を選ぶ．
Face $a$ の一つの critical fiber 上で，puncture $b$ から運んだ meridian
$\beta_b$ の period は
\[
 \int_{\beta_b}\lambda=-M_a\mathbf1_{\{b=f(a)\}}.
\]
\item $d=6$ では Petersen 型の labelled quartet system を選ぶ．
Face $a$ の四つの external critical basins の，puncture から運んだ cycles の
period は全て零とする．内部 annulus の cycle は，二つの端でそれぞれ
$0$ と $-M_a$ の period を持つものとする．この後者の条件は，
実際の torus cycle の Euler framing relation と外側の零 period から従う場合には
独立な条件でない．
\end{enumerate}
Normal face の regular annuli では，これらの標識 cycles の period 差が
向き付けた reduced annulus の面積に等しいとする．これは一つの primitive と
固定 framing を用いた Stokes の同定であり，独立な second-action gauge の
一致を仮定することではない．

このとき，同じ link 葉を用いた Liouville filling は，core の近傍に
作用座標に適合する標準 $K_4$，それぞれ Petersen cone chart を持つ．
各 face と branch の格子は飽和し，全 branch denominator は $1$ である．
元の Liouville domain，その skeleton と等方的 filtration は変えない．
保持 corner の全 torus 葉と cap arcs も保存される．
さらに各 branch star で，その normal action の延長 $q_a$ はこの chart の
ambient integer affine function であり，branch tangent 上の勾配は primitive
で非零である．
\end{proposition}

\begin{proof}
$K_4$ の crossing face $ab$ では，既存の整数行列
\[
 A_{ab}=\begin{pmatrix}-1&\epsilon_{ba}\\
                         \epsilon_{ab}&-1\end{pmatrix},\qquad
 \epsilon_{ab}=\mathbf1_{\{b=f(a)\}}
\]
を使って $(I_a,I_b)^T=A_{ab}(t_a,t_b)^T$，
$F=t_ar_a+t_br_b$ と定める．有向四角形には二辺 cycle がなく，
$\det A_{ab}=1$．仮定の critical periods により各 branch 上では
$F=(-M_a)r_a$ となり，隣接 face maps が一致する．

正値性には比例曲率を使わない．例えば $a\to b$ なら
$t_a=-I_a$，$t_b=-I_a-I_b$ である．$a$ 側の normal face では
$I_a=M_a<0$ が一定で，$I_b$ は負の reference value から
$-M_a$ へ面積に従って厳密単調に増える．従って critical endpoint を除き
$t_a,t_b>0$．$b$ 側では $I_a<0$，$I_b=M_b<0$ から同じ結論を得る．
非選択辺では $t_a=-I_a$，$t_b=-I_b$．保持 corner と丸めでは
両 meridian periods が負なので，どの場合も係数は正である．

Petersen の crossing faces では
$F=-I_ar_v-I_br_w$，内部 faces では
\[
 F=(-I_a-I_\beta)r_A+I_\beta r_B
\]
とする．前者では external periods が負から零へ増え，
後者では $I_\beta$ が $0$ から $-M_a>0$ へ増える．従って両係数は
regular interior で正であり，端点では所定の共通 ray に一致する．
二つの係数への行列の determinant はそれぞれ $1$，$-1$ である．
Cycle framing は実際の overlap で一致するので，これらは単なる抽象 graph の
同定から作った period maps ではない．

各 link edge の parameter を $u$ とし，その action vector を $a(u)$ とする．
Liouville cone 上の作用写像は $(r,u)\mapsto ra(u)$ である．
正則 Lagrangian fibration の作用座標は局所座標だから
$\det(a(u),a'(u))\ne0$．各 face の正の sector 内では，projective action ratio
が厳密単調であり，二つの boundary rays の間を一度だけ走る．
従って各 face map は所定の cone sector の原点近傍へ同相を与える．
有限個の edges と compact な link を用い，原点近傍を共通に縮めればよい．
Ray 上の係数は全 incident faces で同じであり，標準 fan の face lattices と
上の行列は unimodular なので，branch lattices も飽和する．
Normal period の ambient integer linear extension は，$K_4$ では
\eqref{per:k4-lattice-map}，Petersen では補題\ref{per:branch-duality}の
同じ整数行列を使う．$q_a=I_\alpha+\rho_a$ なので，normal action は
ambient integer affine となり，その branch 勾配は primitive である．

補題\ref{ke:anchor-liouville-actions}が，全 radial levels へ同じ関係を
運ぶ．Core point への一様収束と全 singular strata の等方性，skeleton の
連結性は補題\ref{ke:actual-link-cone-filling}の同じ filling から従う．
Corner では invariant primitive の Liouville flow が actual torus action と
可換であり，torus orbits の全 partition を保つ．従って保持 cap arcs を含む
全 corner overlap に余分な葉の接着条件を生じない．
\end{proof}

\begin{lemma}[Normal face の Stokes 総量と必要な target areas]
\label{ke:anchor-stokes-budget}
固定 normal face の circle bundle $E\to P$ を取り，$P$ は $k$ 個の境界を
持つ球面とする．Normal cycle を $\alpha$，各 puncture を正に回る，
指定された実際の lift cycles を $\beta_b$ とする．
Clutching relation と primitive の period を
\[
 \sum_{b=1}^k[\beta_b]=-[\alpha]\quad\hbox{in }H_1(E;\mathbb Z),
 \qquad \int_\alpha\lambda=M<0,\qquad
 c_b=\int_{\beta_b}\lambda
\]
と仮定する．ここで homology relation に非閉形式 $\lambda$ を直接積分しない．
Reduced form を $\omega_P$ とすると
\begin{equation}\label{ke:eq:anchor-area-total}
 \operatorname{Area}(P)=-\sum_bc_b-M.
\end{equation}
従って $K_4$ の target critical periods
$A_b=-M\mathbf1_{\{b=f(a)\}}$ に必要な basin areas は
$a_b=A_b-c_b$ で，その和は厳密に全体面積に等しい．
Petersen の external target $A_b=0$ に必要な外側 basin areas は $a_b=-c_b$，
内部 annulus の target area は $-M$ であり，これらの和も厳密に全体面積に等しい．
\end{lemma}

\begin{proof}
$H^2(P;\mathbb Z)=0$ なので bundle を自明化し，周期 $1$ の座標 $\theta$ と
global section を選ぶ．Circle invariance と $M$ の一定性から
$\lambda=M\,d\theta+\pi^*\beta$ と書け，$d\beta=\omega_P$．
$\beta_b$ の section に対する整数 winding を $n_b$ とすれば，
clutching relation より $\sum_bn_b=-1$ である．
正の puncture circle $\gamma_b$ の向きは $P$ の境界向きと反対なので
\[
 \int_P\omega_P=-\sum_b\int_{\gamma_b}\beta
 =-\sum_b(c_b-Mn_b)=-\sum_bc_b-M.
\]
同じ固定 framing を critical cycle まで運ぶと，annulus の Stokes は
$a_b=A_b-c_b$ を与える．各指定値を代入すると残りの結論を得る．
\end{proof}

\begin{remark}[Actual 適用で確認する部分と残る輸送条件]
一次 normal bundle の Euler 数 $+1$ と puncture-compatible framings は，
上の clutching relation の標準的な実現である．本稿の actual component に
適用する場合には，component graph gauge と circle conjugacies がこの同じ
meridian marking を運ぶことを確認する必要がある．
補題\ref{ke:component-meridian-normalization}は meridian period の厳密な
規格化を与えるが，任意に選んだ residual cycle の framing を無条件に同定しない．

保持集合が puncture collars のみで，その内側境界でも同じ meridian period
$c'_b<0$ なら，その collar の面積は $c'_b-c_b$ である．
上の非負 critical target $A_b$ について
$a_b-\operatorname{Area}(K_b)=A_b-c'_b>0$ となり，正の面積の余裕を得る．
この限定された場合には，補題\ref{ng8:parametric-moser}と同じ相対 density--Moser
argument により，指定された marked class 内で anchor の葉を構成できる．
ただし全 source annuli など，追加の保持集合 $K$ が補集合を分断する場合には，
この総量一致と各 basin の余裕だけでは十分でない．
各補集合成分への面積割当てと，critical regions との incidence を保つ
有限輸送条件を別途満たす必要がある．

上の命題は，その実際の anchor 葉が得られた後の core cone filling と
action chart を与える．Neck の反対側の critical lines，残余 primitive の
定数差，normal framing，source 周囲の全保持領域を含む marked interpolation，
および全 open overlaps の whole-fiber 一致までは与えない．
これらを\eqref{ke:eq:anchor-intercept}又はその gauge付きの式と合わせて
確認した後に，\texttt{ke:moment-relative-replacement}を用いる．
\end{remark}

\begin{remark}[Irrational period による具体的な障害]
固定した anchor の $A/(\delta_a-\rho_a)$ が無理数なら，その Liouville filling
の branch は整数作用格子に対して有理方向を持たない．
整数可逆な作用座標変換は有理方向を有理方向へ，無理方向を無理方向へ写すので，
この filling に作用適合 matroidal core chart を与えることはできない．
これは Liouville cone construction が自動的に atlas を与えない具体的理由である．
枝を含む regular annulus の内部で正の面積を小さく移すことで，end collars と
normal moment を保ったまま critical period を連続に変えられる場合，この
無理数条件を満たす anchor foliation は存在する．その小変更の可否も，
保持集合の補集合内に対応する可動領域がある場合に限定する．
Critical action が直線であることと，その直線が整数作用格子に適合することは
異なる．一般的な cone-extension 原理や action-angle 原理の新規性は主張しない．
\end{remark}

\input{sections/paper/mixed_global_obstruction.tex}
\input{sections/paper/mixed_all_denominator_obstruction.tex}
\input{sections/paper/invariants_examples.tex}
\begin{thebibliography}{99}

\bibitem{AAK}
M. Abouzaid, D. Auroux and L. Katzarkov,
\emph{Lagrangian fibrations on blowups of toric varieties and mirror symmetry for hypersurfaces},
Publications Mathématiques de l'IHÉS \textbf{123} (2016), 199--282.
\url{https://doi.org/10.1007/s10240-016-0081-9}.

\bibitem{AbbondandoloMajerMS}
A. Abbondandolo and P. Majer,
\emph{Stable foliations and CW-structure induced by a Morse--Smale gradient-like flow},
arXiv:2003.07134v2 (2020).
Introduction, pp.~1--2 (closure relations), and Theorem~1.
\url{https://arxiv.org/abs/2003.07134v2}.

\bibitem{BSV}
M. S. Borman, N. Sheridan and U. Varolgunes,
\emph{Quantum cohomology as a deformation of symplectic cohomology},
Journal of Fixed Point Theory and Applications \textbf{24} (2022), article 48.
Remark~2.2, Definition~2.4, Propositions~2.5, 2.7, 2.10 and Remark~2.11,
published pp.~28--31 (wrapping numbers, commuting Hamiltonians and adapted primitive).
\url{https://doi.org/10.1007/s11784-022-00965-6};
\url{https://arxiv.org/abs/2108.08391v2}.

\bibitem{CEFlexible}
K. Cieliebak and Y. Eliashberg,
\emph{Flexible Weinstein manifolds},
in \emph{Symplectic, Poisson and Noncommutative Geometry},
MSRI Publications \textbf{62} (2014), 1--42.
\S4A, Lemma~4.1 (holonomy deformation).
\url{https://library.slmath.org/books/Book62/files/eliashberg.pdf}.

\bibitem{CGblowup}
P. Chakravarthy and Y. Groman,
\emph{Almost toric fibrations on symplectic blow ups},
arXiv:2510.00994v1 (2025).
Theorem~1.1, Definition~2.2, \S2.4, Propositions~3.1, 3.6, Lemma~3.7 and \S4.
\url{https://arxiv.org/abs/2510.00994v1}.

\bibitem{CGK3}
P. Chakravarthy and Y. Groman,
\emph{Almost toric fibrations on K3 surfaces via degenerations},
Advances in Mathematics \textbf{502} (2026), Part~B, 111138.
Definition~1.2, Theorems~1.6, 1.8, Corollary~1.9;
Lemma~2.13 is cited in the revised arXiv version of 28 May 2026.
\url{https://doi.org/10.1016/j.aim.2026.111138};
\url{https://arxiv.org/abs/2502.04304v3}.

\bibitem{CohenSuciuBoundary}
D. C. Cohen and A. I. Suciu,
\emph{Boundary manifolds of projective hypersurfaces},
Advances in Mathematics \textbf{206} (2006), no.~2, 538--566.
Proposition~2.5, Corollary~2.6 and \S6.9 (numbering in arXiv:math/0502506v2,
PDF pages~6--7 and 28).
\url{https://doi.org/10.1016/j.aim.2005.10.003};
\url{https://arxiv.org/abs/math/0502506v2}.

\bibitem{DMT}
J.-P. Dufour, P. Molino and A. Toulet,
\emph{Classification des systèmes intégrables en dimension 2 et invariants des modèles de Fomenko},
Comptes Rendus de l'Académie des Sciences, Paris, Série~I, Mathématique
\textbf{318} (1994), no.~10, 949--952.
An English translation is available at
\url{https://webpages.csus.edu/mjv/DufourMolinoToulet.pdf}.

\bibitem{DMVMT}
T.-C. Dinh, G. Marinescu and D.-V. Vu,
\emph{Moser--Trudinger inequalities and complex Monge--Amp\`ere equation},
Ann. Sc. Norm. Super. Pisa Cl. Sci. (5) \textbf{24} (2023), 927--954.
Theorem~1.2 and the complex Sobolev norm, pp.~928--929.
\url{https://www.mi.uni-koeln.de/~gmarines/PREPRINTS/Dinh_Marinescu_Vu_ASNP_2023.pdf}.

\bibitem{DHJH}
S. Dwivedi, J. Herman, L. C. Jeffrey and T. van den Hurk,
\emph{Hamiltonian Group Actions and Equivariant Cohomology},
SpringerBriefs in Mathematics, Springer, Cham, 2019.
Theorem~3.2 and Corollary~3.4, pp.~17--20 (symplectic neighborhood);
Chapter~10 (Duistermaat--Heckman).
\url{https://doi.org/10.1007/978-3-030-27227-2}.

\bibitem{EliashbergWeinsteinRevisited}
Y. Eliashberg,
\emph{Weinstein manifolds revisited},
arXiv:1707.03442v2 (2017).
\S1, pp.~2--3 (exact Moser transport and isotropic stratification).
\url{https://arxiv.org/abs/1707.03442v2}.

\bibitem{Evans}
J. Evans,
\emph{Lectures on Lagrangian Torus Fibrations},
Cambridge University Press, Cambridge, 2023.
Theorem~1.42, p.~16; Lemmas~2.10, 2.15, 2.25 and Remark~2.17,
pp.~24--29 in arXiv:2110.08643v4 (action-angle coordinates, action periods and local fibred equivalence).
\url{https://arxiv.org/abs/2110.08643v4}.

\bibitem{EvansNotes}
J. Evans,
\emph{Lagrangian torus fibrations (notes in progress)},
8 February 2019.
Lemmas~6.1.5, 6.2.1 and \S6.3.
The numbering here refers to this preliminary set of notes, not to \cite{Evans}.
\url{https://jde27.uk/misc/ltf.pdf}.

\bibitem{EM}
J. D. Evans and M. Mauri,
\emph{Constructing local models for Lagrangian torus fibrations},
Annales Henri Lebesgue \textbf{4} (2021), 537--570.
Lemma~2.16, pp.~543--544; Remarks~2.18--2.19, p.~544;
Lemma~2.22 and Remark~2.23, p.~545;
Lemma~3.1 and \S3.1, pp.~546--547; Remark~2.7;
Assumption~1, Propositions~6.11--6.12 and Theorem~6.14
(published numbering, pp.~561, 565--566).
\url{https://doi.org/10.5802/ahl.80};
\url{https://arxiv.org/abs/1905.09229v3}.

\bibitem{EvansUrzua}
J. D. Evans and G. Urz\'ua,
\emph{Antiflips, mutations, and unbounded symplectic embeddings of rational homology balls},
Ann. Inst. Fourier \textbf{71} (2021), no.~5, 1807--1843.
\url{https://doi.org/10.5802/aif.3429}.

\bibitem{BHS}
J. Brendel, J. Hauber and J. Schmitz,
\emph{Semi-Local Exotic Lagrangian Tori in Dimension Four},
Ann. Inst. Fourier, online first (2026), 48 pp.
\url{https://doi.org/10.5802/aif.3789};
\url{https://arxiv.org/abs/2403.00408}.

\bibitem{FBP}
J. Fernández de Bobadilla and T. Pełka,
\emph{Fibrations by Lagrangian tori for maximal Calabi--Yau degenerations and beyond},
Communications on Pure and Applied Mathematics (2026), e70060,
published online 8 July 2026,
\url{https://doi.org/10.1002/cpa.70060}.
Theorem~6.14(b) (partial Lagrangian fibration),
Theorem~6.14(c) (whole-fiber GH convergence for the constructed Kähler metrics when proper),
Propositions~5.6, 6.13, Example~5.8 and Remarks~5.9, 7.4 are cited from the final preprint,
arXiv:2312.13248v4, 9 July 2026.
\url{https://arxiv.org/abs/2312.13248v4}.

\bibitem{Fleming}
A. Fleming,
\emph{Kähler packings of projective complex manifolds},
European Journal of Mathematics \textbf{7} (2021), 1382--1400.
Lemmas~4.4--4.5, Proposition~4.6 and Construction~4.7, pp.~1394--1396.
本稿では容量記号を移さず，例外曲線の実面積を指定した radial 構成を直接用いる．
\url{https://doi.org/10.1007/s40879-021-00491-8};
\url{https://arxiv.org/abs/1905.03140}.

\bibitem{FordFulkerson1956}
L. R. Ford, Jr. and D. R. Fulkerson,
\emph{Maximal flow through a network},
Canadian Journal of Mathematics \textbf{8} (1956), 399--404.
Theorem~1 (minimal cut theorem), p.~400; proof pp.~400--402.
\url{https://doi.org/10.4153/CJM-1956-045-5}.

\bibitem{Forstneric}
F. Forstnerič,
\emph{Euclidean domains in complex manifolds},
Journal of Mathematical Analysis and Applications \textbf{506} (2022), no.~1, 125660.
The Introduction recalls the holomorphic tubular-neighborhood theorem of Docquier--Grauert.
\url{https://doi.org/10.1016/j.jmaa.2021.125660};
\url{https://arxiv.org/abs/2104.03051v3}.

\bibitem{GollaMartelli}
M. Golla and B. Martelli,
\emph{Pair of pants decomposition of 4-manifolds},
Algebraic \& Geometric Topology \textbf{17} (2017), no.~3, 1407--1444.
\url{https://doi.org/10.2140/agt.2017.17.1407};
\url{https://arxiv.org/abs/1503.05839}.

\bibitem{GV}
Y. Groman and U. Varolgunes,
\emph{Locality of relative symplectic cohomology for complete embeddings},
Compositio Mathematica \textbf{159} (2023), no.~12, 2551--2637.
\S7.3, Lemma~7.19 and Proposition~7.21.
\url{https://doi.org/10.1112/S0010437X23007492};
\url{https://arxiv.org/abs/2110.08891v4}.

\bibitem{GrossSiebert}
M. Gross and B. Siebert,
\emph{From real affine geometry to complex geometry},
Annals of Mathematics (2) \textbf{174} (2011), no.~3, 1301--1428.
\url{https://doi.org/10.4007/annals.2011.174.3.1}.

\bibitem{Hatcher}
A. Hatcher,
\emph{Algebraic Topology},
Cambridge University Press, Cambridge, 2002.
\S\S0.2, 2.2 and 4.2 (CW complexes, cellular homology and homotopy equivalences).
\url{https://pi.math.cornell.edu/~hatcher/AT/ATpage.html}.

\bibitem{Hartshorne}
R. Hartshorne,
\emph{Algebraic Geometry},
Graduate Texts in Mathematics \textbf{52}, Springer, New York, 1977.
Chapter~V, Theorem~1.10 (the Nakai--Moishezon criterion for surfaces).
\url{https://doi.org/10.1007/978-1-4757-3849-0}.

\bibitem{Hattori}
A. Hattori,
\emph{Topology of $\mathbb C^n$ minus a finite number of affine hyperplanes in general position},
Journal of the Faculty of Science, University of Tokyo, Section~IA, Mathematics
\textbf{22} (1975), no.~2, 205--219.
\url{https://repository.dl.itc.u-tokyo.ac.jp/records/39761}.

\bibitem{IKM}
A. Izosimov, B. Khesin and M. Mousavi,
\emph{Coadjoint orbits of symplectic diffeomorphisms of surfaces and ideal hydrodynamics},
Annales de l'Institut Fourier \textbf{66} (2016), no.~6, 2385--2433.
Theorems~3.11 and 3.17 (measured Reeb graphs and weighted Reeb graphs).
Lemma~4.8, p.~2421, the proof of Theorem~4.6, p.~2423,
and \S5, pp.~2424--2425 (leaf periods and regular-leaf cycles; published pagination).
\url{https://doi.org/10.5802/aif.3066};
\url{https://arxiv.org/abs/1504.05629}.

\bibitem{LeungSymington}
N. C. Leung and M. Symington,
\emph{Almost toric symplectic four-manifolds},
Journal of Symplectic Geometry \textbf{8} (2010), no.~2, 143--187.
Theorem~2.1 (pp.~148--149) and Table~1 (p.~144), published version.
\url{https://doi.org/10.4310/JSG.2010.v8.n2.a2};
\url{https://arxiv.org/abs/math/0312165}.
\href{https://archive.ymsc.tsinghua.edu.cn/pacm_download/60/558-2010_Leung_Symington_Almost_Toric_Sympl_Mfd_JSG.pdf}{掲載版PDF}.

\bibitem{LeungWan}
N. C. Leung and T. Y. H. Wan,
\emph{Calabi--Yau components in general type hypersurfaces},
Journal of Differential Geometry \textbf{83} (2009), no.~1, 43--74.
\S\S1 and 3.1.
\url{https://doi.org/10.4310/jdg/1253804351};
\url{https://arxiv.org/abs/0807.1784}.

\bibitem{Maksymenko}
S. Maksymenko,
\emph{Path-components of Morse mappings spaces of surfaces},
Commentarii Mathematici Helvetici \textbf{80} (2005), no.~3, 655--690.
Main Theorem (p.~658), Lemma~4.1 (p.~665), and \S6.1, Lemma~6.4 (p.~668).
\url{https://doi.org/10.4171/CMH/30};
\url{https://arxiv.org/abs/math/9910085}.

\bibitem{MatessiPants}
D. Matessi,
\emph{Lagrangian pairs of pants},
International Mathematics Research Notices \textbf{2021} (2021), no.~15, 11306--11356.
Theorem~1.1 and its scope, pp.~1--2, and the standard cotangent form, p.~6,
in arXiv:1802.02993v1; the planar tropical-curve case is proved in this paper.
\url{https://doi.org/10.1093/imrn/rnz126};
\url{https://arxiv.org/abs/1802.02993v1}.

\bibitem{Matessi}
D. Matessi,
\emph{Lagrangian submanifolds from tropical hypersurfaces},
International Journal of Mathematics \textbf{32} (2021), no.~7, 2150046, 63~pp.
Theorem~1.1 and \S1.2, pp.~1--3 in arXiv:1804.01469v1:
Lagrangian submanifold lifts of tropical hypersurfaces in $\RR^3$,
not a fibration of a compact hypersurface.
\url{https://arxiv.org/abs/1804.01469v1}.

\bibitem{IvanSmithTorus}
I. Smith,
\emph{Torus fibrations on symplectic four-manifolds},
Turkish Journal of Mathematics \textbf{25} (2001), no.~1, 69--96.
\url{https://journals.tubitak.gov.tr/math/vol25/iss1/4/}.
分類の仮定は connected homologically essential Lagrangian tori と
tame singularities である（原論文のabstractと定理の射程を区別する）．

\bibitem{Mikhalkin}
G. Mikhalkin,
\emph{Decomposition into pairs-of-pants for complex algebraic hypersurfaces},
Topology \textbf{43} (2004), no.~5, 1035--1065.
Definition~7, pp.~13--14, Theorems~1 and~1$'$, pp.~17--18,
and \S4.2 in arXiv:math/0205011v3.
The symplectic conclusion supplies a Lagrangian fibre in each maximal cell.
\url{https://doi.org/10.1016/j.top.2003.11.006};
\url{https://arxiv.org/abs/math/0205011v3}.

\bibitem{PapadimaSuciu}
S. Papadima and A. I. Suciu,
\emph{Higher homotopy groups of complements of complex hyperplane arrangements},
Advances in Mathematics \textbf{165} (2002), 71--100.
\S1.1, p.~72 (Hattori's theorem for generic arrangements).
\url{https://doi.org/10.1006/aima.2001.2023};
\url{https://arxiv.org/abs/math/0002251}.

\bibitem{Pascaleff}
J. Pascaleff,
\emph{On the symplectic cohomology of log Calabi--Yau surfaces},
Geometry \& Topology \textbf{23} (2019), no.~6, 2701--2792.
\S4.2.
\url{https://doi.org/10.2140/gt.2019.23.2701};
\url{https://arxiv.org/abs/1304.5298v3}.

\bibitem{Ruan}
W.-D. Ruan,
\emph{Lagrangian torus fibration of quintic hypersurfaces. I. Fermat quintic case},
in \emph{Winter School on Mirror Symmetry, Vector Bundles and Lagrangian Submanifolds}
(Cambridge, MA, 1999), AMS/IP Studies in Advanced Mathematics \textbf{23},
American Mathematical Society, Providence, RI, 2001, 297--332.
\S2, pp.~6--7 (Fubini--Study form), Theorem~3.1, p.~15,
and \S6, pp.~36--37 (piecewise smooth fibration),
in arXiv:math/9904012v2; the hypersurfaces are Calabi--Yau threefolds.
\url{https://arxiv.org/abs/math/9904012v2}.

\bibitem{RSTZ}
H. Ruddat, N. Sibilla, D. Treumann and E. Zaslow,
\emph{Skeleta of affine hypersurfaces},
Geometry \& Topology \textbf{18} (2014), no.~3, 1343--1395.
\url{https://doi.org/10.2140/gt.2014.18.1343};
\url{https://arxiv.org/abs/1211.5263}.

\bibitem{StacksAmpleness}
The Stacks Project Authors,
\emph{The Stacks Project}, 2026.
Lemma~30.17.2, Tag~0B5V (ampleness under finite surjective pullback);
Lemma~37.50.3, Tag~0D2S (relative ampleness near an ample fiber);
Lemma~71.15.3, Tag~0D3A (the algebraic-space version).
\url{https://stacks.math.columbia.edu/tag/0B5V};
\url{https://stacks.math.columbia.edu/tag/0D2S};
\url{https://stacks.math.columbia.edu/tag/0D3A}.

\bibitem{VuNgoc}
S. Vũ Ngọc,
\emph{On semi-global invariants for focus--focus singularities},
Topology \textbf{42} (2003), no.~2, 365--380.
Theorem~2.1, Proposition~3.1, Definition~3.1, Remark~3.2,
Lemma~4.1 and \S\S4, 6--7.
確認した arXiv:math/0208245v1 の PDF では，定理は p.~4，正則化作用は pp.~6--8，
座標の自由度は pp.~9--11，実現構成は pp.~14--16，``Exact version'' は p.~18．
\url{https://doi.org/10.1016/S0040-9383(01)00026-X};
\url{https://arxiv.org/abs/math/0208245}.

\bibitem{ZungFocus}
N.~T. Zung,
\emph{A note on focus-focus singularities},
Differential Geometry and its Applications \textbf{7} (1997), no.~2, 123--130.
\S2 (local quadratic type), \S3, Proposition~3 and Corollary~1 (compact leaf and monodromy),
\S4 (singular reduction).
確認した arXiv PDF の p.~4 の訂正脚注は，位相的同値から微分同相を結論する旧記述を訂正している．
本稿では一般の smooth foliated classification の入力には用いない．
\url{https://doi.org/10.1016/S0926-2245(96)00042-3};
\url{https://arxiv.org/abs/math/0110147}.

\bibitem{Zhang}
Y. Zhang,
\emph{Pair-of-Pants Decompositions of 4-Manifolds Diffeomorphic to General Type Hypersurfaces},
Journal of Mathematical Study \textbf{58} (2025), no.~2, 211--252.
Figure~1.1, p.~2, Theorem~1.1(ii),(iv), p.~4,
and the quintic example in \S1 are checked in arXiv:2102.10037v1.
The singular boundary fibre is $\Theta\times S^1$, with $\Theta$ the theta graph;
the corresponding passages of the published PDF have not been checked.
\url{https://doi.org/10.4208/jms.v58n2.25.06};
\url{https://arxiv.org/abs/2102.10037v1}.

\bibitem{Zhou}
P. Zhou,
\emph{Lagrangian skeleta of hypersurfaces in $(\mathbb C^*)^n$},
Selecta Mathematica (N.S.) \textbf{26} (2020), no.~2, article~26, 33~pp.
\url{https://arxiv.org/abs/1803.00320}.

% Bibliography for the SNC theorem and subsequent applications.

\bibitem{DGMS}
P. Deligne, P. Griffiths, J. Morgan and D. Sullivan,
\emph{Real homotopy theory of K\"ahler manifolds},
Inventiones Mathematicae \textbf{29} (1975), no.~3, 245--274.
Lemma~(5.11), pp.~266--267, and Remark~(5.14).
\url{https://doi.org/10.1007/BF01389853};
\url{https://www.math.stonybrook.edu/~dennis/publications/PDF/DS-pub-0024.pdf}.

\bibitem{TMZ}
M. Farajzadeh Tehrani, M. McLean and A. Zinger,
\emph{Normal crossings singularities for symplectic topology},
Advances in Mathematics \textbf{339} (2018), 672--748.
Equations~(2.9)--(2.10), Definitions~2.8, 2.9, 2.11, 2.12 and Theorem~2.13,
pp.~13--17 in arXiv:1410.0609v3 (normalization and deformation allowed by regularization).
\url{https://doi.org/10.1016/j.aim.2018.09.035};
\url{https://arxiv.org/abs/1410.0609v3}.

\bibitem{FeichtnerSturmfels}
Eva Maria Feichtner and Bernd Sturmfels,
\emph{Matroid polytopes, nested sets and Bergman fans},
Portugaliae Mathematica (N.S.) \textbf{62} (2005), no.~4, 437--468.
\href{https://arxiv.org/abs/math/0411260}{arXiv:math/0411260}.

\bibitem{ArdilaKlivans}
Federico Ardila and Caroline J. Klivans,
\emph{The Bergman complex of a matroid and phylogenetic trees},
Journal of Combinatorial Theory, Series B \textbf{96} (2006), no.~1, 38--49.
\href{https://doi.org/10.1016/j.jctb.2005.06.004}{doi:10.1016/j.jctb.2005.06.004}.
\href{https://arxiv.org/abs/math/0311370}{arXiv:math/0311370}.

\bibitem{Shaw}
Kristin M. Shaw,
\emph{A tropical intersection product in matroidal fans},
SIAM Journal on Discrete Mathematics \textbf{27} (2013), no.~1, 459--491.
\href{https://doi.org/10.1137/110850141}{doi:10.1137/110850141}.
\href{https://arxiv.org/abs/1010.3967v2}{arXiv:1010.3967v2}.

\bibitem{KapranovChow}
Mikhail M. Kapranov,
\emph{Chow quotients of Grassmannians. I},
in \emph{I.~M.~Gelfand Seminar},
Advances in Soviet Mathematics, vol.~16, part~2,
American Mathematical Society, Providence, RI, 1993, pp.~29--110.
\href{https://arxiv.org/abs/alg-geom/9210002}{arXiv:alg-geom/9210002}.

\bibitem{Aksnes}
E. Aksnes,
\emph{Tropicalization of curve arrangement complements and arroids},
Mathematica Scandinavica \textbf{131} (2025), no.~2.
\url{https://doi.org/10.7146/math.scand.a-156674};
\url{https://arxiv.org/abs/2404.14380v3}.

\bibitem{GKM}
A. Gathmann, M. Kerber and H. Markwig,
\emph{Tropical fans and the moduli spaces of tropical curves},
Compositio Mathematica \textbf{145} (2009), no.~1, 173--195.
\url{https://doi.org/10.1112/S0010437X08003837};
\url{https://arxiv.org/abs/0708.2268v2}.

\bibitem{MichalakReeb}
\L{}. P. Michalak,
\emph{Realization of a graph as the Reeb graph of a Morse function on a manifold},
Topological Methods in Nonlinear Analysis \textbf{52} (2018), no.~2, 749--762.
Definition~2.2 and Theorem~5.6.
\url{https://doi.org/10.12775/TMNA.2018.029};
\url{https://arxiv.org/abs/1805.06727v2}.

% Add this entry inside the main thebibliography environment.

\bibitem{ShawSurfaces}
Kristin Shaw,
\emph{Tropical surfaces},
arXiv:1506.07407v1 (2015), 43 pp.
Definition~3.1, pp.~15--16 (ambient integer affine transitions);
\S3.2 (boundary), \S4.2 (tropical summation).
\url{https://arxiv.org/abs/1506.07407v1};
\url{https://arxiv.org/pdf/1506.07407}.

\bibitem{JellRauShaw}
Philipp Jell, Johannes Rau, and Kristin Shaw,
\emph{Lefschetz $(1,1)$-theorem in tropical geometry},
\'Epijournal de G\'eom\'etrie Alg\'ebrique \textbf{2} (2018), Article~11, 27 pp.
Definitions~2.1, 3.1 and 5.1 (rational polyhedral spaces, the affine sheaf, and matroidal charts).
\url{https://doi.org/10.46298/epiga.2018.volume2.4126};
\url{https://epiga.episciences.org/4981/pdf};
\url{https://arxiv.org/abs/1711.07900}.

\bibitem{LRSChern}
Luc\'ia L\'opez de Medrano, Felipe Rinc\'on, and Kris Shaw,
\emph{Chern Classes of Tropical Manifolds},
arXiv:2309.00229v2 (2023), 30 pp.
Definitions~2.2--2.4, pp.~4--5 (extended integer affine maps, tropical manifold atlas, and boundary).
\url{https://arxiv.org/abs/2309.00229v2};
\url{https://arxiv.org/pdf/2309.00229}.

\bibitem{CartwrightComplexes}
Dustin Cartwright,
\emph{Tropical complexes},
Manuscripta Mathematica \textbf{163} (2020), 227--261.
Proposition~2.1; Lemma~2.2; Definitions~2.5, 2.8;
Proposition~2.10; Construction~3.2 and Remark~3.4.
The linked arXiv version is v4 (2019).
\url{https://doi.org/10.1007/s00229-019-01148-3};
\url{https://arxiv.org/abs/1308.3813v4};
\url{https://arxiv.org/pdf/1308.3813}.

% must remain explicit; they do not prove the present branched atlas.
% Existing VuNgoc entry can additionally mention Remark 3.1,
% which displays the focus--focus monodromy matrix.

\bibitem{ShawWerner}
K. Shaw and A. Werner,
\emph{On the birational geometry of matroids},
Combinatorial Theory \textbf{3} (2023), no.~2, Paper No.~17.
Theorem~7.2 (parallel connections and product Bergman fans).
\url{https://doi.org/10.5070/C63261996};
\url{https://arxiv.org/abs/2207.13639v1}.

\bibitem{AllermannRau}
L. Allermann and J. Rau,
\emph{First steps in tropical intersection theory},
Mathematische Zeitschrift \textbf{264} (2010), no.~3, 633--670.
Definition~3.4, Remark~3.5, Proposition~3.7;
Definition~6.1 and Construction~6.4.
\url{https://doi.org/10.1007/s00209-009-0483-1};
\url{https://arxiv.org/abs/0709.3705v3}.

\bibitem{SymingtonFour}
M. Symington,
\emph{Four dimensions from two in symplectic topology},
arXiv:math/0210033v1 (2002).
Proposition~2.15, \S4.3, equation~(4.10), and Proposition~4.14.
\url{https://arxiv.org/abs/math/0210033v1}.

\bibitem{QuinticArithmetic}
\emph{五次曲面の整数アフィン接着と高分母支持グラフ}，
未発表草稿，第25稿，2026年10月28日．
本稿の companion 論文．


\bibitem{RuddatAnnouncement2026}
H. Ruddat,
\emph{Lagrangian torus fibrations of Calabi--Yau hypersurfaces},
Institut Mittag--Leffler, seminar abstract, 30 June 2026.
Joint work with C. Y. Mak, D. Matessi and I. Zharkov, as stated in the abstract.
Announcement; accessed 28 October 2026.
\url{https://www.mittag-leffler.se/seminar/helge-ruddat-tba/}.

\bibitem{ZharkovSoftCells2026}
I. Zharkov,
\emph{Phase Tropical hypersurfaces and their cousins},
ICMS--Sofia Geometry Seminar, 15 April 2026; announcement posted 8 April 2026.
Seminar abstract; accessed 28 October 2026.
\href{https://icms.bg/phase-tropical-hypersurfaces-and-their-cousins-geometry-seminar-talk-by-ilia-zharkov/}{ICMS--Sofia: seminar announcement}.


\bibitem{VenugopalanWoodward2026}
S. Venugopalan and C. T. Woodward,
\emph{Tropical disk potential for almost toric manifolds},
arXiv:2604.03161v1, 3 April 2026.
Definitions~1.1--1.2, Theorem~1.3, and the discussion of the nonmonotone potential in \S1.
\url{https://arxiv.org/abs/2604.03161v1}.

\bibitem{ZhangNegativeKE}
Y. Zhang,
\emph{Collapsing of negative Kähler--Einstein metrics},
Mathematical Research Letters \textbf{22} (2015), no.~6, 1843--1869.
\url{https://doi.org/10.4310/MRL.2015.v22.n6.a16}.
Theorems~1.1, 2.4, 2.6 in \url{https://arxiv.org/abs/1505.04728v2}.

\bibitem{SongNegativeKE}
J. Song,
\emph{Degeneration of Kähler--Einstein manifolds of negative scalar curvature},
arXiv:1706.01518v1 (2017), Theorem~1.2.
\url{https://arxiv.org/abs/1706.01518}.

\bibitem{MandelNegativeKE}
H. Mandel,
\emph{Degenerations of negative Kähler--Einstein surfaces},
Journal of the London Mathematical Society \textbf{109} (2024), no.~1, e12818.
\url{https://doi.org/10.1112/jlms.12818}.

\bibitem{FuHeinJiangCusp}
X. Fu, H.-J. Hein and X. Jiang,
\emph{A continuous cusp closing process for negative Kähler--Einstein metrics},
Geometric and Functional Analysis \textbf{35} (2025), 542--632.
\S1.1, Main Theorem.
\url{https://doi.org/10.1007/s00039-025-00708-y}.

\bibitem{MultiscaleKECurrent}
K. Ueda,
\emph{Multiscale K\"ahler--Einstein Geometry I:
Spatial Infinity, Metrized Dual Complexes, and Volume Asymptotics},
現行未公表原稿（2026年11月22日確認）．
\texttt{thm:intro-metric}，\texttt{cor:arrangement-cone}．
\url{https://github.com/uedakazushi/multiscale-kaehler-einstein/blob/main/papers/multiscale-kaehler-einstein.tex}.

\bibitem{KECubicalCurrent}
K. Ueda,
\emph{多重スケール Kähler--Einstein 幾何 V：超平面和への退化，立方体極限，および直径},
現行未公表原稿（2026年11月22日，v1.11）．定理 \texttt{thm:degree-d-limit}，
仮定 \texttt{hyp:CPd}, \texttt{hyp:PBEd}．
\url{https://github.com/uedakazushi/kaehler-einstein-degeneration/blob/main/papers/cubical-gh-limit.tex}.

\bibitem{RuanHMinimal}
W.-D. Ruan,
\emph{H-minimal Lagrangian fibrations in Kähler manifolds and minimal
Lagrangian vanishing tori in Kähler--Einstein manifolds},
arXiv:math/0309177v2 (2003).
Proposition~3.1, p.~7, Lemma~3.3, p.~8,
Theorem~3.1, p.~9 (local H-minimal fibration),
and Theorem~4.2, p.~11 (minimal vanishing tori), in the cited arXiv version.
\url{https://arxiv.org/abs/math/0309177v2}.

\bibitem{KELocalCurrent}
K. Ueda,
\emph{多重スケール Kähler--Einstein 幾何 IV：超平面和退化の局所安定性と閉 corner の収縮},
現行未公表原稿（2026年11月22日，v1.10）．
\texttt{pair:prop:bundle-conic}，\texttt{pair:cor:global-primitive}，
\texttt{pair:cor:actual-common-moment}，\texttt{pair:lem:diagonal-fixed-end}，
\texttt{neck:lem:rank-one-circle}，
\texttt{coarse:thm:boundary-to-triple}，\texttt{coarse:thm:full-triple-tensor}，
\texttt{coarse:hyp:TC5}，
\texttt{coarse:prop:density-counterexample}，
\texttt{coarse:lem:actual-buffer-comparison}，
\texttt{coarse:cor:degree-d-volume-lower}，
\texttt{coarse:cor:open-square-tensor}，
\texttt{coarse:thm:full-triple-tensor}，
\texttt{simplex:lem:rank-one-schwarz}，
\texttt{simplex:cor:actual-dominant-product}，
\texttt{simplex:prop:facet-product}，
\texttt{simplex:thm:closed-corner-model}，
\texttt{simplex:cor:actual-closed-corner-metric}．
\href{https://github.com/uedakazushi/kaehler-einstein-degeneration/blob/main/papers/kaehler-einstein-degeneration.tex}{現行ソース}.

\bibitem{SavinSmall}
O.~Savin,
Small perturbation solutions for elliptic equations,
\emph{Comm. Partial Differential Equations} \textbf{32} (2007), 557--578,
Theorem~1.3.
\href{https://doi.org/10.1080/03605300500394405}{doi:10.1080/03605300500394405}.
\href{https://www.math.columbia.edu/~savin/Small_pert.pdf}{Author PDF}.

\bibitem{BabsonBilleraChan}
E.~K.~Babson, L.~J.~Billera and C.~S.~Chan,
\emph{Neighborly cubical spheres and a cubical lower bound conjecture},
Israel Journal of Mathematics \textbf{102} (1997), 297--315.
\href{https://doi.org/10.1007/BF02773804}{doi:10.1007/BF02773804}.
本文確認は1996年4月10日preprint, \S2.3, p.~6,
Proposition~2.2, Corollary~2.4.
\href{https://www.math.ucdavis.edu/~deloera/MISC/LA-BIBLIO/trunk/Billera/BilleraEhrenborgcdindex.pdf}{一次scan PDF}（連結PDF第24頁）．
書誌は\href{https://pi.math.cornell.edu/~billera/research.html}{著者の刊行論文一覧}と照合．

\bibitem{RochonZhang}
F.~Rochon and Z.~Zhang,
\emph{Asymptotics of complete K\"ahler metrics of finite volume on quasiprojective manifolds},
Advances in Mathematics \textbf{231} (2012), no.~5, 2892--2952.
\S5, (5.3)--(5.6), Theorem~5.1 (normal-crossing case).
\url{https://doi.org/10.1016/j.aim.2012.08.005}.
\url{https://arxiv.org/abs/1106.0873v2}.

\bibitem{GGZDiameter}
V.~Guedj, H.~Guenancia and A.~Zeriahi,
\emph{Diameter of K\"ahler currents},
arXiv:2310.20482v1 (2023), Proposition~1.4, p.~5.
\url{https://arxiv.org/abs/2310.20482v1}.

\bibitem{YauSchwarz}
S.-T.~Yau,
\emph{A general Schwarz lemma for K\"ahler manifolds},
American Journal of Mathematics \textbf{100} (1978), no.~1, 197--203.
Theorem~2 and its proof, pp.~199--201.
\href{https://doi.org/10.2307/2373880}{doi:10.2307/2373880}.
\href{https://archive.ymsc.tsinghua.edu.cn/pacm_download/59/11034-Shing-Tung_Yau_30.pdf}{著者登録の掲載版PDF}.

\end{thebibliography}

\end{document}
